% =====================================================================
%  R. Nielsen, OM HINDRINGER OG BETINGELSER FOR DET AANDELIGE LIV I
%  NUTIDEN. Sexten Forelæsninger, holdte ved Universitetet i Christiania
%  September--Oktober 1867. Kjøbenhavn. Forlagt af den Gyldendalske
%  Boghandel (F. Hegel). Trykt hos J. H. Schultz. 1868.
%  4 unnumbered preliminary leaves + 326 pp. + Rettelser.
%
%  Diplomatic transcription from Det Kongelige Bibliotek's digitisation
%  (KB label and barcode on the front pastedown; KB stamps on the title
%  page and its verso), `bibliotek/Nielsen, Rasmus/1868-om-hindringer.pdf`,
%  343 PDF pages, each page one 300 ppi colour JPEG with an embedded text
%  layer, sha256
%  594ccf9f09687286b9e7c0c3dd2b3cedbf23ca8d3deecde3f30def62902bbe76.
%
%  The exact title, from the title page (PDF 6), line by line:
%      Om
%      Hindringer og Betingelser
%      for
%      det aandelige Liv i Nutiden.
%      Sexten Forelæsninger,
%      holdte ved Universitetet i Christiania
%      September—Oktober 1867          [letterspaced]
%      af
%      R. Nielsen.                     [letterspaced]
%      [publisher's device: a monogram in a buckled garter, legend
%       "1770 · SG · JD · EH" (?), not transcribed]
%      Kjøbenhavn.
%      Forlagt af den Gyldendalske Boghandel (F. Hegel).
%      Trykt hos J. H. Schultz.        [small; a KB stamp over it]
%      1868.
%
%  This header is the edition's apparatus. The transcription harness
%  (pagemap.py, ocr.sh, spacing.py, check.py, splice.py, joints.py,
%  layerdiff.py) is not tracked in git and is unrunnable without the scan,
%  so everything an outside reader needs in order to check this text
%  against the scan is recorded here.
%
% ---------------------------------------------------------------------
%  1. PAGE MAP  (verified by eye 22 Sep 2026; see pagemap.py)
% ---------------------------------------------------------------------
%
%      printed 1--326   =   PDF page - 11      (PDF 12--337)
%
%  Uniform. The head of every one of PDF 12--337 was rendered and its
%  folio read on the image: 2--295 and 298--326 in unbroken sequence on
%  PDF 13--306 and 309--337, no leaf missing or doubled, no misprinted
%  folio. Three pages carry no folio: p. 1 (PDF 12, opening of
%  Forelæsning I), p. [296] (PDF 307, BLANK), p. [297] (PDF 308, opening
%  of the Efterskrift). The other fifteen lectures open on fresh pages
%  that do carry their folio.
%
%  Leaves (PDF pages):
%      PDF   1      KB "Digitaliseret af | Digitised by" sheet
%      PDF   2--5   cover, pastedown (KB label), endpaper, blank
%      PDF   6      TITLE PAGE                                [transcribed]
%      PDF   7      title verso, blank (KB stamp)
%      PDF   8      FORORD, one page                           [transcribed]
%      PDF   9      blank
%      PDF  10      DEDICATION                                 [transcribed]
%      PDF  11      blank
%      PDF  12--337 text, pp. 1--326 (p. [296] blank)
%      PDF 338      RETTELSER, unnumbered (p. [327])          [transcribed]
%      PDF 339--341 blank
%      PDF 342--343 back pastedown (pencil shelf notes), back cover
%  The preliminary leaves stand in this copy in the order title, Forord,
%  dedication; that order is kept. There is no Indhold.
%
% ---------------------------------------------------------------------
%  2. STRUCTURE  (from the images, not the OCR)
% ---------------------------------------------------------------------
%
%  Sixteen lectures and an Efterskrift. Each opens on a fresh page with a
%  centred head consisting of the Roman numeral and a full stop ("I.",
%  "II.", ... "XVI.") and NOTHING ELSE -- no title, no argument -- followed
%  by the text, which begins with a large initial capital (set here as an
%  ordinary letter, with a % comment). Lecture V opens inside a quotation,
%  „I skulle“, with the large initial on the I. The head "Efterskrift." is
%  set the same way. Lectures end with a short centred closing rule where
%  the page has room for one (found on pp. 16, 72, 89, 118, 134, 153, 207,
%  255, 278, 295, 326).
%
%      I     pp.   1--16      IX    pp. 135--153
%      II    pp.  17--35      X     pp. 154--175
%      III   pp.  36--53      XI    pp. 176--193
%      IV    pp.  54--72      XII   pp. 194--207
%      V     pp.  73--89      XIII  pp. 208--224
%      VI    pp.  90--104     XIV   pp. 225--255
%      VII   pp. 105--118     XV    pp. 256--278
%      VIII  pp. 119--134     XVI   pp. 279--295   (p. [296] blank)
%                             Efterskrift  pp. [297]--326
%
%  Heads are set \chapter*{I.} etc., exactly as printed. The table of
%  contents entries ("Forelæsning I" ...) and the running heads are
%  EDITORIAL: the book has no Indhold and its only running head is the
%  folio.
%
%  RETTELSER (PDF 338), the book's own errata, four entries:
%      S. 50  Linie 1 f. o.    vente at finde den  læs: vente det
%      - 70   —  1 - -         stridt mod          læs: stridt med
%      - 109  —  3 f. n.       det absolute Selv   læs: det Absolute selv
%      - 174  —  9 og 10 f. o. Saducæer            læs: Sadducæer
%  The text is transcribed AS PRINTED at each of these places, with a %
%  comment citing the Rettelser; the errata leaf is transcribed at the
%  end. In this copy a reader has pencilled the corrections into the
%  margins (p. 50 "det", p. 70 "e"): those marks are not transcribed.
%
% ---------------------------------------------------------------------
%  3. TYPOGRAPHY AND TRANSCRIPTION CONVENTIONS
% ---------------------------------------------------------------------
%
%  Body in ANTIQUA (roman) throughout. aa (not å). Emphasis is by
%  LETTERSPACING -> \emph{} (e.g. p. 1 "Hindringer og Betingelser for det
%  aandelige Liv i Nutiden"). A true ITALIC -> \textit{} is used (a) for
%  French ("Le présent est gros de l'avenir", p. 17; p. 34; p. 182;
%  "Mécanique céléste", p. 207, so printed); (b) for SOME Latin tags
%  ("Credo ut intelligam", p. 32; "Legant in parietibus ...", p. 124;
%  "Equus Caballus", p. 157; "Si fractus illabatur orbis ...", p. 211;
%  "in abstracto", "concordia discors", "horribile dictu", pp. 238--254;
%  p. 272); and (c) for emphasis in Danish, chiefly in Lectures VII, IX,
%  XI and XV (e.g. p. 112 "Gud", "hellig"; p. 148 "Du skal!"; the
%  numbered heads 1)--3), pp. 140, 146, 151; p. 190). Other Latin
%  ("nervus rerum gerendarum", p. 5; "descendit ad inferos", p. 151;
%  "exceptio fori", p. 319) and German (p. 32, p. 287) are set in the
%  same ROMAN as the body and are NOT marked, unless letterspaced. Each
%  occurrence follows its own page.
%  Faint raised marks between words are risen spaces (quads inked in the
%  forme): they are not typed, and are noted in a % comment at the site. Greek is set in Greek type (p. 22 "αὐτοὺς δὲ
%  ἑλώρια τεῦχε κύνεσσιν", pp. 63, 94) and is typed in Unicode Greek with
%  its accents and breathings, as printed; the Plato on p. 320 is printed
%  WITHOUT accents (smooth breathings only) and is so transcribed.
%  The name is printed "Bröchner" (ö) on p. 111; each occurrence of any
%  name is transcribed as printed.
%
%  Quotation marks „...“, low-high. Long rule (pause, and ranges such as
%  "Iliad. I, 1—5") -> ---. The id-est mark ɔ:, where printed, is typed
%  U+0254. No footnotes were found at set-up (no footnote rule on any
%  page, no note marks in the layer); should one appear, \thefootnote is
%  defined to give *), **), ***) in Nielsen's usual style, reset per
%  printed page by \setcounter{footnote}{0}% after each page marker.
%  Line-break hyphens are resolved (Cau\-%). Signature marks at the page
%  foot ("2", "1*"), the folio, library stamps and the reader's pencil
%  marks are not transcribed.
%
% ---------------------------------------------------------------------
%  4. WORD-ACCURACY
% ---------------------------------------------------------------------
%
%  Word-accuracy: every batch diffed against the embedded (KB) text layer,
%  re-ordered into printed lines (layerdiff.py -> ocrdiff.py --ocr), before
%  splicing. 28 batches (26 body + front matter + Rettelser), 767
%  candidates, 2 transcription errors corrected, 0 unresolved. See
%  TRANSCRIPTION-PLAYBOOK.md §3 step 4. Because the batch agents drafted
%  with the layer in hand, that diff is not wholly independent of them;
%  each agent therefore also re-read every page of its range against the
%  image (11 errors fixed there before the diff).
%  Then a full SECOND COLLATION of all 326 pages and the four unpaginated
%  leaves at the image (300 dpi, 600 dpi for glyphs, NO OCR consulted), by
%  28 readers who had not transcribed, 22 Sep 2026: 9 errors found and
%  corrected, 0.03 per page -- a silently corrected printer's error
%  (p. 324 Universitetsaeret), two accents (p. 88 Rénan not Rènan,
%  p. 202 Nødvendighedéns), three risen spaces that had been typed as
%  punctuation (p. 36 an apostrophe, p. 41 a semicolon, p. 135 a comma),
%  a spacing "Absolut - Absolute" (p. 141), a missed paragraph break
%  (p. 228), four points for three (p. 325). No dropped or added word
%  was found in any batch.
%  Sensitivity control: six errors of the kinds this corpus has suffered
%  (a dropped "ikke", a dropped "vel", "Charakteer" modernised,
%  "begribeligt" losing its -t, a semicolon made a comma, an added "og")
%  were planted in pp. 105--118 without the collator's knowledge; all six
%  were found and reverted, and nothing else was changed.
%  After splicing, joints.py removed 9 spurious paragraph breaks at batch
%  seams. Items the collators could not settle are kept as transcribed
%  with a "% AUDIT doubtful" comment at the site (pp. 54, 84, 103, 108,
%  114, 211, 299, 300, 302, 323, 324).
%
% ---------------------------------------------------------------------
%  5. PRINTER'S DEFECTS LOG
% ---------------------------------------------------------------------
%
%  All transcribed as printed, with a % comment at the site; the four
%  Rettelser places are listed in §2. Printer's errors (not in the
%  Rettelser): p. 23 "Ro et nyde" (et for at); p. 46 "Horledes";
%  p. 55 "denné"; p. 63 "Σεύς" (Σ for Ζ); p. 83 "Virkeligked", broken
%  "paa"; p. 88 "Rénan"; p. 124 "Mennskebud"; p. 143 displaced type in
%  "den Viden"; p. 149 "hehøver", "Fordring." (point for comma?);
%  p. 165 no stop after "med"; p. 202 "forkjellige", "Nødvendighedéns";
%  p. 207 "céléste"; p. 211 "au" (turned n); p. 228 "videnskalige";
%  p. 262 "af gjøre"; p. 267 "communication", "Dictionaire"; p. 277 no
%  stop after "forsvinder"; p. 292 "Váaben"; p. 297 "Vidnetsbyrdets";
%  p. 310 "ndrykket"; p. 319 "ndvikle"; p. 320 "Toen"; p. 324
%  "Universitetsaeret". Author's references as printed: p. 100 "XXIII"
%  (Gen. XVIII), p. 132 "1 Moseb.", p. 168 "Psalm. IV", p. 170 "1 Sam.",
%  p. 172 "Joh. XIX.".
%  QUOTATION BALANCE: „ 559, “ 561, balance -2, entirely from logged
%  defects: opening mark missing pp. 66, 97, 176, 300 (-4); closing mark
%  missing pp. 200, 217 (+2).
%
% =====================================================================
\documentclass[11pt]{book}
\usepackage[utf8]{inputenc}
\usepackage[T1]{fontenc}
\usepackage{libertinus}
\usepackage{libertinust1math}
\usepackage{textalpha}
\usepackage{graphicx}
\DeclareUnicodeCharacter{0254}{\reflectbox{c}}
\usepackage{amsmath}
\usepackage[danish]{babel}
\usepackage{fancyhdr}
\setlength{\headheight}{28pt}
\addtolength{\topmargin}{-13.5pt}
\pagestyle{fancy}
\fancyhf{}
\fancyhead[LE,RO]{\thepage}
\fancyhead[RE]{\textit{Om Hindringer og Betingelser}}
\fancyhead[LO]{\textit{\leftmark}}
\renewcommand{\thefootnote}{\ifcase\value{footnote}\or *)\or **)\or ***)\else
  \arabic{footnote})\fi}
\usepackage{setspace}
\setstretch{1.3}
\usepackage{marginnote}
\newcommand{\opage}[1]{\marginnote{\footnotesize\textit{[#1]}}}
\newcommand{\apage}[1]{\marginnote{\footnotesize\textit{[#1]}}}

\begin{document}

\frontmatter
% --- front matter: title page (PDF 6), Forord (PDF 8), dedication (PDF 10),
% in the order of the leaves in this copy. Unpaginated: no page markers.
% --- PDF 6: title page ---
\begin{titlepage}
\centering
Om\par
{\Large Hindringer og Betingelser}\par
for\par
{\LARGE det aandelige Liv i Nutiden.}\par
\bigskip
{\large Sexten Forelæsninger,}\par
holdte ved Universitetet i Christiania\par
\emph{September---Oktober 1867}\par
\medskip
af\par
\emph{R. Nielsen.}\par
\vfill
% [publisher's device: Gyldendal's circular badge, "1770 · SG · JD · FH" round a GB monogram]
\vfill
Kjøbenhavn.\par
Forlagt af den Gyldendalske Boghandel (F. Hegel).\par
{\small Trykt hos J. H. Schultz.}\par
1868.\par
% (a circular library stamp over "Boghandel" / "hos J. H. Schultz" is not transcribed)
\end{titlepage}

% --- PDF 8: Forord ---
\chapter*{Forord.}
% the head "Forord." is letterspaced; a short rule stands below it
\addcontentsline{toc}{chapter}{Forord}

Endskjøndt disse Forelæsninger i deres nærværende
Form ikke i Eet og Alt kunne siges at være en lige\-%
frem Gjengivelse af det mundtlige Ord, har jeg dog
ved at afbenytte deels egne Optegnelser deels udfør\-%
ligere Referater i norske Blade søgt at komme det
oprindelig frie Foredrag saa nær som muligt. Her og
der ere imidlertid Begrebsanalyserne lidt videre udførte,
og enkelte Partier, navnlig i de to næstsidste Foredrag,
indskudte, ligesom jeg ogsaa har fundet mig opfordret
til at imødegaae adskillige Betænkeligheder og epilo\-%
giske Mistydninger i en særegen Efterskrift.

\medskip
\noindent\hspace*{2em}Kjøbenhavn den 17de Januar 1868.

\begin{flushright}
R. Nielsen.
\end{flushright}

% --- PDF 10: dedication ---
\cleardoublepage
\thispagestyle{empty}
\vspace*{\fill}
\begin{center}
Mine norske\par
\bigskip
{\Large Tilhørere og Tilhørerinder}\par
% wide display type, but not letterspaced (compared at 300 dpi with the line below)
\medskip
{\small tilegnes}\par
\medskip
{\large disse Forelæsninger}\par
\bigskip
\emph{med dyb Erkjendtlighed}\par
\bigskip
af\par
\bigskip
Forfatteren.\par
\end{center}
\vspace*{\fill}
\cleardoublepage

\tableofcontents

\mainmatter

\chapter*{I.}
\addcontentsline{toc}{chapter}{Forelæsning I}
\markboth{Forelæsning I}{Forelæsning I}
\label{ch:i}
% pp. 1--16
% --- p. 1 ---
\opage{1}\setcounter{footnote}{0}%
% [large initial E]
En Tanke, en stor og frugtbar Tanke, har samlet her
en talrig Hæderskreds. Det er Lærere og Studerende
ved denne Høiskole, gamle Norges, unge Høiskole, som
i skjøn Forening have sat Tanken i Bevægelse, den
sikkert følgerige Fremtidstanke, at et nærmere viden\-%
skabeligt Samqvem imellem de nordiske Universiteter
maa særlig bidrage til gjensidig Forhøielse af vort aande\-%
lige Liv, Vækkelse af aandige Kræfter. Saa udgik Ind\-%
bydelsen, ogsaa til mig; og da jeg nu, greben af den
skjønne Tanke, ret maatte paaskjønne den sjeldne Hæder,
at det ved en saa betydningsfuld Anledning blev mig
forundt at fremkomme med mine Bidrag og føre Ordet
i en Forsamling som denne: da vidste jeg ikke at vælge
nogen anden eller mere passende Opgave end et Forsøg
til en almindelig Fremstilling af \emph{Hindringer og Be\-%
tingelser for det aandelige Liv i Nutiden}.

Ja, i Nutiden har det aandelige Liv tilvisse eien\-%
dommelige, gunstige Betingelser, men møder just derfor
ogsaa store eiendommelige Hindringer. See vi bort fra
Hindringerne for udelukkende at fæste vort Blik paa
% --- p. 2 ---
\opage{2}\setcounter{footnote}{0}%
Betingelserne, da aabne sig lyse Udsigter; af de lyse
Udsigter fødes et glimtende, flygtigt Haab, men et flygtigt
Haab kan bedaare. Glemme vi Betingelserne og fortabe
os i en Betragtning af de mangfoldige Hindringer, da
blive Udsigterne mørke; af de mørke Udsigter fremkom\-%
mer Mismod; men Mismod er som en Taage for Øiet.
Vi ville da ikke hengive os til eensidig Overveielse enten
af Betingelserne alene, saa at vi derover glemme Hin\-%
dringerne, eller af Hindringerne alene, saa at vi oversee
Betingelserne; nei, vi skulle i al Besindighed søge at veie
Hindringer mod Betingelser, Betingelser mod Hindringer,
saa kommer der Ligevægt.

Det er et gammelt Ord, at den, der vil Øiemedet,
maa ville Midlerne; men Ordet har en dobbelt Mening.
Det kan forstaaes saaledes, at Intet kan komme frem i
Øiemedet, som ikke ligger i Midlerne, at det virkelig\-%
gjorte Øiemed er Midlernes Resultat, saa at det først
og sidst kommer an paa Midlerne: det er Forstandens
Opfattelse, det er Klogskabens Raad, og Viisdommen
selv kan ikke opgive Klogskaben. „Thi“ --- siger det
hellige Ord --- „hvo iblandt Eder, som vil bygge et
Taarn, sidder ikke først og beregner Bekostningen, om
han haver, hvad der hører til at fuldføre det med; at
ikke, naar han faaer lagt Grundvold og ei kan fuldende
det, Alle, som see det, skulle begynde at spotte ham og
sige: dette Menneske begyndte at bygge og kunde ikke
fuldende det.“ Men det kan da ogsaa forstaaes saaledes,
at Midlerne ere det Underordnede, Øiemedet det Over\-%
ordnede; det sande selvgyldige Maal har i evig Forstand
Midlerne i sin Magt: først Øiemedet, det første Afgjø\-%
rende, det Ene Fornødne, og Betingelserne, Midlerne,
de endelige Fornødenheder, komme som af sig selv! dette
er Aandens Fordring. Aandens Rige er i høieste For\-%
stand Retfærdighedens Rige, Guds Rige. I Guds, Ret\-%
% --- p. 3 ---
\opage{3}\setcounter{footnote}{0}%
færdighedens Rige maa Klogskaben tjene Viisdommen,
Bekymringen for det Timelige forsvinde i Stræben mod
det Evige; thi --- siger det hellige Ord --- „søger først
Guds Rige og hans Retfærdighed, saa skulle og alle
disse (timelige) Ting tillægges Eder.“

Den angivne Vexelbestemmelse af Midler og Øiemed
kan nu tjene til at opklare os Forholdet mellem Hin\-%
dringer og Betingelser for det aandelige Liv. I virksom
Forening med det overgribende, altbestemmende Øiemed
ere de endelige Midler lutter gunstige Betingelser; det
Sandselige tjener det Aandelige derved, at det gjennem\-%
trænges, beaandes og bevæges af Aanden. Men tages
Betingelserne i endelig Mening som et Indbegreb af
Midler, der ved sig selv i samlet Sum skulle udgjøre
Øiemedet, da staae de Aanden imod og blive for Aands\-%
livet lutter Snarer, lutter Skranker, lutter Hindringer.
Fra dette Synspunkt kan det da ikke falde vanskeligt at
charakterisere Nutiden.

Nutiden er i særegen Forstand at kalde Klogskabens
Tid, en Tid, der seer paa Midlerne og arbeider med Be\-%
tingelserne. Herfra har denne Tid paa een Gang sin Lys- og
% [suspended compound: "Lys-" ends the printed line, "og" begins the next]
sin Skyggeside. Vor Tidsalder kalder sig med Rette en
Videnskabernes, Oplysningens og Culturens Alder; Cul\-%
turens Blik er et Blik for Betingelser; den praktiske
Culturbevidsthed forstaaer tilgavns at erobre Virkelig\-%
heden, forøge Midlerne og bringe dem i Anvendelse.
Hvad har den nyere Tid ikke allerede naaet i denne
Retning? Naturvidenskabernes Fremskridt, disse sindrige
Opfindelser, disse storartede, i Sandhed forbausende Op\-%
dagelser, hvilken Frugtbarhed paa Midler, hvilken Rig\-%
dom af Betingelser! At Skibene gaae imod Vind og
Strøm; at Damphesten, lig en fnysende Leviathan, farer
gjennem Landene fra Sted til Sted og bringer de uhyre
Afstande til at forsvinde; at Tankernes Ilbud udsendes
% --- p. 4 ---
\opage{4}\setcounter{footnote}{0}%
og naaer i Tankens Nu til en anden Verden: dette og
Mere vidner tilstrækkelig om Snillets Magt, om Forsk\-%
ningens Frugt; ja, det forskende Blik er et Blik for
Betingelser.

Den levende Sands for Betingelser og det Betingede
indeholder naturlig en Fristelse til at oversee det Abso\-%
lute, det Ubetingede, altsaa en Fristelse til at miskjende
Aanden. At miskjende Aanden er at drage den saaledes
ned under naturbundne Vilkaar, at Frihedens Love uden
videre falde sammen med Naturnødvendighedens Love.
Hvad der gjælder om det udvortes Naturlige, vil Cultur\-%
bevidstheden kun altfor gjerne overføre paa det Indvortes
og Aandelige. Ved skærpet Blik for Betingelser opdage
vi Naturens Love; ved Naturens Love erkjende vi Natu\-%
rens Kræfter; ved Indsigt i Naturens Kræfter blive vi
mere og mere istand til at beherske Materien og vise
os som Herrer i Naturens Rige. Det ligger da nær at
spørge, om ikke noget Lignende kunde skee ogsaa med
det Aandelige? „Skulde det dog ikke omsider kunne
lykkes vort forskende Blik at opdage ogsaa Aandens
Love, ved Aandens Love at bestemme og maale Aandens
Kræfter, ved Aandens Kræfter at magte Aanden selv og
hæve os til Selvherskere i Aandens Rige, forstaaer sig,
paa Betingelser?“ Her er Anmasselsen. Det er Men\-%
nesket umuligt at opkaste sig til Aandens Herre; det er
tvertimod hans høieste Bestemmelse at være dens Tjener.

At Nutiden miskjender Aandens Ret, bliver i al
Ubestemthed udtrykt derved, at Tidsalderen siges at
være for materialistisk. Hvori bestaaer dens Materialisme?
Den kan da umuligt bestaae deri, at den anseer Materien
for noget Virkeligt, thi deri har den jo Ret; den kan da
heller ikke sættes deri, at den lægger Vægt paa materielle
Goder, arbeider paa at forøge de materielle Betingelser; thi
deri gjør den visselig Ret. Nei, Tidens Materialisme bestaaer
% --- p. 5 ---
\opage{5}\setcounter{footnote}{0}%
just deri, at den maaler det Evige med Timelighedens
Maal, gjør det Absolute relativt og drager det Ubetin\-%
gede ned under endelige Betingelser. For lutter Midler
er den ikke istand til at see Maalet. Hvilken eensidig
Vægt den nærværende Tid lægger paa Midlerne, sees
blandt Andet af den umaadelige Iver, hvormed den i
Smaat og i Stort, ad alle Veie og paa alle Maader, stræber
at erhverve sig Midlernes Middel, ikke det selvgyldige,
men det sølvgyldige Øiemed (nervus rerum geren\-%
darum). Nutidens Magt er i særlig Forstand Capitaler\-%
nes Magt. Og dog er Nutidens Charakteer ikke egentlig
Gjerrighed. Gjerrighedens Særkjende er, at den gjør
Guldet til Maal, alt Andet til Middel; men den nærvæ\-%
rende Tid gjør kun Guldet til høieste Maal, forsaavidt
den gjør det til høieste Middel. Midlernes Tid er Inter\-%
essernes Tid; Alt har sin Interesse, forsaavidt det har
sin relative Værdi; Alt har sin relative Værdi, forsaavidt
det kan betale sig. Kan nu Friheden betale sig? Hvad
koster Krigen? Hvormange Procent giver Retfærdighe\-%
den, naar den sættes paa Rente? Dersom virkelig Alt
skal kunne betale sig, maa jo ogsaa, det ligger i Sagens
Natur, Alt kunne betales. Hvormeget betaler man saa
for en Mands Ære, en Qvindes Dyd? Ved at overveie,
hvor conseqvent disse Spørgsmaal fremgaae af den hele
Grundanskuelse, og hvor de føre hen, vil man idetmindste
kunne finde det forklarligt, at Digtere og dybere Naturer
ansee Interessernes Guldalder for Aandens Jernalder,
mistvivle om at faae det Ubetingede med i en Verden af
lutter Betingelser og klage mismodig over Tidens aande\-%
lige Slaphed. Det er som et Sjælens Suk i Digterens
Klage, at Slægten ældes, at Ungdommen alt er gammel,
Charaktererne blege, at de travle i Frihed og Selvstæn\-%
dighed interesserede Individer allerhelst gaae med Vin\-%
den, overladende til Skibene at gaae imod Vinden. Der
% --- p. 6 ---
\opage{6}\setcounter{footnote}{0}%
er et Sjælens Suk i Digterens Klage, at Nutidens Aand
er mat og kold, at Aandens Kulde i Livets Høst er
Vinterens Forbud. Der er et Sjælens Suk i Digterens
Klage, at Skabelsens Dage forkortes, at Lyset viger og
Skyggerne voxe, at „Solen daler dybere og dybere, at
Nætterne blive længere og længere, og de kolde mørke
Magter træde mere og mere ind i Livet.“

Dog, hvad hjælper det at klage; selvisk Forstemt\-%
hed, overvældende Mismod forraade jo ogsaa Krænkelse
af Aanden. Ere Hindringerne store, nu vel, saa ere Be\-%
tingelserne desto rigere. Thi Hindringer og Betingelser
forholde sig til hinanden paa en ganske anden Maade i
den aandelige end i den sandselige Verden. Jo mere
sandseligt Modsætningen opfattes, desto stærkere ville de
modsatte Led optræde imod hinanden og udelukke hin\-%
anden; jo mere aandelig derimod Opfattelsen er, desto
klarere vil det indsees, at de modsatte Led forholde sig enten
som Følge til Grund eller som Magt til Afmagt. I sandselig
Mening gjælder det vel, at Hindring ikke er Betingelse,
Betingelse ikke Hindring: jo stærkere Hindringer, desto
svagere Betingelser, og omvendt; men aandelig forstaaet
gjælder det, at Hindringen i Mulighed er Betingelse,
Betingelsen Hindring: jo større Hindringer, desto væsent\-%
ligere Betingelser, og omvendt.

Vi kunne oplyse dette ved nogle Exempler. Man
sætte blot Modgang mod Medgang, Liden mod Handlen,
Frygt mod Haab, og Forskjellen mellem sandselig og
aandelig Forstaaelse vil være iøinefaldende. Sandselig
seet, staaer Modgang imod Medgang som Armod imod
Rigdom, Sygdom imod Sundhed, Sorg imod Glæde.
Hvor skarpt udelukkende Modsætningen er, fremlyser
umiddelbart deraf, at en Formerelse af det ene Led nød\-%
vendig forudsætter en Formindskelse af det andet: jo
mere Modgang, desto mindre Medgang, jo større Sorg,
% --- p. 7 ---
\opage{7}\setcounter{footnote}{0}%
desto mindre Glæde. I Aandens Verden er det ander\-%
ledes. Her er Medgang ikke uden videre en Formind\-%
skelse af Modgang, Rigdom af Fattigdom, Sygdom af
Sundhed, Sorg af Glæde; nei Forholdet er i al Oprinde\-%
lighed tvertimod dette, at forstærket Modgang bliver
Grundlag og Forudsætning for en tilsvarende forhøiet
Medgang; Aandens Magt aabenbarer sig just deri, at
den forvandler Modgang til Medgang, grundlægger en
aandelig Medgang paa sandselig Modgang, en aandelig
Rigdom paa sandselig Armod, en aandelig Sundhed paa
sandselig Sygdom, en aandelig Glæde paa sandselig Sorg.
Saa lyder det hellige Ord: „Salige I, som sørge; I
skulle husvales. Salige I, som hungre; I skulle mættes.
Salige I, som græde; thi I skulle lee.“ Vistnok er
Sorgen som Sorg ikke selv Glæde, thi Glæden er i
Husvalelsen; men Husvalelsens Glæde er betinget af
Sorgen. Her er da ikke blot Tale om en nu\-%
værende Modgang og en tilkommende Medgang, en nu\-%
værende Sorg og en tilkommende Glæde; men, naar vi
see nærmere til, tillige om en nærværende, af Modgang
betinget Medgang, en nærværende, af Sorgen selv betinget
Glæde. Det er saa langt fra, at den saaledes nærvæ\-%
rende Medgang skulde være bleven formindsket, fordi
den falder sammen med Modgang, at den tilsvarende
Glæde tvertimod er hævet til sit Høieste, til Salighed.
Klarere kan da Aanden ikke bevise sin Magt ved at
forvandle Hindringer til Betingelser, end naar den saa
vidunderlig omdanner Modgang til Medgang, at det endog
bliver saligt at sørge, saligt at hungre, saligt at græde,
saligt at lide.

For Aanden, vi mene den absolut guddommelige
Aand, fordi et aandeligt Liv uden den ligesaalidt er
muligt i Nutiden, som det nogensinde har været det i
Fortiden --- for Aanden selv er Intet i den naturlige
% --- p. 8 ---
\opage{8}\setcounter{footnote}{0}%
Verden ligefrem enten Hindring eller Betingelse; men Alt
kommer an paa en Forvandling. Derfor kan det Gud\-%
dommelige i en vis Forstand bruge Alt og Alle. I Aan\-%
dens Rige er Ingen for svag, naar Aanden forvandler
ham hans Svaghed til Styrke, og Ingen for stærk, naar
Aanden forvandler hans Styrke til Svaghed. Her kan
fra dette Synspunkt ikke være Tale om nogen væsentlig
Forskjel imellem gode og onde, beleilige og ubeleilige
Tider; enhver Tid er i Grunden lige god for det Øie\-%
med, der frembringer sine Midler, altid beleilig for det
Ubetingede, der selv skaber sine Betingelser. Aanden
er i sin Virken ingenlunde ligegyldig mod Betingelser,
ikke betingelsesløs, heller ikke skaber den Betingelserne
ligefrem af Intet; nei den danner dem af naturlige For\-%
udsætninger, den danner dem ved en Forvandling. Hvad
der i de naturlige Tilstande med de naturlige Vilkaar
lader sig forvandle, bliver, hvormeget det end i det Ud\-%
vortes tager sig ud som Hindring, dog til Betingelse;
hvad der forbliver i sin Naturlighed og ikke lader sig
forvandle, bliver, om det for en sandselig Betragtning
seer ud som den allergunstigste Betingelse, dog væsentlig
en Hindring.

Men, vil man maaskee indvende, hvis der i den
naturlige Verden virkelig er Noget, som den guddomme\-%
lige Aand ikke kan forvandle, Noget, der kan gjøre den
Modstand, sætte den en Skranke og stille sig som Hin\-%
dring, saa er Aanden jo ikke almægtig; altsaa heller ikke
i Sandhed guddommelig. Er Aanden i Sandhed gud\-%
dommelig, altsaa almægtig, formaaer den overalt at hæve
enhver Modstand, gjennembryde enhver Skranke og skabe
Betingelser af lutter naturlige Hindringer: hvorfor kom\-%
mer den da i Menneskehedens Historie ikke med lige
Styrke til enhver Slægt, hvorfor er den ikke lige mægtig
i enhver Tidsalder? Hvis den behøvede Tjenere, der
% --- p. 9 ---
\opage{9}\setcounter{footnote}{0}%
kunde gaae forud og berede den Vei; hvis den nu og
da maatte ligge stille og vente paa gunstige Betingelser,
ligesom en Skipper ligger stille og venter paa gunstig
Vind, saa kunde vi forstaae det. Men naar den selv
maa beaande sine Tjenere og selv bringe Betingelserne
med, er det da ikke ubegribeligt, at den ikke altid og
overalt vil virke med lige kraftig Nærværelse? Her maa da
vel i Virkeligheden være Noget til, der kan sætte Aanden
Skranker, Noget, der kan berede den alvorlige Hindrin\-%
ger og bevirke, at den, skjøndt almægtig, synes at være
afmægtig. Vi svare: der er virkelig et saadant Noget
% [a reader has pencil-underlined "der er virkelig et saadant Noget" and marked the margin; not transcribed]
til! Men dette Noget er ikke til udenfor Aanden; nei,
det er for Aanden selv, i dens eget Rige: det ligger i Fri\-%
heden, og Aandens Væsen er Frihed.

Fatte vi Friheden i kort Begreb som Villiens Een\-%
hed af Viisdom og Magt, da fatte vi strax, at den evig
frie Aand maa være hævet over mørk Naturnødvendighed
og blind Vilkaarlighed. Fri i sig selv vil Aanden virke
paa Frihedens Vilkaar i frie Organer; saadanne frie Or\-%
ganer ere de menneskelige Villier. Kom det nu alene
an paa at bruge Magten og herske som Despot, da
vilde der ganske vist Intet være hverken i Himlen eller
paa Jorden, der formaaede at sætte den almægtige Aand
Skranker; men ved at herske despotisk over frie Villier,
vilde Frihedens Aand modsige sig selv. Her viser sig
da det Særsyn, at det, som i Skabningens Rige maa sy\-%
nes at være det Svageste af Alt, den menneskelige Villie,
virkelig formaaer, hvad ikke de haarde Klipper, ikke de
mægtige Fjelde, ikke de høie Stjerner formaae, at berede
den guddommelige Villie Hindringer, sætte Almagten
selv en Skranke. Det maa ved flygtig, overfladisk Be\-%
tragtning forekomme os ganske urimeligt. Thi hvad er
dog den menneskelige Villie, naar vi see bort fra Drift
og Tilbøielighed? Den pure Villen, den saakaldte rene
% --- p. 10 ---
\opage{10}\setcounter{footnote}{0}%
Selvbestemmen er jo i sig selv et tomt Skyggevæsen, et
ideelt Phantom: hvad kan vel være mere skyggeagtigt
end den sig villende Villen, indholdsløsere end den sig
bestemmende Bestemmen? Og dog er det bogstavelig
sandt, at den rene Selvbestemmen er i hele Verden det
Eneste, der formaaer at sætte en Hindring for Aanden.
Eller hvis Nogen mener, at Hindringen ikke oprindelig
udgaaer fra Villien i Mennesket, han sige os da, hvorfra
den udgaaer. Vel vide vi, at vor Villie i Virkeligheden
kun kan bestemme sig til Noget, forsaavidt den bestem\-%
mes af Noget; men vi vide jo ogsaa, at Villien umulig
kunde blive fri, dersom den ikke oprindelig kunde be\-%
stemme sig selv til at lade sig bestemme. Sagen er ikke
saa indviklet som det ved første Øiekast kunde synes;
den klarer sig strax, naar vi betænke, hvad Forskjel det
gjør, om Villien bestemmes af Driften eller af Aanden.
Den naturlige Villie er i Begyndelsen Eet med den na\-%
turlige Drift; sandselig bevæget af sandselig Drift, griber
Villien efter det Sandselige. Bestemmes nu Driften nær\-%
mere som Tilbøielighed, og stiger Tilbøieligheden til Li\-%
denskab, saa bliver den sandselige Villie stærk ved den
stærke Lidenskab; men den af Lidenskaben beherskede
Villie er med al sin sandselige Styrke ufri. Hvorfor ufri?
Fordi den i sin Naturlighed ikke selv har bestemt sig
for Lidenskaben, men er uden videre besat, opflammet
af Lidenskaben. Det er en let Sag for Tilbøieligheder
og Lidenskaber at faae den naturlige Villie i deres Magt,
fordi de, om vi saa maae sige, kunne komme bag paa
den, slaae ned paa den, gribe den og bevæge den de\-%
spotisk. De stærke Lidenskaber ere den sandselige
Villies Despoter. Allerede heraf skjønnes, hvi Aanden
ikke som Driften kan magte Villien. Frihedens Aand
vil frigjøre Villien; men Villien kan kun frigjøres derved,
at den selv inderlig vil Friheden. Ligeoverfor Drift og
% --- p. 11 ---
\opage{11}\setcounter{footnote}{0}%
Tilbøielighed formaaer Villien Intet, naar den overlades til
sig selv; hvad den formaaer, formaaer den ene og alene
ved Aandens Bistand. Hvis Aandens Bistand var umid\-%
delbar som de Tilskyndelser, der udgaae fra Driften, da
vilde det ligefrem kunne komme til en af Villien uaf\-%
hængig Strid mellem den frie Aand og den natur\-%
bundne Drift; fik saa Driften Overhaand og rev Villien
med sig, da var det ikke Villiens, men Aandens Skyld:
den guddommelige Aand havde da været for svag. Har et
Menneske Lov til at see Sagen fra dette Synspunkt, kan
han let undskylde sig: „Hvad kan jeg gjøre for, at jeg
er født med sandselige Drifter og urene Begjærligheder?
jeg kan da umulig forandre min egen Natur. Vil Aan\-%
den have Begjærlighederne udrensede, maa den selv ud\-%
rense dem; vil den frigjøre min Villie, maa den selv
gribe og bevæge den. Jeg vil jo gjerne lade mig fri\-%
gjøre; min Villie gjør aldeles ingen Modstand; Aanden
kan for den Sags Skyld tage den, vende den, dreie den
og gjøre ved den, hvad den vil, naar den blot vil gjøre
Noget. Men naar den guddommelige Aand selv enten
ikke kan eller ikke vil umiddelbart gjøre Noget, hvad
skal jeg saa kunne gjøre?“ Her er en uhyre Misfor\-%
staaelse. Dersom Aanden virkelig skulde bevæge Villien
ganske umiddelbart, i Lighed med Driften, saa maatte
den jo begynde Villiens Frigjørelse med at gjøre den
ufri; hvilken Modsigelse! Villiens første Frihedsact,
Frugten af dens allerførste Beaandelse, er netop den for\-%
høiede Selvbestemmelse, hvormed den inderlig attraaer
og, i Strid mod Driften, giver sig hen til Aandens Ind\-%
virkninger. Idet saa disse Indvirkninger modtages og
tilegnes af den ved Tilegnelsen sig frigjørende Villie, bliver
Forholdet ingenlunde dette, at Villien i Selvadskillelse
% [p. 11: a small speck above the second e of "Selvadskillelse" (layer reads "Selvadskillélse"); at 600 dpi not a positive accent, read as e]
fra Aanden skulde have nogen Kraft til nogensomhelst
selvstændig Virken; tvertimod! den kan kun bevæge sig
% --- p. 12 ---
\opage{12}\setcounter{footnote}{0}%
efter Aanden, idet den med Frihed bevæges af Aanden.
Frihedens bevægende Kraft er ene og alene Aandens
Kraft; men det beroer paa Villien selv, om den i Inder\-%
lighed, i sand Oprindelighed, vil tilegne sig denne Kraft.
Hvis den ikke vil, da „bedrøver den Aanden“; da er
dens Afmagt en negativ Magt, der, sørgeligt nok, sætter
Almagten selv en Skranke.

Og hvad er saa det praktiske Resultat af de her an\-%
stillede Betragtninger? Hvad kræver Aanden af det
med Virkelighedens Betingelser arbeidende Menneske;
hvad er det vel for en Grundfordring, den særlig stiller
til denne nærværende Tid? Er vor Tidsalder verdslig
fremfor alle foregaaende Tider? Er det Verdensforsa\-%
gelse i store udadgaaende Demonstrationer, middelalder\-%
lige Tilstande med Munkevæsen, Beden, Fasten, kirkelige
Processioner og klosterlig gode Gjerninger, Aanden for\-%
drer for at forvandle de naturlige Hindringer til gun\-%
stige Betingelser? Ingenlunde! Aandsfordringen er ikke og
kan ifølge sit Princip ikke være nogen endelig Fordring paa
nogetsomhelst endeligt Udvortes. Efter sit Udvortes kan
den nærværende Tid vel siges at være, om ikke bedre, saa
dog idetmindste ligesaa god som nogen forudgaaende.
Nei, den Grundfordring, Aanden gjør til denne som til
alle andre Tider, er først og sidst en Inderlighedens For\-%
dring; den fordrer Tro, Tro paa dens evig frie, frigjø\-%
rende Magt; den fordrer Villie, Villie til at gjøre sig fri
ved at lade sig frigjøre, kort og bestemt: den fordrer
Alvor.

Enhver endelig bestemt Fordring kan Nutiden med
sine urolige Reflexioner og sin næsten feberagtig over\-%
spændte Kritik let omskrive, halvere og, især naar For\-%
dringen opstilles enkelt, let protestere. Men een For\-%
dring er der, den ikke kan omskrive, ikke halvere, ikke
protestere, og det er Fordringen Alvor!

% --- p. 13 ---
\opage{13}\setcounter{footnote}{0}%
Denne Fordring er streng nok, men den er ikke for
streng. Der kan maaskee gives en Slægt, der roser sig
af at have opgivet det evige Liv for desto kraftigere at
leve det timelige; men der gives dog vist ikke en
Slægt, der for Alvor vilde rose sig af at have opgivet
al Alvor. Der kan maaskee fødes en Generation, der
sætter en Ære i at mangle Sands for Gud og hans hel\-%
lige Ord; men skulde der vel kunne fødes en Genera\-%
tion, der vilde sætte sin Ære i at mangle Alvor? Naar
da Aanden i dette Nærværende kun fordrer væsentlig og
sand Alvor, saa fordrer den jo i Grunden ikke Andet af
Tidsalderen, end hvad Tidsalderen nødvendigt maa fordre af
sig selv: her maa da være en Forstaaelse mulig. At gjøre
Alvor af Noget er jo, som man siger, at tage sig Noget
nær, at lægge sig det paa Sinde; ved Alvoren sættes der
Eenhed af Sag og Sind. Men heraf følger igjen, at
Sindet, der optager Sagen i sig og gjør sig til Eet med
den, selv igjen maa blive af samme Væsen som Sagen.
Er Sagen uvæsentlig, i sig selv ubetydelig, da er Sindet
usselt; kun et usselt Sind kan gjøre Væsen af det Ube\-%
tydelige. Er Sagen ikke ligefrem uvæsentlig, men dog
af endelig Natur, saa giver Sindet sig til Priis for det
Endelige, endeliggjør sit Væsen ved at tage sig Sagen
nær og give sig ganske hen i den. Det Endelige er
ifølge sin Natur foranderligt; men Alvoren fordrer jo, at
Sindet i Forhold til Sagen skal være sig selv ligt, skal
midt i det Foranderlige vide sig hævet over Forandrin\-%
gerne. Interessernes og Midlernes Verden er ubestandig.
Den alvorlige Mand maa være beredt paa Omskiftelser;
hans Frygt og Haab bestemmes ved, at han altid kan
vente sig det Modsatte, efter Modgang Medgang, efter
Vinding Tab og saa fremdeles. Her passe de kloge
Ord: Sperat infestis, metuit secundis alteram sortem
% --- p. 14 ---
\opage{14}\setcounter{footnote}{0}%
bene præparatum pectus. See, det er nu Klogskabens
Alvor.

Men Klogskabens Alvor kan umulig være sand Al\-%
vor; den lider af en medfødt Skrøbelighed, den slæber
paa en tung Modsigelse, den Modsigelse, at Sindet paa
een Gang skal baade give sig hen og dog ikke give sig
hen i det endelige Væsen. Ganske vist har det Ende\-%
lige sin Interesse som Betingelse, som Middel. Men
bliver det i den Grad Alvor med Midlerne, at Maalet
udelukkes, saa er Selvbedraget aabenbart. Betingelsen
gribes; men den har kun Betydning for Noget, som der\-%
ved skulde betinges. Hvad er nu dette for et Noget?
I Endeligheden selv en ny Betingelse! Og den nye Be\-%
tingelse? Har kun Betydning for et Noget, som derved
skulde betinges! Betingelsens egentlige Værdi falder i
det Uendelige udenfor Betingelsen selv, ligesom Midler\-%
nes Værdi falder udenfor Midlerne. Den, der nu paa
disse Vilkaar vil leve af sine Midler og mætte sin Sjæl
med lutter Betingelser, kan rigtignok siges at være en
alvorlig Mand, alvorlig, ja som Tantalus, naar han vil
plukke Frugterne og drikke af Strømmen. Den Alvor,
der gjør en Tantalus gammel og graa, kan jo nok gjøre
en heel Slægt gammel før Tiden. Er det virkelig sandt,
at Menneskeheden ældes, da kan det ikke være Naturens
Skyld, thi om Naturen gjælder, hvad Digteren siger:
„hver Old og hvert Aar har sin egen, sin blomstrende Vaar“.
Endnu mindre kan det være Aandens Skyld, thi Aanden
som Aand har Evighedskilder, og det Evige kan ikke
ældes. Skylden maa da falde paa den falske Alvor, In\-%
teressernes, Midlernes, Betingelsernes endelige Alvor.

Har falsk Alvor gjort en Slægt gammel, da er det at
haabe, at sand Alvor kan gjøre den ung igjen. Sand Alvor
er en Aandens Gave, men paa Frihedens Vilkaar. Be\-%
tingelsen er, at den beaandede Villie vil modtage, bevare
% --- p. 15 ---
\opage{15}\setcounter{footnote}{0}%
og bruge Gaven. Den sande Alvor lægger afgjørende
Vægt paa Øiemedet, men saaledes, at den med det Samme
lægger Vægt paa Midlerne. At sværme for tomme,
Virkeligheden overflyvende Idealer er ikke Alvor. Idea\-%
lernes Alvor er, at de ville virkeliggjøres. Til Virkelig\-%
gjørelsen udkræves Anstrengelse, Udholdenhed, Bestan\-%
dighed. Men Udholdenhed og Bestandighed under store
Anstrengelser svare jo netop til væsentlig Alvor. Alvor
er en Sjælens Kraft, der holder ud i Prøvelser, en Sin\-%
dets Grundstemning, der er sikkret ved fast Beslutning;
Alvor er Selvets Fornyelse, Charakterens Indvielse, Vil\-%
liens Daab, at den døbes med Aand. Derfor er sand
Alvor ikke blot Noget, hvormed det aandelige Liv skal
fuldendes og sluttes, men ogsaa Noget, hvormed det skal
begyndes. Alvor er ikke for de Gamle alene, men ogsaa
for de Unge. En begavet Ungdom har væsentlig Alvor,
fordi den har væsentlig Aand. Vel er der noget Kjøligt,
Skjult og Strengt i Alvorens Dybder, ligesom der er
noget Kjøligt i Havets Dybder; men ligesaa vel som det
dybe Hav kan bestraales af Solen, saa kan og Alvoren
selv omstraales af Glæden. Dersom Sands for Livets
Alvor ikke oprindelig hørte med til Ungdommens rige
Anlæg, hvordan vilde det da see ud med Fremtiden?
Hvad gavnede saa al vor Stræben med Opdragelse, Un\-%
derviisning og Dannelse, naar det, der i sin Tid skulde
anvendes og bruges i Alvor, ikke blev opfattet og for\-%
staaet i Alvor? Hvad vilde, for at holde os til hvad
der nu ligger saa nær, en Bestræbelse for at bringe
rigere Afvexling ind i de akademiske Foredrag og ind\-%
lede nøiere Samqvem mellem de forskjellige Universiteter
vel gavne, dersom ikke den studerende Ungdom selv var
levende interesseret i Alt, hvad der fremmer og forhøier
det aandelige Liv.

Det er da med levende Glæde og et godt Haab, vi tør
% --- p. 16 ---
\opage{16}\setcounter{footnote}{0}%
prøve paa at iværksætte vort Forehavende, fordi vi er\-%
kjende, at de Yngre have ligesaa megen, lad os kun sige
det, nok saa megen Andeel i dets Iværksættelse som de
Ældre. Det er de Unge, der have villet, at nu skulde
det være Alvor, og hvad der saaledes villes i Alvor, vil,
med Styrelsens Bistand, ogsaa blive udført med Alvor.
En Ungdom, der lytter til, naar ideelle Strenge røres, og
selv virker med til, at aandelige Formaal fremmes, en
akademisk Ungdom med klare Øine, vaagen Sands og
foraarsfriske Tanker vil snart forstaae hvad det gjælder
om, naar det gjælder om Valg af Midler og Maal, om
Hindringers Forvandling ved Aandens Magt, om Fri\-%
hedens Seir i Livets Strid: den vil forstaae det i Alvor.

\begin{center}\rule{3cm}{0.4pt}\end{center}

\chapter*{II.}
\addcontentsline{toc}{chapter}{Forelæsning II}
\markboth{Forelæsning II}{Forelæsning II}
\label{ch:ii}
% pp. 17--35
% --- p. 17 ---
\opage{17}\setcounter{footnote}{0}%
% [large initial D]
Der fortælles om Alexander, Philip af Macedoniens
Søn, at han som Yngling græd, naar hans Fader havde
udført en stor Bedrift. Det var Beundringens Taarer,
Ynglingen græd, thi han saae med sønlig Ærefrygt op
til Faderen; han saae i ham den berømmelige Helt.
Men det var ogsaa Misundelsens Taarer; „thi“, sagde
han, „min Fader udfører alt det Store, hvad bliver der
tilovers for mig?“

Dette saa charakteristiske Forhold imellem Ynglin\-%
gen Alexander og hans Fader kan i en vis Forstand ---
fra een Side idetmindste --- tjene til at belyse Historiens
paa hinanden følgende Aldere. \textit{Le présent est gros de
l'avenir}, siger Philosophen Leibnitz, og med Rette; den
nærmest følgende Tid er, om jeg saa maa udtrykke mig,
altid en Søn af den foregaaende. Men har den nærvæ\-%
rende Tidsalder den foregaaende til Fader, altsaa, for at
blive i Billedet, den næstforegaaende til Bedstefader, saa
bliver det, naar vi skue tilbage paa Slægtleddenes
lange Række, paa Oldefader, Tipoldefader, Tip-Tip\ldots{}
o. s. v., en mægtig Stamtavle. Den nulevende Slægt er
en adelig Slægt med mange Ahner og store Fædre; naar
da Nutiden vender sig og betragter Fortiden, dens Be\-%
drifter og Helte, maa den uvilkaarlig udbryde: „Mine
% --- p. 18 ---
\opage{18}\setcounter{footnote}{0}%
Fædre have udført det Store!“ og kunde da let fristes til
at spørge med Alexander: „Hvad bliver der tilovers
for mig?“

Det er, alvorlig seet, ikke ungdommelig Ærgjær\-%
righed, ikke Misundelse, endnu mindre tom Forfæn\-%
gelighed, der aander i et saadant Spørgsmaal; det er
tvertimod stille, dyb Bekymring; det er aandelig
Næringssorg. Det hjælper ikke, at Nogen vil sige til
os: „Hvor ville I hen med Eders Spørgen om det Store?
Have I som Alexander de store Kræfter, velan, saa
ville I ogsaa som Alexander selv vide at udføre store
Bedrifter; hvis ikke, da gjøre I bedst i at blive ved
Jorden, lære at krybe, indtil I kunne gaae, og tage til
Takke med det Mindre.“ Denne Afviisning er ubeføiet.
Vi ville for vort eget Vedkommende gjerne nøies med
det Mindre; men for at have det Mindre med Aand og
Sandhed, maae vi nødvendig hjælpes ved det Store.
Forsvinder det Store saaledes fra en given Samtid, at
det ikke mere har Nærværelsens Magt, saa bliver det
en aandelig Qval at leve i det Mindre. Vel maa det
Forbigangne afgive Forudsætninger for det Nærværende;
men det Nærværende sygner hen, naar det udelukkende skal
tære paa en Forbigangenhed; hver Tidsalder maa have
Næringskilder i sig selv: det Nærværende lever af et
Nærværende og det Mindre af det Større. Det gaaer med
Bevægelserne i Verdenshistorien som i Verdenshavet. Der\-%
som ikke Vind og Storm og Ebbe og Flod og stærke
Strømme uophørlig bevægede Vandene i det store Hav,
hvordan vilde det da see ud med Bugter og Fjorde?

Det er et sandt Ord, at Ingen kan leve, medmindre
han har Noget at leve af og Noget at leve for. Sandse\-%
ligt kunne de to Fordringer adskilles, men ikke aande\-%
ligt. I materiel Henseende er det første Nødvendige
unegtelig dette, at den, der skal leve, maa have Noget
% --- p. 19 ---
\opage{19}\setcounter{footnote}{0}%
at leve af. Ja denne Fordrings Paatrængenhed er endog
saa iøinefaldende, at Mange tage den for den eneste
gyldige og tænke ved sig selv: havde vi kun Nok at
leve af, vi skulde sagtens finde paa Noget at leve for.
Men i aandelig Forstand ere de to Fordringer uad\-%
skillelige: hvad et Menneske aandelig lever for, det lever
han ogsaa af. En Kunstner f. Ex., der udelukkende
lever for sin Kunst, lever jo ogsaa, vistnok ikke altid
materielt men dog ideelt, af sin kunstneriske Frembrin\-%
gelse, har i denne Frembringelse sin Aandsføde, sit Liv,
sin Glæde. Den, der ubetinget lever for Troen, lever
aandeligt af Troen; den, der lever for Viden, af Viden
o. s. v. Det gjælder i det Ideelle ligesaa vel som i det
Materielle om Driftscapitaler.

I Handelsverdenen, overhovedet i Arbeidernes, de
borgerlige Interessers praktiske Verden, vaages der over,
at de store Capitaler ikke spredes og fordeles i lutter
Skillemynt. Det er ligesaa meget i de Ubemidledes som i de
Bemidledes Interesse, at Capitalerne bevares. Saaledes
ogsaa i Historien. Der maa Aandscapitaler til, at enhver
Tidsalder kan have Noget at leve for og at leve af;
der maa i kritiske Vendepunkter gjøres store Indsatser.
De store Indsatser betegne Historiens Epocher; Capitalens
Anvendelse og Fordeling i det Større og Mindre, indtil
den har givet sit Udbytte og er brugt op, er betegnende
for en Periode. Jo større Epochen er, desto større er
den indsatte Livscapital og desto længere vil ogsaa den
tilsvarende Periode være. Den absolute Aandsepoche
er i Religionens Sprog det Vendepunkt, der kaldes Ti\-%
dens Fylde, det Vendepunkt, da ikke blot en af de
store Evighedstanker, men den Evige selv, han, der
siger: „før Abraham blev, er jeg“, aabenbaret i Kjød og
Blod, traadte ind i Tiden. Seet fra dette Høidepunkt er
Livets Ord den aandelige Grundcapital, hvoraf Slægterne
% --- p. 20 ---
\opage{20}\setcounter{footnote}{0}%
skulle leve; for denne Epoche ere alle følgende Tider at
betragte som en eneste, fortløbende Periode.

Anderledes forholder det sig derimod med de rela\-%
tive Epocher; til disse høre naturligviis kun relative, i
Forhold til den større eller mindre capitale Indsats læn\-%
gere eller kortere Perioder. Saadanne relative Epocher
have vi i de ideelle Gjennembrud, paa hvilke Adskillel\-%
sen mellem Oldtiden, Middelalderen og den nyere Tid
oprindelig beroer; de kunne, hvad Aandslivet angaaer, i
det Hele siges at repræsentere: Folke-Aanden, Kirke-%
Aanden, Verdens-Aanden. Det er i denne Sammenhæng
nødvendigt at besinde sig paa Adskillelsen mellem de
relative Former, hvorunder Aanden virker i Menneske\-%
livet, og den guddommelige Aand selv. I Culturhistorien
indvirker Aanden ikke umiddelbart paa Menneskets Tan\-%
ker og Kræfter; her er et indre Mellemværende mellem
Aanden absolut og den menneskelige Selvbevidsthed, et
Mellemvæsen, hvori Aanden hver Gang indlægger et
særeget, ved en særegen Form bestemt Indhold; dette
Mellemvæsen kalde vi Ideen. Kun idet Menneskets Be\-%
vidsthed i Sind og Tanker, i Villen og Virken tilegner
sig Ideen, kan Menneske-Aanden faae sit eiendommelige
Præg.

Ved at udhæve Forskjellen mellem Idee og Aand
kunne vi strax charakterisere Modsætningen mellem
Aabenbaring og Mythologie. I Aabenbaringen er Ideen
saa fuldkommen gjennemsigtig, at det i mange Tilfælde
seer ud, som om den ganske var forsvunden. Gud taler
til Mennesket, underviser Mennesket, stiller sig imod
Mennesket som Villie mod Villie, Aand imod Aand.
Paa Mythologiens Trin forholder Sagen sig omvendt;
her er Ideen som Idee bleven uigjennemsigtig for Men\-%
neske-Aanden; Beviset herfor er, at Ideerne selv i bil\-%
ledlig Form opstilles som Guder: de mythiske Guder
% --- p. 21 ---
\opage{21}\setcounter{footnote}{0}%
ere personificerede Ideer. Modsætningen mellem Jøder
og Grækere er i denne Henseende charakteristisk; Jøde\-%
folket er Aandens, Grækerne Ideernes Folk. Hos Jø\-%
derne gaaer Alt ud paa at fyldestgjøre Aanden; derfor
kunne Kunstens og Videnskabens humane Ideer ikke
komme til Udvikling; end ikke Staten fik som Stat Lov
til at udvikle sig af sin egen Idee: den theokratiske Stat
er Aandens Magt umiddelbart underlagt. Anderledes
hos Grækerne. Her er det saa langt fra, at Aanden
overvælder Ideen, at Livets Indhold i alle Retninger
tvertimod udvikler sig af Ideen, formes og individualiseres
til et eiendommeligt Idee-Indhold. Heri ligge de stærke
Impulser til Udvikling af Kunst og Videnskab, Politik
og rationel Ethik; men heraf forklares da ogsaa den
Grundmangel, at det Overnaturlige overalt smelter sam\-%
men med det Naturlige, Villien med Driften, Friheden
med Nødvendigheden paa en saadan Maade, at ikke en
eneste Aandens Fordring kommer til at gjælde i absolut
Ubetingethed.

Som Oldtidens mest begavede Culturfolk afgive
Grækerne et træffende Exempel paa, hvad en Folke-Aand
efter sin Oprindelighed betyder, hvilke indre Grændser
den har, hvorledes den kæmper med Hindringer og Be\-%
tingelser. Hos ethvert begavet Folk rører sig fra Begyn\-%
delsen af en Kreds af særegne naturkraftige Ideer. Hvor
betegnende slige Ideer ere for den folkelige Eiendomme\-%
lighed, falder i Øinene, naar vi sammenligne den græske
Mythologie med den nordiske. Zeus og Odin, Frigga
og Here, Helheim og Hades, Olympen og Valhalla,
hvilken Forskjel! Forskjellen er dobbelt; deels udad\-%
gaaende, forsaavidt Gudernes Skikkelse, Individualitet,
Bedrifter og Skjæbne udmales anderledes af den græske
Phantasie end af den nordiske; deels indadgaaende, for\-%
saavidt enhver af de forskjellige Gudeverdener forudsæt\-%
% --- p. 22 ---
\opage{22}\setcounter{footnote}{0}%
ter en egen Grad af Selvbevidsthed, en egen Verdens- og
Livsanskuelse hos det paagjældende Folk. De græske
Guder ere Skjønhedsguder, de nordiske Kampguder.
Dette maa naturligviis ikke forstaaes, som om de græske
Guder ikke kunde kæmpe, og de nordiske ikke være
skjønne; men Sagen er, at de græske have deres fuld\-%
komneste Aabenbarelser i Skjønheden, de nordiske i
Kampen. Hvor indgribende denne Modsætning er ved
Bestemmelsen af de tvende Folks eiendommelige Aands\-%
retning og Rolle i Verdenshistorien, kan oplyses ved
et Par Træk.

At en Guddoms Aabenbarelse fuldendes i Skjønhe\-%
den, vil med andre Ord sige, at dens Eiendommelighed,
dens ideale Selvstændighed væsentlig er i æsthetisk Indi\-%
vidualitet. Men naar Individualiteten i sit umiddelbare
Phænomen er høieste Udtryk for Gudernes Selvstændig\-%
hed, da maa den ogsaa være det for Menneskenes; thi
det Guddommelige er Forbillede for det Menneskelige.
Homer vil, siger han (Iliad. I, 1---5), lade Gudinden be\-%
synge Achills Vrede, der sendte de mange Heltes Sjæle
til Hades, men gjorde dem selv til Bytte for Hundene
% Greek type in the print; set from the image at 300 ppi.
(αὐτοὺς δὲ ἑλώρια τεῦχε κύνεσσιν). Vistnok maae vi
ved dette „dem selv“ uvilkaarlig tænke paa Ligene; men
det Charakteristiske i Udtrykket er dog dette, at det,
der betegner Heltene selv, ikke er deres Sjæle, men deres
Legemer. Den legemlige Form er nemlig Heltens æsthe\-%
tiske Individualitet; at have Individualitet, er fælles for
Guder og Helte. Men nu kommer Forskjellen: Guderne
ere udødelige, Menneskene dødelige; ved Adskillelsen
mellem udødelige og dødelige Slægter har Skjæbnen sat en
tragisk Skranke mellem Guder og Helte. Forladte af
den legemlige Individualitet friste de blodløse Skygger
i Hades en ynkelig Tilværelse. Det hedder om Hektors
Død (Iliad. XXII, 362): „Ud af Lemmerne fløi hans
% --- p. 23 ---
\opage{23}\setcounter{footnote}{0}%
Sjæl og iilte til Hades, jamrende over den Lod at skilles
fra Kjækhed og Ungdom.“ Anderledes med Heltene i
det gamle Norden; der ventede dem som Einheriar i
Valhalla en anden Lod, og hvorfor? Fordi Odin var
saa forskjellig fra Zeus, Valhallas Guder i Sind og
Tanker forskjellige fra Olympens. Odin og Thor satte
ikke det personlige Væsen i skjøn Individualitet men i
den Villieskraft, der lever af Kamp og voxer i Kampen;
det er en saadan Villieskraft, som Heltene have tilfælles
med Guderne, der adler Bedriften og overvinder Døden.
At Døden seirer, at Lyset slukkes, og Sjælene skjælve i
Skyggeriget, er ikke Skjæbnens Skyld, Skylden er i
Feigheden og den svage Villie. Derfor er Thor „de
rædde Helheims Mænd saa gram“, og tiltaler dem haardt:
„I usselige Daarer, som frygted Saar og Død! Nu Hel
Jer evig saarer med Qval og bitter Nød.“

Af Forskjellen mellem den æsthetiske Individualitet
og den kæmpende Charakteer i de mythiske Forbilleder
flyder endnu en anden Conseqvens, der ret er beteg\-%
nende for den Rolle, ethvert af de nævnte Folkeslag
skulde spille i Historien. Olympens „salige Guder“ have
jo vistnok, det ligger i Individualitetens Natur, ikke kunnet
trænge igjennem med den klare fuldendte Form og
faae Bugt med de gamle, halvt formløse, af Chaos op\-%
rundne Elementarguder, uden efter store Omvæltninger
og haarde Kampe; men medens Aserne, Valhallas villie\-%
stærke Kampguder, maae føre en stadig Krig med Jet\-%
terne, en Krig, der begynder med Ymers Drab og ender
med Ragnarokur, blive de olympiske Individualiteter i for\-%
holdsviis kort Tid færdige med den hele Gudekrig og faae
saaledes Leilighed til i plastisk Skjønhed og classisk Ro et
% printer's defect: "et" printed for "at" (infinitive marker); transcribed as printed.
nyde al Livets Lyst med Idealitetens salige Glæde: „Hjertet
ei savnede her et rundeligt Maaltid, eller den yndige Lyra,
som sloges af Phoibos Apollon, eller de vexlende Qvad
% --- p. 24 ---
\opage{24}\setcounter{footnote}{0}%
af Musernes deilige Stemme.“ Medens Nordboerne med
deres guddommelige Førere ere saaledes indhvirvlede i
den uendelige Kamp, at de først kunde begynde at
komme til Ro i den fremrykkede Middelalder og paa
Christendommens Grundlag udvikle nogen Civilisation,
naae Grækerne gjennem forholdsviis kortvarige Kriser og
storartede Katastropher det Hviletidens Vendepunkt, da
Ideerne klares, og Stoffet formes, saa at ikke blot Statens
men ogsaa Tankens og Skjønhedens Love opfyldes. Det
er nu let at see, at Folke-Aanden bestemmer Folkets
virkelige Liv, at dets historiske Stigen, Culmination og
Dalen paa indvortes Maade betinges af de i Aanden selv
arbeidende Ideer. Den græske Historie er i særegen
Forstand at kalde katastrophisk. Det Katastrophiske egner
sig just for en alsidig Udvikling af den æsthetiske Indi\-%
vidualitet og giver os iøinefaldende Exempler paa, hvor\-%
ledes en Folke-Aand først beviser sin Kraft ved at for\-%
vandle Hindringer til Betingelser, og saa igjen forraader
sin Svaghed, naar Betingelserne vende sig om og for\-%
vandles til Hindringer. Den græske Individualitet er i sin
heroiske Umiddelbarhed stridbar; allerede i Kampen for
Troja er Striden mellem Achill og Agamemnon af stor
katastrophisk Virkning. Fik nu de stridbare Individu\-%
aliteter Lov til at tumle sig frit imellem hverandre, da
vilde Staten, dette almene Væsen, umulig kunne dannes
og komme til Raadighed; men ligesom Platos Gud ved
et Blik paa Chaos oprindelig saae, at Orden er bedre end
Uorden, saaledes saae det græske Blik, at Uorden i lige
Grad stred saavel mod det Skjønne som mod det Gode.
Men idet Individerne bøiede sig under Staten som under
en Guddom, bøiede de sig dog ikke saa slavisk, at In\-%
dividualitetens Præg udslettedes; tvertimod! Individuali\-%
teterne gik just op i Staten saaledes, at Statens Væsen
derved selv blev individualiseret. Istedenfor een eneste
% --- p. 25 ---
\opage{25}\setcounter{footnote}{0}%
almindelig græsk Stat danner der sig fra Begyndelsen
af en Fleerhed af individuelle, ifølge deres umiddelbare,
altsaa ligeledes æsthetiske, Individualitet indbyrdes stridige
Stater.

Af de rige Individualiteters Grund fremsprudle Ideer\-%
nes Kilder; her er da allevegne i det Udvortes som i
det Indvortes, i Colonierne som i Moderstaterne Impul\-%
ser, Kræfter, Elementer nok i Bevægelse. Skulde der
imidlertid komme noget Heelt og Stort ud af denne
brogede Mangfoldighed, Noget, hvorpaa Skjønhedens
Aand kunde sætte sit Præg, altsaa Noget, som i Sand\-%
hed var af een Støbning: da maatte der fremfor Alt til\-%
veiebringes et uhyre Tryk; der maatte en Betingelse til,
der havde Noget at sige. Og Betingelsen indtraf; men
i det Øieblik den indtraf, tog den sig rigtignok ingen\-%
lunde ud som en Betingelse, nei, den tog sig tvertimod
ud som en knusende Hindring. Det er Xerxes, der med
sine Hundredetusinder fik Opgaven, den skjønne histo\-%
riske Opgave at skaffe Betingelsen for den græske Folke-%
Aands rige, vidunderlig stærke Udvikling tilveie. Og
Xerxes maa vel ruste sig og haste; thi det haster med
Betingelsen. I Bugten ved Salamis har Kiv og Tvedragt
allerede beredt Themistokles saa mange Hindringer, at
han kun veed at gjøre Eet: forvandle disse Hindringer
selv til een eneste stor, afgjørende Betingelse. Det er jo
naturligt, at Individualitet og Frihed maae svare til hin\-%
anden. Den æsthetisk stærke Individualitet opflammes
af Frihedsbegeistring, naar Faren er udvortes overhæn\-%
gende; i Begeistringens Øieblik har den en uhyre Spænd\-%
kraft; men det gjælder ogsaa om at benytte Øieblikket;
og de græske Helte benyttede Øieblikket: Marathon,
Thermopylæ, Salamis og Platæa tale herom til alle Slæg\-%
ter. Den græske Frihedskamp var en uhyre Katastrophe;
det Slag, der skulde have knust Individualiteten, bevir\-%
% --- p. 26 ---
\opage{26}\setcounter{footnote}{0}%
kede en Reisning, hvorved den kom til at udfolde Kræf\-%
ter, der har forbauset Verden og overrasket den selv.
Det er i denne Kamp, at Capitalen indsættes, den Livs\-%
% line-end break "Livs-/capital" (lower-case c; p. 19 prints "Livscapital"): read as one word.
capital, der skulde give store aandelige, uberegnelige
Renter.

Det græske Statsliv, den græske Poesi og fremfor
Alt den græske Kunst skulde sikkert aldrig have naaet
det Høidepunkt, hvori en reen, streng og fuldendt Form
paa saa oprindelig enfoldig Maade oplives og fyldes af
Ideens Indhold, at den sildigste Tid vil see derpaa med
Forundring, hvis ikke den katastrophiske Frihedskamp
med sine Rædsler og Rystelser havde bragt Hindringer
med, som Aanden kunde forvandle til Betingelser. Det
var, siges der, et græsk Ord, at den ikke kunde blive
salig, som ikke havde seet Zeus i Olympia. Hvad var
der da at see i Olympia? Den forklarede Individualitet
i mandig Ungdomskraft, i Magt og Majestæt, i Ro og
Høihed, den „Gudernes og Menneskenes Fader“, hvis
Lynstraale havde ramt Perserne, som den fordum ramte
Titanerne. Vi læse hos Homer om det hellige Tegn,
hvormed Zeus stadfæster sit Løfte til Thetis (Il. I, 528):
„Talt, da nikkede Zeus med de sortblaa Bryn i det
samme, og fra udødelig Isse faldt Drottens ambrosiske
Lokker bølgende ned, og høien Olymp han kom til at
skjælve.“ Ved at læse dette faae vi en Anelse om det
Gudesyn, der i Begeistringens Øieblik maa have fore\-%
svævet Phidias; men for os --- der da ogsaa haabe at
blive salige uden et saadant Syn --- er det kun som en
Anelse; for den græske Kunstner var det derimod vir\-%
kelig og levende Nærværelse. Det var den græske Folke-%
Aand, der aabenbarede sig i Synet, og det med en
Klarhed og Umiddelbarhed, som vidnede om, at den nu
virkelig var kommen til fuld Bevidsthed om sin guddom\-%
melige Kraft. Hvis vi ikke turde forudsætte dette, hvor
% --- p. 27 ---
\opage{27}\setcounter{footnote}{0}%
skulde vi da søge Betingelsen for hvad Historien kalder
„den store Kunst“.

Paa den katastrophiske Udvikling med den hurtige
Opblomstring fulgte saa under en ligeledes katastrophisk
Splidagtighed (den peloponesiske Krig), en tilsvarende
hurtig Dalen, Afblomstring og Hensygnen. Da Demosthenes
tordnede mod Philip, var Lynstraalen i Zeuses Haand
allerede mat; Aandskræfterne vare udtømte, thi medens
Stater og Stammer svækkede hinanden i indbyrdes Stridigheder,
udtømte den til de classiske Former bundne
Idee sin oprindelige Productivitet. Det var en tragisk
Conflict imellem det Ethiske og det Politiske, hvori den
græske Frihed gik under: ethisk havde vel Demosthenes
Ret imod Philip; men at Philip, politisk seet, havde Ret
imod Demosthenes, beviste hans Søn Alexander.

I Middelalderen bliver Folke-Aanden ikke saameget
gjennemtrængt og væsentlig gjenfødt, som meget mere
umiddelbart overvældet og systematisk behersket af Kirke-%
Aanden. Hvad er nu Kirke-Aanden? Ingenlunde Eet
med den oprindelige Menighedsaand, men en afledet og
dog mægtig Væsenhed, en asketisk, ved Dogmer og Traditioner
disciplineret Hersker-Aand, en græsk-romersk
Aand, udgangen, ikke fra Faderen og Sønnen, men fra
Bisper og Concilier, Eremiter og Munke, Cardinaler og
Paver. I Kirke-Aanden er den umiddelbare Folke-Aand
nedsat til Moment, det Naturlige kuet og forvandlet til
Grundlag. Den principielle Modsætning imellem Kjødet
og Aanden, det Profane og det Hellige, det Jordiske og
det Himmelske, det Nærværende og det Tilkommende,
griber Følelsen, antænder Phantasien og slaaer i alle Retninger,
paa alle Punkter og under alle Former ud i
mægtig, anskuelig Dualisme. Imellem de strengeste, mest
ubetingede Aandsfordringer paa den ene og Verdsligheden
med de sandselige Drifter, de raae Lyster og stærke Li\-%
% --- p. 28 ---
\opage{28}\setcounter{footnote}{0}%
denskaber paa den anden Side dukke Ideerne op paa Ny,
og en uhyre Gjæring indtræder. Ved Ideernes Gjæring
bliver det Guddommelige verdsligt, det Verdslige guddommeligt,
og dog vedblive de modsatte Poler trods
gjensidig Tiltrækning at frastøde hinanden. Kirken bliver
en Stat, og Staten bliver kirkelig, og dog vedblive de,
Kirken som Gudsriget og Staten som Verdensriget,
under uophørlig Strid og Forsoning at modvirke hinanden.
Den for Middelalderens Dualisme saa charakteristiske
Brydning mellem Keiser og Pave begynder fra det
Øieblik af, da Leo den Tredie har salvet og kronet Carl
den Store. Det varer ikke længe, før Brydningen bliver
saa stærk, at Samfundets Piller rystes. Medens det
carolingiske Statssystem splittedes, og den mangesidige
Vexlen af Privatherredømme og personlig Afhængighed,
der charakteriserer Hierarchie og Lehnsvæsen, forvirrede
Retsbegreberne og gav Individerne til Priis for Vold og
Vilkaarlighed, stormede Normannerne fra Nord, Magyarerne
fra Øst, Saracenerne fra Syd, og fuldendte Tidens
Nød og Elendighed. Folkene, hvis Øie den nye Tro
havde opladt for Fredens og Retfærdighedens Velsignelser,
og som just havde gjort det første Indblik i Civilisationen,
betragtede Tilstanden i Religionens Lys, bleve grebne af
Rædsel for Verdens Forfængelighed og ventede Dommedag.

Hvormed skal nu alt dette ende? Ikke med Verdens
Undergang men med en Civilisationens Seier, en
Culturens Fremgang: de stærke Gjæringer ere de Fødselsveer,
der bringe en ny Tid med en ny Skikkelse for
Lyset. Gjæringen er en uophørlig Omsætning af Hindringer
i Betingelser og Betingelser i Hindringer. Gjæringen
selv kan ikke frembringe noget Virkeligt; men
den kan, om jeg saa maa sige, ælte og behørigt tilberede
Mulighederne. Det, der ventes paa, er et stort Gjennem\-%
% --- p. 29 ---
\opage{29}\setcounter{footnote}{0}%
brud, en Moderidee, som ved Kirkens Aand undfanger
og føder de mange Døttre. Med denne Idee bliver saa
Indsatsen gjort; dens Medgift er den Livscapital, hvorpaa
Middelalderen længe har samlet, og hvoraf den længe
skal leve. Vendepunktet er, at der viser sig et Aandens
Maal, der i sig kan optage og omfatte alle andre Formaal,
at der udfindes en Gjerning, som kan omslutte og
indbefatte alle andre Gjerninger, at der vækkes en Begeistring
saa rig og stærk, at dens Virkninger kunne
udstrække sig til Alt, hvad der drager Næring af det
Ideelle, kort og godt: det er den hellige Grav, der skal
erobres.

Korstogenes Idee kan med Føie kaldes Middelalderens
Grundidee; det Maal, som herved blev sat, var paa
een Gang saa jordisk, at det kunde gribes med Hænder,
og dog saa overjordisk, at det umiddelbart stod som en
Port til Himlen. Hvormange gode Gjerninger, der end
kunde opregnes, saa var der dog nu blandt de Christne
i Vesten Ingen, der tog feil, naar Spørgsmaalet var om
den første og største. Ved Korstogene blev Kirkens
Sag først for Alvor en Folkesag: ved Christi Kors mod
Muhameds Maane vilde Vestens frie Mænd maale sig i
Kamp med Østens Slaver; ved Korstogene kom Keiseren
og Fyrsterne først i den, efter Kirkens Anskuelse, rette
Stilling til Paven; ved Korstogene fik Munken og Ridderen
i den Grad fælles Opgave, at Munken blev Ridder
og Ridderen Munk; ved Korstogene blev det ellers saa
particulariserede, i Provindser, Stænder, Stæder og Corporationer
deelte Samfundsliv gjennemstrømmet af een
Følelse, bevæget af een Tanke. Eventyr og Andagt,
beregnende Egennytte og mystisk Fromhed, ærgjerrig
Higen og romantisk Længsel, udadstræbende Sandselighed
og asketisk Inderlighed, kort: Tilskyndelser af det allerforskjelligste
Udspring, Bestræbelser udrundne af de
% --- p. 30 ---
\opage{30}\setcounter{footnote}{0}%
mest ueensartede Kilder fløde her sammen i een eneste
uhyre Strøm. Der skal vist ikke kunne nævnes nogen
Aandsretning, nogen enten ideel eller materiel Fremtoning
af Middelalderens Liv, være sig i Handel eller i Videnskab,
i Industri eller i Poesi, i religiøs Ceremonie eller i
ridderligt Galanteri, som jo middelbart eller umiddelbart
stod i Forbindelse med Korstogene.

Korstogenes Idee er en Dualismens Idee, en Idee,
der udfolder Middelalderens Væsen. Da Korstogenes
Tid var endt, fik Paven Feber, det var forbi med hans
sunde Constitution, hans Magt var væsentlig brudt. Det
er ikke uden Betydning, at Dante just begynder sin Nedfart
til Helvede og dermed sin Dommedag over Middelalderen
i det samme Aar, da Bonifacius den Ottende
samlede de Andægtige skareviis om Apostlenes Grave
og ved at indstifte Jubelaaret grundlagde Afladshandelen
i det Store. Ligesom Korstog og Korstog og atter
Korstog var et Kirke-Aandens Tegn paa den kraftig
kæmpende Middelalder, saaledes bliver nu, efter at St.
Peters Stjerne har en Tid lang staaet i Zenith, Aflad og
Aflad og atter Aflad et tydeligt Tegn paa Middelalderens
Synken.

At erhverve sig Syndernes Forladelse ved at kæmpe
som Korsfarer er dog altid at sætte Livet ind for at
vinde Livet; men at tiltuske sig Syndernes Forladelse
for endelige Værdier, ja for Pølser og Skinker og rede
Penge, det er mere end Mysticisme, det er Materialismen
selv, den frækkeste, der er seet i Historien.

Hvorledes Hindringer ved Aandens Magt forvandle
sig til Betingelser, sees nu tydeligt nok paa Reformationen.
Det var ikke fordi det Gode og Ufordærvede i Pavekirken
endnu havde Overvægt over det Slette og Fordærvelige,
at Reformationstanken omsider slog igjennem;
nei, det var tvertimod fordi Fordærvelsens Udslet var saa
% --- p. 31 ---
\opage{31}\setcounter{footnote}{0}%
stærk, saa iøinefaldende, at den ved intet hierarkisk Kunstgreb
lod sig skjule. Hvad det vil sige, at en Centralidee
optager, befrugter og bevæger et System af relative Ideer,
at det aandelige Liv vækkes, gjenfødes og fremmes, idet
Aanden gjennem Øiemedet tildanner Midlerne, lære vi af
Reformationens Gang. Thi vel er det sandt, at de reformatoriske
Bevægelser i det Hele betegne en Reisning af
Nationaliteterne, en Reaction mod det fremmede Aag, en
Staternes Opstand mod det clericale Despoti og mod
Udsugelserne fra Rom; men dog staaer det fast, at denne
nationale og statsøkonomiske Reisning ikke vilde have
ført til noget afgjørende Resultat, endnu mindre have
faaet nogen for Aandslivet blivende Betydning, dersom
den ikke var bleven optaget som relativt Moment i den
absolut religiøse Grundtanke. Vel er det sandt, at de
intellectuelle Vaaben, hvormed Reformatorerne kæmpede,
for en stor Deel hentedes fra Humanismens og den gjenopvakte
Videnskabeligheds velforsynede Arsenaler; men
dog staaer det fast, at hverken de classiske Studier eller
de humanistisk-kritiske Reflexioner vilde ved sig selv have
været tilstrækkelige til at styre Tidsalderens Tankegang,
allermindst til at udfylde Huulheden og væde Livstræets
Rødder med Strømme af et rigt Væld. Vel er det sandt, at
Reformationen hverken vilde eller kunde give noget i
religiøs Forstand oprindelig Nyt, den kunde og vilde kun
udrense det Skadelige ved at afsløre Mistydninger og
fjerne vilkaarlige Tilsætninger; men dog staaer det fast,
at Reformationstanken var en af Aanden selv befrugtet
og just derfor frugtbar Tanke, saa frugtbar og stærk, at den
% "en Slægt": a raised spacing mark inks after "en" (not an apostrophe; not transcribed)
løftede en Slægt ved Begeistringens Kraft og forfriskede
Blodet i alle dens Aarer. Med nye Vækkelser af heroisk
Opoffrelse, med nye Tanker i nye Ideer blev saa Capitalen
indsat, den Livscapital, hvis rige Midler vilde have
givet Folkelivet et varigt Opsving, hvis disse Midler ikke
% --- p. 32 ---
\opage{32}\setcounter{footnote}{0}%
i sørgeligt Maal vare blevne offrede paa dogmatiske
Formler, skolastiske Lærebygninger og absolutistiske
Statskirker, indtil Capitalen omsider svandt hen og fortæredes
i Trediveaarskrigen.

Med Reformationen kom Troestanken vel forsaavidt
til sin Ret, som den fik sit Evangelium renset og sin
Mellemregning opgjort med Middelalderen; men en klar
Udvikling af Troesprincipet som et i sig eiendommeligt
Erkjendelsesprincip blev endnu ikke gjennemført. Midt
i de endeløse Stridigheder om Skriftens og Aandens
% "Credo ut intelligam" is set in true italic (Latin, but italic here, unlike pp. 5, 82, 151)
Vidnesbyrd bevarede det skolastiske \textit{Credo ut intelligam}
endnu sin hele Tvetydighed, en Tvetydighed, der ikke
blot har havt sine dogmatiske, men har endnu den Dag
idag sine ethiske, vidtforgrenede Conseqvenser.

Efterhaanden som Reformationens oprindelige Idee
fordunkledes, efterhaanden som den principielle Modsætning
mellem Guds- og Verdensriget, en Modsætning, om hvis
% Luther's German is set in roman, as the body
Forqvakling Luther selv udbryder: „der lebendige Teufel
hört ja nie auf, mir diese zwei Reichen in einander zu
kochen und zu brauen“, blev i Existensen selv mere og
mere udvisket, viste det sig ogsaa, at de store Reformatorer,
efter Tidsalderens aandelige Vilkaar, havde maattet
lade det rationelle Spørgsmaal om Tro og Viden, den
største af alle intellectuelle Opgaver, henstaae uopløst.
Følgen heraf var, at medens man i den orthodox protestantiske
Dogmatik stadig tog Fornuften fangen under
Troens Lydighed, begyndte man i den protestantiske
Politik mere og mere at tage Troen fangen under Fornuftens
Lydighed. For en menneskelig Ordens Skyld,
altsaa for Fornuftens Skyld, gjorde man kirkelig Rettroenhed
til et statsborgerligt Anliggende. At Fyrster
og Regjeringer, forvexlende deres humane Respect for
Evangeliets hellige Sag med levende Tro paa Evangeliet
selv, toge sig for at gjøre Statsstyrelsen kirkelig og sørge
% --- p. 33 ---
\opage{33}\setcounter{footnote}{0}%
for Undersaatternes Salighed ved kirkelige Tvangslove,
seer vel ved et flygtigt Øiekast ud, som var det en
praktisk Gjennemførelse af den Grundsætning, at tage
Fornuften fangen og lade Troen herske; men det er et
Sandsebedrag. Politikens Verden er en Verden af endelige
Interesser; i de endelige Interessers Verden maae,
det følger af sig selv, de endelige Hensyn herske; ved
at drages ned i en herskende Endelighed bliver Troen
som Tro, ikke Maal, men Middel. At benytte Troen
som Middel er at tage den tilfange. Naar det da i de
orthodoxe Statskirker seer ud, som om Troen sad i
Høisædet og førte Regeringens Scepter: da er det ikke
den frie Tro, der hersker over en verdslig Fornuft; nei,
det er den fangne Tro, hvis tilsyneladende Ophøielse er
saameget mere fornedrende, som den ikke engang faaer
Lov til at fremvise sit Fangenskab, men skal, til Skalkeskjul
for den i Virkeligheden herskende Fornuft, figurere
som Herskerinde.

Den skjulte Modsigelse blev allerede afsløret under
Trediveaarskrigens videre Gang, særlig ved den westphalske
Fred; dog endnu kun for de Klarsynede, de
egentlig Indviede, Diplomaterne. Men Diplomater maae
kunne tie; den diplomatiske Hemmelighed blev først
efterhaanden røbet, idet de mere Oplyste kom til den
Erkjendelse, at det kun kan være verdslige Maal, der
umiddelbart fremmes ved verdslige Midler. Saa forandredes
Tænkemaaden, og, uden at man just kan paavise
noget endeligt Hvorfor, var Religionskrigenes Periode tilende.
Eftersom de religiøse Anliggender mere og mere
trængtes i Baggrunden, trængte naturligviis Fornuftfordringerne,
de theoretiske saavel som de praktiske, sig
frem i Forgrunden. Havde Troen i Form af objectiv,
blind Autoritet, forhen villet tyrannisere Fornuften, saa
naaede den emanciperede Fornuft gjennem Philosophie,
% --- p. 34 ---
\opage{34}\setcounter{footnote}{0}%
Kritik og Politik under stedse stærkere Brydninger med
det Bestaaende dog omsider det Vendepunkt, da Frihedstanken
brød igjennem: det epochegjørende Gjennembrud
var Revolutionen.

Revolutionsaanden staaer i afgjørende Strid med
Kirke-Aanden, de ere hinandens svorne Fjender. Den
revolutionære Frihedsaand er, saa store Lidenskaber den
end i sin Tid vakte hos Folkene, dog ingenlunde nogen
folkelig Aand. Hvad den har tilfælles med de naturlige
Folke-Aander er vel det, at den identificerer Ideen med
Guddommen; men medens Folke-Aanden i sin første
Naturlighed grunder i mythiske Forudsætninger og sygner
hen, naar dens udvalgte Guder blegne, er Revolutionsaanden
derimod saa reflecteret, at den gjør Verdensfornuften
selv til Guddom; det er en kosmopolitisk Aand,
der veed med sig selv, at den hverken kan eller vil
tilstræbe andre Formaal end saadanne, der tilhøre Betingelsernes
endelig-uendelige Verden. Revolutionens Aand
er den til Selvbevidsthed vaagnede Verdensaand, en praktisk
Aand, hvis Lys- og Skyggesider fornemmelig maae
søges i Politiken.

Seet fra Lyssiden er denne politiske Aand en Menneskerettighedernes,
altsaa en Humanitetens Aand; dens
Løsen er Frihed og Lighed; dens Alvor er, at den, om
end først efter mange Kampe, gjennem mange Hindringer,
ad mange Omveie, vil sætte Menneskerettighederne i
Kraft og gjøre den borgerlige Frihed til en virkelig
Sandhed. Det er et charakteristisk Træk, at da den
constituerende Forsamling i Paris (1789) discuterede
Rettighederne, og Mirabeau, utaalmodig over at Discussionerne
antoge et for philosophisk Sving og trak sig i
Langdrag, udbrød: \textit{N'employez pas le mot de droits,
mais dites: Dans l'intérêt de tous, il a été déclaré \dots},
da holdt Forsamlingen med Instinctets Sikkerhed fast
% --- p. 35 ---
\opage{35}\setcounter{footnote}{0}%
paa, at det netop var Rettigheder, altsaa udtrykkelig
Principfordringer, der skulde sættes igjennem.

Men just i Principfordringerne sees en uhyre Strid
og dog mærkværdig Overeensstemmelse imellem Revolutionens
og Reformationens Aand. Kan Maalet naaes,
og Menneskerettighederne sættes i Kraft, dersom det
religiøse Princip fra 1530 forkastes? Vi svare ubetinget:
Nei! Kan Maalet naaes, og Menneskerettighederne sættes
i Kraft, dersom det politiske Princip fra 1789 forkastes?
Vi svare med fuld Overbeviisning: Nei! Naar nu desuagtet
de to Principer ere saa indbyrdes ueensartede, at
en umiddelbar og ligefrem Overeenskomst imellem dem
hidtil har viist sig at være en Umulighed, da er der tilvisse
stillet den nærmeste Fremtid et saa betydningsfuldt Problem,
en saavel for det timelige som for det aandelige Liv afgjørende
Opgave, at en bekymret Nutid, som med Beundring
seer, at Fædrene have udrettet det Store, ikke
tvivlraadig skal vakle og spørge: Hvad bliver der tilovers
for mig?

\begin{center}\rule{3cm}{0.4pt}\end{center}

\chapter*{III.}
\addcontentsline{toc}{chapter}{Forelæsning III}
\markboth{Forelæsning III}{Forelæsning III}
\label{ch:iii}
% pp. 36--53
% --- p. 36 ---
\opage{36}\setcounter{footnote}{0}%
% [large initial N]
Naar Krigsrygterne høres, bliver Børsen urolig; naar
Nordlysene tændes, bliver Magneten urolig; naar Aands\-%
elementerne røres, og Tankerne vækkes til Strid, bliver
Sindet uroligt. Ligesom der kan tales om et Kri\-%
gens Uveir, om et magnetisk Uveir, saaledes kan der
tales, og det med Grund, om et Aandens Uveir. Be\-%
tragte vi Himlen, naar det trækker op til Torden, see
% "Himlen,": a faint raised vertical mark in the thin space before the comma is a risen space/quad (no apostrophe head or tail at 600 ppi); not typed
vi Skyerne bevæge sig, nogle med Vinden, andre, som
det synes, imod Vinden: hvorfra komme disse modsatte
Bevægelser? De komme fra Tiltrækninger af modsatte
Poler, fra modsatte Strømme i lavere og høiere Luftlag.
Udbruddet selv overrasker os desto mere, jo mere For\-%
beredelsen er foregaaet skjult og stille. Paa tilsvarende
Maade med det aandelige Uveir! Forberedelsen skeer
derved, at modsatte Ideer, ueensartede Tanker, under
gjensidig Frastøden og Tiltrækning bevæge sig med sti\-%
gende Spænding i modsatte Strømme gjennem lavere
og høiere Regioner. Bevægelserne kunne i Begyndelsen
være saa skjulte, det hele Røre saa dunkelt ulmende, at
Samvittigheden døser, og den umiddelbare Opmærksom\-%
hed Intet mærker. Men under den tiltagende Spænding
fornemmes en Uro; hvad der forhen yttrede sig som flyg\-%
tige Paastande, glimtende Indfald, abstracte, den praktiske
% --- p. 37 ---
\opage{37}\setcounter{footnote}{0}%
Verden ligegyldige Forestillinger, bliver efterhaanden
paatrængende, udbreder sig stedse videre og faaer mere
Vægt; endda er det Hele, udvortes seet, som en Gaade.
Endelig naaes det kritiske Vendepunkt; de kæmpende
Ideer mødes i Strid: der skeer et Sammenstød, der tæn\-%
des en Gnist, der blinker et Lyn, og en Zittren farer
igjennem Sjælene.

Med Reformationens religiøse og Revolutionens po\-%
litiske Gjennembrud til Forudsætning er den nærværende
Tidsalder idelig bevæget i to diametralt modsatte Ret\-%
ninger, underkastet to aandige Magter, af hvilke den ene
% "aandige" as printed (clear at 300 ppi); sense perhaps "aandelige"
uophørlig modarbeider og bekæmper den anden. Vel
ere Brydningerne ikke til enhver Tid lige følelige; tverti\-%
mod! der gives Mellemtider, i hvilke den aandelige Strid
næsten ganske døer hen, og en stiltiende Overeenskomst
synes at bebude en, om just ikke evig, saa dog langvarig
Fred. Men hvor lidet man kan stole paa den Slags
Fred, maa de foregaaende Tidsalderes Historie med dens
Epocher og Perioder dog vel tilstrækkelig have lært os.
Det gaaer med de hvilende Tilstande i Historien, som
det gaaer i vulkanske Egne. Vulkanen er i Ro; vel
ikke uden al Virksomhed, thi Dampe og Røgsøiler minde
jo om, at der er Noget, der rører sig derinde. Men
Svælget er lukket og Krateret selv med Lava og Aske\-%
hobe en Fordybning, man kan betræde. Saa seer den
hele Omgivelse fredelig, indbydende og betryggende ud:
Lavaen størkner, Viinranken voxer af den varme Grund,
og Menneskene bygge ved Bjergets Fod og fortælle om
gamle Dage, da Jorden revnede, og Krateret spyede Ild.
Dog, det er jo længe siden, Udbrudstiderne ere forbi;
nu har det ingen Fare: saaledes har en fredelig Nærvæ\-%
relse, om den end gaaer svanger med det Forfærdelige,
altid for den umiddelbare Sands noget Bedaarende, ja
noget Bedøvende. Thi det ulmer derinde; det brænder
% --- p. 38 ---
\opage{38}\setcounter{footnote}{0}%
i Dybet: pludselig, før Nogen aner det, høres et Bulder,
og Jorden skjælver, og Flammen slaaer ud, og Kata\-%
strophen er der igjen med alle sine Rædsler.

Det gaaer med de rolige Tilstande i Aandens som
i Naturens Verden; hvad Roligheden og Ligevægten be\-%
tyder, og hvorlænge den vil vare, beroer paa, hvorledes,
under hvilke Forudsætninger og paa hvilke Betingelser
den er tilveiebragt. Den Ligevægt, der for nærværende
Tid tilsyneladende finder Sted imellem Grundsætningerne
fra 1789 og Troessætningerne fra 1530 og i det Hele
mellem de religiøse og de humane Bestræbelser i den
civiliserede Verden, kan med Rette kaldes en ustadig
Ligevægt. Den er ingenlunde fremgaaet af nogen sand
gjensidig Forstaaelse imellem de stridende Magter. Langt
fra! Naar afgjørende Spørgsmaal reises, staae Princi\-%
perne endnu i dette Aarhundrede ligesaa skarpt og ufor\-%
ligeligt imod hinanden som i Slutningen af det forrige.
Hvorledes er den tilsyneladende Ligevægt da tilveie\-%
bragt? Ved en forandret Stemning, ved en endelig, af
praktiske Hensyn ledet Overeenskomst: Betingelsernes,
Midlernes, Interessernes Tid forstaaer sig fortræffeligt paa
praktiske Hensyn og endelige Overeenskomster. Revo\-%
lutionen gik for vidt; dens blodige Conseqvenser viste sig
snart som oprørende Inconseqvenser. Den væltede Thro\-%
nerne, fordi den nedbrød Altrene; det var afskyeligt:
saa vaktes Sandsen igjen for den foragtede Religion;
Trangen til Fædrenes Tro faldt sammen med Trang til
borgerlig Orden: saa maatte Krigshærenes Fyrste selv
indrømme, at et Folk umulig kan leve uden positiv Re\-%
ligion. Hermed havde Revolutionsmagten fra sin Side
anerkjendt Religionen, dog, det forstaaer sig, ikke som
ubetinget og selvgyldigt Maal, men som et høiere politisk
Middel. Paa den anden Side havde Religionens For\-%
kæmpere under vedholdende, skjøndt ingenlunde seierrig
% --- p. 39 ---
\opage{39}\setcounter{footnote}{0}%
Strid mod „den sunde Fornuft“ og den negative Viden
omsider lært, at der maatte gives en borgerlig, af kirkelig
Troesbekjendelse uafhængig Ret og Stilling i Staten, og
fandt sig nu nærmest henviste til ved Hjælp af Følelse,
Phantasie og Romantik at virke paa Gemytterne. Efter
nogle voldsomme Optrin, hvorunder en fanatisk Reaction,
for at slukke de sidste revolutionære Blus, havde givet
den kirkelige Rasen frit Løb, fik endelig Besindigheden
saa megen Overvægt, at de stridende Parter fandt det
raadeligt at tolerere hinanden.

Hvad der i Mellemtiden alt er gjort for atter at
tænde Troens Lys og vække den indslumrede Menighed,
skulle vi med fuld Anerkjendelse paaskjønne; men vi
maae udtrykkelig tilføie, at hvad der i denne Henseende
er naaet, det er naaet, ikke ad videnskabelig Vei, men
alene ved religiøse Midler. Vel have Theologerne paa
deres Viis bestræbt sig for, deels at afværge Videnska\-%
bens Angreb paa Religionen, deels at mægle mellem Tro
og Viden; men at disse Bestræbelser, langt fra at gavne
Religionens Sag, have været den til ubodelig Skade,
derfor har det sidste Aarhundredes Litteratur leveret et
Beviis, der ikke kan modsiges.

Det er ikke Forholdet imellem Kirke og Stat eller
imellem Religion og Politik, nei det er Forholdet imellem
Religion og Cultur, nærmere bestemt, imellem Religion
og Videnskab, imellem den troende og den vidende Be\-%
vidsthed, hvori Knuden oprindelig er knyttet. Som
Bidrag til en Belysning af det Cardinalpunkt, hvorom
Nutidens aandelige Liv med dets Hindringer og Betin\-%
gelser, seet fra det intellectuelle Synspunkt, maa siges
at bevæge sig, skulle vi i faa Træk angive den principielle
Modsætning imellem den troende og den vidende Bevidst\-%
hed, imellem den Bevidsthed, der vil tilegne sig Evan\-%
% --- p. 40 ---
\opage{40}\setcounter{footnote}{0}%
geliets Aand, og den Bevidsthed, der fyldes og bevæges
af Verdensaanden.

Spørge vi os selv, hvori Menneskets Fortrin for
Dyret oprindelig maa søges, da ligger Svaret nær: det
maa søges i Bevidstheden. Dyrets Selvfornemmelse taber
sig i det Sandselige, det Øieblikkelige; dets Vilkaarlig\-%
hed ledes af Instinct og Drift; dets Forstand og Hukom\-%
melse er i den Grad optagen af Sandsebilleder og vex\-%
lende Forestillinger, at det umulig kan faae fat paa sig
selv, gjøre sit Væsen til Gjenstand for særegen Anskuelse
eller, med andre Ord, komme til at anstille nogensom\-%
helst Selvbetragtning. Kunde en Hest f. Ex., eller Hun\-%
den med de kloge Øine blot en eneste Gang tænke ved
sig selv eller sige til sig selv i sit eget Sprog: \textit{jeg er et
Dyr}, saa vilde et saadant Dyr have et Jeg, og Grænd\-%
sen mellem Dyr og Menneske være udslettet. Menne\-%
skets Bevidsthed er i Modsætning til Dyrets udtrykkelig
Selvbevidsthed; heri ligger, at Mennesket har, hvad Dyret
ikke har, et frit, personligt Selv, at Mennesket er, hvad
Dyret ikke er, et aandeligt Væsen, et Jeg. Det er nu
let at forstaae, at intet Menneske kan leve med en aldeles
tom Bevidsthed: Bevidstheden maa opfyldes og bevæges
af et Indhold; som Indholdet er, saaledes er Bevidstheden.
Det Menneske, hvis Bevidsthed udfyldes af sandselige
Forestillinger, sandselige Begjæringer og timelige Inter\-%
esser, er et raat, sandseligt og aandløst Menneske. Først
idet Bevidstheden optager og tilegner sig et Indhold af
væsentlig, almeengyldig Natur, et eiendommeligt Idee-%
Indhold, saaledes at Selvet kan siges at leve deri, at leve
deraf og at leve derfor, bliver Mennesket Aand. Ideen
for sig er endnu ikke Aand, saalidt som Selvet for sig
uden Idee er Aand: Aanden som Aand er Ideens Een\-%
hed med Selvet. Heraf følger, at Aand og Bevidsthed
maae bestemme hinanden gjensidig: som Aanden er,
% --- p. 41 ---
\opage{41}\setcounter{footnote}{0}%
saaledes er Bevidstheden; som Bevidstheden er, saaledes
er Aanden.

Støttende os til disse Forudsætninger skulle vi nu
see, om vi kunne faae Modsætningen imellem Tidsalde\-%
rens ideale Magter, Verdensaanden og Evangeliets Aand,
tilstrækkelig opklaret. Lader os begynde med Ver\-%
densaanden.

Idet Bevidstheden opfyldes af Ideen, være sig Sta\-%
tens, Videnskabens, Kunstens eller hvilkensomhelst ind\-%
% "Statens,": a faint raised vertical mark above the comma is a risen space/quad, not a semicolon dot (a true semicolon on this page, l. 1, has a solid round dot); not typed. Both witnesses read ";"
holdsrig Idee, frigjøres Selvet fra sin naturlige Indskræn\-%
kethed, sin Privategoisme, mættes med Væsen, udvides
og opløftes. Derfor have alle ædle Naturer en umis\-%
kjendelig Trang til at give sig hen i Ideetilskyndelser
og leve for noget Ideelt. Men naar den almene Idee\-%
tilfredsstillelse uden videre skal gjælde for absolut Til\-%
fredsstillelse af Selvet, saa kommer Modsigelsen. Selvet
er timeligt, Ideen evig; Eenheden af det timelige Selv
og den evige Idee er en tvetydig Eenhed, og Verdens\-%
aanden som Følge heraf en tvetydig Aand. Først seer
det ud, som skulde det personlige Selv ved at indaande
Ideen blive evigt; men da de organisk-physiske Betin\-%
gelser, hvortil dets individuelle Tilværelse væsentlig er
knyttet, ingenlunde sikkres ved Ideebevægelsen, men
tvertimod udsættes for Fare: saa ender det med, at den
Individualitet, der har vovet det Yderste i Ideens Tje\-%
neste, ligesom opgives og forraades af Ideen, uagtet Op\-%
offrelsen, der medfører Individets Undergang i det tra\-%
giske Øieblik, har noget for det selv Opløftende og
uendeligt Tilfredsstillende. Modsigelsen er, at den høieste
Tilfredsstillelse af Selvet ender i det Selvløse; og Tve\-%
tydigheden er, at denne Selvopløsning tilfredsstiller Selvet.
Homer har allerede seet det Typiske. At Hektor, der
kæmper saa drabelig for at frelse sit Fædreland ved at
frelse Staden, kæmper for en Idee, er da klart nok;
% --- p. 42 ---
\opage{42}\setcounter{footnote}{0}%
ligesom det ogsaa er klart, at han midt i den forfærde\-%
lige Kamp føler en høiere Tilfredsstillelse: i den time\-%
lige Strid nyder han Ideens Evighed. Saa kommer det
sidste store Øieblik; men i det store Øieblik, da han i
Kamp med Achill skal vove det Yderste, da staaer han,
udvortes seet, forladt af Guderne og bedaaret af Athene,
med andre Ord: priisgiven og forraadt af Ideen. „Vee
mig! til Døden forsand nu Himmelens Guder mig kalde!
Fuldt jeg troede, den kjække Deiphobos stod ved min
Side; men han er hist i Staden, og mig har Athene be\-%
daaret“ (Iliad. XXII). Det kan næsten oprøre En at
være Vidne til Gudernes Ligegyldighed og Partiskhed
mod den ædle Hektor, især maa Athenes Adfærd fore\-%
komme os uværdig. Men saadanne ere nu de ideale
Magter, de mægtig raadende Verdensideer; allerede før
Homers Tid endte de Legen med at give Selvet til Priis,
og de gjøre det Samme endnu den Dag idag. Og dog,
lader os ikke krænke de høie Magter, der er i Verdens\-%
guderne noget Guddommeligt, som ikke bedrager. Hektor
forstaaer det; han forstaaer det heroisk, netop saaledes
som det skal forstaaes i det store tragiske Øieblik. Gu\-%
derne unddrage ham deres Bistand i det Udvortes, thi
hans Maal er sat, og hans Time er kommen. Hektor
forstaaer det: „Nu er min Redning umulig, i forrige
Dage mig frelste Zeus og Phoibos Apollon, hans Søn,
som begge mig fordum huldt beskjermed i Fare, men
nu indhenter mig Skjæbnen.“ Ja, Skjæbnen er den
sidste, altformaaende Verdensmagt; Helten føler det;
men fortvivler han da over Skjæbnen? Nei, thi han bøier
sig med Resignation under Skjæbnen. Spotter han
maaskee Guderne? Nei, thi han beder til Guderne:
„Lad mig da ei uden Daad, uden Roes her gange til\-%
grunde! Lad en Bedrift mig øve, som vordende Slægter
vil mindes!“ Og Guderne høre hans Bøn. Det Gud\-%
% --- p. 43 ---
\opage{43}\setcounter{footnote}{0}%
dommelige er ham uendelig nær; thi det er i ham selv,
det styrker ham i Striden! „Talende saa han af Balg
uddrog sit hvæssede Slagsværd, som over Hoften var
spændt, en stærk og drabelig Glavind, stormed saa frem
som en Ørn, der flyvende høit gjennem Luften oppe fra
sortnende Skyer slaaer ned paa Marken og griber enten
et nyfødt Lam eller og en ængstelig Hare; saa frem\-%
stormede Hektor og svang sit hvæssede Slagsværd.“
Alt paa Homers Tid var altsaa Høidepunktet naaet, det
heroiske Høidepunkt, hvortil Verdensaanden formaaer at
løfte det sig for Ideen offrende, sig i det selvløse Væsen
opløsende Selv: det var mythisk paa Homers Tid, og
det er mythisk endnu den Dag idag.

Skal Selvet frelses derved, at Bevidstheden bevæges,
fyldes og gjennemstrømmes af Livets Idee, da maa Livets
Idee ikke tabe sig i det selvløse Almene, men i Aanden
være Eet med det evige, almægtige Selv, det Selv, der
er Livets Ord, det Ord, der fra Begyndelsen af er Eet
med Gud. „I Begyndelsen var Ordet, og Ordet var
hos Gud, og Ordet var Gud. Det var i Begyndelsen
hos Gud. Alle Ting ere blevne ved det, og uden det
er ikke een Ting bleven til af det, som er. I det var
Livet, og Livet var Menneskens Lys“ (Joh. 1, 1). Saa
% "Menneskens" as printed (clear at 300 ppi); the Danish Bible has "Menneskenes"
taler Evangeliets Aand, ikke for at fremhjælpe en specu\-%
lativ Viden, men for at bestemme Troens Princip og
klare sin Modsætning til Verdensaanden.

Evangeliets Aand fornegter ikke den virkelige Verden
med de virkelige Love; tvertimod! Den gjør jo netop
den virkelige Verden med de strenge Love til en Op\-%
dragelsesanstalt for den menneskelige Aand, en Selvud\-%
viklingsskole for Selvet. Han, Livets Ord, om hvem det
siges, at han kom til sit Eget, og at hans Egne annam\-%
mede ham ikke, han kom ikke med fromme Phantasier
og taagede Forestillinger om en Verden uden Love; nei,
% --- p. 44 ---
\opage{44}\setcounter{footnote}{0}%
han kom til Lovenes Verden, at han, der evig er Eet
med Lovenes Ophav, omsider kunde, saavel i det Natur\-%
lige som i det Aandelige, faae Selvhedslovene opfyldte.
„Hvo, som vil frelse sit Liv, skal miste det; og hvo, som
mister sit Liv for min Skyld, skal frelse det.“ Ved
disse Ord bliver Evangeliets Grundlov sat i skarp Mod\-%
sætning til Verdensloven. Ifølge Verdensloven gjælder
det vel, at den, der krysteragtigt vil frelse sit Liv, skal
miste det; thi kun ved at vove Livet for en Idee, kan
Mennesket vinde et Liv, der er værd at leve. Saa sætter
den Kjække Livet paa Spil, og Krigeren synger i
\textit{Wallensteins Lager}: „Und setzet Ihr nicht das Leben
ein, nie wird Euch das Leben gewonnen sein“. Men
den til Loven knyttede Forjættelse: „hvo, som mister
sit Liv for min Skyld, skal frelse det“, kan Verdens\-%
aanden ikke udtale; i det Øieblik den vil prøve derpaa,
bliver den tvetydig og mythisk. Hvad har den vel ogsaa
at byde den Helt, der vover det Afgjørende? Kun det
gamle Dilemma: „enten Cæsar eller Intet“! Det er vel
værd at mærke sig disse Alternativer. Ideen kan paa givne
Vilkaar yde Selvet den høieste selviske Tilfredsstillelse;
paa Selvtilfredsstillelsens heroiske Høidepunkt er da Hel\-%
ten en Cæsar: fra denne Side seer det ud, som om Ver\-%
densloven virkelig var en oprindelig Selvhedslov. Men
det er et Skin; Forherligelsen af det heroiske Selv har
et Tilfældighedens Præg, er afhængig af Lykken, af
Cæsars Lykke. Imod Lykken staaer Ulykken; og den
mod det enkelte Selv ligegyldige Verdenslov er ligesaa
indgaaet med Ulykken som med Lykken. Verdensloven
er Forandringernes, Udvorteshedens, Omstændigheder\-%
nes, de mødende Tilfældes, de vexlende Betingelsers,
med eet Ord: Andethedens Lov; og naar den vender
denne Side frem, saa trykker Skjæbnen, saa blegner
Helten, saa bliver en Cæsar selv til et Intet. Ved sin
% --- p. 45 ---
\opage{45}\setcounter{footnote}{0}%
Opfyldelse idelig svævende mellem Selvhed og Andethed,
mellem Selvet og det Selvløse, er Verdensloven tvetydig
som Verdensaanden.

Men i Evangeliets Lov er der ingen Tvetydighed;
den er fra Først til Sidst en absolut afgjørende Selv\-%
hedslov og som saadan en Frihedslov. Evangeliets Lov
fritager Ingen fra Virkelighedens Byrder, men den for\-%
vandler de udvortes Hindringer til Betingelser for Selvets
Udvikling og gjør Byrderne lette; den fritager Ingen
fra Smerte og Død, men den giver Taalmod i Smerte
og Seier over Døden: det ligger i dens Væsen som ube\-%
tinget Selvhedslov. Gjennem sin egen Negation, sin
indre Modstand, sit fjendtlige Andet faaer Selvet tilfulde
sig selv i sin Magt; gjennem Forkrænkelighed indgaaer
det til Uforkrænkelighed; derfor burde det ham, der op\-%
fylder Loven, Frihedens fuldkomne Lov, gjennem Lidelse
og Død at indgaae til Herlighed. Hvor Selvet seirer
under Verdensloven, der seirer det egoistisk, thi dersom
det opoffres, saa forsvinder det i det Selvløse og seirer
altsaa ikke: det kan umulig baade opoffres og dog seire.
Naar Selvet derimod seirer under Evangeliets Lov, saa
seirer det just ved Selvopoffrelse. Selvopoffrelsen er saa
langt fra at tilintetgjøre det frie Selv, at den tvertimod
styrker, luttrer og forherliger det: saa guddommelig er
Selvhedsloven, efter at Forbilledet, det evige Selv, der ene
har kunnet opfylde den, har faaet den sat i Kraft.

„Den tragiske Helt“, siger en stor Forfatter, „han
behøver Taarer, og han fordrer Taarer“. Rigtigt! thi
han behøver Medlidenhed og fordrer Medlidenhed: hans
Liden udtrykker det Almene, at Selvets Opløsning er
dets Forklarelse. Han fordrer Taarer. Rigtigt! thi han
vækker Veemod. Selvets Opløsning er dets Forsvinden:
det forsvinder for de Efterlevende, og det forsvinder for
sig selv; den sidste Forklaring er det veemodige Resultat,
% --- p. 46 ---
\opage{46}\setcounter{footnote}{0}%
den i Stemning formildede Modsigelse, at den Forkla\-%
rede ikke mere er sig selv, men et Andet, et aandigt
Efterskin af sig selv, et Minde. Anderledes med Troens
Forbillede; han afviser alle tragiske Taarer, thi han af\-%
viser Medlidenheden. Hans Liden truer ham ikke med
Selvopløsning, den stadfæster kun Selvets Overlegenhed;
just fordi Lidelsernes Vei er hans Vei til Herlighed,
kan han sige til de grædende Qvinder: „I Jerusalems
Døttre! græder ikke over mig, men græder over Eder
selv og over Eders Børn.“

At nu to saa modsatte Principer som Evangeliet og
Verdensaanden maae til enhver Tid bevæge Slægten i
modsatte Retninger, er en Selvfølge. Det er ikke blot
Middelalderen, der har været dualistisk; ogsaa den nyere
Tid er, og enhver følgende Tidsalder vil fremdeles ved\-%
blive at være, dualistisk. Horledes skulde Dualismen vel
% printer's defect: "Horledes" as printed (for Hvorledes; no v in the image)
kunne forsvinde? Det kan da ikke skee derved, at Ver\-%
densaanden med sine Ideer, sine Virkelighedsopgaver og
Interesser forsvinder; thi den forsvinder først ved Ver\-%
dens Undergang. Saa maatte det da skee derved, at
Evangeliets Aand med sit guddommelige Forbillede, sine
Naadegaver og Forjættelser forsvandt fra Jorden; men
saa maatte hvert menneskeligt Selv da først have lært at
trodse Frihedens Lov og foretrække Døden for Livet.
Nei, Tingen er ikke den, at Dualismen skal forsvinde;
men Tingen er, at hver Tidsalder maa opfatte og be\-%
stemme Dualismens Charakteer paa sin eiendommelige
Maade. I Nutiden er Modsætningen tilsløret, det Duali\-%
stiske afdæmpet. Det kommer deraf, at Verdensbevidst\-%
heden med Interesserne for det Nærværende er traadt i
Forgrunden. Den Sandhed, som Evangeliet stiltiende
underforstaaer, men som Middelalderen ikke ret kunde
faae frem, at en Leven for det Tilkommende er uadskil\-%
lelig fra en kraftig Leven i det Nærværende, at Udvik\-%
% --- p. 47 ---
\opage{47}\setcounter{footnote}{0}%
lingen af de humane Ideer væsentlig hører med i det
virkelige Aandens Rige, har vor Tidsalder givet en saa\-%
dan Overvægt, at det Himmelske er nær ved at gaae op
i det Jordiske, som et Primtal gaaer op i sit Produkt.
Skulle det Himmelske og det Jordiske, Evangeliets Aand
og Verdensaanden holde sig i indbyrdes Ueensartethed,
vedblive at danne Extremer uden dog at adskilles ved
en gabende Kløft, da maae Extremerne have en Midte;
en saadan Midte er netop Folke-Aanden.

Hvorledes Folke-Aand og Folkeliv bestemmes ved
hinanden, er allerede omhandlet i al Almindelighed; her
skulle vi endnu forsøge, om end blot foreløbigt, at be\-%
stemme det Folkelige lidt nærmere. Hvad er det Folke\-%
lige? Det er det Grundmenneskelige i naturkraftig Eien\-%
dommelighed. Det høres ofte nok i den nyere Tid, at
Mennesket er et Fornuftvæsen; sjeldnere lægges der be\-%
hørigt Eftertryk paa, at Mennesket ifølge sin Natur op\-%
rindelig er et religiøst Væsen; og dog maae ikke blot
begge Grundbestemmelser udledes deraf, at Mennesket
er et selvbevidst Væsen, et personligt Selv, men den
første maa endogsaa kunne betragtes som afledet af den
sidste. Som religiøst Væsen staaer Mennesket gjennem
Verden i Forhold til Gud; for det personlige Selv er
Gudsforholdet det oprindelige. Som Fornuftvæsen staaer
Mennesket i Forhold til den naturlige Tingenes Orden,
saaledes at grundig Indsigt i Verdensforholdene bliver
ham en væsentlig Opgave. Erkjendelsen af Gudsforhol\-%
det er en religiøs Erkjendelse og beroer paa Tro; Er\-%
kjendelse af Verdensideen og af Tingenes Natur er en
Fornufterkjendelse og beroer paa Viden. Mennesket
har da af Naturen som Menneske, altsaa som folkeligt
Individ, inderlig, væsentlig Trang baade til Tro og Vi\-%
den, saasandt det af Naturen har inderlig, væsentlig
Trang baade til Gudserkjendelse og til Verdenserkjen\-%
% --- p. 48 ---
\opage{48}\setcounter{footnote}{0}%
delse. Selvbevidstheden er, og maa ifølge sin Natur
nødvendig være, en i sig fordoblet Bevidsthed: Bevidst\-%
hed om Andet, og Bevidsthed om sig selv. Som sit
Andet har Selvbevidstheden i absolut Forstand Gud, i
relativ Forstand, Verden; efter sin Grundfordobling vil
den menneskelige Bevidsthed da med Hensyn hertil ad\-%
skille sig i Gudsbevidsthed og Verdensbevidsthed. Hvad
der saaledes gjælder om hvert folkeligt Individ, naar vi
see paa de naturlige Anlæg og den tilsvarende Trang, det
maa uimodsigelig gjælde om Folket som Folk. Ethvert
Folk begynder da sit aandelige Liv med en fordoblet
Grundbevidsthed, et fordoblet Anlæg, en fordoblet Trang,
Trang til religiøs og til fornuftig Selvudvikling, Trang til
Sjælefrelse og Trang til Cultur.

Dersom en folkelig Aandsudvikling nu enten var,
eller i den virkelige Historie kunde være, normal, da
skulde det vel vise sig, hvorledes det religiøse Liv og
Culturlivet, adskilte ifølge Grundfordoblingen i to absolut
ueensartede Bevidsthedskredse, formedelst den personlige
Selvbevidstheds Eenhed med sig selv, gjennem alle For\-%
doblinger udviklede sig i harmonisk Vexelvirkning; men
da ingen Folkehistorie er normal, maae vi føre Beviset
omvendt. Folkelivet er i sin første aandskraftige Be\-%
gyndelse overveiende religiøst; først efterhaanden som
Fornuftudviklingen og dermed Culturen skrider frem,
faaer Verdensbevidstheden Overvægt over Gudsbevidst\-%
heden. Hvorlænge Religionen nu vil formaae at veie op
mod den stigende Cultur, skal netop vise sig at beroe
paa den Dybde, hvori Modsætningsforholdet fra Begyn\-%
delsen af er opfattet, altsaa paa den dualistiske Grund\-%
fordobling. Forskjellen mellem Grækerne og Jøderne
er i denne Henseende charakteristisk.

Det er træffende sagt om de græske Guder, at de
i al deres Overnaturlighed og overmenneskelige Salighed
% --- p. 49 ---
\opage{49}\setcounter{footnote}{0}%
vare altfor naturlige, altfor menneskelige. Heri ligger,
at Gudsbevidstheden hos Grækerne faldt for umiddelbart
sammen med Verdensbevidstheden. Den græske Tro var
uden Dybde, uden Inderlighed og ethisk Alvor. Jo mere
Gudeskikkelserne, saa at sige, trængte paa for at faae
Aabenbarelserne æsthetisk fuldendte, desto mere forvand\-%
lede sig da ogsaa den religiøse Tro til æsthetisk Tro,
til en ved Phantasieanskuelse bestemt æsthetisk Stemning.
Vi see det hos Homer. Den olympisk huuslige Scene
(Iliad. I), hvori Heres Stiklerier bringer Zeus i en saa
huusfaderlig Harnisk, at baade Mennesker og Guder maae
briste i Latter, er vistnok ingenlunde lagt an paa at
give Guderne --- hvad en Lucian langt senere gjorde ---
til Priis for Spot og Foragt; nei, den er et Beviis paa,
at Digteren, uden at bekymre sig om den religiøse Ære\-%
frygt, har benyttet sit poetiske Motiv og forvandlet Reli\-%
gionen til Poesi. Medens Religionen saaledes hos de
Dannede efterhaanden blev forvandlet til Poesi, sank den
umiddelbare, folkelig religiøse Tro ned til raa, phantastisk
Overtro, og Philosophien fik i Fornuftens Navn Ret til
at anvende Kritik og lade Viden opløse Troen. Uagtet
Sokratesses Domfældelse vel maatte gjøre Plato forsigtig,
kan han dog ikke fortie, at de homeriske Guder bør
udelukkes af Republiken.

Medens den religiøse Udvikling hos Grækerne, endog
temmelig tidlig, blev hindret og tilbagetrængt af Cultu\-%
ren, saa blev Culturudviklingen omvendt hos Jøderne i
flere Henseender hindret og tilbagetrængt af Religionen.
Gik den græske Gudsbevidsthed umiddelbart op i Ver\-%
densbevidsthed, saa gik den jødiske Verdensbevidsthed
til Gjengjæld umiddelbart op i Gudsbevidsthed. Ogsaa
hos Jøderne savne vi den gjennemgaaende Modsætning
mellem Gudsriget og Verdensriget, hvortil den oprinde\-%
lige Bevidsthedsfordobling conseqvent maa føre; men
% --- p. 50 ---
\opage{50}\setcounter{footnote}{0}%
rigtignok saaledes, som vi maatte vente at finde den hos
% Rettelser (PDF 338): læs "vente det" (for "vente at finde den"); the print is
% transcribed. A reader has struck "at finde den" in pencil and written "det"
% above it in this copy -- not transcribed.
det udvalgte Folk, paa Aabenbaringens første Trin, hvor
det guddommelige Eftertryk falder paa Religionen. I
det jødiske Theokratie er Forskjellen mellem religiøs og
borgerlig Lovgivning ophævet; Davidsriget er et Guds\-%
rige, Propheternes Moral en religiøs Politik. Troen paa
Jehovas Gjengjældelse ved timelige Straffe og timelige
Belønninger svarer hertil. Følgen heraf var, at Aands\-%
livet med sit frie Indvortes mere og mere tabte sig i et
krystalliseret Udvortes. Efter Exilet, da de prophetiske
Røster omsider forstummede, var det religiøse Indhold
forstenet i Traditioner; Folket, uden Troes-Inderlighed
og uden Cultur, lededes af Pharisæer og Skriftkloge.

Kaste vi et Blik paa Romerfolket, da faae vi ret en
uhyggelig Forestilling om, hvad en endelig Sammenblan\-%
ding af Religiøsitet og Verdslighed har at betyde. De ro\-%
merske Guder ere i Ordets laveste, mest prosaiske Betyd\-%
ning Verdensguder; disse, i sig selv aandløse, paa smaaligt
beregnede Ceremonier og endeløse Formularer smaaligt
opmærksomme Forstandsguder have intet andet Maal
end en grændseløs Udvidelse af Romerrigets Magt; intet
Andet at tænke paa, end hvorledes de kunne være de
romerske Borgere til Nytte og Gavn ved Agerdyrkning,
Fædrift, Faareavl o. s. v. At en saa officiel Religiøsitet
omsider maatte blive sig selv komisk, naar den blev sig
sin Modsigelse bevidst, var naturligt. Tilsidst kunde jo
to Augurer, trods al deres romerske Værdighed, ikke
møde hinanden uden at lee. Hvoraf loe Augurerne?
Det var en Bevidsthedsfordobling; den var kun formel,
og havde aldrig været reel.

Og hvad er Resultatet af alt dette? At Folkelivet
gaaer under, naar Religionen undergraves! Ved den
livskraftige Begyndelse, da alle Folke-Aandens Ele\-%
menter ere i fuld Virksomhed, er Religionen en aan\-%
% --- p. 51 ---
\opage{51}\setcounter{footnote}{0}%
delig Magt; det Hellige vogtes med Ærefrygt: de
% "de er Alvor" as printed (both witnesses and the image agree; sense may want
% "det"). A real word, so not logged as a defect -- noted only.
er Alvor med Gudernes Dyrkelse. Men med den voxende
Cultur foregaaer der en Adskillelse i Samfundet; Ad\-%
skillelsen er paa een Gang borgerlig og religiøs. Fra den
borgerlige Side viser Adskillelsen sig deri, at den store
Hob mere og mere udelukkes fra Deeltagelse i Stats\-%
styrelsen og de offentlige Anliggender. Udelukkelsen er
ikke vilkaarlig; den følger af Forholdenes Natur. Det
er Forskjel i Livsstilling, i Formue, i Begavelse, og
fremfor Alt, Forskjel i Dannelse, hvorpaa det kommer
an. Culturudviklingen grundlægges aldrig, og kan aldrig
grundlægges, ligefrem paa Religionen; det er en Ud\-%
vikling, ikke af Gudsbevidstheden, men af Verdensbe\-%
vidstheden; det er Reflexion og Kritik, Viden og Viden\-%
skab i mangfoldige, deels theoretiske, deels praktiske
Forgreninger og Conseqvenser, der oprindelig afgive
Culturlivets ideelle Grundlag. Med dette Grundlag
bliver Mængden aldrig fortrolig: Lyset skinner den i
Øinene, men den blendes af Skinnet og kan ikke see.
Skulde Mængden følge med, da maatte dens Øie klares
af et andet Lys; da maatte en religiøs Oplysning med
tilsvarende Aandsdannelse træde til og holde Skridt med
Culturen. Naar dette nu ikke skeer, saa synker den
store Hob ned i Raahed, Uvidenhed og Overtro; saa
lever Massen paa Landet i Tilstand af Almue, i de
større Stæder i Tilstand af Pøbel, og Folkeudviklingen
er standset. Staten kan organiseres og udvides, Sæderne
forfines, Videnskab, Kunst og Literatur gjøre beundrings\-%
værdige Fremskridt: det hjælper ikke; Folkets Skjæbne
er afgjort, Livets Puls overskaaren.

Det er et tragisk Skuespil, enhver af de store Old\-%
tidsstater frembyder, idet de nærme sig Culturudvik\-%
lingens Maal: i Midten en lille Kreds af virkeligt Dan\-%
nede, saadanne, som ved Opdragelse og Selvudvikling
% --- p. 52 ---
\opage{52}\setcounter{footnote}{0}%
virkelig have tilegnet sig Culturen; udenom dem en
Fleerhed af Halvdannede, der ere med i Civilisationen
paa anden og tredie Haand, og udenom disse igjen den
tankeløse Mængde, hvis Raahed altid danner det Extrem,
der har Overforfinelsen til Modsætning. Man tænke sig
Rom paa Augusts Tid, de romerske Riddere, de romerske
Sæder, den romerske Luxus, og saa en romersk Pøbel,
hungrig efter Brød og Skuespil (panes et circenses), og
man faaer et Indtryk af, hvad det vil sige, at en Folke-%
Aand kan udslukkes, og et Folk gaae tilgrunde.

Folkedannelsens Princip er væsentlig religiøst. Kun
i Religionens Lys kan Folket som Folk skride frem paa
Culturens Vei og deelagtiggjøres i Civilisationen; ud\-%
slukkes dette Lys, saa sidder Folket i Mørke. Natur\-%
religionerne forsvinde, den ene efter den anden; Old\-%
tidens Lys slukkes, det ene efter det andet: kun eet
Lys udslukkes ikke, og det er Evangeliets Lys.

Det er betegnende for Christendommen, at den fra
første Begyndelse af anmelder sig som en almindelig
Folkesag. Christi Fødsel forkyndes af Englene som en
stor Glæde, der skal vederfares alt Folket; Oldingen
Simeon med Christusbarnet paa sine Arme priser Gud,
at hans Øine have seet Frelsen, der er beredt for alle
Folk, Lyset, der skal oplyse Hedningerne. Christi Op\-%
træden er i alleregentligste Forstand folkelig. Han taler
i Synagogen, han lærer fra Skibet, han prædiker paa
Bjerget, og allevegne strømme Skarer til ham, og Folket
samles om ham og hører ham gjerne, thi „han lærte dem
som den, der havde Myndighed og ikke som de Skrift\-%
kloge“. Hele Jesu offentlige Liv er offret til at vække og
forberede Folket. Men saa stærk er igjen Dualismen af
Helligt og Syndigt, at Offeret ender med et Sonoffer. Først
da Sonofferet er bragt, og han, der opoffrede sig, indgangen
til sin Herlighed, udsendes Aanden, og Menigheden stiftes.
% --- p. 53 ---
\opage{53}\setcounter{footnote}{0}%
Eenheden af det Religiøse og det Folkelige har
just Grundfordoblingen, den mest afgjørende Adskillelse
imellem Guds- og Verdensriget, Gudsbevidstheden og
Verdensbevidstheden, til Forudsætning. Just fordi Chri\-%
stendommen er absolut ueensartet med alt Verdsligt, kan
den paa aandig Maade optage og bruge alt Verdsligt;
den spørger ikke om, hvad Verdensideerne gjælde som
relative Maal i sig selv; thi for det absolute Formaal
ere de alle Midler. Hvorfor kan det, at være Fisker og
det at være Agerdyrker, folkelig forenes med det at være
Christen? Fordi Fiskeri og Agerdyrkning ikke blot ere
relativt, men absolut ueensartede med Christendommen
selv. Af samme Grund kan da ogsaa det, at være
Kunstner, og det, at være Videnskabsdyrker, folkelig
forenes med det at være Christen. Skulde den Vi\-%
dende have Grund til at klage over, at hans Bevidsthed
halveres, fordi de theoretiske Erkjendelser i Videnskaben
ikke falde sammen med de religiøse Erkjendelser i Troen,
da vilde Agerdyrkeren jo ogsaa kunne klage over, at
de agronomiske Erkjendelser ikke lade sig udvikle af
Evangeliet.

Paa Grundfordoblingens Vilkaar vil den religiøse
Oplysning altsaa kunne skride frem med Culturoplys\-%
ningen, og begge forene sig til en sand Folkeoplysning.
At Fordoblingen maa gjøre strengere Fordringer til
Aandslivet, end den formelle Eenhed, er naturligt; at
de strengere Fordringer maae kunne ægge Culturbe\-%
vidstheden til Modstand, og denne Modstand tilveie\-%
bringe en Spænding, som tilsidst udlader sig i et
Aandens Uveir, er muligt: vi skulle i det nærmest Føl\-%
gende see, hvorpaa Spændingen beroer, og hvorledes
% "hvorledes": a speck over the e in this copy, not an accent
den stiger.

\begin{center}\rule{3cm}{0.4pt}\end{center}

\chapter*{IV.}
\addcontentsline{toc}{chapter}{Forelæsning IV}
\markboth{Forelæsning IV}{Forelæsning IV}
\label{ch:iv}
% pp. 54--72
% --- p. 54 ---
\opage{54}\setcounter{footnote}{0}%
% [large initial D]
Det er Verdensaanden med dens Tankeindskydelser
og Ideebevægelser i Nutidens Culturliv, vi skulle forsøge
at charakterisere i nogle faa, men store, anskuelige Træk.
Vi skulle betragte denne Aand, ikke fra Lyssiden, men
fra Skyggesiden; thi Lyssiden er allerede tilstrækkelig aner\-%
kjendt. At Cultur og Oplysning, Dannelse og Videnskab
ere væsentlige Goder, indseer ethvert fornuftigt Menneske,
og hvad der tilfulde er anerkjendt af Enhver, behøver
ikke vor Lov eller Anpriisning. Men den mørke
Side! Verdensaanden, ikke som Culturens og Frem\-%
skridtets Aand., % printer's defect: a stray full stop (or broken sort) stands between "Aand" and the comma; transcribed as printed
% AUDIT doubtful: at 600 dpi the dot sits a little above the baseline, tight against the comma head -- a full stop is possible, but it may be a speck or damaged sort; reading "Aand," rejected only for want of certainty
men som Blendværkets, de kraftige
Vildfarelsers, den skjæbnesvangre Nødvendigheds Aand,
er dette Væsen nu ogsaa tilstrækkelig gjennemskuet, er
det Culturbevidstheden fuldkommen gjennemsigtigt, eller
skulde her i Dybet ikke være Dunkelheder, til hvis Op\-%
klaring der kunde behøves et stærkere Lys? Her skal
naturligviis ikke
tales enkeltviis om nogen særegen Ugudelig\-%
hed hos denne Tids enkelte Mennesker --- at ivre mod Fri\-%
tænkeres Udskeielser og Gudsfornegteres Bespottelser maae
vi overlade til Bodsprædikanterne --- her skal i det Hele
ikke tales om Individernes Vantro, men om Vantroens
Aand i Individerne, en Aand, der forpester Livsluften,
afsætter Giftknuder, og fra de giftige Knudepunkter bre\-%
% --- p. 55 ---
\opage{55}\setcounter{footnote}{0}%
der sig ud og angriber det Ædleste og Bedste i Sam\-%
fundet. Det er fornemmelig Videnskaben, denné Aand % printer's defect: "denné" printed with an acute accent (clear at 300 ppi); transcribed as printed
misbruger; gjennem Naturlære, Historie og Psychologie
søger den ved misledende Kritik at forvanske Troes\-%
sandheden; gjennem Viden og Videnskab aander den
Gift i Sjælene. Men saa er det maaskee Videnskabens
Skyld? Ingenlunde! Videnskaben er reen og ædel og
god, naar den faaer Lov at gjælde for, hvad den op\-%
rindelig er, i objectiv videnskabelig Charakteer. At kaste
Skygge paa Videnskaben sømmer sig kun lidet for en
akademisk Lærer, og vilde staae i Strid med Sandheden:
ingen Green af Videnskaben er enten farlig eller for\-%
dærvelig i og for sig. Kun en udlevet Papisme kan
i fortvivlet Kamp mod de stærke lysbringende Ideer
falde paa at lyse Forbandelse over vor Tids viden\-%
skabelige Fremskridt. Nei, det er hverken Videnskaberne
eller de vidende Individualiteter, hvorom Talen er; Talen
er om et Bevægende, der forvender Ideen og giver den
en falsk Retning. Talen er om et negativt, upersonligt-%
halvpersonligt Uvæsen, der gjør, at Videnskabens Skin
bedaarer den Vidende; Talen er om en Tænken uden
Tænker, en Villen uden Villie; det seer noget mystisk
ud, hvad mon det vel kan være? Vi behøve just ikke at
sænke os ned i den gudsfornegtende Villies inderste Mørke
eller fortabe os i Grublerier over, hvad Skriften kalder
Satans Dybheder, for at faae fat paa vort mystiske Uvæsen;
vi holde os til det Nærmere, det mindre Umenneskelige
og mere Fattelige, og betegne da Verdensaanden, seet
fra Skyggesiden, som den gøglende Tidsaand. For at
faae Syn paa Tidsaanden maae vi særlig rette vort Blik
mod Culturbevidstheden og den i Tidsalderen herskende
Grundidee. Midlernes, Betingelsernes, Interessernes Tid
er ikke ideeforladt. Vel er den i det Praktiske fristet
til at glemme det Absolute, det Ubetingede, men theore\-%
% --- p. 56 ---
\opage{56}\setcounter{footnote}{0}%
tisk glemmer den det ikke; Culturtiden kan aldrig
glemme, at dens klareste, reneste Lys udgaaer fra Viden
og Videnskab; i Videnskaben har den sit Absolute:
Videns Idee er herskende Grundidee. At nu Videns\-%
ideen kan siges at være baade halvpersonlig og uper\-%
sonlig, er let at indsee. Dersom der ingen Vidende var,
kunde der jo heller ikke være nogen Viden: Videns\-%
ideen maa bæres og bevæges af vidende Individualiteter.
Men de vidende Individualiteter maae da ogsaa med det
Samme finde sig saaledes opfyldte, gjennemtrængte og
bevægede af Ideen, at det ikke bliver deres egne Tanker,
der subjectivt bestemme Ideens Væsen, men omvendt Ideens
Væsen med de nødvendige Love, de evige Sandheder, der
objectivt bestemme Tankerne. Men „de evige Sandheder“
og almeengyldige Love ere jo upersonlige: Grundbe\-%
vægelsen er altsaa den, at det Upersonlige bestemmer
det Personlige, at det Selvløse faaer Magt med Selvet.
I denne Svæven mellem Selvet og det Selvløse bliver
Selvreflexionen tilsidst usikker paa sig selv. Den Vidende
tænker „de evige Sandheder“ og er ved denne Tænken
nedsænket i det Almene og ved det Almene bevæget til
intellectuel Resignation. Men den Vidende existerer til\-%
lige i en Verden af endelige, virkelige Interesser. Inter\-%
essen er selvisk; for Interessernes Skyld er den existerende
Tænker paa det Eftertrykkeligste opfordret til at hævde
sin personlige Ret og lægge selvisk Vægt paa sit eget
Selv. Paa een Gang uendelig resigneret og dog endelig
interesseret, er altsaa den vidende Individualitet baade
upersonlig og dog personlig bevæget: hvor ligger nu den
absolute Sandhed?

Den absolut personlige Sandhed ligger udenfor
Videns Grændse og kan alene gribes i Troen. Paa disse
Vilkaar maatte den Vidende da absolut erkjende et Sand\-%
hedsprincip, der er høiere end Viden; men just det, der i
% --- p. 57 ---
\opage{57}\setcounter{footnote}{0}%
dybeste Forstand alene formaaer at tilfredsstille Selvets
Fordring, betragter den Vidende nu som en Modsigelse
i Selvets Væsen, en fortvivlet Selvmodsigelse: her er
Bedraget. Fra dette Bedrag maae de umiddelbart existen\-%
tielle Fristelser og al den Tankeforvirring, hvormed den
til Grundideen svarende Tidsaand blender og bedaarer
Culturmennesket, væsentlig udledes.

Det er ingenlunde Ideen som Idee, der bedaarer
den vidende Individualitet; det er tvertimod Indi\-%
vidualiteten, som ved at forgude sin Viden og tage
Ideen forfængelig blender og bedaarer sig selv. Cultur\-%
menneskets absolut personlige Krav kan og skal Viden\-%
skaben ikke tilfredsstille; fra Vidensforgudelsens objective
Standpunkt maa et saadant Krav ganske rigtig afvises. Men
et absolut personligt Krav kan Personligheden dog ikke
afvise; hvergang det trænges tilbage, indstiller det sig paany.
Det absolut Personlige er Frihedens Ideal; men Videns
„evige Sandheder“ hvile i Nødvendigheden. Fangen af
Nødvendigheden maa Culturmennesket da uden Ophør
hige efter Friheden. Den Tidens Aand, der blender
ham, er netop af den Grund at kalde en ustadig, vægel\-%
sindet Aand, at den idelig svæver mellem to Poler,
mellem Frihedens lyse og Nødvendighedens dunkle
Region; for at skjule Bedraget og fuldende sit Gøglespil
forelsker den sig først i Frihedens Lys og vender sig
bort fra den mørke Nødvendighed, men saa skal Frihe\-%
den igjen være Skin, og Nødvendigheden selv Oplysningens
Kilde.

Vi mindes Asa Loke fra de gamle Myther. Loke
er af Jettenatur, men rigt begavet, fuld af Vid og
Snille, deilig at see paa, naar han vender sig mod
Lyset, saa at Odin selv i Tidernes Morgen har indgaaet
Fostbroderskab med ham. Skjøndt Jette føler han sig
tiltrukken af Aserne, han lader sig bruge i Lysets Tjeneste;
% --- p. 58 ---
\opage{58}\setcounter{footnote}{0}%
men egentlig hjemme i Lysregionen bliver han aldrig,
dertil er hans indre Væsen for mørkt. Seet fra Skygge\-%
siden er han hæslig, og denne Hæslighed vidner netop
om hans Jettenatur. Saa har han med Angerbod avlet
Fenrisulven, Midgardsormen og Hel, og staaer bestandig
i Pagt med Mørkets Rige. Han mangler hverken Liv\-%
lighed eller Kritik; men hans Livlighed er løsagtig, og
hans ondskabsfulde Reflexioner fortære hans Villie.
Ved disse og lignende Træk kan Loke jo vel vise sig
som et Billede paa Tidsaanden, men dog kun paa Tids\-%
aanden i al Almindelighed som den falske, tvetydige Ver\-%
densaand. Vor Opgave er udtrykkelig den, at bestemme
Tidsaanden i det Nærværende ved at vise, hvorledes den
ud fra en mystificeret Viden som fra en skyggende Bag\-%
grund driver sit Gøglespil.

At begynde med lutter Løgn, lader sig ikke
gjøre; en reen Løgn er desuden en Modsigelse, thi i
enhver Løgn er der altid nogle Qvintin Sandhed: Løgn
er forqvaklet Sandhed, en Brøk Sandhed, der udgiver
sig for et heelt Tal. Enhver intellectuel Mystification
maa begynde med noget Utvetydigt for at indføre det
Tvetydige. Det er f. Ex. en utvetydig Sætning, at den,
der har sagt A, ogsaa maa sige B; paa denne Sætning
falder nu Eftertrykket, naar Verdensaanden ved Hjælp
af Viden skal til at forvirre Culturen. Har man ind\-%
rømmet en Forudsætning, saa tvinges man, enten man
vil eller ikke --- hvis man da ikke vil staae som en tan\-%
keløs Person i Selvmodsigelse --- til at indrømme dens
Conseqvenser. Det er fuldkommen sandt og utvetydigt,
at den, som siger A, maa sige B; men glemmer man
over dette Utvetydige at undersøge, om Forudsætningerne
ere eensartede eller ueensartede med de deraf dragne
Conseqvenser, saa kommer Tvetydigheden. Sætningen
er kun gyldig paa den Betingelse, at Forudsætninger
% --- p. 59 ---
\opage{59}\setcounter{footnote}{0}%
og Conseqvenser høre til samme Sphære: fra Forud\-%
sætninger i een Sphære kan man umulig udlede Conse\-%
qvenser, der skulle gjælde i en anden. Der er saaledes
Noget, der hedder Harmonie, saavel i Tonernes som i
Farvernes Kreds; sætter nu Nogen Forudsætningen i
Tonernes Harmonie, og udleder deraf Conseqvenser for
Farverne, saa bliver han Sophist, og dog kan man sige:
Harmonie er altid Harmonie. Man kan tale om Stenenes
Haardhed og om Hjerternes Haardhed, og man kan sige
i al Almindelighed: Haardhed er altid Haardhed; men,
man kan ikke fra Bestemmelser om Stenenes Haardhed
udlede Conseqvenser for Hjerternes Haardhed eller om\-%
vendt. Man kan tale om Lys i Aanden og om Lys i
Naturen, og man kan sige, at Begrebet Lys er altid Lys;
men man kan ikke af det naturlige Lys udlede Conse\-%
qvenser for det aandelige. Her er altsaa Spidsfindigheden:
inden man vel veed deraf, dreies Alting rundt; Conse\-%
qvenser i een Sphære skydes ind under Forudsætninger
i en anden. Gaaer man ind paa en saadan Logik, saa
lader man sig narre, og naar Sophisteriet først er kom\-%
men i Gang, skal det ikke mangle paa Reflexion og
Dialektik for at fuldende Bedraget. Saaledes, naar der
handles om Religion og Videnskab.

Ved religiøse Forudsætninger skal man bøie sig for
religiøse Conseqvenser og ligesaa i Videnskaben for vi\-%
denskabelige; heri er ingen Tvetydighed. Men vender
man det nu saaledes, at man af religiøse Forudsætninger
vil udlede videnskabelige Conseqvenser og omvendt, saa
have vi Sophisteriet. Det er i dette Tidsaandens So\-%
phisterie, at Fordærvelsen ligger; det er dette, som
angriber det personlige Liv fra Roden af; det er dette,
som gjør, at Aandskræfterne aftage, at Charaktererne
blive svagere, at de af Tvivlen svækkede Individua\-%
liteter flagre op mod Troens Lys uden dog at kunne
% --- p. 60 ---
\opage{60}\setcounter{footnote}{0}%
troe. Vi ere, som sagt, forgiftede; Giften er en sophistisk
Blanding af ueensartede, hinanden indbyrdes fortærende
Elementer. Man kan, for at paavise dette, gaae
ud fra Naturvidenskaben, fra Historien, fra Psycho\-%
logien. Vi dvæle i nærværende Sammenhæng ved Na\-%
turvidenskaben.

Naar Tidsaanden gjennem Naturvidenskaben blender
Culturbevidstheden, begynder Sophisteriet omtrent saa\-%
ledes: „Det er vist, at vi Alle uden Hensyn til subjec\-%
tive Ønsker skulle søge efter Sandheden, thi Sandheden
er dog det Høieste, det Evige.“ Heri er ingen
Tvetydighed. „Videre skulle vi erkjende, at hvor der
er Sandhed, der maa der ogsaa være Virkelighed. Vel
vide vi af den høiere Logik, at Virkelighedens Be\-%
greb ikke fuldt udtømmer Sandhedsbegrebet, men hvor
der er fuldstændig Sandhed, maa der nødvendig være en
tilsvarende Virkelighed. Strider Sandheden mod Virke\-%
lighedens Kjendsgjerninger, da mangler der Noget;
Sandheden er da kun halv og i denne Halvhed forsaa\-%
vidt usand.“ Her lægger altsaa Culturbevidstheden, for\-%
saavidt den negativ vender sig mod det Høieste, en
egen Betoning paa Ordet Virkelighed og gjør --- i en
vis Forstand med fuld Ret --- opmærksom paa, at
det, vi kalde Virkelighed, er Noget som vi i det
praktiske Liv maae have ubetinget Respect for. Hvis
vi ikke havde Respect for de virkelige Kjendsgjerninger,
kunde vi jo ikke engang gaae paa egen Haand eller be\-%
væge os frit i Tingenes Verden, vi vilde da f. Ex. ikke
respectere Adskillelsen mellem Fjeld og Afgrund eller
mellem Ild og Vand. Enhver veed, at naar han stikker
Haanden ind i Ilden, maa han brænde sig; han behøver
ikke at gjentage Forsøget mange Gange, nei, et eneste
Forsøg er nok, selv for en Mucius Scævola. Det er kun
skeptisk Forskruethed og speculativ Affectation at ville
% --- p. 61 ---
\opage{61}\setcounter{footnote}{0}%
gjøre Erfaringen usikker ved at lade den beroe paa en
Sum af mange enkelte Tilfælde, eller afhænge af Regler,
der altid tilstede Undtagelser. Ganske vist gives der
Kjendsgjerninger, hvis eiendommelige Naturbeskaffenhed
først efter mangesidige Iagttagelser og omhyggelige Prø\-%
velser kunne sikkres vor Erfaring; men Tvivl om enkelte,
endnu ikke tilstrækkelig prøvede Erfaringssætninger er
væsentlig forskjellig fra Tvivl om Erfaringens hele for\-%
nuftbestemte Grundlag.

De Dybtskuende, som i en idealere Videns Navn
mene at kunne hæve sig over „de empiriske Viden\-%
skabers Stykværk“ ved en Bemærkning som den, at
Naturlærens „saakaldte Love“ kun bestemmes ved In\-%
duction, og at Inductionen som ufuldkommen Slutning,
staaer langt under den kategoriske Syllogisme, vilde
gjøre vel i at betænke sig lidt, førend de gaae videre.
Er det en usikker Slutning, at Massernes Tiltrækning
forholder sig omvendt som Qvadraterne af deres Afstande,
da er det ogsaa en usikker Slutning, at den, der springer
i Vandet, maa blive vaad. Er det en usikker Slutning,
at Ilt og Brint under bestemte Betingelser maae forbinde
sig til Vand, da er det ogsaa en usikker Slutning, at
den, der bliver halshugget, kommer af med Livet. Skulde
nu Nogen for Inductionens Skyld betvivle enten den
mechaniske eller den chemiske Slutnings Almeengyldighed,
lad en saadan dybere Skuende, for Inductionens Skyld,
engang forsøge, om det dog ikke, tiltrods for de blot
empiriske Love, kunde lykkes ham at leve uden Hoved.

For at gjendrive en naturvidenskabelig Induction,
nytter det kun lidt, at man i tom Almindelighed opkaster
Tvivl om Inductionens abstract logiske Gyldighed, nei,
Fordringen, den sande videnskabelige Fordring, er, at
man in concreto paaviser Tilfælde, om end kun et eneste,
der vitterlig staaer i Strid med det ved den be\-%
% --- p. 62 ---
\opage{62}\setcounter{footnote}{0}%
stemte Induction bestemte Resultat. Og naar et saa\-%
dant Tilfælde indtræffer, hvad saa? Ja, saa er det endda
ikke Inductionens almeenvidenskabelige Gyldighed som
Induction, der bliver rokket, men det er en særegen, paa
særegne Tilfælde anvendt Induction, der bliver corrigeret.
Den videnskabelige Erfaringslære har et i sin Art ligesaa
sikkert Grundlag som nogensomhelst apriorisk Fornuft\-%
videnskab. Det er ikke det Uvisse og Svævende i
Naturvidenskabernes Resultater, der skal trøste os i vore
Bekymringer for de religiøse Sandheder, nei, det er
noget Andet; det er for det Første en Afsløring af et
eiendommeligt Bedrag, der ligger i Tidsaanden.

Vi skulle nu see Bedraget, idet vi, for at charak\-%
terisere Tidsaanden, give os i Færd med at udlede reli\-%
giøse Conseqvenser af videnskabelige Forudsætninger
og omvendt videnskabelige Conseqvenser af religiøse
Forudsætninger; vi skulle i nogle Exempler see, hvor\-%
ledes Sophisteriet kommer i Bevægelse.

Et Exempel fra det gamle Testamente! Jehova,
hedder det (1 Mosb. IX.), sætter udtrykkelig Regnbuen i
Skyerne, idet han ved en høitidelig Handling slutter
Pagt med Noah. Den religiøse Forudsætning er nu den,
at Jehova virkelig ved en bestemt Villiesact har sat dette
Tegn, altsaa gjort et Under, og, hvad der jo ogsaa er et
Under, talt til Noah og fortolket ham Tegnet. „Min
Bue vil jeg sætte paa Skyen, og den skal være et Pagtens
Tegn imellem mig og imellem Jorden. Og det skal skee,
naar jeg samler Skyen over Jorden, da skal Buen sees
i Himmelen, og jeg vil mindes min Pagt, som er imellem
mig og imellem Eder og imellem hver levende Sjæl af
alt Kjød, og ikke mere skal Vandet blive til en Flod at
ødelægge alt Kjød; men Buen skal være paa Himmelen,
og jeg skal see den og mindes den evige Pagt imellem
Gud og imellem alt Kjød, som er paa Jorden.“

% --- p. 63 ---
\opage{63}\setcounter{footnote}{0}%
Skulle vi nu af denne religiøse Forudsætning have
Lov til at uddrage en videnskabelig Conseqvens, da
gjøre vi vist bedst i, jo før jo hellere, at kaste Optiken
paa Ilden. Thi Optiken kan ikke og vil ikke paa nogen
Maade see i Regnbuen enten et Under eller et ud\-%
trykkeligt Mindetegn om Guds Pagt med alt Kjød.
Optiken forklarer Regnbuen slet og ret af Lysstraalernes
Brydningsvinkel, altsaa som en naturnødvendig Conse\-%
qvens af naturlige Betingelser.

Man forsøge da at gaae den omvendte Vei; man
gjøre sig det ret tydeligt, hvorledes Regnbuen dannes
under visse Betingelser, og hvorledes vi ikke blot i
Skyerne, men ogsaa i Vandspring kunne see Regnbuer:
saa have vi da her en sikker, videnskabelig Forudsætning.
Skal man nu af denne videnskabelige Forudsætning
drage religiøse Conseqvenser, kommer man i Strid med
Religionen. Thi, sige de Vidende, vi vide, at det, som
siges i det gamle Testamente om Regnbuen, ikke er
noget Virkeligt, men en Phantasie; det gamle Testamente
kjender ikke Regnbuens Dannelse, Betingelser og Love.
Det er en barnlig Aand, ubekjendt med de virkelige
Kjendsgjerninger, som taler i de Ord, der lægges Jehova
i Munden. Men er dette saa, da kan det heller ikke
være sandt, at Jehova selv har talt til Noah om nogen Pagt;
Pagten selv er her ligesaa indbildt som den Forklaring,
der gives af dens Phænomen. Naar den ubekjendte For\-%
fatter af første Mosebog lader Jehova sige: „Og det skal
skee, naar jeg samler Skyen over Jorden“, eller: „naar
jeg fører Skyen frem“, da opfatter han Skysamleren Jehova
ligesaa billedlig, barnlig og mythisk, som Homer opfatter
Skysamleren Zeus (Νεφεληγερέτα Σεύς). % printer's defect: capital sigma Σ printed for zeta (Homeric Ζεύς); the sort shows Σ's pointed middle at 600 dpi. Transcribed as printed. The ρ is the curly-tailed form, typed ρ.
Er det nu i det
Hele betegnende for Tidsalderens aandelige Liv, at Tanke\-%
bevægelser i saa modsatte Strømme gaae igjennem Cultur\-%
bevidstheden frem og tilbage, hvad Under da, om Tvivlene
% --- p. 64 ---
\opage{64}\setcounter{footnote}{0}%
voxe en stille Væxt; hvad Under, at de Anstukne nøies
med Palliativer af æsthetisk opstillet Gudsdyrkelse, med
Lidt af Troen og Lidt af Viden, og føle kun liden
Trang til at indlade sig med dybere religiøse Problemer.

Et andet Exempel! I første Mosebogs første Kapitel
møde vi den høieste, den første af alle religiøse Forud\-%
sætninger i Fortællingen om Skabelsen. „I Begyndelsen
skabte Gud Himmelen og Jorden.“ Opmærksomheden
fæster sig strax paa Jorden: „Og Jorden var øde og
tom“ o. s. v. Om Jorden som Midtpunkt maa Himlen
hvælve sig som en udstrakt Befæstning midt imellem
Vandene. „Og Gud gjorde den udstrakte Befæstning
og adskilte imellem Vandene, som vare nedenfor Be\-%
fæstningen, og Vandene, som vare ovenfor Befæstningen,
og det skete saa.“ Det siges her udtrykkelig, at Himmel\-%
legemerne bleve til for Jordens Skyld. „Og Gud gjorde
de to store Lys, det store Lys til at have Herredømme
om Dagen, og det lille Lys til at have Herredømme om
Natten, og Stjernerne. Og Gud satte dem paa Himmelens
Befæstning til at lyse over Jorden.“

Jorden, seet fra Religionens Standpunkt, er altsaa
ikke blot den physiske Midte i Verdensrummet, men
om vi saa maae sige, Guds Øiesteen, fordi der paa
Jorden er Noget, Gud sætter ubetinget Priis paa, Men\-%
nesket nemlig, der bærer hans eget Billede. Ligesom
Stjernehimlen med Sol og Maane er til for Jordens
Skyld, saaledes er Jorden igjen til for Menneskets Skyld:
Mennesket er ikke blot Jordens Herre, men Skabningens
Krone. Just derfor betragtes det som en Verdenskata\-%
strophe, at Mennesket bukkede under i Fristelsen og
tabte sit Gudsbillede; og just derfor betragtes det som
en ny Skabelse, at Gud selv i sin eenbaarne Søn ved at
paatage sig Kjød og Blod opreiser den faldne Slægt og
giver den det tabte Billede tilbage. Just derfor hedder
% --- p. 65 ---
\opage{65}\setcounter{footnote}{0}%
det, at, naar Christi Rige er bleven forkyndt for alle
Hedninger, saa følger den sidste Katastrophe, Verdens
Undergang og Dommen. Idet Christus forudsiger
Verdens Undergang, siger han, at der skal skee Tegn i
Sol og Maane, at Solen skal formørkes og Stjernerne
falde ned. Dette er Grundanskuelsen i den aabenbarede
Religion, og denne Grundanskuelse er ikke tilfældig.
Alt er fra Først til Sidst saa sammenhængende, at dersom
vi rokke ved een Ting, saa ende vi med at opløse det
Hele: den Verden, som skal være en Skueplads for Guds
egen Menneskebliven, maa være en Centralverden, det
høieste Midtpunkt i Universet.

Skal man nu af denne religiøse Forudsætning uden
videre uddrage videnskabelige Conseqvenser, saa maa
man bryde Staven over Nutidens Physik og Astro\-%
nomie. Thi Astronomien indrømmer ikke, kan ikke og
vil ikke indrømme, at denne Jord er Verdens Midtpunkt.
Tvertimod! hvis man vil høre Videnskabens Stemme,
faaer man et andet Syn paa Verden, et virkeligt Syn
paa det Virkelige.

Hvad Astronomien lærer os om Verdensbygningen, er
ikke løse Meninger og tomme Formodninger, det er virke\-%
lige Kjendsgjerninger og virkelige Love, Noget, der er saa
uimodsigeligt, at man ikke kan slaae det hen i Veiret.
Det gjælder her, at jo mere dannet et Menneske er, desto
større Respect har han for den virkelige Natur og dens
Love. Jo lavere han staaer, desto lettere slaaer han sig
til Ro i den Indbildning, at de Lærde dog ikke vide rigtig
Besked, og at man kun faaer den sande Physik og den
rette Astronomie ved at tage Bibelens Ord bogstaveligt.
Ja, paa Bogstavtroens Dovnepude slumrer man sødelig ind;
men man sover sig ikke til Sandheds Erkjendelse. Hvis
Øiet er sundt og Forstanden vakt til at see, med hvilken
Sikkerhed og Conseqvens Videnskaben gaaer frem ved
% --- p. 66 ---
\opage{66}\setcounter{footnote}{0}%
at bestemme Forholdet mellem Jord og Himmel, vil
man erkjende, at Jorden, langt fra at være Verdens
Midtpunkt, ikke engang er Midtpunktet i den lille Deel
af Verden, som kaldes vort Solsystem. De herhen
hørende Læresætninger blive ikke blot opstillede, de
blive ogsaa begrundede. Den, der ligeover for Viden\-%
skaben paa dens nuværende Standpunkt vil kunne be\-%
tvivle Uimodsigeligheden af en Sætning som den, at det
ikke er Solen, der bevæger sig om Jorden, men Jorden,
der bevæger sig om Solen, at det altsaa er Solen og ikke
Jorden, der er Planetsystemets Centrum, lad ham be\-%
tvivle Uimodsigeligheden af den Sætning, at et hovedløst
Menneske ikke kan leve. At Jorden i vor lille Verden
er en forholdsviis ringe Størrelse ved Siden af Kloder
som Jupiter, Saturn og Uranus, staaer klart for Enhver,
der har gjort et Indblik i Videnskaben. Nu betænke
man hine hellige Ord: „Der skal“ --- for Jordens Skyld ---
„skee Tegn i Sol, Maane og Stjerner: Solen skal for\-%
mørkes, og Maanen ikke give sit Skin, og Stjernerne
falde ned af Himmelen, og Himmelens Kræfter røres.“
Og hvorfor? Fordi Maalet paa Jorden er naaet; fordi
Menneskehedens Levnetsløb er tilende, fordi det Høieste
er forkyndt for alle Folk og enten tilegnet i Tro eller
forkastet i Vantro.

Hvorledes det nu vil gaae, naar man af videnskabe\-%
lige Forudsætninger skal uddrage religiøse Conseqvenser,
kan Enhver let sige sig selv. Der staaer i Aabenbarin\-%
gens Bog (VI, 12---13): Og jeg saae, at det (Lammet)
% Printer's defect: the opening „ before "Og jeg saae" is not printed
% (the closing “ after "Veir." is); transcribed as printed. Quote balance +1 “.
oplod det sjette Segl, og see, der skete et stort Jord\-%
skjælv, og Solen blev sort som en Haarsæk, og Maanen
blev som Blod. Og Himmelens Stjerner faldt ned paa
Jorden, som et Figentræ nedkaster sine umodne Figen,
naar det røres af et stærkt Veir.“ Til Grund for den
Antagelse, at Stjernerne kunne falde til Jorden, ligger
% --- p. 67 ---
\opage{67}\setcounter{footnote}{0}%
aabenbart en dobbelt Forestilling, deels at Stjernerne ere
meget smaa i Sammenligning med Jorden, deels at Af\-%
standen imellem Himlen og Jorden endda just ikke kan
være saa overordentlig stor.

Hos de Gamle fortælles Noget, som vidner om,
hvor lidet den menneskelige Phantasie er i Stand til
at fatte de uhyre Afstande i Universet. Homer (Iliad. I,
586) lader Hephaistos, der saa gjerne vil staae Here
bi imod den vrede Zeus, tiltale sin Moder: „Den olym\-%
piske Drot er slem at bekæmpe. Alt tilforn engang,
da til Hjælp jeg vilde dig komme, tog han ved Foden
mig fat, og smed mig fra Gudernes Tærskel. Dagen
igjennem jeg svæved omkring, men da Solen sig dukked,
faldt paa Lemnos jeg ned, ei synderligt Liv var der i
mig.“ Det gyser i En ved at tænke sig den faldende
Hephaistos svæve en heel Dag imellem Himmel og Jord;
hos Hesiodus bliver denne Svæven endog udstrakt
til ni Dage og ni Nætter, men hvad forslaaer det? Hvor
langt mon der da være fra Jorden til en Stjerne udenfor
Solsystemet, til en Fixstjerne?

Betragte vi Høiden fra en af de nærmeste Fixstjerner,
da maae vi bruge Lysets Hastighed, 40,000 Mile i Se\-%
cundet, til Maalestok og ikke Hephaistoses Fald. Af\-%
standen fra Solen til Jorden, der er 20 Millioner Mile, gjen\-%
nemløbes af Lyset i 8 Minutter og 13 Secunder; men fra
den nærmeste Fixstjerne, Alpha i Centauren, er Lyset
% "3½": printed as a case fraction, superior 1 over inferior 2.
3½ Aar, fra Sirius 14 Aar, fra Polarstjernen 42 Aar
underveis, inden det naaer Jorden. Polarstjernen kunde
altsaa slukkes og dog endnu lyse for os i 42 Aar, efter
at den var udslukket. Det er uhyre Afstande, og dog
svinde de ind til ubetydelige Størrelser, naar vi hæve
vort Blik til Melkeveiens Stjerner, til Stjerner, om hvilke
det uden Overdrivelse kan siges, at de lyse for os med
præadamitisk Lys.
% --- p. 68 ---
\opage{68}\setcounter{footnote}{0}%

Hvis vi, for at føie Tidsaanden, tillade os den
Taabelighed at drage religiøse Conseqvenser af viden\-%
skabelige Forudsætninger, da vil det i Sandhed komme
til at gaae ud over Religionen. Kan den religiøse An\-%
skuelse af Stjernehimlen ikke have sand religiøs Gyl\-%
dighed, med mindre den videnskabeligt stemmer med
den astronomiske, da lader os være oprigtige og sige til
Christus i al Frimodighed: Herre, Du har taget feil, Du
har ikke studeret Din Astronomie; og da Du ikke er
videre sikker i Din Astronomie, saa tør vi heller ikke
ubetinget troe, at Du altid er ganske sikker i Din Theo\-%
logie. Du taler om, at Stjernerne paa Dommens Dag
skulle falde fra Himlen. Hvorledes skulde de falde fra
Himlen, dersom ikke Din Mening var, at de skulde falde
paa Jorden? Men Sirius og Polarstjernen og de endnu
fjernere Sole kunne umuligt falde paa Jorden, da Jorden
er et imod dem forsvindende Punkt.

Lad os forlade Himmelrummet og dvæle ved Jorden
og see, om det da ikke her er tilladt af videnskabelige
Forudsætninger at drage religiøse Conseqvenser, eller
omvendt; det lader saa smukt, at Videnskaben er
lidt religiøs, og Religionen lidt videnskabelig. Aaben\-%
baringens hellige Ord taler om Jordens Skabelse, Talen
er fyndig og kort. Jorden var øde og tom, og der var
Mørke over Afgrundene, og Guds Aand svævede over
Vandene; saa høres Almagtens Ord, og Dag for Dag
skrider Værket frem til Fuldendelse. Efter Skabningens
Fuldendelse, efterat Mennesket, dannet af Jord, er skabt i
Guds Billede, fortælles der saa om Paradiis og den Til\-%
stand, da Fredens og Velsignelsens Aand hvilede over
det Skabte. „Og Gud saae alt det, han havde gjort, og
see, det var saare godt.“ Dette er en sand og væsentlig
religiøs Forudsætning, og den, der undergraver denne
Forudsætning, seet som Princip, undergraver Reli\-%
% --- p. 69 ---
\opage{69}\setcounter{footnote}{0}%
gionen. Men dersom man heraf skal uddrage viden\-%
skabelige Conseqvenser, saa maa man nødvendigt bryde
Staven over Geologien. Geologien er en ny Videnskab,
der som saadan endnu har mange Dunkelheder; men den
skrider dristig frem: den har allerede sin sikkre Methode,
sin Maalestok og sine Correctiver. Ved at bryde Staven
over en saa grundigt forskende Videnskab, vilde vi
komme i afgjørende Strid med den hele Cultur. Om\-%
vendt, skal man fra videnskabelige Forudsætninger drage
religiøse Conseqvenser, altsaa gaae ud fra Læren om
Jordlegemet og det organiske Livs Udvikling paa
Jorden, da kommer man til det Resultat, at det ikke
kan være gaaet saaledes til med Skabelsen, som der fortælles
i første Mosebog. Der maa da være hengaaet lange,
lange Mellemrum, inden den ene Art Skabninger fulgte
paa den anden; Jorden selv maa være undergaaet mange
Omdannelser, inden Mennesket kunde fremtræde paa den.
Og hvad bliver der af Edens Have: den moderne Vi\-%
denskab har jo slet ingen Plads for et Paradiis?
Nei, Middelalderens Litteratur veed Besked om Pa\-%
% "divina comedia": roman, as printed (not italic, not spaced).
radiset, vi see det i Dantes divina comedia. Paa
Dantes Tid var Jerusalem med Zions Bjerg endnu
Centrum for Jorden paa denne Side; til Antipode havde
det saa paa den anden Side Skjærsildens Bjerg, paa
hvis Top Paradiset laae, det saligt jordiske, hvor
% A reader's pencil mark in the outer margin here: not transcribed.
Dante mødte sin Beatrice. Men skal man i vor Tid
drage religiøse Conseqvenser af videnskabelige Forud\-%
sætninger, saa har der aldrig været noget jordisk Para\-%
diis. Eet endnu! Ifølge Skriften er Synden kommen
ind i Verden ved Mennesket og Døden ved Synden;
altsaa er Døden Syndens Sold, og den, som krænker
denne Sætning, krænker det Religiøse. Men skal man
uddrage videnskabelige Conseqvenser, maa man krænke
denne Forudsætning; thi Geologien viser, at mange lavere
% --- p. 70 ---
\opage{70}\setcounter{footnote}{0}%
% Rettelser (PDF 338): læs "stridt med". The print has "mod"; a reader has
% struck the o in pencil and written "e" above it in this copy: not transcribed.
og høiere Dyreslægter have stridt mod hinanden og for\-%
tæret hinanden, længe før Mennesket blev til, at der
altsaa har været Død og Smerte før Mennesket, at For\-%
krænkeligheden ikke kom, fordi Mennesket syndede, men
at Mennesket syndede, fordi det blev til i Forkrænkelighed.
% "Dersom" is flush left (no indent): no paragraph break, as printed.
Dersom det nu virkelig er sandt, at Verdensaanden,
under nærværende Vilkaar nærmere bestemt som Tidsaand,
netop fremkogler sine Illusioner derved, at videnskabelige
og religiøse Lysstraaler under vildsomt krydsende Inter\-%
ferenser danne gøglende Taagebilleder; dersom dette, sige
vi, er det Bedaarende i Tidsalderens Væsen, da maa et
saadant Væsen jo ogsaa blive aabenbart i tilsvarende
Phænomener, og disse Phænomener særligt bidrage
til en Charakteristik af Nutidens aandelige Liv. Vi
opkaste os ikke til Dommere over det Nærværende i
dets umiddelbare Virkelighed; vi betragte kun Phæno\-%
menerne, forsaavidt de med indre Nødvendighed maae
antages at ledsage Tidsaanden.

Her maa da for det Første frembyde sig to Extremer,
to hinanden aandeligt modsatte Tilstande. Jo fjernere
nemlig Individerne staae den videnskabelige Oplysning,
desto mindre ville Videnskabens Resultater indvirke paa
dem, desto lettere ville de kunne fastholde de religiøse,
i Kirken overleverede, i Overleveringen bestemte Sand\-%
heder; det er den Tilstand, hvori det overveiende Fleertal af
et Lands Befolkning naturligt befinder sig. Blandt Almuen
og de Halvdannede, som endnu have bevaret nogen væsentlig
Alvor, vil Aabenbaringens Ord da i det Hele falde
sammen med Kirkens og Bibelens Ord og gjælde med
guddommelig Autoritet. Jo nærmere Individerne derimod
staae det videnskabelige Culturlys, jo alsidigere Dannelsen
er, og jo mere Kritiken er vakt, desto stærkere ville de
frastødes af den hellige Skrifts Vidundere, af dens enfol\-%
dige Ord om Aandens Almagt. Fuldt forvissede om den
% --- p. 71 ---
\opage{71}\setcounter{footnote}{0}%
videnskabelige Verdensbetragtnings ubetingede Ret imod
den bibelske, ville da disse Vidende lade den positive
Religion være Almuens Sag, medens de selv, deels aaben\-%
lyst, deels i Stilhed, af alle Kræfter arbeide paa at under\-%
grave Troens Grundvold hos de Dannede, dem de saa,
naar det endelig forlanges, byde til Vederlag og fuld
Erstatning en vandklar-uklar Tankereligion, en Reli\-%
gion af Hvidt i Hvidt og Graat i Graat, „det Absolute
selv“ i logisk Skyggerids og metaphysisk Dæmring.

Imellem disse Yderligheder bølger den store Masse,
fornemmelig i de store Stæder, den bevægelige Mængde
af mere og mindre Dannede, der udgjøre Tidsalderens
medvidende, paa Alt reflecterende Publikum. Betragtes
et saadant Publikum fra Skyggesiden --- og det er ud\-%
trykkelig Skyggesiden, vi her have for os --- da er dets
Bevægelighed allernærmest at forklare deraf, at Alt
sidder løst i Bevidstheden: Viden sidder løst, og Troen
sidder løst. Vist er her endnu nogen Tro tilbage, et
dunkelt Indtryk af det Mystiske, en praktisk Respect for
det kirkeligt Bestaaende. Men, det forstaaer sig, de Dan\-%
nedes Religion maa jo være en „perfectibel Religion“, de
Dannedes Præst en perfectibel Mand, den dannede Tro
revideret af Viden; saa gjælder det at træffe det bedste
Arrangement og faae en kirkelig Orden slaaet fast.
Men at slaae Noget fast, hvor Grunden er løs, har sine
Vanskeligheder; saa er Publikum bekymret, og dets
Førere bekymrede, de Ulærde bekymrede for Viden, de
Lærde for Troen: Ja, her er alskens Bekymringer; Be\-%
kymringer i Publikum og for Publikum, Bekymringer
for det Opbyggelige i Bekymringer for det Videnskabe\-%
lige, Bekymringer nok for „Silkepræster“ og speculative
Theologer!

Men i disse Bekymringer er der ingen Alvor. Og
hvorfor ikke? Fordi det ligger i Sagens Natur, at der
% --- p. 72 ---
\opage{72}\setcounter{footnote}{0}%
ingen kan være. Naar Hindringerne for Nutidens aande\-%
lige Liv oprindelig henføres til en intellectuel Forvirring,
en uklar Blanding af religiøse og videnskabelige Betragtnin\-%
ger, da maa Alvoren svare dertil; skal den virkelig svare
dertil, maa dens Grundopgave være af intellectuel Natur.
% "ville": a speck over the final e (read "villé" by tesseract) is not an accent.
Tidsalderen maa ville fæste sit Blik paa de første Hindringer
for i dem at opdage de første Betingelser; den maa ville
tage Modsætningen mellem det Religiøse og det Viden\-%
skabelige skarpt i Øiesyn. Om Kløften mellem Tro og
Viden, Religion og Videnskab end er saa stor, at den
seer ud som en gabende Afgrund: det har dog ingen
Fare, naar Tænkningen blot kan vækkes til Alvor. Cul\-%
turbevidstheden kan, naar den vil, forstaae, at her trods Mod\-%
sætningen er Noget, som den Troende og den Vidende maae
have tilfælles, og det er en sand Overbeviisning. Cultur\-%
bevidstheden kan, naar den vil, forstaae, at sand Over\-%
beviisning umulig kan hverken for Troens eller for Videns
Vedkommende danne sig uden klar og grundig Tænk\-%
ning, at altsaa Troen, saasandt den er et selvstændigt
Erkjendelsesprincip, maa have Tænkningen med, lige saa
vel som Viden. Den Dannede maa, naar han vil, kunne
indsee, at Tænkningen, om den end i Forhold til de ab\-%
solut ueensartede Principer bevæger sig forskjelligt i de
forskjellige Kredse, dog overalt vil kunne forstaaes af
Tænkning, og den tænkende Bevidsthed forstaae sig selv.
% "Her er" is flush left (no indent): no paragraph break, as printed.
Her er i de mange Hindringer Mulighed til ligesaa mange
Betingelser; her er Reflexioner nok, og Aandens Opgave
kræver Reflexion; her er Dannelse nok, og Aandens Op\-%
gave kræver Dannelse; her er Frihed nok, og Aandens
Opgave kræver Frihed; men Eet er endnu fornødent:
paa dette Ene vil Alt beroe, og det er --- Alvor.

\begin{center}\rule{3cm}{0.4pt}\end{center}

\chapter*{V.}
\addcontentsline{toc}{chapter}{Forelæsning V}
\markboth{Forelæsning V}{Forelæsning V}
\label{ch:v}
% pp. 73--89
% --- p. 73 ---
\opage{73}\setcounter{footnote}{0}%
% [large initial I]
„I skulle“, siger Christus, „erkjende Sandheden, og Sand\-%
heden skal frigjøre Eder.“ Den moderne Viden
opfatter og tilegner sig disse Ord paa en ret eiendommelig
Maade, og under denne Tilegnelse anbringer Verdens\-%
aanden sine Tvetydigheder. Sandheden erkjende vi ved
at prøve Forskjellen mellem Sandt og Usandt; For\-%
skjellen mellem Sandt og Usandt prøves ved Tænkningen:
herimod er Intet at indvende; men nu kommer For\-%
dreielsen. Tænkningen, siger man, er kun grundig, klar
og reen, forsaavidt den er videnskabelig. At erkjende
Sandheden vil altsaa sige at erkjende i Videnskaben, hvad
der er Gyldigt og Ugyldigt i Religionen. Det Religiøse,
der udholder Prøven i Videnskaben, kan anerkjendes,
det, som ikke bestaaer Prøven, maa forkastes. Efter denne
Opfattelse er altsaa Videns Lys det høieste Lys. Viden\-%
skaben sidder paa Pilati Domstol; Religionen er stævnet,
og Dommeren spørger: hvad er Sandhed? Ja nu forstaaer
Culturbevidstheden, hvad det vil sige: „I skulle erkjende
Sandheden, og Sandheden skal frigjøre Eder“; det vil
jo sige: „I skulle erkjende Videnskaben, og Videnskaben
skal frigjøre Eder.“ At Videnskaben, videnskabelig
seet, virker, har virket og fremdeles vil virke rensende
og frigjørende paa Sind og Tanke, skal ikke blot ind\-%
% --- p. 74 ---
\opage{74}\setcounter{footnote}{0}%
rømmes, men, for dog saavidt muligt at afværge taabe\-%
lige Misforstaaelser, atter og atter indskærpes. Men
Frigjørelse og Frigjørelse ere to Ting. Et er Bevidst\-%
hedens Frigjørelse fra falske Meninger og indgroede
Fordomme; et Andet er Sjælens Frigjørelse fra For\-%
krænkelighedens Baand til et evigt Liv; Et er den viden\-%
skabelige, et Andet den personlige Sandhed. Den Sand\-%
hed, Christus taler om, er ikke Videnskabens almeen\-%
gyldige, upersonlige, fornuftnødvendige Sandhed, men den
personlige, frie, frigjørende Sandhed. Idet han taler om
Sandheden, taler han af sig selv om sig selv og siger
udtrykkelig: „Jeg er Veien, Sandheden og Livet“. Her
er Knudepunktet, som ingen Videnskab er istand til at
løse, det Knudepunkt, at en existerende Personlighed er
Sandheden, ja, at den evige Sandhed selv er denne i
Kjød og Blod existerende Personlighed. Lad kun den
Videnskab, der med et betegnende Udtryk kalder sig
Videnskabernes Videnskab, Philosophien, anstrenge sig
af alle Kræfter, den vil i det Høieste drive det til at
udvikle Personlighedens og Sandhedens Begreb og Be\-%
grebernes Eenhed i Ideen. Forvexlende Personlighedens
almene Idee med Personligheden selv, kan man saa maaskee
indbilde sig at have grebet det Absolute i det rene Begreb
og fremmet Sandheden ved at sige i Christi Navn: „I skulle
erkjende Philosophien, og Philosophien skal frigjøre Eder.“
Men er det Knudens Løsning? Det er jo en stor Mis\-%
forstaaelse. Vil Philosophien i Videns Navn underkaste
Evangeliet en Kritik, opløse den historiske Christus og
forvandle Gudmenneskets Existens til en Idee, det skal
den have Lov til; ihvorvel man, hvad Frugtbarheden
angaaer, vel kunde fristes til den Bemærkning, at en
Langhalm, der er tærsket saa tidt, neppe vil give nogen
Kjerne. Men, hvad Philosophien paa ingen Maade har
Lov til, er at udgive sin kritiske Opløsning af Christi
% --- p. 75 ---
\opage{75}\setcounter{footnote}{0}%
personlige Existens for en Løsning af Christendommens
Gaade. At Christus, ifølge den evangeliske Opfattelse ---
og denne maa dog vel gjælde for den oprindeligt reli\-%
giøse --- ikke er kommen til Verden for at hjælpe Philo\-%
sophien til et almindeligt Begreb om Personlighedens
Idee, men for at hengive sit Liv, sit virkelige, personlige
Liv til Verdens Frelse, har han tydelig beviist i Ord
saavel som i Gjerning. Vi høre ham i Johannes Evan\-%
gelium (6te Kap.) sige: „Jeg er det levende Brød, som
kom ned af Himmelen; om Nogen æder af dette Brød,
han skal leve til evig Tid; og det Brød, som jeg vil
give, er mit Kjød, hvilket jeg vil give for Verdens Liv.“
Dette er nu ganske vist en haard Tale, saavel for Vidende
som for Uvidende. Og dog lægger den Talende selv
afgjørende Vægt paa sit Ord. Thi da Jøderne kivedes
indbyrdes om, hvorledes han kunde give dem sit Kjød
at æde, sagde han til dem: „Sandelig, sandelig siger jeg
Eder: dersom I ikke æde Menneskens Søns Kjød og
drikke hans Blod, have I ikke Livet i Eder.“ Det er
ikke Brødets Begreb, men det virkelige Brød, der mætter
den Hungrige, og det er ikke Personlighedens Idee, men
den virkelige Personlighed, der frelser de Troende. I
Christi Tale gjøres der ikke Forskjel imellem Vidende
og Ikke-Vidende, thi for begge Slags Tilhørere kan
Talen være lige forargelig; men der gjøres udtrykkelig
Forskjel imellem Troende og Ikke-Troende: Kun for
de Ikke-Troende er den forargelig. „Forarger dette
Eder? Hvad om I da faae at see, at Menneskens Søn
farer op, hvor han var før? Det er Aanden, som levende\-%
gjør; de Ord, som jeg taler til Eder, ere Aand og ere
Liv. Men der er Nogle iblandt Eder, som ikke troe.“

Men, siger man, lad saa være, at Christus giver sig
selv hen, ikke som Personlighedens Idee, men som virkelig
Personlighed, til Frelse for de Troende, saa vil han jo dog,
% --- p. 76 ---
\opage{76}\setcounter{footnote}{0}%
som vi høre af hans Tale, ikke paa sandselig Maade til\-%
egnes af kjødædende Disciple og bloddrikkende Tilhængere.
Nei, hans Ord, der kun ere Liv, forsaavidt de ere Aand,
maae tilegnes i Aand og Sandhed. Men naar det dog
tilsidst kommer an paa Aanden og en aandelig Tilegnelse
af Christus, saa maa jo Aandens Indhold klares ved
Tænkning, og ved Tænkning gaaer Videnskabens Dør
op, og see, saa ere vi dybt inde i Viden! Ja, vi
ere dybt inde i Misforstaaelsen, i fuld Gang med
at forblande Aabenbaringens Aand med Viden\-%
skabens. Tidsaandens Vidende dømme forresten
om Aabenbaringens Aand, hvad de ville; men Eet
maae de uden kritisk Overanstrengelse kunne indsee,
og det er, at denne Aand i Væren og Virken, i Kom\-%
men og Svinden ikke har mindste Lighed med Viden\-%
skabens Aand. Hvorledes kom Aanden paa Pintsefesten
over Apostlene? „Og der Pintsefestens Dag var kommen
% "de": a speck over the e (read "dé" by the layer) is not an accent.
(Ap. G. c. 2), vare de alle eendrægtige tilsammen. Og
der kom pludselig en Lyd af Himmelen, som af et frem\-%
farende vældigt Veir, og fyldte det ganske Huus, hvor
de sadde.“ Hvis en Spotter her tillod sig den Bemærk\-%
ning, at en Aand, der kommer med en saadan Blæst,
ikke kan være videnskabelig: paa hvem faldt saa Spotten
tilbage? Ikke paa dem, der adskille Tro og Viden!
„Og der viste sig for dem Tunger, som af Ild, der for\-%
deelte sig, og satte sig paa Enhver af dem. Og de bleve
alle opfyldte af den Hellig Aand og begyndte at tale
med andre Tungemaal, eftersom Aanden gav dem at
tale.“ Det er nødvendigt, at en Aands Eiendommelighed
maa charakteriseres ved dens eiendommelige Virkemaade.
Men at den sproggivende Aands Virkemaade her er,
ikke blot noget forskjellig, men uendelig forskjellig fra
den videnskabelige, vil Ingen vove at negte. Vi ville til
yderligere Bekræftelse høre af Apostelen Peters egen
% --- p. 77 ---
\opage{77}\setcounter{footnote}{0}%
Mund, hvorledes Aanden kommer, og hvorledes den vir\-%
ker. Hvorledes kommer den da, og hvorledes virker den?
„Og det skal skee i de sidste Dage, siger Gud, da vil
jeg udgyde af min Aand over alt Kjød; og Eders Søn\-%
ner og Eders Døttre skulle prophetere, og de Unge
blandt Eder skulle see Syner, og de Gamle have Drømme“.
Er her i slige Udtalelser mindste Antydning af, at den
apostoliske Aand nogensinde vil kunne identificere sig
med Videnskabens Aand? Vistnok vil Christus tilegnes
i Aand og Sandhed, men Tilegnelsen er ikke videnska\-%
% A small dot on the baseline before "tvertimod" (a stray mark or inked quad) is not transcribed.
belig: tvertimod siger han i Matthæi 2det Kap. de be\-%
kjendte Ord: „Jeg priser Dig, Fader, Himmelens og
Jordens Herre! at Du har skjult dette for de Vise og
Forstandige og aabenbaret det for de Umyndige.“ En
stærkere Protest mod Videnskabens Opfatning kunde
han ikke nedlægge, end naar han udtrykkelig siger, at
kun de Umyndige forstaae ham, og at de Vise og For\-%
standige slet ikke kunne forstaae ham. Modsætningen
imellem de Umyndige og de Vise er ikke en Modsæt\-%
ning imellem de Tankeløse og de Tænkende; Christus
har aldrig sagt, at han ikke kunde forstaaes af de Tæn\-%
kende, men kun af de Tankeløse; hvorledes skulde han
vel kunne sige til de Tankeløse: „I skulle erkjende Sand\-%
heden, og Sandheden skal frigjøre Eder“? Nei, ved
de Umyndige forstaaer han de Enfoldige, der begynde
med Tro, hvis Erkjendelse altsaa fra Først til Sidst maa
blive en Troeserkjendelse; ved de Vise og Forstandige
tænker han paa de Vidende; og saa gjælder hans Ord
bogstaveligt, at den himmelske Fader har skjult Evan\-%
geliets Sandhed for de Vidende og aabenbaret den for
de Troende.

At Christus ikke vil tilegnes af Vidende, men af Troende,
sees ogsaa paa hans Apostle. Vi maae dog vel indrømme, at
Apostlene have tilegnet sig Christi personlige Liv paa en
% --- p. 78 ---
\opage{78}\setcounter{footnote}{0}%
sand, men ikke paa en videnskabelig Maade. De have
ikke ved at tænke over hans Personlighed klaret sig
Personlighedens objective, almeengyldige Begreb. Vist\-%
nok er Paulus en Tænker, en af de største i Verden,
men hans Tanker ere ikke videnskabelige. Det er en
ligefrem Mishandling af den store Tænker, at man ud\-%
byder hans Lærebegreb i Form af et videnskabeligt
System, tager, om jeg saa maa sige, Lærebegrebet med
op paa Cathedret, men Apostelen selv kommer ikke
paa Cathedret. Om en videnskabelig Tænker, en Philo\-%
soph f. Ex., gjælder det, at naar man prøver hans
Viden, er det ligegyldigt med hans Person. Men
naar man prøver en Apostels Lære, prøver man med
det Samme hans Liv. Apostelens Tanker ere Per\-%
sonlighedstanker, hans saakaldte Lærebegreb en per\-%
sonlig Forklaring af hans Tro, hans Liv i Christus,
hans Opoffrelse for Menigheden. Hvorledes Aanden,
hvor den virker i Kraft, uden Ophør forvandler Hin\-%
dringer til Betingelser, medens Verden omvendt frem\-%
% A raised tick after "Hindringer" (no opening quotation mark precedes it) is taken as a stray mark, not transcribed.
stiller disse Betingelser som lutter Hindringer, kan, som
vi allerede have seet, udvikles af Troestanken, og Troes\-%
tanken udvikles af Troens Princip, uden at det dermed
er afgjort, at den, der saaledes udvikler Troestanken,
selv er en Troende. Anderledes med en Apostel; han
udvikler Troestanken, idet han forklarer sit Liv. Vi
vælge et Exempel hos Paulus: „I ere“, siger han til de
„opblæste“ Corinthier (1 Cor. 4) „allerede blevne mætte, I
ere allerede blevne rige, I ere blevne Herrer foruden os,
og gid I vare blevne Herrer, paa det og vi kunde blive
det med Eder. Thi mig synes, at Gud haver frem\-%
stillet os Apostle som de Ringeste, som han antvordede
til Døden; thi vi ere blevne et Skuespil for Verden,
baade for Engle og Mennesker. Vi ere Daarer for
Christi Skyld, men I ere Kloge i Christo, vi skrøbelige,
% --- p. 79 ---
\opage{79}\setcounter{footnote}{0}%
men I stærke, I herlige, men vi foragtede. . . . Overskjel\-%
dede velsigne vi, forfulgte taale vi, bespottede formane
vi, vi ere blevne som Udskud i Verden, Alles Skovisk
indtil nu.“ See, det er en Deel, en meget væsentlig
Deel af „det paulinske Lærebegreb“; men et saadant
Lærebegreb lader sig ikke filtrere til videnskabelig Sy\-%
stematik, med mindre man bortfiltrerer den paulinske
Existens. Men naar man har faaet den Eenhed af
Tænken og Handlen, der charakteriserer den aposto\-%
liske Existens, videnskabeligt renset ud, hvad saa? Ja,
saa er Lærebegrebet maaskee theologisk, men det er
ikke længer apostolisk.

Et andet Kjendemærke, hvorved vi let kunne gjøre
Skilsmisse mellem den videnskabelige og den religiøse
Tænkning, er Miraklet. Det er et charakteristisk Mærke
for Aabenbaringen: uden Mirakel ingen Aabenbaring og
uden Aabenbaring ingen Tro. Tages Miraklet bort eller
skydes det tilside, saa udviskes Aabenbaringens Charak\-%
teer, saa svækkes Troen. Videnskaben sige om Miraklet,
hvad den vil; den kan ikke negte, at Miraklet er Reli\-%
gionens charakteristiske Mærke. At ville forklare og
begrunde Religionen videnskabeligt, uden at ville for\-%
klare og begrunde Miraklet videnskabeligt, er Nonsens:
gives der en Religionsvidenskab, saa maa der ogsaa
gives en Mirakelvidenskab. Men Miraklet er Videns
Skranke, det er Videnskabens Skandalon. Videnskaben
kan, naar den gjøres til Dommer, negte Undergjernin\-%
gerne, og dermed Religionen, det er der Mening i; men
den kan ikke erkjende Religionen og udviske Miraklet.
Christi Personlighed er et Mirakel. Det er ikke nok at
søge hans Eiendommelighed deri, at han har gjort
mange Mirakler: nei, Christus er Underet selv, Aaben\-%
baringens store Mirakel. Hermed stemmer det conse\-%
qvent, at den Aand, der skal forklare Christus, kommer
% --- p. 80 ---
\opage{80}\setcounter{footnote}{0}%
ved et Mirakel, at Troen, hvori Aandens Gaver tilegnes,
gives ved et Mirakel; fra alle Sider støder Viden her
paa sin oprindelige Skranke, sit Skandalon, nemlig
Miraklet.

Man har jo rigtignok villet tage Sagen fra en anden
Side, idet man har sagt, at Religionens Mirakel vel ikke
optages i verdslig Viden, men at man just af den
Grund har gjort Forskjel mellem verdslig og hellig
Videnskab. Som hellig Videnskab anmasser Theologien
sig ikke at ville begribe Miraklerne, thi Miraklerne ere
ubegribelige; nei, ikke ved at begribe, men ved at for\-%
udsætte det Ubegribelige og fremhæve Aabenbaringens
og Troens Kjendsgjerninger i skarpe Træk, vil Theolo\-%
gien levere en Troeslære i videnskabelig Form og for\-%
sone Tro med Viden.

Dette Standpunkt er vistnok valgt i den bedste
Mening, i den troskyldige Forudsætning, at det engang
skulde lykkes paa Troens Grundlag at reise en ny vi\-%
denskabelig Bygning. Men det vil aldrig lykkes. Alle\-%
rede det, at her lægges et Grundlag, bestaaende af over\-%
naturlige Kjendsgjerninger, altsaa af uvidenskabelige Forud\-%
sætninger, og at man paa dette Grundlag vil opføre en
Lærebygning af noget Videnskabeligt, allerede det maa,
som i vort Foregaaende er viist, vække Betænkelighed.

Der er sagt, at den troende Dogmatiker forholder
sig til de hellige Kjendsgjerninger paa samme Maade
som Naturforskeren til Naturens Kjendsgjerninger; hvilken
uhyre Misforstaaelse! De naturlige Kjendsgjerninger vise
sig, naar Spørgsmaalet er om Viden og Videnskab, som
ganske ueensartede med de hellige. Naturens Kjends\-%
gjerninger ere underlagte Naturlovene; her haves det
Almeengyldige, det Fornuftnødvendige, det, som gjør
Prøvelse og Kritik mulig. Hvad et saadant fornuftnød\-%
vendigt, almeengyldigt Grundlag betyder, kan sees end\-%
% --- p. 81 ---
\opage{81}\setcounter{footnote}{0}%
og i de materielle Former. Spørges der f. Ex. om hvad
Granit og Flint er, da kunne disse Begreber med til\-%
svarende Kjendsgjerninger fastsættes saa bestemt, at
Enhver, som har sund Forstand, Fordannelse og Sands
for videnskabelige Undersøgelser, han være forresten
Jøde, Græker, Muhamedaner, uden Hensyn til folke\-%
lige Forskjelligheder og Tro, bliver tvungen til at ind\-%
rømme Begrebernes Gyldighed: „Dette maa være Granit,
dette maa være Flint“. Saaledes ogsaa, naar der spørges
om Svovlsyre, Kulsyre o. s. v. Seer man derimod hen
til saadanne overnaturlige Kjendsgjerninger, som at for\-%
andre Vand til Viin, at vandre paa Havet, at opvække
Døde, da er det umuligt at fremstille slige Kjendsgjer\-%
ninger saaledes, at de antage Præg af noget Almeen\-%
gyldigt eller Uimodsigeligt. Grækere, Jøder, romersk
Katholske, Protestanter see jo de hellige Phænomener med
forskjellige Øine. Dem, som troe, kan Videnskaben ikke
gjøre vantroe; dem, som ikke troe, kan den ikke gjøre
troende: her er imellem Tro og Vantro en personlig
Modsætning, som Videnskaben ikke magter. Den væ\-%
sentlige Forskjel mellem de naturlige og de hellige
Kjendsgjerninger sees i denne Sammenhæng fra en ny Side.
De naturlige Kjendsgjerninger ere nærværende. Natur\-%
gjenstandene foreligge til stedse fornyet Prøvelse, thi Na\-%
turens Rige er et altid nærværende Rige. Anderledes
med de hellige Kjendsgjerninger; de kunne ikke mere
blive Gjenstand for fornyet Selvsyn. Han, som gjorde
% A raised tick after "nærværende" (cf. p. 78 "Hindringer") is taken as a stray mark, not transcribed.
Miraklerne, er jo ikke mere sandselig nærværende;
Undergjerningernes Tid er forbi: „Der skeer“, siger
man, „i vore Dage ingen Mirakler“. Aabenbaringens
overnaturlige Kjendsgjerninger høre til Forbigangen\-%
heden. De ere ikke at sammenligne med Naturens
Kjendsgjerninger, men snarere med Historiens. Sæt\-%
ningen forandres altsaa derhen, at den troende Dogma\-%
% --- p. 82 ---
\opage{82}\setcounter{footnote}{0}%
tiker forholder sig til de hellige Kjendsgjerninger paa
samme Maade som Historikeren til de historiske Kjends\-%
gjerninger. Mon dette vil gaae? Vel forudsætter Er\-%
kjendelsen af Historiens Kjendsgjerninger en historisk
Tro; men Evangeliets Kjendsgjerninger fordre udtryk\-%
kelig en religiøs Tro, og imellem den religiøse (fides
religiosa) og den almene historiske Tro (fides historica)
er et Svælg befæstet. De historiske Kjendsgjerninger
underkaster Videnskaben en Kritik; tør vi nu ogsaa,
spørger Culturbevidstheden, anvende Kritik paa de hel\-%
lige Kjendsgjerninger? Ja eller Nei? Hvis der svares
Nei, er det jo øieblikkelig klart, at den hellige Histories
Kjendsgjerninger faae en ganske anden Stilling for Be\-%
vidstheden end Kjendsgjerningerne i den profane Hi\-%
storie. Skal Bibelkritik holdes ude ved et dogmatisk
Magtsprog, saa maae vi nødvendigt vende tilbage til den
gamle Lære, at ethvert Ord, enhver Tøddel i den hel\-%
lige Skrift er indblæst af Gud, at Testamenterne, det
gamle saavel som det nye, ere i Eet og Alt ufeilbare
Aabenbaringsdokumenter. Den videnskabelige Reflexion
maa pludselig standse, hvorved? Ved Autoritet --- blind
Autoritet! Og dog skal den være videnskabelig! Iste\-%
denfor Grunde, fuldgyldige, videnskabelige Grunde be\-%
gynder vor Dogmatiker sin „hellige Videnskab“ med at
paakalde den blinde Autoritet. Lad os betragte den
blinde Autoritet lidt nærmere! Muhamedanerne f. Ex.
have ogsaa deres Aabenbaringsdokument, deres hellige
Kjendsgjerninger, deres Autoritet. Foruden Koranen an\-%
erkjende Sunniterne en levende Tradition og i denne
Tradition en Mængde guddommelige Eventyr om Pro\-%
pheten. Han blev, fortælles der, engang haanet af de
Vantroe, fordi han ikke strax kunde bevise sin guddom\-%
melige Sendelse ved et Mirakel. Længe taug den For\-%
haanede sagtmodig; men omsider slog hans Time. Halv\-%
% --- p. 83 ---
\opage{83}\setcounter{footnote}{0}%
maanen stod just skinnende blank paa Himmelhvæl\-%
vingen; han kaldte paa den, og den faldt med en Sølv\-%
klang til Jorden. Saa tog han den op og bar den under
sin Arm, som man bærer en Chapeaubas. Da dette var
skeet, kastede han den op i Luften; den fløi da til
% "paa": the final a is a broken sort, only a fragment inked.
Himlen igjen og lyste, som før, paa Hvælvingen. Det
er en hellig Kjendsgjerning; vee den Sunnit, der ikke
kan troe paa et saadant Mirakel. Ved blind Autoritet
kan Kritiken jo vistnok standses; men, spørge vi, kan
en saadan Standsning betragtes som Indledning til et
videnskabeligt Lærebegreb? Kan blind Autoritet forenes
med Christi Ord: „I skulle erkjende Sandheden, og
Sandheden skal frigjøre Eder“? Man maatte da fortolke
dem saaledes: „I skulle erkjende den blinde Autoritet,
og den blinde Autoritet skal frigjøre Eder“.

Uden videnskabelig Kritik kan der ikke være Tale
om Videnskab. Der indrømmes altsaa Kritik, forstaaer
sig, med Maade, Varsomhed, Ærbødighed. Men her
gjælder Sætningen: har man sagt A, maa man ogsaa
sige B. Skal Kritik anvendes, maa den anvendes con\-%
seqvent, ikke blot paa naturlige Begivenheder, men ogsaa
paa Mirakler og Spaadomme. Er det videnskabelige
Synspunkt først slaaet fast; er man først godtroende
gaaet ind paa den grundfalske Forudsætning, at de hel\-%
lige Kjendsgjerninger skulle sees i samme Lys, have
% "Virkeligked": printer's error (k for h), as printed; read "Virkelighed".
deres Virkeligked i samme Sphære og godtgjøre deres
Realitet paa samme Maade som de profane: da siger
den personlige Sandhed Godnat, thi nu skal det Aaben\-%
barede mønstres paa Vrangen. Det er „den moderne
Videnskab“, efter at den ved Theologiens Bistand har
vundet sikkert Fodfæste i de religiøse Erkjendelsers
Kreds, en let Sag at vise, hvorledes Miraklerne udgjøre
et Væv af Modsigelser, hvorledes Spaadommene sigte
paa noget ganske Andet end deres foregivne Opfyldelse,
% --- p. 84 ---
\opage{84}\setcounter{footnote}{0}%
og hvorledes de bibelske Beretninger ikke blot stride
mod andre historiske Beretninger, men endog mod hver\-%
andre indbyrdes.

Vi vælge som Exempel det Mirakel, hvorved Chri\-%
stus opvækker Lazarus, et Mirakel, der skal bevise
Sandheden af hans Ord: „Jeg er Opstandelsen og
Livet“. Det Indtryk, Miraklet gjør paa de mange
Vidner, er høist forskjelligt, idet Nogle ære Christus i
Tro, medens Andre gaae som hemmelige Angivere til
Pharisæerne. Og hvad gjøre saa Pharisæerne? „Der\-%
% "(Joh. 11)": the numeral is blotted in this copy (two upright strokes under an ink smudge; layer "tø", tesseract "W"); read 11 from the passage quoted, Joh. 11,47--53. Doubtful.
% AUDIT doubtful: at 600 dpi the two heavy strokes look like a reader's hand-written "11" laid over the printed numeral; a printed bowl like a "9" shows under the second stroke. Printed reading may be "9" or "19"; kept 11.
for“, hedder det (Joh. 11), „forsamlede de Ypperste-%
Præster og Pharisæerne Raadet, og sagde: hvad gjøre
vi? thi dette Menneske gjør mange Tegn. Dersom vi
lade ham saaledes blive ved, ville Alle troe paa ham;
og Romerne skulle komme, og tage baade vort Land og
Folk . . . Fra den Dag af raadsloge de derfor om at
ihjelslaae ham.“ Men hvad er dog dette for en Psycho\-%
logie? Ypperste-Præsterne og det ganske Raad begynde
jo som Troende med den Forudsætning, at Jesus virke\-%
lig „gjør mange Tegn“, ja saa vidunderlige Gjerninger,
at dersom han fik Lov at blive saaledes ved, saa vilde
„Alle troe paa ham“. Lazari Opvækkelse staaer altsaa
for disse Skriftkloges Bevidsthed, ikke som et Bedrag,
et Skinmirakel, udført af en forslagen Gøgler, men som
en for Mange vitterlig Kjendsgjerning, et Vidunder, det
ikke kan falde dem ind at negte. Men, hvordan kan det
da i al Verden falde saa kloge Mænd ind at ville ihjel\-%
slaae et Menneske, der kan opvække de Døde? Antages
et Menneske virkelig at have Magt til at opvække de
Døde, da behøves der ikke religiøs Tro, men kun lidt
sund Menneskeforstand, for at indsee, hvor unyttigt, ja
hvor meningsløst det er at stræbe et saadant Menneske
efter Livet.

Pharisæerne og de Skriftkloge, vil man maaskee
% --- p. 85 ---
\opage{85}\setcounter{footnote}{0}%
sige, troede kun historisk, men kunde ikke troe religiøst,
thi Gud havde forhærdet deres Hjerte. For at faae
endog kun et Videnskabsglimt ud af alt dette, maatte
man da kunne forklare, hvormeget af Forhærdelsen
der tilhører disse Skriftkloge, og hvormeget der kommer
fra Gud. Men gjøres der Alvor af en saadan Opgave,
maa dog vel Enhver kunne see, at Hjerternes Forhær\-%
delse ikke vil gaae op i noget videnskabeligt Begreb.
Hvilken Vrimmel af Modsigelser kan Videnskabskritiken
ikke her paavise! Det hedder i Jacobs Brev (1, 13---14):
„Ingen sige, naar han fristes; jeg fristes af Gud; thi Gud
fristes ikke af det Onde, men han selv frister heller
Ingen“. Men at forhærde et Menneskes Hjerte er jo
dog at friste det til det Onde. „Derfor“, siger Christus,
„taler jeg til dem“ (det jødiske Folk) „ved Lignelser,
fordi de, skjøndt seende, dog ikke see, og hørende dog
ikke høre, og forstaae ikke heller. Og i dem fuld\-%
kommes Esaias Spaadom, som siger: med Eders Øren
skulle I høre, og ikke forstaae, og see med Eders Øine
og ikke fatte.“ Theologien vil udligne Modsigelserne
saaledes, at Forklaringen paa een Gang skal fyldestgjøre
Troen og tilfredsstille Viden; den vil i al Tjenstagtighed
tjene to Herrer, men faaer Utak af dem begge. Uop\-%
hørlig forvexlende Troens Psychologie med en viden\-%
skabelig Bevidsthedslære, opnaaer den kun at berede den
moderne Videnskab en let Seier og drive dens negative
Kritik til det Yderste.

Ved at ville forsvare de hellige Kjendsgjerninger
videnskabeligt, bidrager Theologien til at vedligeholde
det Bedrag, at disse Kjendsgjerninger for Alvor kunne
trues af videnskabelige Angreb. Men kunne de hellige
Kjendsgjerninger trues af videnskabelige Angreb, da
ligger heri en Indrømmelse af, at de muligviis engang
kunne falde for slige Angreb. Den moderne Videnskab
% --- p. 86 ---
\opage{86}\setcounter{footnote}{0}%
har taget Indrømmelsen til Indtægt og svaret med trium\-%
pherende Spot: „De ere faldne!“

Der gaaer igjennem den religiøse Psychologie en
egen Række Reflexionsfordoblinger; om disse Reflexions\-%
fordoblinger med deres umiddelbare Udslag i tilsyne\-%
ladende Modsigelser har Videnskaben intet Begreb, og
kan --- saalænge Troens egen Tænkning endnu ikke har
faaet Lov til klar og alsidig Udvikling af sine egne Er\-%
kjendelser, men hver Gang den vil begynde paa en
systematisk Selvudvikling, øieblikkelig dreies om i det
Videnskabelige --- umulig have noget Begreb. For at
antyde, hvortil her sigtes, sammenligne vi et Par Steder
af det gamle Testamente. I anden Samuels Bog (XXIV)
berettes Følgende: „Og Herrens Vrede blev atter op\-%
tændt mod Israel, og han tilskyndede David imod dem,
idet han sagde: Gaa, tæl Israel og Juda!“ Og Kon\-%
gen gjør udtrykkelig, som Jehova har befalet ham, og
byder Joab, som var Høvding for Hæren, at bereise alle
Israels Stammer fra Dan til Beersaba og tælle Folket.
Ikke desto mindre faaer David, just som Folketællingen
er endt, en Forestilling om, at han har gjort noget Ondt.
„Men“, hedder det, „da slog Davids Hjerte ham, efter
at han havde talt Folket, og David sagde til Herren:
Jeg har høilig syndet i hvad jeg har gjort, men nu,
Herre, tilgiv dog din Tjeners Brøde, thi jeg har handlet
meget daarligt!“ Altsaa har David, ved ligefrem at
handle efter Jehovas udtrykkelige Befaling, begaaet en
Synd; altsaa har Gud udtrykkelig fristet David til det
Onde. Dersom det ikke er det, her siges med rene Ord,
hvad er det da her siges? Den theologiske Fortolkning
beviser saa gjerne sin Videnskabelighed ved at samle
Parallelsteder, og lade „Bibelen fortolke sig selv“ ved
at lade det ene Sted forklare det andet. Parallelstedet
til anden Samuel (XXIV) er, som bekjendt, første
% --- p. 87 ---
\opage{87}\setcounter{footnote}{0}%
Krønikerne (XXI), der begynder saaledes: „Satan stod
Israel imod og tilskyndede David til at tælle Israel, og
David sagde til Joab“ o. s. v. Man vende sig herme\-%
neutisk og vride sig exegetisk, hvordan man vil, saa
staaer der udtrykkelig paa det ene Sted, at Satan gjør
det, hvorom det paa det andet Sted hedder, at Gud selv
gjør det. Og at Gud (And. Sam.) ikke lader Satan
friste David, men gjør det selv umiddelbart, fremlyser
tydelig nok af de utvetydige Ord: Gaa, tæl Israel og
Juda!

Kaste vi et Blik paa den videnskabelig kritiske Op\-%
fattelse af Spaadommene, gaaer det ikke Theologien et
Haar bedre. Fritænkeren kritiserer, og Theologen kri\-%
tiserer, og jo flere Lys Kritiken tænder, desto blegere
bliver Religionen. Man har omhyggelig opregnet og
classificeret de Spaadomme i det gamle Testamente, som
gaae i Opfyldelse i det nye, men sete igjennem Kriti\-%
kens Briller falde de aldeles hen; thi hvad Propheterne
sige, er noget Andet, end hvad Evangelisterne lade dem
sige. Et eneste Exempel! I Matthæi Evangelium (II),
hvor Flugten til Ægypten fortælles, hedder det, at dette
skete, for at det skulde fuldkommes, som er talet af
Herren ved Propheten, som siger: „Jeg kaldte min Søn
ud af Ægypten.“ Det er Hoseas (XI, 1), hvortil her
sigtes. Men hos Hoseas staaer: „Der Israel var ung,
da havde jeg ham kjær; og jeg kaldte min Søn af
Ægypten.“ For Kritiken er her ingen Spaadom; thi
ved den Søn, hvorom Propheten taler, mener han ikke
Jesus, født af Maria, men Israels Folk. Nu faaer man
rigtignok den aandrige Forklaring, at Israels Folk er et
Forbillede paa Christus, men aandrige Omsvøb have
vi nok af. Hvad Kritiken seer ved første Øiekast, er,
at Prophetens Tanke gaaer i en Retning, Evangelistens
i en anden. Propheten taler om et Forbigangent, men
% --- p. 88 ---
\opage{88}\setcounter{footnote}{0}%
Evangelisten opfatter ham, som om han havde talt om
et Tilkommende. Der er i Prophetens Tale ingen Spaa\-%
dom, og dog vil Evangelisten i den see en Spaadoms
Opfyldelse.

Hvad er nu Følgen af disse videnskabelige Forsøg?
% "Rénan": an accent is printed over the e, as transcribed; the name has none. AUDIT (600 dpi): the inked stroke is heavy at top right and tapers down-left, i.e. an acute, not a grave; a faint grey scratch crossing it from upper left is not ink of the type.
Man kjender Navne som Strausz, Rénan og fl. Det er
vitterligt, at Strausz med videnskabelig Grundighed og
Kritik har opløst den hele evangeliske Historie og gjort
Christi Liv paa Jorden til en Mythe. Det er Verdens\-%
aanden, som gjennem ham fuldender sine nedbrydende
Conseqvenser; den gamle Theologie har længe sagt A, og
nu kommer Strausz tilsidst og siger B. Ved Hjælp af den
strausziske Kritik bliver man da ikke alene fri for de
mange videnskabelige Modsigelser, hvormed Religionen
hidtil har været beheftet, men man bliver videnskabeligt
fri for Religionen selv; saa er Kritiken sand, og Sandhe\-%
den kritisk; saa tolke vi Christi Ord i Tidens Aand:
„I skulle erkjende Kritiken, og Kritiken skal frigjøre
Eder.“

Dersom det nu var Opgaven, det Maal, vi her havde
sat os, at opklare Troens Psychologie og udvikle de
religiøse Begreber af Religionens eget Princip, saa maatte
vi vel i nærværende Sammenhæng være forpligtede til at
angive, hvorledes de udpegede Misgreb Punkt for Punkt
lode sig hæve. Men da vi ikke have foresat os at med\-%
dele enten Troespsychologie eller Religionsphilosophie i
udførlige Foredrag, men derimod i al Almindelighed at
tale om Hindringer og Betingelser for det aandelige Liv:
saa vil en Henpegen paa de Mishandlinger, Religionen
maa lide af den ved Theologien indførte Videnskabelig\-%
hed, allerede kunne afgive Bidrag. Det er, efter hvad
vi alt have seet, ingenlunde en umiddelbar Begeistring
for det Ideale, men i det Høieste en intellectuel Stræben
efter Sandheds Erkjendelse, der kan siges at betegne
% --- p. 89 ---
\opage{89}\setcounter{footnote}{0}%
Tidsalderens aandelige Fortrin. Men en Tidsalder, hvis
Hovedstyrke er intellectuel, maa først og fornemmelig
arbeide paa at klare sig ud af intellectuelle Forvirringer.
Vi sige det frit og med fuld Overbeviisning, at den første
og største Hindring for Aandslivet i den nærværende
Tid er en theologisk Forqvakling af de hellige Kjends\-%
gjerninger, en dogmatisk Sammenblanden af religiøse og
videnskabelige Erkjendelser, et kritisk udhulet men offi\-%
cielt afstivet Myndighedsvæsen, der støtter sig til Troen,
naar det skorter paa Viden, og værger sig med Viden,
hvor det glipper for Troen. Vi sige det frit og med
fuld Overbeviisning, at den første og største af vor intel\-%
lectuelle Tids mange intellectuelle Opgaver er, saasandt
den vil bevare og fremme sit aandelige Liv, den Opgave:
at faae Religionen, jeg mener den aabenbarede --- ikke ved
et Orthodoxiens Magtsprog, ikke ved ekstatiske Væk\-%
kelser og mystiske Visioner, men ved Tænkning, en til
% A raised dot after "til" is a stray mark, not transcribed.
Troen svarende, grundig, klar, oplysende Tænkning ---
een Gang for alle udfriet og frelst af Videnskabens
Kløer. Indtil denne Opgave grundig er løst, maae
Hindringer og Betingelser vedblive at brydes i dump
Gjæring; men med Opgavens Løsning vil Sandheden
seire ved Aandens Magt, ikke den upersonlige, men den
personlige; ikke den kolde, der grunder i Nødvendighe\-%
den, men den levende, der lyser i Friheden, just den,
der høres i Christi Ord: „I skulle erkjende Sandheden,
og Sandheden skal frigjøre Eder.“

\begin{center}\rule{3cm}{0.4pt}\end{center}

\chapter*{VI.}
\addcontentsline{toc}{chapter}{Forelæsning VI}
\markboth{Forelæsning VI}{Forelæsning VI}
\label{ch:vi}
% pp. 90--104
% --- p. 90 ---
\opage{90}\setcounter{footnote}{0}%
% [large initial C]
Culturbevidsthedens Strid mod Evangeliets Aand er
allerede charakteriseret fra to Synspunkter, forsaavidt
nemlig Tankebevægelsen gaaer gjennem Naturlæren og
igjennem Historien. Det staaer nu tilbage at opfatte og
charakterisere denne Strid fra Aandslærens eget Stand\-%
punkt. Her er Kamppladsen indvortes; Striden naaer
sit afgjørende Vendepunkt, idet den ender som en Strid
i Bevidstheden selv. Spørge vi Culturmennesket, om
han, i Conseqvens af den negative Verdensaand, vil for\-%
negte det Høieste, Idealet, det Evige, svarer han Nei;
spørge vi ham, om han vil fornegte Guddommen selv,
svarer han Nei; men tilføie vi endnu det Spørgsmaal,
om han da vil troe paa Gud og hans aabenbarede Ord,
svarer han ogsaa Nei. Culturbevidstheden opfatter nemlig
sin Guddom, ikke som almægtig villende Personlighed, men
som Videns absolute Idee. I Troen antages Gud for et
personligt Væsen; men, siger den Vidende, dersom Gud
virkelig skal tænkes at være et personligt Væsen, saa maa
han jo ogsaa tænkes at være et indskrænket Væsen. „I til\-%
lægge Gud“, siger Fichte til de Troende, „Personlighed
og Bevidsthed. Hvad forstaae I da ved Personlighed
og Bevidsthed? Dog vel det, som I have fundet i Eder
selv, og betegnet med dette Navn? Men at I paa ingen
% --- p. 91 ---
\opage{91}\setcounter{footnote}{0}%
Maade tænke eller kunne tænke Eder dette uden Be\-%
grændsning og Endelighed, det kan den ringeste Op\-%
mærksomhed ved Eders Construction af dette Begreb
lære Eder. Ved at tillægge dette Væsen hiint Prædicat
gjøre I det altsaa til et endeligt, til et Væsen, der er
Eder ligt, og I have ikke, som I vilde, tænkt Gud, men
kun mangfoldiggjort Eder selv i Tænkningen.“ Den
absolute Aands Selvbevidsthed er vel, efter Hegel, Reli\-%
gionen; men hvis Nogen heraf slutter, at den speculative
Philosophie virkelig tillægger Gud en evig, i sig bestaa\-%
ende Selvbevidsthed, da vil han komme paa andre Tan\-%
ker, naar han ved at trænge dybere ind i Systemet lærer
at indsee Betydningen af de mærkværdige Ord, at „Men\-%
neskets Viden af Gud er Guds egen Viden af sig selv.“
Fra dette philosophiske Høidepunkt vil da Culturbevidst\-%
heden fortsætte sin intellectuelle Kamp for at udrydde
Troen.

Og hvorfor vil Culturbevidstheden udrydde Troen?
Dersom dens Standpunkt er gyldigt, da er Bestræbelsen
ikke blot conseqvent, men ædel og god. Her er saavel
udadtil som indadtil Aandsphænomener nok, der kunne
forarge den Oplyste. Der er af en Tænker, der ivrede
for Troen, etsteds opkastet det Spørgsmaal: „Hvis det,
vi forstaae ved at være Christen, virkelig er det at være
Christen, hvad er saa Gud i Himlene?“ Og Svaret gives:
„Han er det latterligste Væsen, der nogensinde har levet,
hans Ord den latterligste Bog, der nogensinde er kom\-%
men for Lyset: at sætte (som han jo gjør i sit Ord)
Himmel og Jord i Bevægelse, at true med Helvede, med
evige Straffe, for at opnaae det, vi forstaae ved at være
Christne (og vi ere jo sande Christne), nei noget saa
Latterligt er aldrig forekommet.“ Den skærende Mod\-%
sigelse imellem Troesidealerne og det virkelige Liv, der
charakteriserer Midlernes, Betingelsernes, Interessernes
% --- p. 92 ---
\opage{92}\setcounter{footnote}{0}%
Tidsalder, staaer ligesaa klar for den vidende Kritiker
som for den troende Tænker. Naar den troende Tænker,
saaret af Modsigelsens Braad, udbryder: „Den rædsomste
Art Gudsbespottelse er den, Christenheden forskylder:
at forvandle Aandens Gud til et latterligt Vrøvl; og den
aandløseste Art Gudsdyrkelse, aandløsere end Alt hvad
der findes og fandtes i Hedenskabet, aandløsere end det:
som Gud at tilbede en Steen, en Oxe, et Insekt, aand\-%
løsere end Alt er: at tilbede under Navn af Gud --- et
Vrøvlehoved!“ naar, sige vi, slige Anskrig lyde fra de
Troendes Leir, da svarer den negative Kritik med Ja
og Amen.

Den Tidens Kjendsgjerning, at Troen har tabt sin
begeistrende Kraft, opfatter Culturmennesket ikke som
Tegn paa et aandeligt Tilbageskridt, men tvertimod som
Beviis paa et væsentligt Fremskridt. Grunden til at
Troens Kraft aftager, er ifølge hans Overbeviisning den,
at Videns Lys tiltager. Det er, siger han, ikke blot
Troens Gjenstande, de hellige Kjendsgjerninger, der op\-%
løse sig i det Fabelagtige og tabe sig i det Mythiske;
nei, Troen selv er af fabelagtig Natur, den troende Be\-%
vidsthed er væsentlig mythisk. Vi skulle nu høre den
nærmere Forklaring.

Hvorledes den mythiske Bevidsthed kommer til at
opfatte de forskjellige Guddomme, sees hos Homer.
Allerede i Iliadens første Sang, hvor Achilles og Aga\-%
memnon ere i den bekjendte Strid med hinanden, ind\-%
træder et for en Gude-Aabenbarelse betegnende Vende\-%
punkt. Achills Lidenskab er opflammet; han er saa
rasende, at han kalder Agamemnon feig og bruger de
stærkeste Haansord; der er alle Tegn til, at han vil for\-%
løbe sig. „Under det laadene Bryst tvivlraadigen banked
hans Hjerte, enten han skulde fra Lænd flux drage sit
hvæssede Glavind, splitte den stimlende Hob, og hugge
% --- p. 93 ---
\opage{93}\setcounter{footnote}{0}%
ham ned den Atreide, eller betvinge sit hidsige Mod og
dæmpe sin Harme“; da skeer Aabenbaringen. „Just
som Sværdet af Balgen han trak, kom Pallas Athene
ned fra Himlen... Bag Achilles hun treen, og trak i
hans lysgule Lokker, synlig alene for ham, af de Øvrige
saae hende Ingen.“ Den Samtale, som nu føres mellem
Achilles og Gudinden, er høist charakteristisk. „Datter
af Aigissvingeren Zeus, hvorfor er du kommen?“ Hvortil
Athene svarer: „Kommen jeg er for at dæmpe din Harm,
ifald du vil lyde“.

Hvor er ethvert Træk i denne Aabenbarelse ikke
charakteristisk! Det, Achilles i det kritiske Øieblik saa
høilig trænger til, er Besindighed, og see, saa staaer Be\-%
sindigheden selv anskuelig for ham i Pallas Athene. Og han
kjender hende, ja han kjender hende strax, hendes „ræd\-%
somme Blik ham funkler imøde“. Besindighed er en
personlig Sag: saa er Gudinden kun synlig for ham, hun
har kun Ærende til ham alene. Og dog overrasker hun
ham og bringer ham til at studse! Ja Besindigheden kom\-%
mer med guddommelig Magt, med en over Lidenskaberne
ophøiet Ro, med Myndighed, med ideal Overlegenhed.
Athene er ingen Allegorie; hun er en Datter af Aigis\-%
svingeren Zeus: hun er en Gudinde. Her have vi da
den mythiske Fordobling; med andre Ord: her have vi
„Troen“. Achill troer paa Gudinden, og just derfor aaben\-%
barer hun sig for ham. Uden at han selv veed deraf,
henfører han den i hans eget Sind opvaagnende Besin\-%
dighed til et guddommeligt, udenfor ham existerende
Væsen, og saa kommer dette Væsen „ned fra høien
Olymp“ og bringer ham, hvad der oprindelig er hans
Eget. Her er Bevidsthedsfordoblingen, den Fordobling,
som charakteriserer Troen. Den troende Bevidsthed
stiller Guderne op udenfor sig, og gjør det af indre Nød\-%
vendighed, fordi den i sin dunkle Dybsindighed endnu
% --- p. 94 ---
\opage{94}\setcounter{footnote}{0}%
ikke veed af, at den har det Guddommelige i sig selv.
Det er just Sindets Dunkelhed med den til Dunkelheden
svarende Tunghed og Spænding, der gjør, at den Troende
faaer Aabenbaringen. Den troende Bevidsthed lever i
Følelse og Phantasie-Anskuelse; derfor har den en saa
uimodstaaelig Trang til at fremstille sig sit guddomme\-%
lige Indhold i tilsvarende Billeder; naar nu et af disse
Billeder i slaaende Umiddelbarhed og som med en magisk
Kraft paatrænger sig Bevidstheden, saa bilder Mennesket
sig ind, at det Guddommelige, han seer i sit eget Syn,
er en virkelig Guds Aabenbarelse: paa denne Indbild\-%
ning beroer Troen. Alle Religioner udgive sig oprindelig
for aabenbarede; dette er fra Begyndelsen af ikke et
simpelt Præstebedrag; langt fra! det er en sjælelig, idee\-% doubtful: line-end hyphen read as a break (ideenødvendig); a compound hyphen (idee-nødvendig) cannot be excluded from the image
nødvendig Selvskuffelse; den, der har seet det primitive
Gudesyn, troer virkelig fuldt og fast, at han har faaet
en Aabenbaring. Saa besinder han sig paa sit Syn; han
erindrer sig det, og i hans Erindring bliver det bestandig
tydeligere, rigere, vidunderligere; saa beskriver han sit
Syn og fortæller om Aabenbaringen, og de, som med
Andagt høre Fortællingen, gribes af Forundring, for\-%
nemme i sig et Vidnesbyrd om, at det, der fortælles dem,
maa være guddommelig Sandhed. Disse Andægtige blive
nu de egentligt Troende, idet de, uden selv at have faaet
nogen Aabenbaring --- være sig, at Oprindeligheden i
deres Væsen er for svag; eller at de leve paa en senere
Tid, efterat de primitive Gjennembrud alt ere skete ---
troe paa den Aabenbaring, der er bleven en Anden, en
mere Værdig, tildeel. I de egentlige Troende arbeider
saa Erindringen videre, Gudesagnet (Μῦθος) udvides, % Greek set in sloping Greek type, printed with the curly theta (ϑ) and final ς; typed from the image at 300 dpi
ændres, fuldstændiggjøres; Ideen, der virker i Phanta\-%
siens Form, gjør Mythen snart til et Emne for Poesien;
nu staae vi igjen ved Homer.

For at bevise, at der gives en Aabenbaringsbe\-%
% --- p. 95 ---
\opage{95}\setcounter{footnote}{0}%
vidsthed, der ikke er mythisk, beraaber man sig paa det
sædelig Rene, det sande Ethiske, paa Nødvendigheden
af at grunde det Ethiske paa det Religiøse. Men hvor
utilstrækkelig en saadan Beviisførelse er, kan udtryk\-%
kelig sees af vort anførte Exempel. Eller kan man vel
tænke sig nogen renere Sammensmeltning af det Religiøse
og det Ethiske end netop i de Repliker, her vexles mel\-%
lem Achill og Athene? „Hvorfor est du kommen?“ ...
„Kommen jeg er for at dæmpe din Harm, ifald du vil
lyde, kommen fra Himlen, mig sendte den liliearmede
Here, ligelig elsker hun Eder, og ømt for begge hun
sørger.“ Saa ophøiet er her den religiøs-ethiske Tanke,
at Achill i sin brusende Vrede maa opløftes over sig
selv og mærke sig, at der i Himlen er en guddommelig
Kjærlighed, der ligelig omfatter ham og hans Fjende.
Og denne guddommelige Kjærlighed, hvor er den dog
menneskelig, hvor er den human! Gudinden med „det
rædsomt funklende Blik“ truer ham ikke, tvinger ham
ikke; men giver ham et Raad, ifald han vil lyde. Ly\-%
dighedsfordringen er da heller ikke overdreven streng.
Han skal blot afholde sig fra at bruge Magt, men tør
forresten give sin Harme Luft og sætte Agamemnon
irette med Ord, saameget han lyster. Ved at lyde Gud\-%
indens Raad viser Achilles sig her som et Mønster paa
græsk Fromhed og sand Religiøsitet. „Skjelligt, Gud\-%
inde det er, Eders Ord at tage til Hjerte, om man i
Sind er aldrig saa vred, thi tjenligst er dette; hvo,
som er lydig mod Gudernes Bud, ham høre de gjerne.“

Dette var saa Culturmenneskets Opfattelse af Tro
og Mythologie paa det græske Bevidsthedstrin; hvor\-%
dan faaer han det nu til at see ud med Jehovalæren,
Aabenbaringen i det gamle Testamente? I Grunden paa
selvsamme Maade.

Det er et træffende Ord af en moderne Forfatter,
% --- p. 96 ---
\opage{96}\setcounter{footnote}{0}%
at Mennesket plyndrer sig selv for at udstyre sin Gud.
Dette gjælder fuldt saavel i Jødedommen som i Heden\-%
skabet. Det Bedste hos Grækeren var hans skjønne In\-%
dividualitet. I denne Individualitet var der noget For\-%
krænkeligt, det kunde han beholde for sig selv; men In\-%
dividualitetens uforkrænkelige Væsen, det Guddommelige,
det turde han ikke beholde for sig selv, det tilhørte
alene de salige Guder. Det Bedste hos Jøden er hans
personlig tænkende Jeg, hans ethisk villende Selv. I
dette Jeg er der noget Indskrænket, Uklart og Forgæn\-%
geligt, det tør Mennesket beholde for sig selv; men i
Jeget som Jeg, i Aanden som Aand er der et Krav paa
hellig Viisdom og evig Magt, som Mennesket ikke kan
fyldestgjøre; den uendelige Viisdom, den hellige Villie,
den evige Magt tilkommer alene Guden. Ud fra den
brændende Tornebusk taler Jehova (Anden Moseb. III)
til Moses, og giver sig selv det betegnende Navn: \emph{Jeg
er den, som jeg er}. I dette almægtige Jehova-Jeg
anskuer nu Israels Folk sit aandelige Væsens evige Ideal
paa tilsvarende Maade, som Grækeren anskuer Indivi\-%
dualitetens Ideal i de olympiske Guder.

Vel er her en Forskjel, ja, naar vi see nærmere til,
endog en meget væsentlig Forskjel; men Forskjellen
ligger, videnskabelig seet, ingenlunde deri, at Jehova er
virkelig, medens de græske Guder ere uvirkelige: For\-%
skjellen imellem „Afguderne“ og „den eneste sande Gud“ er
kun en Troesforskjel, en Forskjel, der sættes af den jø\-%
diske Indbildning. Nei, den egentlige Forskjel er den,
der følger af Idealernes Forskjellighed. De græske
Guder ere mange; thi Individualiteten har mange Skik\-%
kelser; Jødernes Gud er kun een; thi Jeget som Jeg er
trods alle Verdensforandringer det samme. De græske Guder
ere, om end for et dødeligt Øie usynlige, dog efter deres
Væsen synlige; de ere Billedguder, sande Kunstguder;
% --- p. 97 ---
\opage{97}\setcounter{footnote}{0}%
thi i Kunstens Form forklares det Individuelle. Jehova
derimod, om han end nu og da paatager sig synlig
Skikkelse, er efter sit Væsen usynlig; thi Tanken som
Tanke, Villien som Villie, Aanden som Aand er i al
Evighed usynlig. Men ligesaa lidt som den troende
Jøde veed af, at han i den ubilledlige Jehova anskuer
Idealet af sit eget aandelige Selv, ligesaa lidt veed den
troende Græker af, at han i sine Billedguder anskuer
Idealet af sin egen menneskelige Individualitet. Skal
den jødiske Gudserkjendelse kaldes Tro, da er den
græske ogsaa Tro; skal den græske Tro kaldes mythisk,
da er den jødiske Tro ogsaa mythisk. Seet i Troen, er
det jo rigtignok Jehova, der har skabt Mennesket i sit
Billede; men, seet i Viden, er det Jøden selv, der har
skabt Jehova i sit Billede. Almagtens, Viisdommens,
Godhedens, Retfærdighedens hellige Gud er en sand
Menneskegud, Verdens Skaber en Skabning af Menne\-%
skets Aand.“ % printer's defect: a closing quotation mark with no matching opening mark (the citation of "en moderne Forfatter" opens unmarked on p. 96); transcribed as printed; quote balance -1

Det er ifølge denne Synsmaade saa langt fra, at
den jødiske Tro skulde være mindre mythisk end den
græske, at det Mythiske i Jehovalæren, endog i mange
paafaldende Træk, træder stærkere frem; det kommer
deraf, at Aabenbaringsphænomenerne i den jødiske My\-%
thologie staae i en langt mere skærende Strid med
Gudsbegrebet selv, end i den græske.

Gud er almægtig; den almægtige Gud er Himmelens
og Jordens Skaber; nu kommer det Mythiske. Thi den
Almægtige skaber ikke som en Almægtig; nei, han gaaer
stykkeviis tilværks; han bærer sig ad som en menneske\-%
lig Værkmester, omtrent som Platos Gud i \emph{Timæus}.
Først taler han med sig selv og overveier, hvad han
vil gjøre; derpaa gjør han det, og saa seer han bagefter
paa det, han har gjort, og seer, at det er godt gjort.
At en saadan Skabelse maa være Gjenstand for Tro,
% --- p. 98 ---
\opage{98}\setcounter{footnote}{0}%
forstaaer sig af sig selv; thi saasnart den gjøres til Gjen\-%
stand for Viden og sees i Begrebet, forraader den sin
Modsigelse og forsvinder. Er Jehova almægtig, saa er
hans Magt absolut; men den absolute Magt kan ligesaa
lidt indlade sig paa endelige Handlinger som paa ende\-%
lige Beslutninger. Der siges, at Jehova vandrede i
Paradisets Have, der Dagen var luftig. Hvor gemyt\-%
ligt! Den almægtige Himlens og Jordens Skaber gjør
sig en Spadseretour i Aftensvalen under Skyggen af de
Træer, som han selv nylig har skabt. At en Zeus en\-%
gang imellem gjør sig en Spadseretour i Lunden, er
dog paa en Maade stemmende med hans erotiske Indi\-%
vidualitet; men at en Allestedsnærværende spadserer, er
en uhyre Modsigelse. Hvorlænge kan Sligt vedblive
at være Gjenstand for Tro? Svar: saalænge Bevidstheden
kan vedblive at være mythisk. At Jødefolket, dette saa
eiendommelige, aandskraftige Oldtidsfolk, har havt en
religiøs Barndom, et Bevidsthedstrin, da slige ved Tro\-%
skyldighed og barnlig Ynde tiltrækkende Fortællinger frem\-%
kom med psychologisk Nødvendighed og just derfor vir\-%
kede med psychologisk Sandhed, er begribeligt. Der er
i disse „det udvalgte Folks“ religiøse Ammestuehistorier
noget saa Dristigt og storslaaet Enfoldigt, at det maa
imponere Phantasien endnu den Dag idag. „Og Gud
Jehova gjorde Adam og hans Hustru Kjortler af Skind“;
man tænke sig: han syede selv deres Klæder!

Hvad der for Oldtidens barnlige Anskuelse har havt
sin psychologiske Sandhed, bliver i vor Tid til psycho\-%
logisk Usandhed, til Affectation og Selvbedrag, naar det,
til Trods for Videnskab og sund Fornuft, under Navn
af „Barnetro“ endnu skal tages ligefrem for Aabenbaring.
Vi ville ikke opholde os ved de mange Overraskelser
og heftige Sindsbevægelser, hvoraf Jehova, forglemmende
sin Alvidenhed og Almagt, lader sig henrive --- naar
% --- p. 99 ---
\opage{99}\setcounter{footnote}{0}%
han f. Ex. (Første Moseb. VI) i den Grad overraskes
af Menneskenes Ondskab, at „det angrede ham, at han
havde dannet Mennesket paa Jorden, og det fortrød ham
i hans Hjerte“ --- vi ville, som sagt, ikke opholde os
ved, hvad der er Alle bekjendt; kun som et Exempel % a dot (speck), not an accent, stands over the e of "ved"
paa hvorvidt den psychologiske Usandhed i vore Dage
kan gaae, anføre vi endnu et eneste Træk. Der for\-%
tælles (Første Moseb. VII), at da Noah med alle Dyrene
var vel kommen ind i Arken, saa „lukkede Jehova selv
efter ham“. Hertil har nu En af Nutidens Barnetroende
tilføiet følgende ligesaa skriftkloge som opbyggelige For\-%
tolkning. „Jehova \emph{lukkede}, thi udenfra kunde Noah
jo ikke mere tilslutte den (Arken), saa at Vandet ikke
kunde trænge ind; ogsaa maatte han jo være sikkret
mod de Ugudeliges Forsøg paa at kunne bryde Arken
op for at komme ind.“ Ja, Nutidens Tro --- er Barnetro!

See vi tilbage paa det her Anførte, hvor skulle vi
da kunne tvivle om, at den jødiske Tro er ligesaa my\-%
thisk som den græske. Eller hvad skulde vel berettige
os til at erklære Athene for en falsk Gudinde, naar vi
anerkjende Jehova for en sand Gud. At hendes Til\-%
værelse staaer i Strid med Naturlovene? Men det gjør
Jehovas Tilværelse jo ogsaa! At hun hos Homer stun\-%
dom lader sig henrive af Lidenskaben? Men det gjør
Jehova hos Moses jo ogsaa! At Jehovas Grundcha\-%
rakteer, naar vi see bort fra hans enkelte Svagheder, er
ethisk? Men Athenes Grundcharakteer er, naar vi see
bort fra hendes enkelte Svagheder, jo ogsaa ethisk!
Hendes Aabenbarelse for Achilles er da idetmindste
ligesaa ethisk som Jehovas Anger over at have skabt
Menneskene. Nei, skal Jehova være en sand Gud, da
lad Athene ogsaa være en sand Gudinde. Vi behøve
da heller ikke for Troens Skyld at frygte for nogen
logisk Modsigelse. Der er, naar vi blot ville troe paa
% --- p. 100 ---
\opage{100}\setcounter{footnote}{0}%
Aabenbaringer og Aabenbarelser, aldeles Intet iveien
for, at Athene og Jehova kunne være to forskjellige
Aabenbarelser af den ene og samme „Himlens og
Jordens almægtige Skaber“. Vi vide jo dog --- vi vide
det i Troen --- at Jehova selv tilligemed to Engle har
i Mamre Lund (Første Moseb. XXIII) spiist tyk Melk % "XXIII" as printed (the Mamre meal is Gen. XVIII)
og sød Melk og Kalvekjød hos Abraham under Træet;
kan den Almægtige nedlade sig til at spise sød Melk og
Kalvekjød, naar han vil aabenbare sig for en Abraham,
hvorfor skulde den samme Almægtige da ikke ogsaa
kunne nedlade sig til at optræde som en Datter af
Aigissvingeren Zeus, naar han vil aabenbare sig for en
Achilles? --- Saavidt Culturmennesket.

At der i den moderne Videns ofte skarpsindige
Angreb paa Aabenbaringstroen ligger en væsentlig Hin\-%
dring for det religiøse Livs Udvikling, skal ikke negtes;
men see vi nærmere til, saa turde det dog vise sig,
at netop denne Hindring, saafremt det kunde blive Alvor
med Løsningen af Nutidens intellectuelle Grundopgave,
lader sig forvandle til en væsentlig Betingelse. Jo skar\-%
pere og mere indtrængende den Tænkning er, der an\-%
griber Troen, desto skarpere og mere indtrængende maa
ogsaa den Tænkning være, hvorved Troen værger sig
mod Angrebet. Skarpt imod Skarpt, saa kommer det
vel til en Afgjørelse! Her er da i den triumpherende
Videnskabs Angreb for det Første en Vilkaarlighed, der
paa det Eftertrykkeligste maa tilbagevises. Her forud\-%
sættes uden videre, at Troen i sig selv er tankeløs, at
al klar, conseqvent og grundig Tænkning udelukkende
tilhører Videnskaben: at anlægge Tænkningens Maale\-%
stok er at anlægge Videnskabens Maalestok. Denne
Forudsætning er grundfalsk. Støttende sig til en grund\-%
falsk Forudsætning begynder den videnskabelige Kritik i
al Ugenerthed med at behandle Troesbegreberne som
% --- p. 101 ---
\opage{101}\setcounter{footnote}{0}%
ufuldbaarne, forqvaklede Vidensbegreber, og kan da
sagtens bevise, at disse Begreber indeholde videnskabe\-%
lige Modsigelser. Med den falske Forudsætning, at Troen
er tankeløs, finder Kritiken det naturligviis under sin Vær\-%
dighed at sætte sig tilbørlig ind i den troende Bevidstheds
eiendommelige Væsen og Selvudvikling, paadutter den
uden Omstændigheder Forpligtelsen til at være vidende
Bevidsthed, og kan da, efterdi Troen som Tro er ab\-%
solut ueensartet med Viden, sagtens bevise, at den tro\-%
ende Bevidsthed, opfattet som vidende, er i Modsigelse
med sig selv. Med den grundfalske Forudsætning, at den
troende Bevidsthed er i oprindelig Modsigelse med sig
selv, er det naturligviis Videnskritiken en let Sag at kari\-%
kere de hellige Kjendsgjerninger og faae Aabenbaringen
slaaet sammen med Mythologien.

Vi ville nu see, om det dog ikke er muligt at oprede
dette Virvar. Det er, som sagt, en grundfalsk Forud\-% the r of "grundfalsk" prints faint/broken (layer: "giundfalsk"); read as r
sætning, at Troen som Tro er tankeløs; vi bevise det
ved at fremstille Troesprincipet som et i sig eiendom\-%
meligt Tankeprincip. At Troesprincipet selv er Tanke\-%
princip, betyder ikke, at Troen begynder med den rene
Tænken af den rene Væren; men det betyder, at den
troende Bevidsthed formaaer at tænke sin egen oprinde\-%
lige Grundhandling, ved denne Tænken at adskille sig
fra den vidende Bevidsthed og ved denne Adskillelse
blive sig selv gjennemsigtig. Udgangspunktet er i det
Godes Idee. Det Godes Idee er Villiens Idee: til det
absolute Gode svarer den absolute Villie, det absolute
Selv. Det absolute Selv er den hellige Gud. Lige\-%
overfor det absolute Selv, staaer saa, med Verden til
Mellemværende, det relative Selv, det Selv, der ikke er
godt, men forpligtet til det Gode, den relative Villie, der
ikke er hellig, men syndig. Med disse Forudsætninger
kan den troende Bevidsthed ikke blot blive sig selv
% --- p. 102 ---
\opage{102}\setcounter{footnote}{0}%
klar, forsaavidt den kan indsee deres absolute Uundvær\-%
lighed for Selvet, men den kan endogsaa gjennemtænke
sit Væsen saa fuldstændigt, at enhver af Forudsætnin\-%
gerne forvandler sig til Moment i Troens Begreb.

Er Troesprincipet først i Troens Begreb opfattet
og bestemt som Tankeprincip, da har det ingen Fare
med Mythefortolkningen og det Mythiske. At opløse
Naturreligionernes Gudebilleder i personificerede Ideer % a faint speck above the first e of "Naturreligionernes" (layer: "Naturréligionernes"); not a clear accent, not transcribed
og paavise de ubevidste Illusioners psychologiske Op\-%
rindelse, er en værdig Opgave for Viden og videnskabelig
Kritik; ved Videns Fremskridt gjør Selverkjendelsen i
denne Henseende da ogsaa et virkeligt Fremskridt: det ari\-%
stoteliske Gudsbegreb er i Sandhed langt at foretrække
for homeriske Gudebilleder. Men ansporet ved dette
Fremskridt er „den moderne Videnskab“, særlig fristet
og forført af Theologien, i sørgelig Maade kommen
til at overskride sin absolute Grændse. Er Zeus ---
saaledes slutter man --- en Personification af Græker\-%
aandens Individualitet, da maa Jehova nødvendigviis
ogsaa være en Personification af Jøde-Aandens stærke,
selvraadige Jeg. Denne Slutning er grundfalsk. Den,
der siger, at Jehova kun er en Personification af Selv\-%
hedens menneskelige Væsen, hvad gjør han? Han gjør
det menneskelige Selv til en Personification af den sub\-%
jective Tænkning, og den subjective Tænkning til en tom
Reflex af organiske Naturkræfter. Ved en saadan Frem\-%
gangsmaade gaaer det ingenlunde opad, men tvertimod
nedad med Selvet og dets Selverkjendelse. „Det gaaer
tilbage, Hakon“; ja, det gaaer tilbage! Thi saamegen Re\-%
flexion maa en reflecterende Tidsalder dog vel kunne opdrive,
at Enhver, som alvorlig vil, ogsaa saare let kan --- uden
speculativ Overanstrengelse --- indsee, at det menneske\-%
lige Selv maa staae og falde med det guddommelige, det
% --- p. 103 ---
\opage{103}\setcounter{footnote}{0}%
afmægtige med det almægtige, det relative med det ab\-%
solute.

Det skal villig indrømmes, at Aabenbaringsmirak\-%
lerne i det gamle og det nye Testamente have en umis\-%
kjendelig Lighed med Naturreligionernes Mythemirakler,
naar de sees og bedømmes fra Videns Standpunkt. Ved
at gjøre Bibelen til et stokdogmatisk Aabenbarings\-%
dokument har Theologien ægget den videnskabelige
Kritik til en, vistnok ofte saarende, men i sine Grund\-%
træk seierrig, uigjendrivelig Paaviisning af denne saa
iøinefaldende Lighed. Videnserkjendelsen er, og maa
ifølge sin eiendommelige Charakteer nødvendig være, en
fornuftig sammenhængende Verdenserkjendelse: fra Ver\-%
denserkjendelsens Synspunkt er et mosaisk Mirakel lige\-%
saa anstødeligt, usandsynligt, urimeligt og selvmodsigende
som et homerisk. Men Mirakelbegrebet er et Troes\-%
begreb, et ubetinget Selvhedsbegreb, og kan aldrig blive
Vidensbegreb; de Modsigelser, „den moderne Kritik“
saa let opdager i Miraklerne, lade sig alle henføre til
en ukritisk Sammenblanding af Troesbegreber og Videns\-%
begreber.

Spørgsmaalet om Forskjellen mellem falske og sande
Mirakler falder oprindelig sammen med Spørgsmaalet
om Forskjellen mellem falsk og sand Aabenbaring,
mellem falsk og sand Tro. Hermed er Mytheprincipet
een Gang for alle udelukket, ikke blot fra det nye, men
ogsaa fra det gamle Testamente. Vel kan der, navnlig
i det gamle Testamentes hellige Sagn, være Tale om
Analogier til det Mythiske, men, ifølge Troesprincipet,
ingenlunde om egentlige Myther. At udvikle Forskjellen
imellem Myther og mythiske Analogier i det Enkelte, er
en særlig Opgave for én religionsphilosophisk Behandling % AUDIT doubtful: a faint short stroke over the e of "én"; could be an accent or a speck like the one over "Naturreligionernes" (p. 102), which is not transcribed
af den hellige Historie; i nærværende Sammenhæng skal
det kun i Almindelighed fremhæves, at Forskjellen maa
% --- p. 104 ---
\opage{104}\setcounter{footnote}{0}%
bestemmes principielt ved en tydelig og skarp Udvikling
af Forholdet imellem Menneskeordet og det guddomme\-%
lige Grundord. Bibelordet har, som det foreligger, saa\-%
vel sin menneskelige som sin guddommelige Side, er
ikke et Gudsord ved Siden af et Menneskeord, men fra
Først til Sidst et gudmenneskeligt Ord, opfattet i Tro,
fremstillet i Tro, overleveret i Tro. Den religiøse Tro
er fra Begyndelsen af en Folketro, det religiøse Ord et
folkeligt Ord. Istedenfor at gjøre Bibelen til et ordret
dicteret Aabenbaringsdokument, maae vi tvertimod er\-%
indre, at Bibelen er en Folkebog ligesaavel som Homer.
Men er Bibelen en Folkebog, hvad Under da, at den
udtrykker sig folkeligt i folkelige Billeder, Billeder, som
fremkom, da Aabenbaringens hellige Lys brød sine
Straaler i Folke-Aandens naturlige Media!

Med et forandret Princip er Alt forandret. Tør
det betragtes som den væsentligste Hindring for Nu\-%
tidens aandelige Liv, at Videnskaben har taget Begrebet
fra Troen og gjort Aabenbaringen tankeløs: da vil En\-%
hver jo nu kunne sige sig selv, hvad Fordringen er, hvis
Aanden skal forvandle en saadan Hindring til Betingelse.

\begin{center}\rule{3cm}{0.4pt}\end{center}

\chapter*{VII.}
\addcontentsline{toc}{chapter}{Forelæsning VII}
\markboth{Forelæsning VII}{Forelæsning VII}
\label{ch:vii}
% pp. 105--118
% --- p. 105 ---
\opage{105}\setcounter{footnote}{0}%
% [large initial D]
Dersom Nogen optræder med det Oplysningens Evange\-%
lium, at Ingen kan troe paa det Hellige, eftersom det
Hellige netop er det, som ikke er, og Ingen kan troe paa
det, som ikke er, da vil han i høieste Grad overraske
det sædelige Samfund. Men faaer han blot Lov at tale
ud, vil hans Livsbetragtning fra mere end een Side vise
sig som et Produkt af en i Nutiden herskende Tanke\-%
gang. „Det Hellige“, siger han, „er et Begreb, der for
en klaret Reflexion opløser sig selv. Hvorfor? Fordi
det Hellige er en dunkel Magt, der altid viger for Dagens
Lys, en mystisk Forestilling, en Phantasiens Gaade, der
vækker Angst. Betragt engang de hellige Steder, Gu\-%
dernes Altere, de hellige Lunde! Seer Du dem ved
Dagens Lys, i menneskeligt Selskab, mærker Du Intet
til Frygt eller Gysen, det vil sige: Du kommer ikke i
Berørelse med det Hellige. Men besøg disse Steder i
Nattens Dunkelhed, alene, ved Maanens Skin, og Du
vil gyse; denne Gysen, vakt ved en dunkel Forestilling
om en dunkel Magt, er en Fornemmen af det Hellige.
Det Hellige trænger til Forhæng; dets Væsen er, som
den tilslørede Gudinde i Sais, en Phantasiens Gaade.
Men den kolde Forstand drager Forhænget tilside; en
dristig Tanke løfter Sløret, og Ynglingen forvinder sin
% --- p. 106 ---
\opage{106}\setcounter{footnote}{0}%
Skræk, og det Hellige forsvinder. Vi Oplysningens ægte
Børn, Revolutionens fribaarne Sønner have lært at agte
de borgerlige Love, de politiske Institutioner, Menneske\-%
nes Færd og deres Idrætter; men, man tilgive os et
dristigt Ord, vi troe ikke paa Syner og Aabenbaringer,
vi bøie ikke Knæ for det Hellige. Det Hellige vil, ifølge
sin mystiske Natur, altid gjælde for et Ubetinget; men
i Betingelsernes, Midlernes, Interessernes Verden er Alt
betinget. At Noget er mig helligt, vil kun sige, at
Noget under visse Betingelser er mig dyrebart, maaskee
endog saa dyrebart, at jeg kunde offre Livet derfor.
Men i Betydning af et skjult Guddommeligt, et mystisk
Ubetinget, erkjender jeg intet Helligt, ingen Følelse, intet
Løfte, ingen Ed, intet Alter, ingen Gud; --- jeg siger,
hvad jeg mener: Det Hellige er det, som ikke er, og
Ingen kan troe paa det, som ikke er.“

Om et saadant Menneske vilde nu hele Samfundet
fælde den strengeste Dom; det vilde forskyde ham og
see i ham en ryggesløs og fræk Person. Men, maae vi
spørge, hvorfor? Er det, fordi han siger noget Saadant,
eller fordi han tænker det, eller for begge Dele? Nær\-%
mest kun for det Første. Det dannede Samfund bryder
Staven over den, der udtaler sig saa aabent og hensyns\-%
løst, thi denne Hensynsløshed er netop en Krænkelse af
Samfundet. Seer man derimod paa det Indre og spør\-%
ger, om man da kan ansee et Menneske med en saadan
Tænkemaade for agtværdig, da vil Samfundet udentvivl
dele sig i to Partier. Det ene Parti vil holde sig til
„det slette Hjerte“, det andet til „det gode Hoved“; ja,
det var vel ikke umuligt, at der imellem dem, der gik
til Yderligheder, kunde opstaae et Mellemparti, som
paatog sig at bevise, hvor vanskeligt det er at forene
et godt Hjerte med et godt Hoved, at de interessanteste
Mennesker som oftest ere mistænkte for at have en lille
% --- p. 107 ---
\opage{107}\setcounter{footnote}{0}%
Feil i Hjertet, medens de elskværdige undertiden be\-%
skyldes for at have den i Hovedet. Men med alt dette
vil det hele Samfund let blive enigt om, at det Forar\-%
gelige ligger i Udtalelsen; medens de Stridende vanske\-%
ligt ville blive enige om, at den væsentlige Frækhed
ligger i Tanken. Og dog er Forskjellen mellem den,
der tænker i Stilhed, og den, der udtaler Tanken, kun
den, at den Talende er oprigtig, idet han er fræk, den
Tause derimod er fræk uden at være oprigtig, altsaa en
Hykler, og som Hykler maaskee langt farligere for Sam\-%
fundet. Religiøst Hykleri i Ordets sædvanlige Forstand
er ingenlunde Nutidens Feil; tvertimod, det er et For\-%
trin, som stigende Oplysning og større Frihed ogsaa i
alle aandelige Anliggender har ført med sig, at hyklet
Fromhed foragtes som Tegn paa en uædel, svag og ussel
Charakteer. Men der gives med Hensyn paa religiøse
og kirkelige Spørgsmaal en politisk Accomodation, der
i sine Følger neppe kan sige sig fri for at begunstige et
vist conventionelt Hykleri. I katholske Lande, hvor
Culturudviklingen har antaget, ikke blot et antikirkeligt,
men endog et antireligiøst Præg, yttrer sig en Stræben ---
ikke udtrykkelig efter at udligne den principielle Strid
imellem det Kirkelige og det Politiske ved at føre den
tilbage til de første Grundspørgsmaal, begynde den reli\-%
giøse Folkeunderviisning forfra og af al Magt værne om
den Aandsfrihed, som dertil udfordres --- men meget mere
efter at udvide, fremme og befæste de reent politiske
Interesser, som om borgerlig Frihed og sand Cultur
virkelig kunde gaae Arm i Arm med religiøs Ufrihed
og middelalderlig Mystik. En saadan Accomodation, der
bringer de Oplyste til at hykle en vis Ærefrygt for, hvad
de i deres stille Tanker maae henregne til Overtro og
Obscurantisme, foraarsager ikke blot endeløse Bryderier
% --- p. 108 ---
\opage{108}\setcounter{footnote}{0}%
i det Udvortes, men, hvad værre er, den udhuler Livs\-%
grunden i det Indvortes.

At Stater, der have gjennemgaaet en politisk Revo\-%
lution uden forud at være blevne forberedte ved nogen
egentlig Reformation af det middelalderlige Kirkevæsen,
maae føle den ulmende Spænding som et trykkende
Fatum, er begribeligt; maaskee kan det store Tryk, naar
Aandens Time slaaer, dog komme til at afgive Betingelser
for en aandelig Reisning, hvorom vi endnu ikke kunne
gjøre os nogen Forestilling. Men i protestantiske Lande,
hvor det religiøse Frihedsprincip forlængst er anerkjendt ---
idetmindste pro forma --- at ville overkalke Charaktererne
i de kirkelige og politiske Skjoldmærker og hutle sig
igjennem ved Accomodation er en med Reformationens
Aand kun lidet stemmende Fremgangsmaade. Det Hel\-%
lige kræver til sin Anerkjendelse saa fuldt som det Pro\-%
fane en grundig fast Overbeviisning; men til en grundig,
fast Overbeviisning kan en reflecterende Tidsalder umulig
naae uden grundig, klar og gjennemført Tænkning.

Men dette er, hvis vi ikke feile, netop en af de
væsentlige Hindringer for det aandelige Liv i Nutiden,
at den religiøse Tænkning hos det store Fleertal af Dan\-%
nede øiensynlig er standset. Man tænker, og det med
en vis Udholdenhed, over mangfoldige Anliggender.
Man tænker over Politik og Industri, over Kunst og
% "Poesi": a faint smudge over the e, not a clear diaeresis; read as Poesi.
% AUDIT doubtful: at 600 dpi a dot is visible over the e (Poësi?); only one dot seen, left as Poesi.
Poesi, man tænker over brændende Spørgsmaal, uden\-%
landske saavelsom indenlandske: man tænker og tænker
over Stort og Smaat, kun ikke over det Religiøse. Paa
Reformationstiden derimod, ja, da blev der tænkt over
religiøse Problemer; det var ikke alene de Lærde, der
tænkte, nei, Folket selv tænkte med, thi da var det Alvor.
Da Folkets Selvtænkning saa efterhaanden hørte op,
fordi den religiøse Tænkning blev en lærd Sag, bleve
Theologerne, om just ikke selvtænkende, saa dog i den
% --- p. 109 ---
\opage{109}\setcounter{footnote}{0}%
officielle Kirke saa godt som officielt enetænkende; thi
nu var Orthodoxien en Videnskab. Derefter begyndte
saa Fritænkerne; men Fritænkerne forargede Folket;
Fritænkeriets Tankegang var i den Grad ufolkelig, at
den, endnu mere end den theologiske, bidrog Sit til at
vække Folkets Mistanke mod al grundig, religiøs Selv\-%
tænkning. Følgerne heraf ere kjendelige nok i den nær\-%
værende Tid.

Vi tør ingenlunde sige, at det enten er Sløvhed eller
Ligegyldighed, der afholder de Dannede fra at tage fat
paa Aandsopgaverne; nei, det er snarere Mismod, Ube\-%
hjælpsomhed og Mangel paa Overskuelse. Man har en
hemmelig Frygt for, at der ikke kommer Noget ud af
al den religiøse Tænkning; at det Hele tilsidst ender med
Opløsning, Gjæring og grændseløs Forvirring. Det er i
alt Fald saa lærd, saa kritisk, speculativt og philosophisk,
at Lægmand umulig kan følge med; nei, saa er det bedre
at holde paa det Bestaaende, følge Strømmen, have sine
egne Tanker for sig selv og troe paa Accomodationen.

Men er det da i Virkeligheden saa vanskeligt at
sætte sig ind i en religiøs Opgave? Vi mene at kunne
overtyde enhver Dannet om det Modsatte, og vælge et
Exempel, der ikke skal siges at ligge for lavt, idet vi
vælge en Bestemmelse af Forholdet mellem Idee og Gud,
for derved at bestemme Begrebet af det Hellige.

Forskjellen mellem Tro og Viden kan i Korthed
angives saaledes: det hellige Selv er Troens Gjenstand;
det Helliges Idee er Videns Indhold. Det hellige Selv
er Gud; i Troens Sprog hedder det: Gud er hellig,
retfærdig, algod, barmhjertig o. s. v. Anderledes i „den
moderne Videnskab“; her lyder Opfattelsen saaledes:
% Rettelser (PDF 338): læs "det Absolute selv". Printed "det absolute Selv"
% (lower-case a, capital S); a reader has pencilled the correction into this
% copy (an A drawn over the a, an s over the S); not transcribed.
Gud, det absolute Selv, kan ikke være hellig; det hel\-%
lige Selv kan ikke være Gud, thi det hellige Selv er en
Modsigelse. Vi finde i Strauszes Dogmatik følgende, for
% --- p. 110 ---
\opage{110}\setcounter{footnote}{0}%
den negative Viden høist charakteristiske Fremstilling:
„Den Definition \ldots{} af den guddommelige Hellighed
(3 Mos. XIX, 2. Ps. VI, 5 ff. Mt. V, 48. XIX, 17.
Eph. IV, 24. 1 Pet. I, 15 f., 1 Joh. I, 5 ff.), at den er
den Egenskab, paa Grund af hvilken Gud vil, hvad han
skal, og skal, hvad han vil, vil man paa den kirkelige
Side vel ikke let finde sig i: og dog udtaler den kun
det, som er Forudsætningen i de almindelige Begrebsbe\-%
stemmelser. Thi enten man nu definerer Guds Hellighed
som den høieste, mangelfrie Reenhed i Gud, der ogsaa
fordrer den skyldige Reenhed af Skabningerne; eller som
den guddommelige Villies Rigtighed, paa Grund af hvil\-%
ken den, i Overeensstemmelse med sin evige Lov, vil
Alt, hvad der er ret og godt; eller endelig som den
Egenskab, ifølge hvilken Gud kun attraaer og billiger
det Gode: saa er der bestandig antaget en moralsk Lov,
efter hvilken den guddommelige Villie bestemmer sig,
men dette Forhold mellem en Lov og en Villie bliver
ellers bestandig udtrykt ved en Skullen \ldots{} Meget rig\-%
tigt siger derfor Weisse: Gud selv vilde ikke kunne
kaldes god, hvis ikke det Onde ogsaa for ham var en
metaphysisk Mulighed. Men for ethvert Væsen, for
hvilket det Onde er muligt, staaer det Gode som en
Skullen, og saaledes træder vor ovenfor givne Definition
af den guddommelige Hellighed her ind i sine Rettighe\-%
der. Men just saaledes opfattet ophæver denne Egenskab,
hvis forresten ikke hiin Mulighed til det Onde og denne
Skullen ere tomme Ord, det Absolutes Begreb \ldots{} den,
som holder fast paa det Absolute, han opløser Hellig\-%
heden, der kun er Noget hos et Væsen, der er stillet i
Relationer; og den, som tager det alvorligt med Hellig\-%
heden, han træder Absoluthedens Idee for nær, da den
fordunkles ved den svageste Skygge af Mulighed til at
være anderledes end den er.“ (D. F. Strausz: Fremstil\-%
% --- p. 111 ---
\opage{111}\setcounter{footnote}{0}%
ling af den christelige Troeslære. Overs. af H. Bröchner.
1. Bind. Kbhvn. 1842. S. 513---515).

Af anførte Exempel vil det kunne indsees, hvor for\-%
skjellig den videnskabelige Tænkning er fra den religiøse,
og hvilken praktisk Betydning det har, at denne Forskjel
bringes til Bevidsthed. Skulde det nemlig afgjøres ved
videnskabelige Undersøgelser, om en hellig, retfærdig,
almægtig Gud kan være til eller ikke, da er det klart,
at Fleertallet af de Dannede ikke vil kunne følge slige
Undersøgelser, endnu mindre det hele Folk. Spørgs\-%
maalet om det Helliges Realitet bliver altsaa et indviklet
Spørgsmaal, hvis Besvarelse Lægfolk maae overlade de
Høilærde, medens de, indtil Spørgsmaalets endelige Be\-%
svarelse indtræder, søge at hjælpe sig igjennem med en
Tro, der i det Høieste kræver lidt Snaksomhed, men
kun saare ubetydelig Hjernevirksomhed, nemlig „Barne\-%
troen“. Saa har man da, ved at gaae til en modsat Yder\-%
lighed, „Barnetroen“ i „Barnehjertet“, en Tro, der for\-%
saavidt er tankeløs, som den ikke er istand til at udvikle
Begreber, men maa holde sig til Rørelser, Talemaader
og Billeder, en Tro, der faaer sit Opbyggelige fra en
Bekjendelse uden Selverkjendelse, en Beaandelse uden
Aandsudvikling, en hjertelig Forvisning uden egentlig
Overbeviisning. Jo enfoldigere, det vil her sige, jo
dummere Individet er, desto lettere vil det paa disse
Vilkaar være ham at troe; Barnetroen er den gemyt\-%
ligste Sag af Verden: af Børn og Barnetro forlanger
Vorherre jo slet ingen Tanker.

Men i en reflecterende Tidsalder forundes det kun
„en lille Flok“ at bevare Troen i det complet Tankeløse.
Hos det større Fleertal maae Tankerne tidligere eller
sildigere komme i Bevægelse; og jo stærkere Bevægelsen
bliver, desto mere paatrængende blive ogsaa Spørgsmaal
som disse: hvordan forholder det sig med Tro og Viden?
% --- p. 112 ---
\opage{112}\setcounter{footnote}{0}%
Er det virkelig i Videnskaben, maaskee i Theologien,
at Troestankerne skulle søge deres sidste og høieste Op\-%
klaring? Skulde der dog ikke kunne gives en sand reli\-%
giøs, af Viden og Videnskab uafhængig Tænkning?
Slige Spørgsmaals Betydning kan enhver Dannet for\-%
staae, ja den kan enhver nogenledes Opvakt forstaae.
Men Enhver, der er istand til at indsee Spørgsmaalenes
Betydning, er ogsaa istand til at fatte deres Besvarelse.

Naar Gud i Videnskaben ikke kan være hellig, fordi
han absolut skal være det graa Absolute: hvad gjør Vi\-%
denskaben da ved Begrebet Hellighed? Den ophøier
det til Idee med Prædicat af guddommelig. Religionen,
siger Feuerbach, forestiller sig Gud som hellig. I Sæt\-%
ningen: Gud er hellig, er \textit{Gud} Subject og \textit{hellig} Prædi\-%
cat; en saadan Sætning udtrykker kun en Forestilling,
men intet Begreb. For at opløse Forestillingen i sit
Begreb, maae vi vende Sætningen om, gjøre Prædicatet
til Subject, og Subjectet til Prædicat. Istedenfor Sæt\-%
ningen: Gud er hellig, faae vi da Sætningen: Helligheden
er guddommelig; ligeledes istedenfor Sætningen: Gud er
retfærdig, Sætningen: Retfærdigheden er guddommelig;
istedenfor Sætningen: Gud er barmhjertig, Sætningen:
Barmhjertigheden er guddommelig o. s. fr.

Hvorledes Videnskaben nu forøvrigt vil bære sig ad
med at fastholde saadanne frit flagrende guddommelige
Ideer som Hellighed, Retfærdighed o. s. v. hver for
sig, eller med at forene dem i een Idee, en Idee,
der da umulig kan være absolut, siden det Absolute selv
hverken kan være helligt eller retfærdigt, lade vi de Vidende
om; her skal det kun fremhæves, at Hellighedsbegrebet
oprindelig er et Troesbegreb, ligesaa uafhængigt af Viden\-%
skaben som Gudsbegrebet. At der er en hellig og ret\-%
% [a reader has underlined "At der er ... Vidnesbyrd i" in pencil; not letterspaced, not transcribed]
færdig Gud til, kan ene og alene afgjøres ved Selvfor\-%
staaelse, ved indre Erfaring, ved Aandens Vidnesbyrd i
% --- p. 113 ---
\opage{113}\setcounter{footnote}{0}%
Samvittigheden, altsaa ved Tro. Ganske vist er det en
væsentlig og videnskabelig Opgave at faae det gjort
mere og mere klart, hvorledes de relative Ideer forholde
sig til det Absolute selv, og fremfor Alt, hvorledes det
Ethiske principielt forholder sig til det Metaphysiske;
men disse Undersøgelser vedkomme ikke Religionen som
Religion, ikke Troen som Tro. For den troende Be\-%
vidsthed staaer det urokkelig fast, at Gud ligesaa vist er
hellig og retfærdig, som at han er almægtig. Denne
Forudsætning er ikke en tankeløs Paastand, men en i
Selverkjendelsen dybt grundet Overbeviisning, en Troes\-%
principet selv iboende Begrebsbestemmelse.

Til at fastholde Troestanken udkræves ikke en ny
Revision af Philosophiens Historie, men en ethisk Con\-%
centration i religiøs Inderlighed: det frie Selv tilbeder
den Hellige; det matte Selv saliggjøres i Metaphysik.

Dersom det hellige Selv ikke er, det vil sige: der\-%
som Gud ikke er hellig, saa er Hellighedens Idee en
Mulighed, der aldrig kan blive virkelig; men en Mulig\-%
hed, der aldrig kan blive til Virkelighed, er en Umulig\-%
hed. Istedenfor at sige: Helligheden er guddommelig,
maae vi drage den videnskabelige Conseqvens, tale reent
ud og sige: Hellighed er et selvmodsigende Begreb,
Hellighed er umulig. Som det gaaer med Helligheden,
saaledes ogsaa med Retfærdigheden: kan Gud ikke være
retfærdig, saa kan Retfærdigheden heller ikke være gud\-%
dommelig; den maa smukt finde sig i at være menneskelig.
Men hvad vil saa dette sige? Det vil sige, at Retfærdig\-%
heden først kan virkeliggjøres, efterhaanden som Menneskene
ved sig selv, enkeltviis og i Samfundet, omsider blive ret\-%
færdige. Naar skeer det? Aldrig! Det hører jo ogsaa med
til „den moderne Videnskabs“ store Opdagelser, at intet
Menneske og intet menneskeligt Samfund kan i nogen\-%
somhelst virkelig Tilstand virkeliggjøre Retfærdigheden.
% --- p. 114 ---
\opage{114}\setcounter{footnote}{0}%
Retfærdighedsideen i uendelig Skullen uden Kunnen, er i
endeløs Splid med sin egen Virkeliggjørelse. For at
dække den grelle Modsigelse har man beraabt sig paa,
at de enkelte Menneskers Retfærdighed vel var ufuldkom\-%
men, men at Enhver dog idetmindste kunde fremstille
sin Side af det retfærdige Væsen, med andre Ord, præ\-%
stere sit Stykke Retfærdighed. Naar saa de mange Bidrag
samledes i det sociale Liv, fik vi, ved at lægge Stykkerne
sammen, et næsten retfærdigt Hele. Det er jo rigtignok
kun en Tilnærmelse! Ja, det er i Ethiken en Tilnær\-%
melse til Idealet af samme Art, som det i Poesien vilde
være en Tilnærmelse til Idealet, hvis man, i Mangel af
en eneste sand, stor Digter, vilde lade en Mængde
jævnt begavede Mennesker, som hvert havde sin Smule
Poesi, arbeide sammen i en æsthetisk Comité for ved
mange smaa Bidrag at frembringe et stort Digt.

Den, der ikke kan indsee, at Forvisningen, den
urokkelige, ubetinget afgjørende Forvisning om Guds
Hellighed og Retfærdighed, er given i og med den til
% [a reader has underlined "Selvconcentration" in pencil and marked the margin; not transcribed]
Troen svarende Selvconcentration, han har, hvormeget
han end forøvrigt kan have tænkt, endnu ikke gjennem\-%
tænkt Troens Princip, og gjør da vel i at gjennemtænke det,
inden han forkaster det. Men dersom Nogen, efter at
han virkelig har gjennemtænkt Principet, og indseet, at
Begrebet af en hellig og retfærdig Gud ganske rigtig
indesluttes i Troestanken som absolut Selvhedstanke,
% AUDIT doubtful: a grave-like mark stands over the o of "absolute" (absòlute?) -- stray type or speck; left as absolute.
bryder af med den Erklæring, at den absolute Selv\-%
hedstanke for ham er uden Realitet, eftersom den fordrer
en Anstrengelse af Selvet, en Selvbestemmelsens Energie,
han ikke kan udholde: da maa han søge Skylden hos sig
selv og ikke misbruge Viden til Skjul for egen Svaghed.

Hvilken Forskjel der er imellem Troens hellige Gud
og Videnskabens Gudsidee, kan da ogsaa sees deraf, at
medens ingen Vidende vil falde paa enten at bede til
% --- p. 115 ---
\opage{115}\setcounter{footnote}{0}%
eller at tilbede sin Videns Idee, er Troens Gud derimod
Gjenstand for Tilbedelse, en Gud, til hvem de Troende
bede. Der er af Johannes Climacus (Uvidenskab. Efterskr.
S. 148) anstillet følgende Overveielse: „Dersom En, der
lever midt i Christendommen, gaaer op i Guds Huus, i
den sande Guds Huus, med den sande Forestilling om
Gud i Viden, og nu beder, men beder i Usandhed; og
naar En lever i et afgudisk Land, men beder med Uen\-%
delighedens hele Lidenskab, skjøndt hans Øie hviler paa
en Afguds Billede: hvor er saa mest Sandhed? Den
Ene beder i Sandhed til Gud, skjøndt han tilbeder en
Afgud; den Anden beder i Usandhed til den sande Gud,
og tilbeder derfor i Sandhed en Afgud.“ Hvor nu Op\-%
gaven er, saaledes som hos anførte Forfatter, at faae
Subjectiviteten og Inderligheden ret eftertrykkelig betonet,
ere slige Overveielser vistnok af stor Værdi; men hvor
det, som i nærværende Forhandling, kommer an paa at
fremhæve det principielle Forhold imellem Tilbedelsen
og Tilbedelsens Gjenstand, har Spørgsmaalet kun en
flygtig og forbigaaende Betydning. At bede i Sandhed
til Gud og dog tilbede en Afgud, er umuligt; thi dersom
Tilbedelsen er sand, kan den kun være sand derved, at
den afgudiske Forestilling forsvinder i Forestillingen om
en sand Gud; men en sand Gudsforestilling udelukker
nødvendig Forestillingen om en Afgud. Ligesaa er det
umuligt i Usandhed at tilbede den sande Gud; thi, naar
Tilbedelsen er i Usandhed, skjuler den sande Gud sit
Aasyn for den Tilbedende og efterlader i sit Sted et
Phantom, en usand Gud.

At Tilbedelsen og dens Gjenstand naturlig svare til
hinanden, og at selv oplyste Culturmennesker ofte føle
Trang til Noget, de kunne tilbede, er stadfæstet ved
mange Erfaringer: har man ikke noget Sandt at tilbede,
nu, saa tilbeder man noget Usandt, og Tilbedelsen bliver
% --- p. 116 ---
\opage{116}\setcounter{footnote}{0}%
derefter; kan man ikke anerkjende den sande Gud, fordi
den sande Gud er kommen i Modsigelse med Metaphy\-%
siken, saa kan man jo anerkjende sit eget Absolute og sætte
dette som metaphysisk Afgud. Afgudsdyrkelsen er ikke
ophørt, fordi man har ophørt med at tilbede Sol, Maane
og Stjerner. Har et Individ ikke Selvvigtighed nok til
at tilbede sig selv, saa kan det jo have sin Glæde af at
tilbede Andre. Det er i en reflecterende Tidsalder ingen\-%
lunde usædvanligt, at det ene Menneske driver Afguderi
med det andet. En Moder f. Ex. driver Afguderi med
sin Søn, naar hun i Sønnen seer et sandt Vidunder af Geni
og Dyd, et Mønstermenneske, som Verden i sin Forblin\-%
delse aldrig vil forstaae at paaskjønne. At en Elsker, idet\-%
mindste i en vis Periode, driver Afguderi med sin Elskede,
er da heller ikke noget Uhørt. Derfor kaldes hun jo
ogsaa Gjenstanden, den Tilbedte, og Lykken er i Til\-%
bedelsen: lykkelig den Tilbedte, ti Gange lykkeligere
den Tilbedende! Men ogsaa her viser det sig, hvor
umuligt det er, at Tilbedelsen kan være sand, naar det
kun er en Afgud, der tilbedes. Der er i alt Afguderi
noget Mythisk; derfor er der ogsaa i al afgudisk Tilbe\-%
delse noget Mythisk. Det Mythiske er tvetydigt, og varer
kun til en Tid. Saa møder man Tilbederen et halvt
Aar efter, men hvor er den Tilbedte? Det er forbi;
det Hele var kun en Mythe. „Det var dog besynderligt,“
siger den Troskyldige, „thi han tilbad hende virkelig; hvil\-%
ken Ustadighed, hvilken Troløshed!“ Det Bedste vilde
nu være, om den Anklagede faldt paa at skyde Skylden
paa Videnskaben og expectorere sig saaledes: „Tilbedel\-%
sens Begreb begynder med Tro, men maa naturligviis
ligesom andre Troesbegreber klare sig i Viden. Jeg
tilbad hende virkelig, det er sandt, thi jeg troede, ja, jeg
troede; jeg var jo endnu dengang kun Begynder i
Videnskaben, og derfor troede jeg; men efterat jeg er
% --- p. 117 ---
\opage{117}\setcounter{footnote}{0}%
bleven fortrolig med Videnskaben, og særlig med Meta\-%
physiken, har jeg lært at indsee, at Tilbedelse er et Be\-%
greb, der ikke kan bestaae med det Absolute, et Begreb,
der halverer min Bevidsthed, et modsigende Begreb, og
derfor har jeg ophørt at tilbede.“

Hvad her klinger som Spot, fordi det udsiges om
en Tilbedelse, hvis Tvetydighed Enhver saa let gjennem\-%
skuer, kan for fuldt Alvor anvendes paa Tilbedelsen af
de mythiske Guder. De have til en given Tid kunnet
være Gjenstand for en dunkel, foreløbig, sandselig Til\-%
bedelse; men Ingen har i Aand og Sandhed nogensinde
tilbedet Zeus; thi en Zeus kan ikke tilbedes i Aand og
Sandhed; det strider imod hans mythiske Natur. Det,
jeg skal tilbede, maa være mit Væsens rene, hellige Op\-%
hav, Personligheden selv i Sandhed, Fuldkommenhed og
Salighed: kun Aandens Gud, der er Aanden selv, kan
tilbedes i Aand og Sandhed.

Opkaster man det Spørgsmaal, hvilken Tidsalder,
hvilken Slægt der maa kaldes den ulykkeligste, saa veed
man jo af Historien, at en Slægt har været tugtet paa en,
en anden paa en anden Maade, en har lidt under Krigens
Svøbe, en anden af Pest, Hunger o. s. v. Der kan være
store Ulykker i Verden; men saalænge det Hellige be\-%
vares, og den lidende Slægt har en Gud at bede til og
at tilbede, er Ulykkens Maal ikke fyldt. Den ulykke\-%
ligste Slægt maa vist være den, som ved en Misviisning
kommer ind paa en Culturvei, der fører den bort fra Troen,
bort fra det Hellige, saa at der Intet er, hvorpaa den
kan troe, Intet, hvortil den kan see op med Forventning
og Haab, Intet, Intet, hvorfor den kan bøie sig med
inderlig Ærefrygt. Lad en saadan Tidsalder forresten
saa gjøre Fremskridt i alle Retninger og nyde alle
Civilisationens timelige Goder; den er og bliver dog den
ulykkeligste. Vi have et Forbillede i Rom under Kei\-%
% --- p. 118 ---
\opage{118}\setcounter{footnote}{0}%
serne. Skulde Nogen i det store Rige kunne kaldes en
lykkelig Mand, da maatte det vel i romersk Forstand
være Keiseren selv. Hvad er det da, der volder ham
Qvaler? Det er ikke Bekymringer for Undersaatternes
Vel, thi om Undersaatternes Vel bekymrede han sig slet
ikke; heller ikke kan det være Mangel paa Magt, thi
han var i menneskelig Forstand almægtig; og dog
lider han, han lider utrolige Qvaler, grusom som en
Nero, rasende som en Caligula, dyrisk som en Vi\-%
tellius, vellystig indtil Vanvid som en Heliogabal!
det er aabenbart, at han lider. Men hvoraf kommer
det? Det kommer deraf, at en saadan Almægtig ikke
har Noget, han kan see op til og betragte med Ære\-%
frygt; det kommer af, at Keiseren er en Gud, og har
som en Gud lært at foragte Alt, endogsaa sig selv.

Dette er Maalestokken for en Ulykke af høiere Art:
at det Hellige forsvinder, at al Tro forsvinder, at der
Intet mere er, hvortil Slægten med Ærefrygt kan see op,
Intet, hvorpaa den kan bygge sit Haab, Intet, den i Aand
og Sandhed kan tilbede.

\begin{center}\rule{3cm}{0.4pt}\end{center}

\chapter*{VIII.}
\addcontentsline{toc}{chapter}{Forelæsning VIII}
\markboth{Forelæsning VIII}{Forelæsning VIII}
\label{ch:viii}
% pp. 119--134
% --- p. 119 ---
\opage{119}\setcounter{footnote}{0}%
% [large initial N]
Naar Talen er om Civilisation og Oplysning, bør Nu\-%
tiden ikke glemme, at der er Forskjel mellem Lys og
Lys, mellem Culturoplysning og Folkeoplysning; de to
Slags Oplysning ere ikke to Sider af samme Sag, men
indbyrdes ueensartede Ting. De forholde sig efter deres
Udspring som Viden og Tro, som Videnskab og Reli\-%
gion; thi Civilisationens Høieste er i Viden og Folke\-%
oplysningens Dybeste er i Tro. Da Sammenblandingen
af Begreberne Tro og Viden, som vi have seet, er en af
de væsentligste Hindringer for det aandelige Liv i Nu\-%
tiden, vil det, for at bevirke deres Adskillelse, være hen\-%
sigtsmæssigt at udhæve Forskjellen imellem Cultur- og
Folkeoplysning noget nærmere. Hvor forskjellig den
egentlige Folkeoplysning er fra hvad vi kalde Culturop\-%
lysning, ville vi bedst kunne see ved skarpt at adskille
to Begreber, som Nutiden kun altfor let sammenblander,
nemlig det Folkelige og det Populære.

Til at opklare Forskjellen mellem det Folkelige og
det Populære vil det være nødvendigt at gjøre sig Rede
for, hvad det er, der bestemmer Forskjellen mellem fol\-%
kelige og populære Skrifter, mellem folkelige og populære
Mænd.

Der gives Skrifter, som egne sig til almindelig Fol\-%
% --- p. 120 ---
\opage{120}\setcounter{footnote}{0}%
kelæsning og med Grund kunne kaldes populære, uden
at de derfor nogensinde kunne blive folkelige; omvendt,
gives der Skrifter, som uagtet Fleertallet af Befolkningen
endnu ikke ret er istand til at læse dem eller finde Smag
i dem, dog ifølge deres Grundcharakteer maae kaldes
folkelige. Man har i Litteraturen en populær Astronomie,
en populær Geologie, Agerdyrkningschemie o. s. v., men
slige Bøger kunne aldrig blive folkelige. Og hvorfor
ikke? Fordi de aldrig kunne meddeles Læsere, der ikke
ere af Faget, i den for Indholdet oprindelige, eiendom\-%
melige Form; thi den oprindelige Form er og maa være
videnskabelig. Det er ikke de videnskabelige Undersø\-%
gelser, men kun deres Resultater, der meddeles. Men
videnskabelige Resultater tabe deres oprindelige Værdi,
naar de løsrives fra de Undersøgelser, hvoraf de ere
fremgaaede, meddeles i en Form, hvori de ikke kunne
begrundes. Det Populære ligger i Formen; men den
populære Form er en lavere, uadæqvat, fra den egentlige
Viden afvigende, altsaa mere eller mindre forhutlet Form.
Istedenfor at lede Læseren til selvstændig Indsigt, vænner
en saadan Meddelelse ham tvertimod til at stole paa
Kyndigeres Autoritet og tage det Meddeelte paa Tro og
Love. At dannes ved Læsning af lutter populære Skrif\-%
ter er ingenlunde at dannes til Selvtænkning. Skulde et
Folk ad Culturoplysningens Vei hæves til selvstændig
Tænkning, da maatte det nødvendig hæves til viden\-%
skabelig Indsigt. Men det er en meningsløs og unaturlig
Fordring. Eller vilde det ikke være en meningsløs og
% a speck after "og" at the line end, not a comma; not transcribed
unaturlig Fordring, om man ventede, at Oplysningen
skulde stige saa høit, at Alle i det hele Folk bleve
Videnskabsmænd? Sæt, at Culturen bragte det saavidt,
at alle Mænd bleve Astronomer, Geognoster, Botanikere,
Zoologer o. s. v., hvad skulde der da blive af Qvinderne?
De skulde vel ogsaa blive videnskabelige. Sæt, at vi
% --- p. 121 ---
\opage{121}\setcounter{footnote}{0}%
engang fik lutter videnskabelige, zoologiske, botaniske,
ja mathematiske Qvinder, hvorledes vilde det da gaae
med Qvindeligheden? Videnskaben kan i Charakteer af
Videnskab aldrig blive folkelig; thi det Videnskabelige
er ikke for Alle, men det Folkelige er og maa være
for Alle.

Hvad er det Folkelige? Det er det Grundmenne\-%
skelige i naturkraftig Eiendommelighed. Hvad heraf
% a small raised tick between "skelige" and "i", apparently a speck; not transcribed
følger for Bestemmelsen af Folkeoplysning og folke\-%
lige Skrifter, er da let at see. Kun den Tænkning, den
Oplysning og Begrebsudvikling, hvis første Forudsætnin\-%
ger og Betingelser levende røre sig i den personlige,
grundmenneskelige Bevidsthed, kan kaldes folkelig. En
Sætning som den, at den rette Linie er den korteste Vei
imellem to Punkter, er vel i en vis Forstand baade folke\-%
lig og populær, og dog er her en væsentlig Forskjel
imellem den folkelige og den populære Opfattelse. Op\-%
fattelsen er folkelig, forsaavidt Sætningen uden kunstig
Abstraction er indlysende for den sunde Sands, som
naar det hedder: den lige Vei er den korteste. Popu\-%
lær bliver derimod Sætningen, naar der udtrykkelig for\-%
dres en, om end blot foreløbig, mathematisk Abstraction,
for at fastholde Linie som Linie, Punkt som Punkt. At
vi her ved det Populære, komme bort fra det Folkelige,
% "Populære," : the comma is a faint tail-less mark at the baseline (the layer reads a comma); alternative rejected: no comma
mærkes øieblikkelig, naar vi spørge om den anvendte Ab\-%
stractions Betydning. Abstractionen viser nemlig ud over
det Umiddelbare og minder om, at den populære Forstaaelse
kun er foreløbig og underordnet; thi det er Videnskaben,
der har anlagt Abstractionen, og det er i Videnskaben
alene, at dens Betydning kan udvikles. Der behøves
kun et Skridt endnu paa Abstractionens Vei, og den
populære Sætning har ophørt at være populær. Ander\-%
ledes med en folkelig Sætning! Vel maa dens Udvikling
ogsaa være betinget af Abstraction; men Abstractionen
% --- p. 122 ---
\opage{122}\setcounter{footnote}{0}%
foregaaer naturlig og som af sig selv: Udviklingen er
levende. At store Helte have store Lidenskaber, er en
Sætning, der ikke blot kan anskueliggjøres i Sagn, Hi\-%
storie og Poesi, og forsaavidt kaldes folkelig; men den
kan tillige opfattes af Tænkningen, udvikles og fremstilles
i de allerfineste psychologiske Afskygninger, og endda
være folkelig; naar Øiet falder i Sagas Speil, bringer
den grundmenneskelige Bevidsthed selv de til Sætningens
Forstaaelse nødvendige Betingelser med: Saxo Gramma\-%
ticus og Snorre Sturlesen ere Folkeskribenter af første
Rang, men de kunne ikke kaldes populære.

Udviklingen af det Folkelige er stadig udsat for to
Misviisninger, eftersom man enten overseer det i Folke\-%
naturen Individuelle og holder fast ved det blot Almeen\-%
menneskelige, eller med Ringeagt for det Almeenmenne\-%
skelige udelukkende seer paa den naturkraftige Eien\-%
dommelighed. Den første Vildfarelse er tilstrækkelig
betegnet ved det 18de Aarhundredes Forstandsoplysning,
dets mangehaande kosmopolitiske og philantropiske Ab\-%
stractioner; Revolutionens „Menneskerettigheder“ hvile paa
et kosmopolitisk Grundlag. Det kosmopolitisk-philantro\-%
piske Forstandsmenneske naaer med sin iøvrigt agtværdige
Iver for Opdragelsesvæsen, Skolevæsen, Menigmands Op\-%
lysning osv. aldrig ud over det blot Populære. Advaret
ved Udskeielser af en tom Forstandsoplysning er man i
vore Dage, under Indvirkning af de sig fra alle Si\-%
der paatrængende Nationalitetsspørgsmaal, fristet til en
modsat Yderlighed, til at klamre sig ved den umid\-%
delbare Eiendommelighed i den Grad, at man uden vi\-%
dere slaaer det Tilfældige sammen med det Væsentlige
og glemmer det Almeenmenneskelige. Istedenfor at vække
og forberede Folket til friere, kraftigere Selvudvikling,
smigrer man dets naturlige Selviskhed og nærer dets
almueagtige Hang til selvklog Indesluttethed, til Ligegyl\-%
% --- p. 123 ---
\opage{123}\setcounter{footnote}{0}%
dighed og uforstandig Ringeagt mod det Lys, Culturen
udbreder.

Skulle de to Sider af det folkelige Væsen, det
Grundmenneskelige og det naturkraftige Eiendommelige,
forenet virke i Aand og Sandhed, da maa det skee ved
en grundig Omdannelse, en Gjenfødelse. Den gjenfø\-%
dende Kraft betinges af, at det aabenbarede Ord tilegnes
i Troen: det Troesbegreb, hvorpaa Eftertrykket her
maa falde, er Begrebet om Guds Billede.

At Mennesket er skabt i Guds Billede er den af\-%
gjørende Grundsætning, hvorfra al Lære om Folk og
Folkelighed maa gaae ud. Heri ligger for det Første
en Henviisning til det Grundmenneskelige; kun idet Men\-%
nesket saaledes fordyber sig i sit eget Væsen, at han
deri erkjender det Almeenmenneskelige, kan han erkjende
det Gudsbilledlige. Men Gudsbilledets Væsen kan ikke
fastholdes i tom Almindelighed; det er et Livsbillede,
det Grundbillede af Menneskets Existens, hvortil dets
sande Individualitet maa svare. Naar Gudsbilledet virke\-%
liggjøres, fremkommer det qvindeligt i Qvinden, mandligt
i Manden; Udfoldningen af dets Indhold giver ikke tom
Eensformighed, men rig Forskjellighed; i Gudsbilledet
bliver ethvert Menneskes Eiendommelighed ikke blot be\-%
varet, men fremmet og forhøiet. Hvad der gjælder om
det enkelte Menneske i Folket, at det nemlig ved at ud\-%
vikle sig efter Gudsbilledet udvikler sin oprindelige In\-%
dividualitet, det Samme gjælder om et heelt Folk; men
naar alle Individualiteter i et Folk udvikles, udvikles
jo Folkeindividualiteten.

Dog, her reiser sig en historisk Betænkelighed:
„Dersom det virkelig er sandt, at det evangeliske Grund\-%
billede efter sit Princip afgiver Forudsætninger og Betin\-%
gelser for Udvikling af Folkenes naturkraftige Eien\-%
dommelighed, da maatte et saadant Udviklingsprincip jo
% --- p. 124 ---
\opage{124}\setcounter{footnote}{0}%
udtrykkelig kunne paavises i folkehistoriske Kjendsgjer\-%
ninger; men er det nu Tilfældet?“ Vi svare: Historiens
Vidnesbyrd er utilstrækkeligt, men det er ikke Principets
Skyld; Skylden er hos dem, der have misbrugt og for\-%
qvaklet Principet.

Evangeliet er fra mere end een Side bleven opfattet
og anvendt paa en saa uevangelisk Maade, at Folket og
det folkelige Liv har maattet lide derunder. Hvad har
man da gjort? Man har benyttet Religion og Tro til at
holde den menige Mand i aandelig Umyndighed gjennem
Aarhundreder. Da Gregor den Store vilde forsvare de
hellige Billeder, beraabte han sig paa Folkets Uvidenhed:
% the Latin sentence is set in true italic (not roman as elsewhere in the book)
\textit{Legant in parietibus, quod non legunt in codicibus.} Den
katholske Kirke har, som bekjendt, ikke forstaaet
det saaledes, at det uvidende Folk kun skulde søge
Opbyggelse ved at see paa Billederne, indtil man
efterhaanden kunde faae det oplært til at læse i Skrif\-%
terne; nei, den har forstaaet det saaledes, at den katholske
Almue endnu i vore Dage har „samme Trang“ til at see
paa Billederne, som den havde i Gregors Dage.

Der er ved Troen et Charakteermærke, som Hierar\-%
chiet ypperligt har vidst at benytte til Fordeel for sig, men
til Skade for Folket; dette Charakteermærke er Lydighed!
Ligesom Tro er uadskillelig fra Tilbedelse, saaledes ere
Tro og Tilbedelse uadskillelige fra Lydighed. At troe paa
Gud er at tilbede ham, at tilbede ham er at tjene ham,
at tjene ham er at lyde ham. Aabenbaringens Ord er et
myndigt Ord, det myndige Ord kræver Lydighed. Fri\-%
stelsen til at slaae Guds Bud sammen med Mennskebud
% printer's defect: "Mennskebud" as printed (the e of "Menneske-" dropped; confirmed at 600 dpi, no visible gap); read "Menneskebud"
er i dette Tilfælde stor. Aabenbaringens myndige Ord lod
sig fortræffelig bruge til at forskaffe Kirkens Styrere
Lydighed, den blinde Lydighed, der gjør det saa be\-%
qvemt at holde Orden i Samfundet. Jo mere ubetinget
Lydigheden er, desto mindre behøve de Lydige selv at
% --- p. 125 ---
\opage{125}\setcounter{footnote}{0}%
tænke. Den første Fordring til en troende Katholik er
blind, ubetinget Underkastelse under Kirkens Autoritet;
den blinde Tro er Eet med den blinde Lydighed. Hvor\-%
for er vel Kjætteri en Dødssynd? Det er ikke de en\-%
kelte Afvigelser i de enkelte Artikler, hvorpaa her lægges
Vægt; nei, Kjætterens Forbrydelse er, at han tør drage
Kirkens Ufeilbarhed i Tvivl, at han ikke underkaster
sig dens Bud i ubetinget Lydighed. Blindheden ændser
intet Lys: Hierarchie og Folkeoplysning ere uforenelige.

Men nu Protestantismen? Ja, Protestantismen! At
Reformationen er folkelig fra Grunden af og begyndte
folkeligt, er vist; men hvorvidt er den kommen? Paa
Rigsdagen i Worms, da den trierske Official faldt Lu\-%
ther i Talen og forlangte et slet og ret Svar, om han
vilde kalde tilbage eller ikke, da lød et Svar, der nu
har lydt over en heel Verden, og Svaret lyder: „Eftersom
da Eders Keiserl. Maj., Chur- og Fyrstel. Naader be\-%
gjære et enfoldigt, slet og ret Svar, saa vil jeg give et,
som hverken har Horn eller Tænder, nemlig saaledes:
Med mindre jeg ved den hellige Skrifts Vidnesbyrd, eller
med aabenbare, klare og indlysende Grunde bliver over\-%
beviist (thi hverken Paverne eller Concilierne troer jeg
alene, da det ligger aabenbart for Dagen, at de ofte have
feilet og modsagt hverandre), saa at jeg ved selve de
Sprog, som jeg har anført og brugt, vorder overtydet,
og min Samvittighed bliver fangen i Guds Ord, hverken
kan eller vil jeg kalde Noget tilbage; thi det er hverken
sikkert eller raadeligt at gjøre Noget mod sin Samvit\-%
tighed. Her staaer jeg, jeg kan ikke Andet, Gud hjælpe
mig. Amen!“ Ved dette Svar er Staven brudt over
den blinde Lydighed; i dette Svar er Reformationens
paa een Gang religiøse og folkelige Princip fuldstændig
udtalt --- i et stort Øieblik, af en stor Personlighed!
Det er Alvor med det Religiøse, thi det gjælder om
% --- p. 126 ---
\opage{126}\setcounter{footnote}{0}%
Guds Ord og en god Samvittighed. Det er Alvor med
det Folkelige; thi det er netop den ved Troen frigjorte
Individualitet, som ved sin mandige Optræden her frem\-%
stiller os det Menneskelige, det Grundmenneskelige i na\-%
turkraftig Eiendommelighed. Den personlige Ret, Lu\-%
ther herved hævder for sig selv, hævder han tillige for
enhver Troende i Folket, og den Forpligtelse, han her\-%
ved paalægger sig selv, paalægger han med det Samme
Enhver i Folket, som deler hans Tro. Skal Troen føre
til personlig Selvstændighed, da maa den ved Aanden
% a reader's pencilled cross in the left margin here; not transcribed
og Ordet føre til personlig Selvudvikling, til den frigjø\-%
rende Sandheds Erkjendelse, og kan da umulig være
tankeløs. Vel er den frigjørende Sandhed en Samvittig\-%
hedssandhed, ikke en Videnskabssandhed; men som
Sandhed maa den dog erkjendes, og for at kunne er\-%
kjendes, maa den kunne tænkes. Men gjælder det om
aabenbare, klare og indlysende Grunde, saa gjælder det
for Alvor om Folkeoplysning og Folkedannelse; da kan
Folket ikke slaae sig til Ro ved almueagtig from Enfol\-%
dighed; saa er det ikke tilstrækkeligt at troe paa Pavens
Helgenbilleder, men heller ikke tilstrækkeligt at bryde
med Paven og hans Latinere for at blive siddende i
Ammestuen med Billedsproget om Barnetroen og snakke
Barnesnak.

Da den orthodoxe Lutherdom forvanskede Luthers
Værk, var det igjen Lydigheden, den blinde Lydighed,
Lydigheden, ikke i Samvittigheden mod Guds Ord, men
Lydighed mod Menneskebud, religiøs Lydighed mod
„den høie christelige Øvrighed“, hvorved Forvanskningen
foregik. Hvorledes Luther i det Hele dømte om Fol\-%
kenes Forpligtelser til at lyde deres retmæssige Fyrster,
er bekjendt, især af hans Udtalelser om Bondekrigen.
Den Raahed, hvormed de Undertrykte reiste sig imod
deres Tyranner, bevirkede, at Reformatorerne maatte
% --- p. 127 ---
\opage{127}\setcounter{footnote}{0}%
lægge en ganske overordentlig Vægt paa den borgerlige
Lydighed. Dette var forsaavidt ganske i sin Orden.
Men inden man vidste deraf, blev Lydighedspligten i
% "Men": a small mark above the e reads as a speck, not an accent (layer "Men"); not transcribed
det Borgerlige slaaet sammen med Lydighedspligten i
det Religiøse, og nu gik Forvanskningen og Misbrugene
ud over alle Grændser.

Det er dog ikke saameget Regjering og Fyrste,
som meget mere Kirken og dens Theologer, der maae
bære Skylden for Folkeudviklingens protestantiske Un\-%
derkuelse. I den orthodoxe Theologie bleve Dogmerne
bestemte med en objectiv formel Nøiagtighed, der ret
gjorde dem skikkede til juridisk Anvendelse. Regjeret
af Theologien, regjerede Fyrsten i egen Person Theo\-%
loger og Kirke. Vi have et charakteristisk Exempel i
Christian den Fjerde. Hans gudelige Anordninger og
borgerlige Love gjenlyde af Bibelsteder og Gudsfrygt,
og medens han ikke sjeldent paa en haard og vilkaarlig
Maade revser de enkelte Lærere, er han selv afhængig
af dogmatiske Kjendelser, endog i den Grad, at han i
sine huuslige Stridigheder med Christine Munk maa
bruge det theologiske Facultet som Orakel.

Theologien medførte endnu en anden for Folkeop\-%
lysningen meget væsentlig Hindring. Som Udviklingen
af det theologiske Lærebegreb var, saaledes vare natur\-%
ligviis ikke blot Prædikenerne men ogsaa de religiøse,
af Regjeringen forordnede Lærebøger. Disse til Almuens
Underviisning indrettede Lærebøger vare nemlig Uddrag
af dogmatiske Systemer, altsaa uvidenskabelige Gjengi\-%
velser af Noget, der skulde gjælde for Videnskab; paa
denne Maade maatte Religionsunderviisningen nødvendig
nedsættes til at blive populær, uden at kunne blive fol\-%
kelig. Luthers lille Catechismus er en folkelig Bog;
men for at nævne et Exempel, der ligger os nær, Balles
Lærebog er ikke en folkelig Bog, den er kun populær.
% --- p. 128 ---
\opage{128}\setcounter{footnote}{0}%
Dens populære Charakteer falder i Øinene, naar vi
see, hvorledes Forfatteren i Læren om Guds Væsen og
Egenskaber til en Begyndelse indlader sig paa at bevise
Guds Tilværelse. Beviset er skrøbeligt, det bærer Mærke
af Datidens overfladiske Philosophie; men selv om Phi\-%
losophien gik fremad og opfattede Sagen grundigere, vilde
det populære Beviis derfor ikke blive bedre. Beroer
Forvisningen om Guds Tilværelse i dybere Forstand
derpaa, at det theologiske Beviis tilfulde prøves, da er
Folket umyndigt; thi det kan ligesaalidt prøve et theo\-%
logisk som et philosophisk Beviis.

Saavidt om Forskjellen mellem populære og folkelige
Skrifter.

Vi skulle nu kaste et Blik paa Forskjellen mellem
populære og folkelige Mænd, en Forskjel, der naturligviis
maa bestemmes efter samme Princip som den oven\-%
nævnte. Man blander saa let de to Begreber Populær
og Folkelig, og dog ere de saa uligeartede, at en Mand
kan være populær uden nogensinde at blive folkelig, og
omvendt være folkelig uden at blive populær. Ligesom
der, hvad det enkelte Menneske angaaer, kan være en
stor Forskjel imellem hans Grundvillie, det, han
efter sin dybere Natur vil og maa ville, og hans øie\-%
blikkelige Villie, det, han under givne Omstændigheder
har Lyst og Tilbøielighed til, saaledes kan der ogsaa,
hvad et heelt Folk angaaer, gjøres Forskjel imellem dets
endelige, af Interesser og Klogskabshensyn ledede Til\-%
bøieligheder og dets Grundvillie. Den Mand, der leder
et Folk til at udføre, hvad det helst vil, uden at spørge
om, hvad det maa ville, han bliver populær, men ikke
folkelig; den derimod, der bevæger Folket til at udføre,
ikke hvad det helst vil, men hvad det maa ville, han
bliver folkelig, men ikke populær: den populære Mand
% --- p. 129 ---
\opage{129}\setcounter{footnote}{0}%
gjør Folket til Mængde, den folkelige Mand gjør Mæng\-%
den til Folk.

Det folkelige Ord er et myndigt Ord, den folkelige
Mand en myndig Mand; men det folkelige Ord har i sin
Myndighed tillige noget Vækkende og Frigjørende: den
folkelige Mand vil vække Folket og føre det til Frihed.
Ved det Folkelige maa der lægges Vægt baade paa det
Myndige og det Frigjørende. Saaledes var Biskop Ab\-%
salon for sin Tid en folkelig Mand, men ikke populær;
først i den nyere Tid har han begyndt at vinde nogen
Popularitet. Kong Gustav Erikson Vasa i Sverrig var
en folkelig, men langtfra nogen populær Mand. Hvor\-%
ledes han opfattede sit Liv, sin Regjering og Stilling til
Folket, har han paa en gribende og veltalende Maade
fremstillet i den sidste Tale, han holdt til Stænderne,
hvor han med Døden for Øie tager Afsked fra Regje\-%
ring og Verden og uforbeholdent aflægger Regnskab
for sit Liv og Brugen af sin kongelige Myndighed. Han
siger blandt Andet: „Jeg veed, at jeg i Manges Øine
har været en haard Konge, men de Tider skulle komme,
da Sverriges Børn vilde rive mig op af Mulde, om
det stod i deres Magt.“ Det var folkelige Ord, men det
var ogsaa myndige Ord. Tidens Vilkaar vare haarde,
og haard maatte den Konge være, som med Gustavs
Midler skulde naae hans Maal. Men Maalet var,--- et Be\-%
frielsens Maal, og i Befrielsen laa det Folkelige.
% "var,---": a comma set close against the rule, with a space before it, as printed

Lad os gaae tilbage og kaste et Blik paa Oldtiden.
Allerede hos Homer findes Forskjellen mellem den fol\-%
kelige Mand og Demagogen, vel ikke udviklet i et Be\-%
greb, men dog beskueliggjort ved Modsætningen mellem
Odysseus og Thersites: Odysseus er folkelig og bryder
sig kun lidt om at være populær; men det, hvorpaa
Thersites sigter, er aabenbart det Populære. Da Aga\-%
memnon (Iliad. II.), for at prøve Folkets Stemning, har
% --- p. 130 ---
\opage{130}\setcounter{footnote}{0}%
aabnet Hæren Udsigt til at komme hjem, strømme Fol\-%
kene ned til Kysten: hvilket Optrin, hvilket Syn! „Som
naar en Kuling, der blæser fra Vest, fremstyrtende hef\-%
tig bølger den frodige Vang, saa Kornet med Vipperne
neier, saa kom Folket i Røre, med gjaldende Skrig de
tilhobe styrted til Skibene hen“. Den stærke Bevægelse
anskueliggjør Hærens uhyre Længsel efter at vende til\-%
bage til Fædrelandet; intet Under, naar man betænker,
hvilken Fristelse der laa i Agamemnons Ord: „Ni af
den høie Kronions kredsrullende Aar er forløbne, frøn\-%
net er Tømret i Skibene alt, og Tougene raadne: mid\-%
lertid sidde jo hist i vort Hjem vore Koner og Smaa\-%
børn, ventende langlig paa os, som her aldeles for Intet
øde vor Kraft paa den Id, for hvilken vi hid ere dragne“.
Agamemnon spiller høit Spil; han tier med det Folke\-%
lige og frister ved det Populære. Odysseus seer med
Harme det larmende Folk forvandle sig til Mængde;
hans Opgave er at standse Bevægelsen og forvandle
Mængden igjen til et Folk; dette kan kun skee ved gud\-%
dommelig Bistand. Ligesom Pallas Athene kom med
en Aabenbaring til Achilles, hvor det gjaldt at bringe
ham personligt til Besindelse, saaledes kommer hun ogsaa
her, hvor det gjælder at bringe hele Folket til Besindelse,
med en Aabenbaring til Odysseus. Nu følger en mester\-%
lig Skildring af, hvorledes Odysseus forstaaer at bøie
sin Stemme, idet han med Mildhed formaner Fyrsterne
og med Strenghed standser de Menige. Saa lykkes det
ham ved sin myndige Optræden at bringe Masserne til
Ro og faae Folket i Tale. „Alle blev siddende nu og
holdt paa Plads sig tilbage; ene den Skvaldrer Thersites
lod høre sin knebbrende Stemme“. Thersites er Dema\-%
gogen i et første raat Udkast. Umiddelbart begynder
han som den Enkelte i Mængden, der gjør Spektakel;
men idet han kommer tilorde, taler han uvilkaarligt i
% --- p. 131 ---
\opage{131}\setcounter{footnote}{0}%
Alles Navn, og melder sig strax som en Fører af
Mængden. Det Hule i hans Væsen giver sig paa en ret
charakteristisk Maade tilkjende derved, at den Tale, han
holder, er en Efterklang af Achills Bebreidelser mod
Agamemnon: „Medens de Andre stride, slide og lide,
har Agamemnon nok af Guld og skjønne Piger“. Talen
stiger til agiterende Pathos: „Feige, forkjælede Folk,
Achaïinder! ei længer Achaier!“ minder om den Achill
% "Achaïinder": the diaeresis on the first i is printed
tilføiede Forhaanelse og ender parodisk: „Achilles er
spag, ei Galde, der boer ham i Livet, ellers, Atreide!
var denne din Haan vist bleven den sidste.“

Demagogens aandelige Grimhed er symbolsk betegnet
ved den legemlige. „Thersites var“, siger Homer, „den grim\-%
meste Mand i Hæren, som drog imod Troia, skjæve var begge
hans Been, og stærkt med det ene han hinked, Ryggen var
puklet, og Skuldrene frem mod Brystet sig bøied, spidst
var hans Hoved foroven, og tyndt sad Haaret om Issen.“
At Folket med de mythiske Guder, paa det mythiske
Bevidsthedstrin, endnu var ufordærvet, tilkjendegives der\-%
ved, at Mængden, da den fik Øinene op, kunde see Ther\-%
siteses Grimhed ligesaavel som Fyrsterne. For Mængdens
Øine og under Mængdens Bifald kunde Odysseus ikke
blot irettesætte den Skvaldrende med haarde Ord, men
endog tilføie ham en corporlig Tugtelse: „Med sin Stav
over Ryg og Skulder han slog ham; vaandt Thersites
sig krymped, og Taarerne brast ham af Øine; under den
gyldene Stav paa hans Ryg en Strime sig hæved blodig
og blaa, da tog han sin Plads med Bæven og Vaanden,
og med et taabeligt Blik han Taarerne visked af Øie.
Hjertelig maatte man lee“ ... Det var en stor Seier, det
Folkelige ved denne Tugtelse vandt over det blot
Populære.

Men lader os ikke glemme, at det netop var ved
Gudernes Bistand, at Seiren blev vunden. Efterhaanden
% --- p. 132 ---
\opage{132}\setcounter{footnote}{0}%
som den politiske Udvikling fortrængte det oprindeligt
Religiøse, forandrede Forholdene sig. Da Folkepartiet
i Athen ved Pisistratidernes Fordrivelse gjenvandt sin
tidligere Indflydelse, indførtes Ostracismen. Ved Ostra\-%
cismen var Mængden berettiget til uden Beviis, endogsaa
uden Beskyldning for nogen Forbrydelse, at fordømme
Fædrelandets mest fortjente Mænd til Landflygtighed;
ved Ostracismen blev det den mest patriotiske Helt
umuligt at virke folkeligt uden at være populær, hvorimod
det, hvad en Aristophanes kan bevidne, blev en politisk
„Pølsemager“ muligt at virke populært, uden at være
folkelig.

Jo dybere Livstanken er, for hvilken der virkes, jo
mere den gjennemtrænger det Grundmenneskelige, desto
stærkere vil ogsaa Modsætningen mellem det Folkelige
og det Populære kunne fremtræde. Vi see det paa
Aabenbaringens Trin; vi see af Fortællingen om Guld\-%
kalven (1 Moseb. XXXII), hvilken Forskjel der i denne
% "1 Moseb." as printed (confirmed at 600 dpi); the golden calf is 2 Mos. XXXII --- author's slip, not corrected
Henseende er mellem Aron og Moses. Moses tøver med
at komme ned fra Bjerget; Folket samler sig om Aron
og siger til ham: „Staa op og gjør os Guder, som kunne
gaae for vort Aasyn; thi denne Moses, Manden, som
førte os op af Ægyptens Land, vide vi ikke, hvad der
er blevet af“. Saa taler Folket, fristet af Kjød og Blod;
det vil affalde fra Jehova, Aandens Gud; det vil affalde
fra sig selv; det er ikke længere et Folk, men en sand\-%
selig raa Hob, en Mængde. Hvad var nu Arons Pligt?
Som folkelig Mand at staae Mængden imod, og see at
komme Folket i Tale! Men i det Sted bøier Præsten
Aron sig under Mængdens Villie og støber Guldkalven.
Da saa Moses kommer ned af Bjerget og skal holde Dom,
udspinder der sig en i høi Grad charakteristisk Ordvexel
mellem ham og Aron. „Hvad har Folket gjort dig, at
du har ført en stor Synd over det?“ Aron: „Ikke blusse
% --- p. 133 ---
\opage{133}\setcounter{footnote}{0}%
min Herres Vrede op, du kjender Folket, at det helder
til det Onde“; med andre Ord: hvad skulde jeg gjøre;
jeg kunde jo ikke sige Mængden imod og gjøre mig
% a small raised tick after "sige", apparently a speck; not transcribed
upopulær. Arons Synd imod Folket er, at han har be\-%
handlet det, ikke som Folk, men som Mængde. At han
selv føler det, høre vi af hans undskyldende Forklaring:
„Og de sagde til mig: gjør os en Gud, som kan gaae
foran os ... Da sagde jeg til dem: hvo som haver Guld,
rive det af sig; og de gave mig det, og jeg kastede det
i Ilden, da kom denne Kalv ud.“ Det Hele foregaaer,
uden at Aron ret veed hvorledes; Kalven kommer ud
af Ilden som ved et Tilfælde: man skulde næsten troe,
at det var gaaet til med Hexeri. Men Moses forstaaer
det anderledes. Ene og med Livsfare stiller han sig ved
Leirens Port, og holder Dom: „Hvo som hører Jehova
til, komme til mig!“ saa lød hans myndige Røst; ja,
den folkelige Mand er en myndig Mand.

% a reader's pencilled note ("Gr"?) in the right margin beside the next line, and a pencilled bracket with a note in the left margin of p. 134; not transcribed
Den folkelige Mand er paa sin Viis en Magthaver
ligesaavel som Fyrsten paa sin Throne og den populære
Mand, der hersker ved Folkegunst; men Forskjellen i
Magt maa bestemmes ved Forskjellen i Myndighed.
Christus siger i Evangeliet (Matth. XX) til sine Disciple:
„I vide, at Folkenes Fyrster herske over dem, og de
Store bruge Myndighed over dem. Men saa skal det
ikke være iblandt Eder; men hvo, som vil være stor
iblandt Eder, han være Eders Tjener; og hvo, som vil
være den Ypperste iblandt Eder, han være Eder under\-%
danig. Ligesom Menneskens Søn er ikke kommen, at
lade sig tjene, men at tjene, og give sit Liv til en Gjen\-%
løsning for Mange.“ I disse Ord er al Magt og Myn\-%
dighed paa Jorden væsentlig charakteriseret. Den po\-%
pulære Mand kan til en given Tid opnaae Magt, men
hans Magt er uden Myndighed; udvortes seet hersker
han vel over Mængden, men i Grunden er det Mængden,
% --- p. 134 ---
\opage{134}\setcounter{footnote}{0}%
der hersker over ham. Som Statens Overhoved har
Fyrsten vel baade Magt og Myndighed; men Magten er
physisk; den paa physisk Magt begrundede Myndighed
er politisk, men ikke i strengere Forstand personlig.
Den personlige Myndighed staaer i ligefremt Forhold til
den personlige Overlegenhed, den personlige Overlegenhed
igjen i ligefremt Forhold til Selvopoffrelsen. Men heraf
følger jo saa, at den høieste folkelige Myndighed maa i
sin inderste Grund nødvendig være religiøs. Folkenes
Frelser er Folkenes Herre, en tornekronet, men dog
guddommelig Majestæt, der alene kan lære dem Lydig\-%
hed. Thi som Magten og Myndigheden er, maa ogsaa
Lydigheden være: mod den physiske Magt er Lydigheden
tvungen Underkastelse, og den høieste Formildelse den,
at Underkastelsen bliver populær; mod den frigjørende
aandelige Magt er Lydigheden fri Underkastelse; en fri
og dog ubetinget Underkastelse er i dybeste Forstand
folkelig, men den er ikke populær.

At den principielle Forskjel imellem det Populære
og det Folkelige grundig erkjendes, er en vigtig Betin\-%
gelse for Nutidens aandelige Liv.

\begin{center}\rule{3cm}{0.4pt}\end{center}

\chapter*{IX.}
\addcontentsline{toc}{chapter}{Forelæsning IX}
\markboth{Forelæsning IX}{Forelæsning IX}
\label{ch:ix}
% pp. 135--153
% --- p. 135 ---
\opage{135}\setcounter{footnote}{0}%
% [large initial M]
Modsætningen mellem Jødedom og Christendom, mellem
Aabenbaringen paa det gamle og det nye Testamentes
Standpunkt, er, og med Rette, bleven bestemt derhen, at
Jødedommen, skjøndt en aabenbaret Religion, dog endnu
kun er Religion for et enkelt Folk, altsaa beheftet med
de Indskrænkninger, som et Folkelivs naturlige Begrænds\-%
ning nødvendig medfører, medens Christendommen ved
at hæve sig over disse Indskrænkninger er bleven, ikke
en Religion blandt flere andre, men Religionen selv i
Aand og Sandhed, ikke Religion for et enkelt Folk, men
% AUDIT: the low mark after "for" (read as a comma by both OCR witnesses) is a thin sliver above the baseline with no head and no descender, unlike every comma on the page (600 dpi): an inked space, not a comma; not typed
for alle Folk, Menneskehedens Religion. Hermed stem\-%
mer det da ogsaa, hvad Paulus siger, at der i Christus
ingen Forskjel er mellem Jøde og Græker, at det Fjendt\-%
lige er forligt, at Christus „gjorde Eet af Begge, og ned\-%
brød Adskillelsens Mellemvæg, da han i sit Kjød afskaf\-%
fede Fjendskabet, Budenes Lov med dens Befalinger, paa
det han ved sig selv kunde skabe de To til eet nyt Men\-%
neske og gjøre Fred“ (Ephes. II). Troen paa Christus
ophæver altsaa den naturlige Forskjel saavel imellem
Folkeslagene indbyrdes som imellem Individerne: „Her
er ikke Jøde eller Græker; her er ikke Træl eller Fri;
her er ikke Mand eller Qvinde; thi I ere alle Een i
Christo Jesu“ (Gal. III). At Paulus her henpeger paa
% --- p. 136 ---
\opage{136}\setcounter{footnote}{0}%
den Troens og Livets Eenhed, hvori alle menneskelige
Følelser og Tanker maae mødes, er klart; men hvorledes
skal Eenheden opfattes? En tom Eenhed kan det umulig
være, saa sandt det virkelig er en Livets Eenhed; thi
Livets Eenhed er uendelig rig paa levendegjørende,
Mangfoldighed og Forskjellighed udfoldende, Kraft.

Hvad er det da for en Eenhed, hvorom Paulus taler?
Det er den Eenhed, som fremkommer, naar Aandskjær\-%
ligheden faaer Magt med det Selviske. Hvor Egoismen
raader, vil den i alle Forhold misbruge de Forskjellig\-%
heder i Gaver og Kræfter, Stand og Vilkaar, som den
naturlige Tingenes Orden medfører. Saaledes vil det
stærkere, af Lykken begunstigede Folk ikke blot under\-%
kue det svagere, men, naar Voldsgjerningen lykkes, endog
smigre sig med at være Styrelsens udvalgte Folk, et Folk,
der kjender sin Mission. Faaer Folkeegoismen nu oven\-%
ikjøbet et Støttepunkt i det Religiøse, saa er der ingen
Grændse for Hadet og Fjendskabet; det er bekjendt,
hvorledes Jøderne paa Christi Tid, i den Indbildning, at
de med „Budenes Lov“ jo vare „det udvalgte Folk“,
hadede Grækere og Romere, og igjen bleve hadede af
disse. Dette Folkenes gjensidige Had og Fjendskab for\-%
svinder, naar det rene, ved Selvopoffrelsen overmægtige
Selv, fortrænger Egoismen. Hvor en uhemmet Egoisme
% a small raised mark (probably an inked space) stands between "Egoismen." and "Hvor"; its identity cannot be read at 300 ppi, so no character is typed
raader i Samfundet, vil den naturlige Forskjel mellem
Mand og Qvinde og den borgerlige Forskjel mellem
Herre og Tjener naturlig medføre, at Manden gjør Qvin\-%
den til Slavinde, og Herren sin Tjener til Træl. Og
Qvinden og Trællen ville begge, uden ædlere Selvfølelse,
opfatte Fornedrelsen egoistisk, bøie sig i Had for den
Stærkeres Ret og i Skjul bag Adskillelsens Mellemvæg
lure paa Hævn. Som Skranke for Fjendskabet og den
lurende Hævn opstilles saa „Budenes Lov“; men Loven
kan brydes: Befrielse er kun i Troen paa ham, som
% --- p. 137 ---
\opage{137}\setcounter{footnote}{0}%
ved at bryde Egoismens Magt nedbrød „Adskillelsens
Mellemvæg“.

Det er de menneskelige Forskjelligheders egoistiske
Opfattelse og selviske Misbrug, om hvis Ophævelse Paulus
taler i de stærke Ord, at der i Christi Rige ikke er Jøde
eller Græker, ikke Træl eller Fri, ikke Mand eller
Qvinde; men at de Alle ere Een. Men at Christendom\-%
men skulde forholde sig ligegyldig til de i naturlige Anlæg
og naturlige Forhold begrundede Forskjelligheder, enten
imellem Individerne indbyrdes eller mellem Folkeslagene,
kan umulig være Pauli Mening. Forstaaer man hans
Ord saaledes, at Evangeliet enten ganske overseer eller
endog lægger an paa at udslette de nationale Forskjel\-%
ligheder, da maa man jo ogsaa antage, at Evangeliet
enten ganske overseer eller endog lægger an paa at ud\-%
slette den væsentlige, ikke alene legemlige, men ogsaa
sjælelige Forskjel mellem Mand og Qvinde. De, der
% [a reader has underlined "De, der forstaae ... Eiendommeligheder" in pencil; not letterspaced in the print; not transcribed]
forstaae Christendommen saaledes, at den som Menne\-%
skehedens Religion kun seer paa det abstract Menneske\-%
lige og kosmopolitisk oversvæver alle folkelige Eiendom\-%
meligheder, misforstaae den ligesaa vist som de, der
gjøre dens Virkninger afhængige af Folkenaturerne i den
Grad, at hiin Adskillelsens Mellemvæg, som Christus
nedbrød, mures op paany.

Evangeliets Aand vil rense, forynge og styrke, med
eet Ord: \textit{gjenføde} det Grundmenneskelige. Hvad er det
Grundmenneskelige? Det er det Personlige! I Aner\-%
kjendelse af Personligheden falder det Religiøse paa en
Maade sammen med det Politiske. Kun som Person
% [a mark over the u of "Kun" and over the o of "Troessætningerne" below is a reader's pencil/speck, not an accent; not transcribed]
faaer et Menneske Borgerret i et jordisk Samfund; kun
som Person faaer det Borgerret i Himlen. Det Spørgs\-%
maal, hvorledes Reformationsprincipet i Troessætningerne
fra 1530 og Revolutionsprincipet i Grundsætningerne fra
% --- p. 138 ---
\opage{138}\setcounter{footnote}{0}%
1789 skulle forliges, falder sammen med Spørgsmaalet
om Personlighedsprincipets Udvikling.

Forskjellen imellem den statsborgerlige og den evan\-%
geliske Opfattelse af Personlighedens Væsen er nu den,
at hiin anerkjender Personligheden i dens umiddelbare,
naturlige Tilværen og fordrer blot en Begrændsning og
Forædling af dens Egoisme; medens Evangeliet derimod
vil gribe Personligheden i dens inderste aandelige Rod,
og forhøie dens Selvstændighed ved at udrense dens
Egoisme. Kan der nu ikke være Tvivl om, hvilken af
de to Opfattelser der er den overordnede, og hvilken der
er den underordnede, saa kan der heller ikke være Tvivl
om, at den statsborgerlige maa underordnes den evange\-%
liske, og det oprindelige Eftertryk falde paa den evan\-%
geliske Opfattelse.

Den evangeliske Opfattelse kan henføres til følgende
Grundbestemmelser: den menneskelige Personlighed har
% [pencil underlining of "den menneskelige Personlighed ... Troesartikler" and a pencilled "13" in the margin; not transcribed]
den guddommelige 1) til sit Ophav, 2) til sit Forbillede,
3) til sin væsentlige Bærer og Bevæger. At udvikle
disse Grundbestemmelser er at udvikle de apostoliske
Troesartikler.

Tankeudviklingen i den rene, almindelige Form maa
imidlertid ikke forvexles med Tænkningen i det virkelige
Troesliv: Et er det at tænke Troen, et Andet at leve i
Troen. Det er i nærværende Sammenhæng ikke vor Opgave
at aflægge en Troesbekjendelse; men det er udtrykkelig vor
Opgave at tænke Troen, for at tænke Personligheden.

Hvilken Forskjel der er imellem Troens og Videns
Tankegang, bliver os ret anskueligt, naar det gjælder om
Udvikling af Personlighedsprincipet. Allerede Stand\-%
punktet med den første Forudsætnings reent menneskelige
Side bliver bestemt paa en ganske anden Maade i Troen
end i Videnskaben. Fordringen er, at Mennesket, for at
vide sig som Person og hævde sin Personlighed, maa
% --- p. 139 ---
\opage{139}\setcounter{footnote}{0}%
lægge Vægt paa sig selv. Hvormegen Vægt tør et Men\-%
neske da i Personlighedens Navn vel lægge paa sit eget
Selv? Dette kan umulig afgjøres ved subjectiv Vilkaar\-%
lighed; thi i saa Fald maatte jo de mest opblæste,
anmassende og indbildske Individer være Personlighedens
værdigste Repræsentanter. Nei, den sande Personligheds\-%
fordring maa bestemmes ved Selvhedens eget Princip.
Men naar Spørgsmaalet lyder: hvilken Selvstændighed
kan der oprindelig tillægges det menneskelige Selv? saa
viser det sig, at Viden og Tro maae gaae hver sine
Veie.

Skal Sagen afgjøres i Viden, vil Selvets principielle
Selvstændighed vise sig at være afhængig af Organismen,
af organisk-physiologiske Naturbetingelser. Men raadføre
vi os grundig med Physiologien, da synker Selvet saa
dybt i Betingelsernes Kreds, at det neppe kan gjøre sig
Haab om at overleve Døden. Maaskee kan det ved
Philosophiens Hjælp hæve sig igjen i Psychologien? Ja
vel kan det hæve sig, men kun til Væsenseenhed med Ideen,
en tvetydig Øieblikseenhed, der snart faaer Navn af en
nærværende Evighed, snart, mere langtrukkent abstract,
af en metaphysisk Udødelighed.

Anderledes i Troen. Det første for Troeserkjendel\-%
sen charakteristiske Kjendemærke er just dette, at Selvets
principielle Selvstændighed ikke gjøres afhængig af orga\-%
niske Betingelser, men ene og alene af Selvets eget Væ\-%
sen, et Væsen, der kommer tilsyne i Selvets Concentra\-%
tion til absolut Selvbestemmelse. Hvorledes denne Selv\-%
bestemmelse, menneskelig seet, maa tage sig ud, kan let
paavises. Det troende Selv er kjendeligt paa sin eien\-%
dommelige Livsfordring; hvad fordrer det? Et evigt
Liv i Aand og Sandhed! Med denne sin ufravigelige
Fordring, paa denne Selvbekræftelsens høieste Spidse,
er det troende Selv let at kjende fra et vidende Subject.
% --- p. 140 ---
\opage{140}\setcounter{footnote}{0}%
Den ufravigelige Fordring er ingen Anmasselse; den er
tvertimod et Udtryk for Sjælens inderste Længsel, dens
allerdybeste Trang. De Bestemmelser, der ikke ved\-%
komme Fordringen selv, vedkomme ikke Troen; de Be\-%
stemmelser derimod, der ere nødvendige til Fordringens
Tilfredsstillelse, høre oprindelig med i Troesbegrebet og
ere afgjørende for Troens Tænkning. Da her altsaa kan
være Tale ligesaa vel om Troesbegreber som om Videns\-%
begreber, vil Modsigelsens Grundsætning ligesaavel
spille en Rolle i Troens som i Videns Tankeudvikling;
kun at Modsigelserne, ifølge de ueensartede Principer,
maae blive ueensartede. Alle Videnskabens Modsigelser
ere, i deres Grundcharakteer, objective Begrebsmodsigel\-%
ser; det er enten et Begreb, der modsiger et andet, eller
et Begreb, der modsiger sig selv. Saaledes er, f. Ex.,
en rund Fiirkant en objectiv væsentlig Begrebsmodsigelse.
Anderledes med Troens Modsigelser; de kjendes, ifølge
Grundcharakteren, alle derpaa, at Selvet ved dem, mid\-%
delbart eller umiddelbart, modsiger sig selv: Troesmod\-%
sigelserne ere i dybeste Forstand personlige Selvmod\-%
sigelser. Som en Modsigelse er, maa ogsaa dens Løs\-%
ning være: Viden kan ligesaa lidt løse Troens, som Troen
kan løse Videns Modsigelser. Hermed er den for Ud\-%
viklingen af Personlighedsprincipet nødvendige Forud\-%
sætning bestemt; vi skulle nu i almindelige Omrids angive
Udviklingens Gang.

% section head, set in italic (numeral included)
\textit{1) Den menneskelige Personlighed har den guddomme\-%
lige til sit Ophav.}

Hvor uenige Troen og Viden maae være om dette
Punkt, er allerede klart, af hvad vi i en tidligere Sam\-%
menhæng have fremhævet. Naar Videnskaben ikke kan
anerkjende Guds Personlighed, fordi Begrebet af en
absolut Personlighed er et sig selv modsigende Begreb:
hvorledes skulde den da kunne indrømme, at den per\-%
% --- p. 141 ---
\opage{141}\setcounter{footnote}{0}%
sonlige Gud kan være den menneskelige Personligheds
Ophav? Nu er det jo ganske vist, at Videnskaben hvert
Øieblik forandrer sit Gudsbegreb: enhver original Philo\-%
sophie, den atheistiske ikke undtagen, har sit eget origi\-%
nale Begreb om Gud. Hos Spinoza er Gud den ene
uendelige Substans, hos Leibnitz Monadernes Monade,
hos Fichte den moralske Verdensorden, hos Schelling ---
idetmindste for en Tid --- det Absolut - Absolute, hos
% "Absolut - Absolute": the hyphen is printed with a space on either side (600 dpi); transcribed as printed
Hegel „Begrebet“, hos Feuerbach Nul; den, der ret vil
erkjende Tidernes philosophiske Gud, maa kjende den
metaphysiske Cours, ligerviis som den, der vil taxere
Aarets Afgrøde, maa kjende Capitelstaxten.

At forlade Philosophien og tye til Theologien er kun
at gjøre Ondt værre; det er at røre Troens og Videns
Modsigelser sammen i Eet og plumre det Hele. De
eneste Elementer, som med nogen Anskuelighed hæve
sig frem af en saadan Sammenrøren, ere en Mængde store
Ord, der allesammen ende paa \textit{isme}: Atheisme, Akosmisme,
Pantheisme, Polytheisme, Monotheisme, Deisme, Theisme
o. dsl. Men Personlighedsbegrebet bliver, hvormange \textit{Ismer}
man end kan drive det til, ikke klarere.

Vi holde os da til Troen. Enhver Bestemmelse i et
Troesbegreb bliver, ifølge Standpunktet, at angive kort
og skarpt, uden alle vidtløftige Omsvøb: Correctivet
ligger i Selvmodsigelsen og dens Løsning. Hvoraf veed
den Troende, at Gud er personlig? Han \textit{veed} det jo
ikke, men han \textit{troer} det! Saa er vel hans Forvisning i
Troen svagere end den, han kunde have erholdt i Viden?
Ingenlunde! den er tvertimod langt stærkere. At negte
Guds Personlighed eller endog blot drage den i Tvivl,
vilde for den Troende være en ligesaa oprørende Selv\-%
modsigelse, som det for den Vidende er en oprørende
Begrebsmodsigelse, at det Runde skulde være fiirkantet.
Har Selvets personlige Fordring Gyldighed, da kan det,
% --- p. 142 ---
\opage{142}\setcounter{footnote}{0}%
trods sin relative Afhængighed af Organisme og organiske
Betingelser, umuligt aflede sin Tilværelse af en selvløs
Natur. Med Naturlivet til Grundlag maa det, for at leve
et evigt Liv, nødvendigt have en Selvhedsspire fra det
evige, absolute Selv, d. e. fra den personlige Gud.

Jo mere levende Selvfølelsen er, desto smerteligere
føles den skærende Modsætning mellem Menneskets Per\-%
sonlighed og det upersonlige Naturvæsen: en afmægtig
Frihed staaer her ligeoverfor en overmægtig Nødvendig\-%
hed. Skal Friheden i Sandhed erkjendes for Altilvæ\-%
relsens Princip, da maa Naturriget, den hele naturlige
Verden være absolut afhængig af Gud, et Værk af hans
frie, almægtige Villie: at troe paa Selvets evige Ophav,
er at troe paa „Himlens og Jordens almægtige Skaber“.

Ved Spørgsmaalet om Natur og Skabelse bliver For\-%
skjellen mellem Troens og Videns Væsen ret iøinefal\-%
dende. For Vidensprincipet er Naturen som Natur det
Første, og Skabelsesbegrebet maa som Følge heraf rette
sig efter Naturbegrebet. Kunde Videnskaben nu ved at
lægge sit Naturbegreb til Grund udvikle, klare og tyde\-%
liggjøre det religiøse Skabelsesbegreb: da vilde den med
det Samme være istand til at udvikle alle de andre reli\-%
giøse Grundbegreber; Theologien vilde da forvandles til
Philosophie, og den videnskabelige Forsoning af Tro og
Viden fuldbyrdes derved, at Troen blev opløst i Viden.
For Troesprincipet derimod er Skabelsen som Skabelse
det Første, og Naturbegrebet maa som Følge heraf rette
sig efter Skabelsesbegrebet. Kunde Troeslæren nu ved
at lægge sit Skabelsesbegreb til Grund udvikle, klare og
tydeliggjøre det videnskabelige Naturbegreb: da vilde den
med det Samme være istand til at udvikle alle de andre
Naturens Grundbegreber; den videnskabelige For\-%
soning af Tro og Viden vilde da være at søge og at
finde i „den christelige Dogmatik“; og hvad saa?
% --- p. 143 ---
\opage{143}\setcounter{footnote}{0}%
Ja, saa havde vi --- det theologiske Universalmonarchie!
Men det Sidste er, ifølge den menneskelige Erkjendelses
Vilkaar ligesaa umuligt som det Første (Jvnfr. „Den gode
Villie som Magt i Videnskaben“ S. 30 o. flg.).

Da nu ingen af Delene er muligt, saa kan For\-%
soningen kun tilveiebringes paa det Vilkaar, at den
menneskelige Viden erkjender sin absolute Grændse, og
at Troen som et i sig selvstændigt Princip faaer Ret til
at udvikle sine egne Begreber.

„Men“ --- indvender man --- „det kan vel for et
% a small raised mark (probably an inked space, as on p. 136) follows "man"; not typed
Øieblik tage sig ud med en saadan Fordobling, som om
det var Noget; men, see vi nøiere til, saa er det dog
ikke Noget. Videnskaben indrømmer aldrig, at der
udenfor dens Enemærker skal kunne gives en i sig sam\-%
menhængende, conseqvent Tænkning; navnlig ville de
klare Hoveder, de rette, skarpt tænkende Metaphysikere,
Ihændehaverne af den rene Philosophie, snart bevise os,
at den menneskelige Viden umulig kan have nogen abso\-%
lut Grændse, altsaa heller ikke indrømme nogen Tilvæ\-%
relsens Gaade, altsaa heller ikke anerkjende Troen for
noget i sig selvstændigt Tankeprincip.“ Vi svare, at
dersom der virkelig opstaaer en Metaphysiker, som til
Verdens store Forbauselse beviser og gjør det vitterligt
for Alle, som kunne tænke, at den menneskelige Viden
ikke har nogen absolut Grændse: da skulle vi heller ikke
undlade pligtskyldigst at blive forbausede. Men indtil
dette skeer, benytte vi Forventningens Mellemtid til at
indprente os følgende Grundfordringer, som den Viden,
% printer's defect: in "den Viden," the "d" stands on the line while "en Viden," has dropped below it (displaced type, 600 dpi); the words are legible and transcribed as read
der erklærer den absolute Grændse for ophævet, uden
Omsvøb maa kunne fyldestgjøre.

En saadan Vidende maa vise, hvorledes han bærer
sig ad med: a) at tænke, ikke blot \textit{successivt}, først den
ene Tanke og saa den anden, men simultant, d. v. s.
utallige tydelige Tanker paa een Gang, thi saaledes tænker
% --- p. 144 ---
\opage{144}\setcounter{footnote}{0}%
en Guddom; b) at gjennemskue Materien og det Stof\-%
agtige saaledes, at han ved sin intellectuelle Magt --- at
begribe er jo at magte --- bliver istand til at forvandle
Begreb til Existens, og Existensen igjen til Begreb,
Hestens Begreb til Hestens Existens, Guldets Begreb til
virkeligt Guld; thi at oversee Hesten og ride paa Be\-%
grebet, beholde Guldet selv og betale med Begrebet, er
kun Udflugter, og slige Udflugter ere en metaphysisk
Guddom uværdige; c) at gjennemskue Eenheden af
Hjernevirksomheden og den tilsvarende Tænkning, af
den bevægende Villie og den ved Villien bevægede Nerve
saa klart og gjøre Vexelvirkningen mellem Sjæl og Le\-%
geme saa indlysende, saa tydelig og begribelig, at endog
Physiologerne opgive den indviklede og besværlige Empirie
for at kaste sig paa Metaphysiken; o. s. v.

Hvad vi her fordre, vil den rene Metaphysiker maaskee
ansee for Bagateller og Smaapillerier; thi naar man først
ved Hjælp af Abstractioner har bundfældet Existensen
som et Begrebets Grums, saa falde Tilværelsens Magt\-%
bestemmelser naturligviis bort af sig selv, og naar Magt\-%
bestemmelserne falde bort, kan en afmægtig Viden jo let
spille Guddom. Eller skulde maaskee --- thi hvad den
rene Metaphysik angaaer, maa man være belavet paa det
Utrolige --- den afmægtige Viden engang ved „fine Be\-%
mærkninger“ kunne drive det saavidt, at den endog
deducerede, siger og skriver: deducerede, ikke Magtens
almene Begreb --- thi Magtens Begreb in abstracto er
ofte nok deduceret --- men Magten selv, den virkelige
Magt, siger og skriver: den virkelige Magt med de vir\-%
kelige Magtbestemmelser? Ja, hvis det engang skeer,
saa skeer der et metaphysisk Mirakel. Man dømme om
et saadant Mirakel, som man vil; men ethvert fornuftigt
Menneske maa dog vel indrømme os, at de opstillede
Fordringer have deres Betydning, og at det dog vilde
% --- p. 145 ---
\opage{145}\setcounter{footnote}{0}%
være ret interessant, om man ved Hjælp af den rene
Viden engang kunde komme efter, hvorledes det i en
speculativ Hjerne gaaer til med de metaphysiske
Mirakler.

Dog, om ogsaa det Forbausende skeer, at nemlig
den rene afmægtige Viden faaer Magt med Virkeligheden
selv og deducerer, ikke Magtens ideale Nødvendighed ---
en saadan Nødvendighed er ofte nok deduceret --- men
Magtens \textit{umiddelbare Udslag} i virkelig Existens; hvad
saa? Saa er der dog een Ting, der i denne Sammen\-%
hæng kan siges med afgjørende Bestemthed, og det er:
at en saadan metaphysisk Deduction vil Folket, og det
er nu netop om Folkets Oplysning vi her tale, aldrig
faae ind i sin Hjerne; og det af den gode naturlige
Grund, at en folkelig Hjerne ikke er metaphysisk. Med
andre Ord: hvis det religiøse Skabelsesbegreb ved et
Afmagtens Magtsprog skal opløses, affiltreres, deduceres,
paralyseres og omkalfatres til Vidensbegreb, saa at det
ikke længere faaer Lov at gjælde for, hvad det op\-%
rindelig er og vil være: et Troesbegreb, da gaaer det
folkelige Lys ud; da maae de Uvidenskabelige lære at troe,
ikke paa det aabenbarede Ord, men paa Metaphysikens
Deductioner. De faae naturligviis ikke Viden selv, men
kun et Skin af Viden i det Populære.

At ombytte Troen paa Aabenbaringen med Troen
paa Metaphysiken er et Misgreb, der vil hindre det
aandelige Liv, ikke blot i Nutiden, men sikkerlig ogsaa
i Fremtiden. Skulle Aandskræfterne vækkes, og det
Grundmenneskelige røre sig i al sin Oprindelighed, da
er Betingelsen, den ufravigelige Betingelse, den, at Troes\-%
principet gjør sin Ret gjældende, at Aabenbaringens
Grundbegreber, Frihedsbegrebet, Skabelsesbegrebet og
først og sidst \textit{Personlighedsbegrebet}, udvikle sig i en Form,
% --- p. 146 ---
\opage{146}\setcounter{footnote}{0}%
der paa een Gang svarer til den religiøse og folkelige
Tænkning. „I skulle“, siger Christus, „alle være Theodi\-%
dakter“, d. v. s. underviste af Gud. Disse Ord kunne vel i
nærværende Sammenhæng siges at indeholde en alvorlig
Paamindelse. Thi saa umuligt det er at faae Aandslivet
i et Folk udviklet og klaret ved populære Efterskin af
Videnskabernes Fakler, ligesaa nødvendigt er det at faae
den folkelige Underviisning grundet paa personlig Indsigt i
Personlighedens Væsen. Personlighedstanken er paa een
Gang en folkelig og religiøs Grundtanke. Jo rigere og
dybere denne Grundtanke udvikler sig i et Folks Bevidsthed,
desto rigere og dybere vil ogsaa den guddommelige Aands
Indvirkninger paa Folke-Aanden være; jo rigere og dybere
enhver Folke-Aand efter sit Maal tilegner sig den guddom\-%
melige Underviisning, desto mere vil Hadet svinde, Adskil\-%
lelsens Mellemvæg gjennembrydes, og den grundmenne\-%
skelige Eenhedstanke seire ved Ham, der kom „at gjøre
Fred“.

% [section head, italic, indented]
\textit{2) Den menneskelige Personlighed har den guddomme\-%
lige til Forbillede.}

Culturens Oplyste ville maaskee af særligt Hensyn
til det store Fleertal i Befolkningen indrømme os, at en
religiøs Udvikling af Personlighedsprincipet vel kunde
have sin Betydning, naar det gjælder at vække Bevidst\-%
heden om det Grundmenneskelige; men, tilføie de, med
den udtrykkelige Indskrænkning, at man holder sig til
den første Artikel, der udtrykker det Almeenreligiøse,
og vogter sig for at røre ved den anden. „Istedenfor
at fremstille det Almene, fremstiller den anden Artikel
jo netop det aldeles Særegne, ja, noget saa Særegent, at
end ikke de Christne indbyrdes have kunnet enes om
Forstaaelsen“. At denne Artikel, istedenfor at være en
almindelig Enighedsartikel, tvertimod er bleven en Stri\-%
dens, Splidens og Adskillelsens Artikel, har efter disse
% --- p. 147 ---
\opage{147}\setcounter{footnote}{0}%
Dannedes Mening sin naturlige Grund deri, at et enkelt
i Kjød og Blod existerende Menneske er blevet opstillet
som Normalmenneske, ja som Gudmenneske gjort til
Gjenstand for Tro og Tilbedelse.

Indvendingen beroer paa en stor Misforstaaelse.
Gjennem den anden Artikel gaaer den samme, den selv\-%
samme personlige Eenhedstanke som igjennem den første.
Det vil ved nærmere Overveielse vise sig, at den første
Artikel først faaer sin fulde Betydning ved at sees i Be\-%
lysning af den anden.

Saae vi i første Artikel Modsætningen mellem Men\-%
neskets indskrænkede Personlighed og Naturmagten, saa
maae vi nu i den anden Artikel besinde os paa en ny
Modsætning, den aandelige Modsætning mellem Egoismen
og den hellige Personlighed. Denne Modsætning viser
udtrykkelig hen til en Modsigelse i Mennesket selv. Det
er Samvittighedens første religiøse Vidnesbyrd, at Men\-%
nesket indtræder som Skyldner i Aandens Rige: den
første Bevidsthed om Aanden som Aand gaaer i Eet
med Bevidsthed om Skyld. Selvmodsigelsen er, at den
virkelige Personlighed staaer i Strid med sit Ideal, at
det Grundmenneskelige er forvansket, at Gudsbilledet er
tabt. Hvordan skal denne Selvmodsigelse hæves? Philo\-%
sophien maa, forsaavidt den da ikke vil fornegte ethvert
Glimt af Aand, være enig med Religionen i, at det vir\-%
kelige, existerende Menneske er i Splid med sit eget
Ideal. Og dog er det et Aandens Krav, at Mennesket
skal dømmes af sit Ideal. Fordringen er retfærdig, men
den er streng: der er noget Gruopvækkende i at see en
Adam Homo paa Synderbænken og høre hans Defensor
sige, at han „appellerer til Idealet“.

Enigheden mellem Religion og Philosophie ophører
imidlertid i det Øieblik, der spørges om Idealets Vir\-%
keliggjørelse. Striden herom egner sig ret til at
% --- p. 148 ---
\opage{148}\setcounter{footnote}{0}%
oplyse den forskjellige Maade, hvorpaa Modsigel\-%
sens Grundsætning anvendes af Troen og af Viden.
Forsaavidt Modsætningen mellem Ideal og Virkelighed
opfattes fra Videnskabens Synspunkt, lægges Eftertrykket
paa Virkelighedsbegrebet, og Spørgsmaalet bliver saa,
om Idealet kan virkeliggjøres. Viser det sig nu, at det
existerende Menneske ikke kan virkeliggjøre sit Ideal,
hvad saa? Ja, saa kan jo Virkeliggjørelsen, strengt
taget, heller ikke fordres af Mennesket. Mennesket i
Kjød og Blod behøver altsaa ikke at tage sig Fordringen
altfor nær; naar det virkelig ikke kan opfylde den, er
det jo lovlig undskyldt. Ved at komme i Modsigelse
med Virkelighedsbegrebet kommer Idealets Begreb i ud\-%
trykkelig Modsigelse med sig selv; thi et Ideal, der
fordrer Virkeliggjørelse uden at kunne virkeliggjøres, er
en Modsigelse. Anderledes i det Religiøse. Idet Mod\-%
sætningen mellem Menneskets personlige Ideal og det i
Kjød og Blod existerende Menneske opfattes fra Reli\-%
gionens Synspunkt, lægges det hele afgjørende Eftertryk
paa den ideale Fordring: Du \textit{skal} virkeliggjøre dit Ideal,
her gjælder ingen Undskyldning, \textit{Du skal!} Istedenfor at
lægge Modsigelsen ind i Idealet, lægger Religionen, for
at hævde Personlighedens Ret, Modsigelsen ind i det
menneskelige Selv og appellerer til Samvittigheden. Vil
nu Mennesket unddrage sig Forpligtelsen, saa føler det
en smertelig Selvmodsigelse; thi Idealet skal, det er
guddommelig Alvor, Idealet skal virkeliggjøres. Aner\-%
kjender Mennesket Forpligtelsen, saa føler det Selvmod\-%
sigelsens dybeste Smerte; thi det kan jo i denne Virke\-%
lighed, med „disse Lemmer“, i dette til Døden dømte
Legeme, umulig opfylde Forpligtelsen. Det er en Smerte,
som skulde Personligheden sprænges: hvad kan da Men\-%
nesket? Det kan i sin Smerte forstaae Apostelens Klage:
% --- p. 149 ---
\opage{149}\setcounter{footnote}{0}%
„Jeg elendige Menneske! hvo skal frie mig fra dette
Dødens Legeme?“

Naar det virkeliggjorte Ideal, det i Kjød og Blod
existerende Forbillede betragtes i Videns Lys, siger Kri\-%
tiken conseqvent: et saadant Idealmenneske er umuligt,
følgelig er det uvirkeligt. Naar det evangeliske Forbil\-%
lede betragtes i Troens Lys, siger Kritiken conseqvent:
et saadant Idealmenneske er nødvendigt, følgelig er det
virkeligt.

At Viden bøier sig under den endelige Virkeligheds
% Printed "Fordring." with a point where the sense wants a comma
% (a comma whose tail did not print?); transcribed as printed.
Magt som under en af Personlighedens ideale Fordring.
uafhængig Naturmagt, er jo ganske i sin Orden, efterdi
den, ifølge sit Princip, har maattet sætte Naturbegrebet
over Skabelsesbegrebet: Videnskaben kan ikke anerkjende
den anden Artikel, fordi den ikke har kunnet anerkjende
den første. Med Troen forholder det sig omvendt. Idet
% Printed "hehøver" (h for b; the first letter is an open-arched h,
% compare the closed b of "betingelser" directly below); read "behøver".
Troestanken bøier sig for Idealets hellige Magt, hehøver
den ikke at bøie af for sandselige Naturbetingelser.
Ved at sætte Skabelsesbegrebet over Naturbegrebet har
den gjort de sandselige Naturbetingelser afhængige af
den Almægtiges personlige Villie: ved at anerkjende den
første Artikel har den banet sig Vei til at anerkjende
den anden.

Gjennemgaae vi den anden Artikel Led for Led, da
overbevise vi os om, at ethvert Led netop svarer til
Personlighedsprincipets ideale Fordring. Skal Forløse\-%
rens Person være Gjenstand for Viden, saa er den hele
Artikel en phantastisk Sammensætning af indbyrdes mod\-%
stridende Elementer: „Undfangen af Aanden og født af
Jomfruen“, og dog „piint under Pontius Pilatus“; „korsfæ\-%
stet, død og begraven“, og dog „opstanden fra de Døde og
opfaren til Himlen“ \ldots{} hvilke Urimeligheder! Bliver For\-%
løserens Person derimod opfattet som Gjenstand for
% --- p. 150 ---
\opage{150}\setcounter{footnote}{0}%
Tro, saa følger det ene Led med indre Conseqvens af
det andet.

Som den, der i et virkeligt Menneskeliv skal opfylde
Personlighedens hellige Lov, maa Forbilledet have en
fuldkommen reen og stærk, af Hellighedens Aand bevæ\-%
get Villie; men for at et virkeligt Menneske skal kunne
have en reen og stærk og hellig Villie, maa han und\-%
fanges af Aanden og fødes i Reenhed. „Det er urime\-%
ligt, umuligt, i Modsigelse med al Virkelighed!“ raaber
den Vidende. Men Troen frygter ikke Videns Modsi\-%
gelser; den ændser ikke og kan, ifølge sit Princip, ikke
ændse de Modsigelser, der udspringe af Hensyn til den
materielle Virkelighed; den har i sit Skabelsesbegreb et
sikkert Værn mod alle slige Modsigelser. Derimod
frygter Troen for een Modsigelse: det er den personlige
Selvmodsigelse, thi i den er Døden. Men nu vilde det
jo være en Modsigelse i Personlighedens Væsen, dersom
dens ideale Fordringer ikke lode sig virkeliggjøre, og
en ligesaa skærende Modsigelse, at et Menneske kunde
virkeliggjøre disse Fordringer uden en reen og hellig
Villie, og da endelig ogsaa en Modsigelse, at et Menne\-%
ske skulde kunne have en reen og hellig Villie, skjøndt
han var, som David siger i sin Angers Smerte, „undfangen
i Misgjerning og født i Synd“.

Men nu Idealets virkelige Liv i en sandselig, virkelig
Verden? Ja, er det ikke mærkværdigt, at det indholds\-%
rigeste Levnetsløb kan fattes og fremstilles i saa faa og
enfoldige Ord? „Piint under Pontius Pilatus, korsfæstet,
død og begraven“; i disse Ord er Idealets virkelige Liv
sammenfattet; thi Ordene anskueliggjøre hans Lidelse,
og Idealets Liv paa Jorden var Liden og Lidelse fra Først
til Sidst. Striden mellem Ideal og Virkelighed er altsaa
hævet i Lidelsen: Selvmodsigelsen svandt, da Lidelsen
blev sat i dens Sted.

% --- p. 151 ---
\opage{151}\setcounter{footnote}{0}%
Hvad Nedfarten til Hades (descendit ad inferos)
angaaer, kan dette Led ret egne sig til at vise,
hvad der i Religionslæren kommer ud af at gjøre
Christus til Gjenstand for Viden. Theologernes Stridig\-%
heder om hans overnaturlige Fødsel, hans to Naturer,
hans to Villier, hans Forening af Afmagt med Almagt
paa Korset ere vel i det Hele forargelige og meningsløse;
men naar Striden dreier sig om „Nedfarten til Helvede“,
bliver Huulheden i den videnskabelige Phantasering mod\-%
bydelig. Opfattes Forbilledet derimod som Troens Gjen\-%
stand, da er hans Nærværelse i de Dødes Rige en An\-%
skueliggjørelse af hans virkelige Død, hans Seier over
Døden og Frelsen for Dødens Fanger.

Hvor fuldstændig Personlighedens Ideal har seiret over
Død og Forkrænkelighed, fremstilles paa Aabenbaringens
anskuelige Maade ved den personlige Opstandelse og den
personlige Himmelfart. Den, der opstod, er personlig den
Samme som den, der blev begraven; den, der ophøiedes,
er den Samme, personlig den Samme, som den, der var
i Fornedrelsen. Idealets Livsløb er et fuldstændigt Kreds\-%
løb; saa forstaae vi hans Ord, naar han siger om sig
selv: „Jeg udgik fra Faderen og kom til Verden; jeg
forlader Verden igjen og gaaer til Faderen“.

Det er Aandens Lov, at de virkelige Personligheder
skulle dømmes af deres sande Ideal; Personlighedens
virkeliggjorte Ideal er altsaa Menneskehedens Dommer:
„Han skal komme igjen at dømme Levende og Døde“.

% [section head, italic, indented]
\textit{3) Den menneskelige Personlighed har den guddommelige
til sin væsentlige Bærer og Bevæger.}

Af den forskjellige Maade, hvorpaa Forholdet imel\-%
lem Virkeligheden og Idealet opfattes i Viden og i Troen,
følger med Nødvendighed, at Aanden i Videnskaben maa
være en anden end i Religionen. Aandsbegrebet bestem\-%
mes ved Idealets Begreb: for at forstaae den anden
% --- p. 152 ---
\opage{152}\setcounter{footnote}{0}%
Artikel, maae vi lære at forstaae den tredie. Da Viden\-%
skaben, ifølge sit Princip, ikke formaaer at erkjende det
virkeliggjorte Personlighedsideal, kan den af samme
Grund heller ikke erkjende den guddommelige Person\-%
lighed i Aanden selv. For den Vidende er Aanden kun
en Eenhed af Ideen og Bevidstheden; for den Troende
er Ideen derimod et Mellemvæsen imellem hans eget
menneskelige og det i ham virkende guddommelige Jeg;
den Vidende forholder sig gjennem Ideen til sig selv;
den Troende forholder sig gjennem Ideen (Sandhedens,
Hellighedens, Retfærdighedens Idee o. s. v.) til Gud.

At det guddommelige Selv har virkeliggjort Person\-%
lighedens Ideal paa Jorden, kunde umuligt komme de
endelige Personligheder til Gode, dersom ikke Selvet
fra Gud tillige vilde virke til, at Idealet relativt kan
virkeliggjøres i ethvert Menneske. Ved blot at op\-%
fatte Forbilledet som historisk Person ere vi endnu
saa langt fra at faae Selvmodsigelsen løst, at vi
meget mere faae den forstærket. Naar Personligheds\-%
idealet er virkeliggjort i dette ene Menneske, saa staaer
dette ene Menneske jo for Anskuelsen som den, der
hævder de ideale Fordringers uafviselige Gyldighed; han
opfylder Selvhedens Lov, og Selvhedsloven skal opfyldes.
„Mener ikke, at jeg er kommen at opløse Loven eller
Propheterne: jeg er ikke kommen at opløse, men at fuld\-%
komme. Thi sandelig siger jeg Eder: indtil Himmelen
og Jorden forgaaer, skal end ikke det mindste Bogstav
eller en Tøddel forgaae af Loven, indtil det skeer alt\-%
sammen.“ Men ved paa denne Maade at faae Idealet virke\-%
liggjort i den fuldkomne Personlighed have de ufuldkomne
Personligheder faaet det virkelige Ideal udenfor sig, og staae
da virkelig udenfor Idealet. Skulde de i en saadan Stilling
dømmes af Idealet, vilde Dommen blive en Fordømmelse.
Modsigelsen kan kun løses derved, at det guddommelige
% --- p. 153 ---
\opage{153}\setcounter{footnote}{0}%
Selv ikke indskrænker sig til en udvortes Fremstilling af
det virkeliggjorte Ideal, men udtrykkelig sætter den ud\-%
vortes Virkeliggjørelse som Forudsætning for den Virk\-%
somhed, hvormed det som helliggjørende Aand i hvert
enkelt Menneske gjentager og former det oprindelige
Forbillede. Idet de endelige, ufuldkomne Personligheder
paavirkes og bevæges af den ene og samme Personlig\-%
hedens Aand, forsamles Menigheden, og Forbilledet for\-%
klares. Uden det virkelige Forbillede kan Aanden ikke
udvikle det Grundmenneskelige i de enkelte Individer;
thi Udviklingen af det Grundmenneskelige falder sammen
med Forbilledets Tilegnelse. Derfor siger Christus til
Disciplene (Joh. XVI): „Jeg haver endnu Meget at
sige Eder; men nu kunne I ikke bære det. Men naar
han den Sandheds Aand kommer, skal han veilede Eder
til al Sandhed; thi han skal ikke tale af sig selv. Han
skal herliggjøre mig; thi han skal tage af mit og for\-%
kynde Eder.“

Det vil af disse Grundtræk sees, at den humane
Gjennemførelse af Personlighedens Ret væsentlig maa
beroe paa en religiøs Udvikling af Personlighedsprincipet.
Til Culturoplysningen i Staten svarer den folkelige
Oplysning i Menigheden: en folkelig Stat, der ved sine
Tvangslove værner om den borgerlige Personlighed, for\-%
udsætter en Folkekirke, som ved Frihedens Love og
Aandens Midler opdrager og danner den religiøse Per\-%
sonlighed. En videnskabelig Forqvakling af Folkelivets
religiøse Principer er en væsentlig Hindring for det
aandelige Liv; men Hindringen kan forvandles til Be\-%
tingelse, naar Personlighedens Aand faaer Magt med
Culturoplysningens upersonlige, halvpersonlige Væsen.

\begin{center}\rule{3cm}{0.4pt}\end{center}

\chapter*{X.}
\addcontentsline{toc}{chapter}{Forelæsning X}
\markboth{Forelæsning X}{Forelæsning X}
\label{ch:x}
% pp. 154--175
% --- p. 154 ---
\opage{154}\setcounter{footnote}{0}%
% [large initial N]
Naar Folkeoplysningen engang er naaet saa vidt, at den
virkelig kan maale sig med en alsidig, rigt udviklet
Civilisation, naar Lys staaer mod Lys, Troens Lys mod
Culturens Lys, saa ville Straalerne fra begge Sider til\-%
bagekastes ved forskjellig Modstand og bryde sig i Vir\-%
kelighedens forskjelligartede Media; under Brydningen
vil sees et mangefarvet Lys og i det mangefarvede Lys
mangfoldige Livsbilleder. Men vil denne Brydning altid
være harmonisk? Nei, ikke altid. Ville disse Lys- og
Livsbilleder vidne om et eneste sammenhængende, i sig
conseqvent Levnetsløb? Ingenlunde! Blandt de mange
Brydningsmedia er der særlig et, der ifølge sin Natur
maa frembyde et vanskeligt Mellemværende, og dette
Brydningsmedium er den hellige Historie.
% [a reader has underlined "hellige Historie" in pencil; not transcribed]

Den hellige Historie har det tilfælles med al Historie,
at dens Indhold skal gjengive det Virkelige; men den
er deri forskjellig fra al anden Historie, at det Virkelige,
den gjengiver, paa uigjennemskuelig Maade forener det
Overnaturlige med det Naturlige. Det Virkelige har
altsaa her optaget Noget, der synes at staae i Strid med
Virkelighedens Begreb, og dog skal dets Gyldighed
tænkes. Virkelighedsbegrebet er stillet i Forgrunden;
i dette Begreb er Maalestokken given; vi skulle nu see,
% --- p. 155 ---
\opage{155}\setcounter{footnote}{0}%
hvor forskjelligt den anvendes, eftersom Tanken træder
i Civilisationens eller i Folkeoplysningens Tjeneste.

Fra Civilisationens, fra den negative Kritiks Stand\-%
punkt, vil Virkeligheden naturligviis blive at bedømme
efter de almeengjældende uforanderlige Naturlove. Civi\-%
lisationen bevæger sig efter Viden; i Viden er Skabelses\-%
begrebet saaledes underordnet Naturbegrebet, at det Na\-%
turlige ikke kan optage det Overnaturlige. Fra Folke\-%
oplysningens Standpunkt bliver Virkeligheden derimod
at bedømme efter Personlighedens og Frihedens Lov;
Folkeoplysningen bevæger sig efter Troen: i Troen er
Naturbegrebet saaledes underordnet Skabelsesbegrebet,
at det Naturlige maa have sin Grund i det Overnatur\-%
lige. Saa staaer da igjen Tro imod Viden som Princip
mod Princip. Som Principet er, maa Kritiken være;
har Viden sin Kritik, da maa Troen ogsaa have sin. Saa
staaer da Kritik mod Kritik, Videnskritik mod Troes\-%
kritik: ved Kritikens vexlende Lys maa den hellige
Historie skifte Farve. Det Lys, der udgaaer fra Troes\-%
kritiken, vil overalt opklare Forskjellen mellem Mytho\-%
logie og Aabenbaring, mellem Mythe og hellig Historie.
Det Lys, der udgaaer fra Videnskritiken, vil overalt op\-%
klare Ligheden mellem Mythologie og Aabenbaring,
mellem Mythe og hellig Historie. Vi ville begynde med
Videnskritiken.

Grundsætningen er, at hvad der i en Historie, den
være profan eller hellig, ligner det Mythiske i alle væsent\-%
lige Træk, det maa være mythisk. Nu er f. Ex. Ligheden
imellem en talende Hest og en talende Aseninde saa
umiskjendelig, at dersom Hesten er mythisk, saa maa
Aseninden ogsaa være det. At den ene Historie findes
hos Homer, den anden i Bibelen, kan ikke berettige os
til at ansee den sidste for en Aabenbaring og den første
for en Mythe; en saadan Fremgangsmaade vilde jo være
% --- p. 156 ---
\opage{156}\setcounter{footnote}{0}%
ligesaa ukritisk og vilkaarlig, som hvis en Historiker
faldt paa at antage Noget i et Dokument for usandt,
alene fordi det var skrevet med græske Bogstaver, og
Noget for sandt, naar det blot var skrevet med hebraiske.
Opgaven er at bestemme Grundcharakteren saavel af
den bibelske som af den homeriske Fortælling. I Homers
Iliade (XIX) beskrives, hvorledes Achilles, efterat Pa\-%
troklos er falden, og Forsoningen med Agamemnon
foregaaet, ruster sig til Kamp og bestiger Stridskarmen.
Han tiltaler sine Heste med disse Ord: „Xanthos
og Balios! I navnkundige Føl af Podarge, glemmer
nu ei som sidst at føre fra Val Eders Herre hjem
i Behold til Achaiernes Leir, naar Kampen er ud\-%
stridt!“ Det er Alvor med denne Samtale, thi den Ta\-%
lende faaer et alvorligt Svar. Xanthos, „den flyvende
Hest under Aaget“, bøier sig dybt og svarer med Suk:
„Hjem i Behold skal vi bringe dig end, du stærke Peleide!
Alt dog stunder det stærkt ad din Død; vor Skyld er
det ikke; skyldig er kuns den vældige Gud og den
mægtige Skjæbne.“ Saaledes taler Xanthos, og hvorledes
er den kommen til at tale? „Mæle den fik af den lilie\-%
armede Here.“ Og hvorledes virker Talen paa Achilles?
Det overrasker ham ikke i mindste Maade, at hans Hest
giver sig til at tale, ja endog at tale som den, der er
indviet i Skjæbnens Hemmeligheder; tvertimod, det er,
som om det faldt af sig selv. Xanthos er dybt rørt,
Achilles er ogsaa dybt rørt og svarer veemodig:
„Xanthos, hvi spaaer du min Død? Bebudelsen kan
du dig spare.“ Nu sige vi Alle: Det er skjønt, det er
rørende, men det er kun en Mythe. At et opvakt Barn
kan tage en saadan Historie for Virkelighed, er meget
naturligt; men at et voxent Menneske med sund Forstand
kan blive saa barnligt, stikke saa dybt i „Barnetroen“,
at det ogsaa tager Historien for virkelig, er i allerhøieste
% --- p. 157 ---
\opage{157}\setcounter{footnote}{0}%
Grad unaturligt. Ingen tvivler om, at den Hest, som
Here kan give Talen, og Erinnyerne berøve Mælet, ikke
hører ind under Kategorien \textit{Equus Caballus}, der har sin
Plads i Zoologien, men slet og ret ind under Mythologien.
% [Equus Caballus: italic in the print, unlike the book's other Latin]

Vi have betragtet Achills Hest, lader os nu betragte
Bileams Aseninde! Historien findes i fjerde Mosebog
(XXII). Israels Børn, hedder det, leirede sig paa Moabi\-%
ternes slette Marker, paa hiin Side Jordanen mod Jericho.
Balak, Zipors Søn, Moabiternes Konge, seer denne Magt,
seer i den en indbrydende Fjende og forfærdes; i sin
Forfærdelse veed han ikke bedre Raad end at sende sine
Ældste til Propheten Bileam og ved Spaadoms Løn for\-%
maae ham til at forbande de Fremmede. Men om Nat\-%
ten aabenbarer Jehova sig for Bileam og siger udtryk\-%
kelig: „Du skal ikke forbande Folket; thi det er vel\-%
signet“. Og Bileam viser ingen Ulydighed; tvertimod,
han erklærer, at om ogsaa Balak vilde give ham hans
Huus fuldt af Sølv og Guld, kunde han ikke handle
ulydigt mod Herren og forbande Israel. Man trængte
ind paa ham, der hengik igjen en Nat, og Jehova aaben\-%
barede sig atter for Propheten, men istedenfor at holde
ham tilbage siger han nu: „Dersom de Mænd ere komne
for at kalde Dig, da staae op, gak med dem; men dog,
det Ord, som jeg vil tale til Dig, det skal Du gjøre.“
Bileam vaagner og gjør, som Jehova har sagt, og følger
med „de Moabiters Fyrster“. Men paa Veien optændes
Jehovas Vrede over, at Propheten følger med; hvad er
Jehova bleven vred over? Det veed man ikke. Her
anføres ikke det Mindste, der kunde tyde paa en
Forklaring. Vil man have en Forklaring, maa man selv
digte den. Texten siger os kun, at Jehova er bleven
vred, og i sin Vrede sender Engelen, som skal standse
Bileam. Vreden er meget stor; Propheten er i den Grad
falden i Unaade, at han end ikke findes værdig til at see
% --- p. 158 ---
\opage{158}\setcounter{footnote}{0}%
Engelen, hvorimod Aseninden findes værdig. Engelen
staaer i Huulveien ved Viingaarden, hvor der var Steen\-%
gjærde paa begge Sider, og da Aseninden seer Engelen
og trykker sig ind til Væggen, saa fortørnes Prophe\-%
ten og slaaer hende tre Gange; nu giver Herren Asen\-%
inden Mæle, og den mælende Aseninde spørger ganske
sagtmodig: „Hvi slaaer du mig, hvad har jeg gjort dig?“
Istedenfor at forbauses over denne Tiltale, hører Bileam
paa den med samme Naturlighed, som Achilles hørte paa
Xanthos, og siger: „Havde jeg et Sværd, vilde jeg slaae
Dig ihjel.“

Saa slaaende er Ligheden imellem Historien om
den talende Aseninde og Historien om den talende Hest,
at dersom de begge to fandtes hos Homer, saa maatte
de, forstaaer sig, begge være mythiske, og dersom de
begge to fandtes i Bibelen, saa maatte de, betragtede fra
Troens Standpunkt naturligviis, begge være Aabenbaringer.
Her seer man da, hvad der kommer ud af, at lade Ska\-%
belsesbegrebet fortrænge Naturbegrebet og anerkjende en
Virkelighed, hvori det saakaldte Overnaturlige faaer Lov
at blande sig i det Naturlige. Vistnok ville de Troende
beraabe sig paa, at det kun er Here, der gjør Miraklet
med Hesten, medens det derimod er Jehova, der gjør
Miraklet med Aseninden, at Heres Mirakel kun tilhører
Indbildningen ligesom hun selv, medens Jehovas Mirakel
er et Tegn paa den Magt, han virkelig har som Himlens
og Jordens almægtige Skaber; kort: man vil forsvare
Miraklets Virkelighed ved at beraabe sig paa den virke\-%
lige Almagt. Men hvad kommer der nu ud af at lade
Almagtsbegrebet i ubegrændset Vilkaarlighed fortære Na\-%
turbegrebet? Vi lære det af Asenindens Ord til den
forbittrede Prophet. „Og Aseninden sagde til Bileam:
er jeg ikke din Aseninde, som du har redet paa fra den
Tid, du var til, indtil denne Dag; mon jeg nogensinde
% --- p. 159 ---
\opage{159}\setcounter{footnote}{0}%
havde Vane at gjøre saa imod dig?“ Hvem er nu
egentlig her det talende Subject? Det kan da umulig
være enten Engelen eller Jehova, der taler saaledes,
at det lyder for Bileam, som om Talen kom fra Asen\-%
inden: et saadant mirakuløst Bugtaleri vilde jo være
Guddommen uværdigt. Det maa da være Aseninden,
der taler. Men hvad er dog dette for et Mirakel? Det
er et Mirakel, hvorved en umælende Dyresjæl i et Nu
% [a reader's pencil stroke in the margin beside the two lines above; not transcribed]
bliver omskabt til en fornuftig, selvbevidst Sjæl; og det
saaledes, at denne selvbevidste Sjæl i et eneste Øieblik
udvikles i den Grad, at den kan Sproget, det Sprog,
Bileam som Barn har været længere Tid om at lære; et
Mirakel, hvorved Dyresjælen ikke blot bliver omskabt
fra umælende til mælende, altsaa til en ganske anden
Sjæl, et ganske andet Væsen, men ovenikjøbet saaledes
omskabt, at den trods Omskabelsen dog vedbliver at
være det samme Væsen, den samme Sjæl. Det Menne\-%
ske, paa hvis Hoved Haarene ikke reise sig, naar han
skal til at anerkjende en saadan Virkelighed, hans Hoved
maa enten have tabt Haaret eller --- Forstanden;
thi Synet er grueligt. Den mælende Aseninde er virke\-%
lig; den umælende er ogsaa virkelig: den ene skyder
den anden ud, og dog ere disse to virkelige Dyr i Vir\-%
keligheden eet og samme Dyr; det kan man kalde
Virkelighed!

For at det imidlertid ikke skal faae det Udseende,
% [a speck above "at" in the scan looks like an accent; read as dirt]
som om Videnskritiken, om vi saa maae sige, bed sig fast
i en enkelt, dunkel, for Aabenbaringshistorien i det Hele
mindre væsentlig Fortælling, ville vi til Belysning af den
mirakuløse Virkelighed vælge et Par Exempler, der staae
i en saa inderlig organisk Sammenhæng med Aabenba\-%
ringsøkonomien i dens Heelhed, at de ikke kunne tages
bort, uden at det Hele falder sammen; vi mene Sagnet
om Paradiis og Evangeliet om Christi Opstandelse.

% --- p. 160 ---
\opage{160}\setcounter{footnote}{0}%
Sagnet om Paradiis maa fra Troens Synspunkt op\-%
fattes som en sand Historie og som sand Historie lyde
paa det Virkelige. Der maa altsaa i det virkelige Pa\-%
radiis have været et virkeligt Livets Træ, et vir\-%
keligt Kundskabens Træ, en virkelig Slange o. s. v. o. s. v.
Adam og Eva maae have levet som fuldvoxne Mennesker
i hellig Uskyldighed og fuldkommen Lyksalighed før
Faldet, og paa Faldet maa der være fulgt en virkelig
Forbandelse af hele den virkelige Jord. Vi skulle nu
prøve, hvad det er for en Virkelighed. Hvad der i alle
Henseender ligner det Mythiske, det maa være mythisk;
ud fra denne Grundsætning begynde vi igjen Under\-%
søgelsen. At Geologien ikke har nogen Plads for et
virkeligt Paradiis, er jo vistnok et stærkt Vidnesbyrd
imod dets Virkelighed; imidlertid er det i nærværende
Sammenhæng ikke saameget Naturlærens Kjendsgjer\-%
ninger som meget mere Virkelighedens eget Begreb,
hvorpaa Eftertrykket falder. Saalidt som naturlige Træer,
f. Ex. Æble- og Pæretræer, kunne tænkes at voxe i en
overnaturlig, d. v. s. overjordisk, Jordbund, ligesaa lidt
kunne saadanne overnaturlige Træer som Livets Træ
og Kundskabens Træ paa Godt og Ondt tænkes at voxe
i en naturlig Jordbund. Spørge vi nu, hvordan den
virkelige Jordbund maa have været i Paradiis, saa bliver
Svaret meget tvivlsomt. Det siges udtrykkelig, at Gud
Jehova tog Mennesket og satte ham i Edens Have til
at dyrke og bevogte den. Den Deel af den virkelige
Have, som Adam skulde dyrke og bevogte, tør da vel
antages at have været Virkelighedens naturlige Deel;
hvorimod den Afdeling af Haven, hvori Livets Træ og
Kundskabens Træ voxte, neppe har trængt til nogen
Dyrkelse og altsaa vel maa have udgjort den over\-%
naturlige Deel.

At Adam virkelig skal dyrke og bevogte en virkelig
% --- p. 161 ---
\opage{161}\setcounter{footnote}{0}%
Have, tyder som sagt paa, at Adam er et virkeligt og
naturligt Menneske. Men hvordan skal nu et naturligt
Menneske tænkes at kunne plukke Frugter af et over\-%
naturligt Træ, ja, ikke alene plukke dem, men endogsaa
spise dem? Og dog lykkes det desværre kun altfor godt
først Eva og derpaa Adam i deres nøgne Naturlighed
at plukke og spise det Overnaturlige; det lykkes
dem i Virkeligheden saa godt med Kundskabens Træ,
at Jehova ordentlig bliver bange for, at det ogsaa skal
lykkes dem med Livets Træ. „Og Gud Jehova sagde:
See Mennesket er blevet som Een af os til at kjende
Godt og Ondt, og nu, blot han ei udstrækker sin Haand
og tager ogsaa af Livets Træ og æder og lever evindelig“.
For at dette kan forebygges, bliver Adam med sin Hustru da
virkelig udjaget af det virkelige Paradiis, og Veien til
Livets Træ bevogtet af Cheruberne. Før deres Uddri\-%
velse af Paradiset maae vore første Forældre altsaa i
den første Virkelighed have havt en vidunderlig Blan\-%
ding af Naturligt og Overnaturligt, saa at de først, efter
at være blevne uddrevne, ere blevne rigtig naturlige.

Gaae vi over til Dyreverdenen og spørge Zoologen,
hvad han synes om det Stykke af den hellige Historie,
der omhandler Dyrene i Paradiis, vil han uden Tvivl
svare: Ikke synderlig godt. Der siges om Dyrene, at
de vare skabte af alle Arter, og at de alle bleve førte
til Adam, der gav dem Navne. Blandt de paradisiske
Dyr, der saaledes bleve førte til Adam, maae da ogsaa
Rovdyrene have været. I den paradisiske Tilstand kunde
de dog umulig gaae paa Rov; det vilde jo i Virkelighe\-%
den have forstyrret den virkelige Lyksalighed. Men,
spørge vi, hvad er et Rovdyr, der ikke gaaer paa Rov?
Det er en rund Fiirkant. Det ligger i Rovdyrenes zoolo\-%
giske Charakteer, i deres Kløer, Tænder, Fordøielsesred\-%
skaber o. s. v., at de absolut maae være kjødædende; men
% --- p. 162 ---
\opage{162}\setcounter{footnote}{0}%
hvordan skulde kjødædende og bloddrikkende Dyr kunne
leve, naar de ikke turde gaae paa Rov? Skulde det
maaskee være overnaturligt Kjød, hvoraf de levede?
Men saa maatte det jo ogsaa have været overnaturlige
Rovdyr. At Løver, Tigre, Leoparder o. s. v. endnu ikke
gik paa Rov, vil sige, at Løven endnu ikke var Løve,
Tigeren ikke Tiger o. s. v. Rovdyrarterne vare altsaa
ikke blevne til; og dog siger Fortællingen, at Gud alle\-%
rede havde skabt alle „de vilde Dyr paa Marken“.

Endnu værre bliver det, naar man gaaer til Slangen.
Med den talende Slange har Theologien virkelig taget
det altfor let, idet den snart uden videre har gjort
den til Diævelen selv, „den gamle Slange“, snart mere
speculativt til „det kosmiske Princip“. Hvad det kos\-%
miske Princip angaaer, da hører der mere end menne\-%
skelig Forstand til at indsee, hvorledes et saa meta\-%
physisk Væsen som et Princip skal kunne enten for\-%
vandle sig til en Slange eller tale igjennem en Slange.
Og hvad Djævelen angaaer, da indeholder Fortællingen
ikke mindste Antydning af, at Slangen er besat af
Diævelen.
% [Diævelen (twice) and Djævelen (once) on this page, as printed]

Naar Slangen taler til Qvinden, er det i en ganske
anden Forstand Slangen selv, end naar Xanthos taler
til Achilles, eller Aseninden til Bileam. Hesten og
Aseninden ere umælende Dyr, der kun ved et Mirakel
kunde blive talende. Men saaledes ikke med Slangen;
thi „Slangen var trædskere end alle Dyr paa Marken,
som Gud Jehova havde gjort“. Slangen hører til
Dyrene paa Marken, altsaa til de Umælende; men ---
saa frugtbar er den Opfindelse at underordne Natur\-%
begrebet under Skabelsesbegrebet --- skjøndt efter sin
Naturlighed hørende til de umælende Dyr, har Slangen
dog fra Begyndelsen af vidst at forene det Overnaturlige
med det Virkelige saa snildt, at den er bleven talende
% --- p. 163 ---
\opage{163}\setcounter{footnote}{0}%
ligesaavel som Mennesket. Skulde Nogen endnu nære
Tvivl om, at det efter Fortællingens Mening virkelig
er Slangen, der taler, og taler af sit Eget, da høre
han blot Jehovas Forbandelse over Slangen: „Efterdi
du gjorde dette, da vær forbandet fremfor alt Qvæget
og fremfor alle vilde Dyr paa Marken; du skal gaae
paa din Bug og æde Støv alle dit Livs Dage“. Slangen
har altsaa ved sin Handling paadraget sig Jehovas Vrede
og Straf, er altsaa tilregnelig, har altsaa handlet med
Frihed og Overlæg. Dersom dette ikke er mythisk,
hvad er da mythisk? Kan Nogen vel lukke Øinene for
Ligheden mellem Slangens Brøde og Adams og Evas
Brøde? Eva straffes, thi hun er tilregnelig og har mis\-%
brugt sin Frihed; Adam straffes, thi han er tilregnelig
og har misbrugt sin Frihed. Slangen straffes; den er
da ogsaa tilregnelig og har misbrugt sin Frihed. Men
den Bevidsthed, som, hvad Frihed og Tilregnelighed an\-%
gaaer, kan sætte en Slange paa lige Trin med et Men\-%
neske, er ikke en Bevidsthed, der fatter det Virkelige,
men en mythisk Bevidsthed.

Hvad den mythiske Bevidsthed i al Naivetet har
fremstillet som overnaturlig Kjendsgjerning, kan, naar
Indholdet er af dybere ethisk Betydning, endnu bevare
sin Gyldighed længe efter, at Culturoplysningen har
bortforklaret det Overnaturlige; Phænomenerne antages
vel ikke længere for at være virkelige, men æres og
bruges som hellige Symboler. Istedenfor en overnaturlig
Virkelighed faae vi en religiøs Symbolik, og nu er det
Folkelærernes høiere Aandsopgave --- en Opgave, til
hvis Løsning de jo ogsaa værdig forberede sig ved at
studere „den hellige Theologie“ --- at anvende denne
Symbolik paa en saa tvetydig, snild og sindrig Maade,
at Menigmand ikke mærker Uraad, men tager det Sym\-%
bolske for fuld Virkelighed. Saaledes kan Sagnet om
% --- p. 164 ---
\opage{164}\setcounter{footnote}{0}%
Paradiis med Held anvendes ved Folkeunderviisning som
Symbol, og dog tage sig ud som hellig Historie; saa\-%
ledes kan Mythen om Christi Opstandelse endnu mangen
en Paaskemorgen gjøre Virkning som historisk Begiven\-%
hed og dog udlægges som et Symbol paa Livets Seier
over Døden, Lysets Seier over Mørket, Sandhedens
Seier over Løgnen. Men at Historien om Christi Op\-%
standelse ligesaa vist er en Mythe som Historien om
Paradiset, vil Videns-Kritiken let kunne bevise.

Vi ville ikke her opholde os ved, hvad Strausz tilstræk\-%
kelig har paaviist, at de evangeliske Beretninger i mange
væsentlige Punkter modsige hinanden indbyrdes; vi ville
--- til yderligere Prøvelse af den Virkelighed, Troen
faaer ud af at lade Skabelsesbegrebet fortrænge Natur\-%
begrebet --- holde os til en enkelt Beretning, for at see,
hvor modsigende en saadan er i sig selv.
% [a reader's pencil bracket in the margin beside "--- til yderligere ... begrebet"; not transcribed]

Vi vælge Beretningen hos Lukas (XXIV, 36---43).
De to Disciple, der mødte den Opstandne paa Veien
til Emaus, ere allerede vendte tilbage til Jerusalem,
hvor de finde „de Elleve og dem, som vare med dem“,
forsamlede. Just som de Forsamlede ere fordybede i
Samtale om det Vidunderlige, der var skeet, staaer
Jesus selv pludselig midt iblandt dem. „Da forfærdedes de
og betoges af Frygt, og meente, at de saae en Aand“!
Det første umiddelbare Indtryk er som af et Aandesyn.
Til Disciplenes store Overraskelse bliver dette Aande\-%
syn imidlertid til den Korsfæstede selv i legemlig Virke\-%
lighed, i „Kjød og Been“. Til endnu yderligere Beviis
for, at det Overnaturlige i denne Virkelighed er forenet
med det Naturlige, gjør den sig Aabenbarende dem det
Spørgsmaal: „Have I her noget at æde?“ De give ham
et Stykke af en stegt Fisk og af en Honningkage. „Og
han tog det, og aad det for deres Øine“. At Kritiken
maa ansee en saadan Aabenbarelse for umulig, har slet
% --- p. 165 ---
\opage{165}\setcounter{footnote}{0}%
og ret sin Grund deri, at den maa finde det i sig selv
Modsigende umuligt. At den Opstandne fremstiller sig
for Disciplene som et overnaturligt Væsen, en Aand, er
jo vel i formel Overeensstemmelse med Opstandelsens
Begreb, forsaavidt Opstandelsen er et Mirakel, og den,
hvis Liv beroer paa et Mirakel, maa saavel i legemlig
som i sjælelig Henseende være et mirakuløst Væsen.
Men, hvorledes skal det nu lade sig tænke, at et mi\-%
rakuløst Væsen kan foretage en saa naturlig Handling
som den, at spise „et Stykke af en stegt Fisk og af en
Honningkage“? Det maa vel tænkes i Troen paa samme
Maade, som det i Troen maa tænkes, at Jehova, Himlens
og Jordens almægtige Skaber, virkelig har spiist virkelig
Tykmelk og virkeligt Kalvekjød hos Abraham under
Træet i Mamre Lund; det maa vel tænkes at skee ved
et Mirakel. Men at spise er en naturlig Handling, og
intet Mirakel; det, som intet Mirakel er, skeer altsaa ved
et Mirakel: det Overnaturlige fortærer det Naturlige;
men just ved at fortære det Naturlige viser det Over\-%
naturlige sig som det Uvirkelige, det Mythiske.

Dette er Videnskritikens videnskabelige Resultat.
Vi have ladet den Vidende udtale sig saa alsidigt, ud\-%
førligt og uforbeholdent, at Ingen, som er istand til at
forstaae hans Udtalelser, vil kunne tage feil af Grund\-%
tanken; det Lys, hvormed han lyser op, er Culturens
Lys, og han sætter ikke Lyset under en Skjæppe. Saa\-%
længe Folket imidlertid, navnlig hvad den store Mængde
angaaer, kan holdes i en vis almueagtig Uvidenhed, vil
en saadan Culturkritik enten slet ikke blive bemærket,
eller ogsaa vil Kritikeren blive betragtet som et ugude\-%
ligt Menneske, det ikke er værd at indlade sig med
% printer's defect: no full stop printed after "med"
Men eftersom Folkeoplysningen skrider frem, maa Sagen
forandre sig. De Oplyste ville indsee, at Videnskritikens
Angreb paa Religionen ikke bestaae i kaade, vilkaarlige
% --- p. 166 ---
\opage{166}\setcounter{footnote}{0}%
Indfald, men fremgaae med indre Nødvendighed af en
bestemt Grundtanke, et eiendommeligt Princip. Striden
om Religionens Gyldighed maa --- det er en uafviselig
Betingelse for det aandelige Liv --- blive et Principspørgs\-%
maal. Efterhaanden som Culturprincipet klarer sig i
store Træk for den folkelige Bevidsthed, maa da ogsaa
Troesprincipet klare sig: imod Videnskritiken stille vi
Troeskritiken.

Eftertrykket falder, hvad allerede er antydet, paa
Forskjellen mellem Mythologie og Aabenbaring, mellem
Mythe og hellig Historie. Denne Forskjel kan ikke af\-%
gjørende bestemmes ved endelig Sammenligning af Phæ\-%
nomenerne; thi Aabenbaringsphænomenet har en skuf\-%
fende Lighed med Mythen; i Phænomenet som Phænomen
staaer en talende Hest paa lige Trin med en talende
Aseninde. Den væsentlige Dom afhænger af Principet.

At Videnskritiken forkaster ethvert Aabenbarings\-%
tegn, er ganske i sin Orden, thi Aabenbaringstegnet er
jo netop en phænomenal Afbrydelse af Tingenes natur\-%
lige Sammenhæng; hvis det ikke fremtraadte saaledes,
kunde det umulig være Tegn. Men naar den viden\-%
skabelige Kritiker bilder sig ind, at han ved at gjennem\-%
tænke Aabenbaringstegnet kan opløse det i en Modsigelse,
og fremstille det som et meningsløst Brud paa Naturens
Orden, da overskrider han Videns absolute Grændse.

Aabenbaringstegnet har aldeles Intet med den viden\-%
skabelige Opfattelse af Natursammenhængen at gjøre;
det er og bliver et religiøst Mellemværende mellem den
sig i Tegnet aabenbarende Gud og den Troende, der
modtager Aabenbaringen. Der kan fra Fortællingen om
den talende Aseninde ikke drages nogensomhelst natur\-%
videnskabelig Slutning enten til Angreb paa eller til
Forsvar for Phænomenet. Hvad er Phænomenet af en
talende Aseninde? Et Almagtens Tegn! Hvad er Phæ\-%
% --- p. 167 ---
\opage{167}\setcounter{footnote}{0}%
nomenet af en brændende Tornebusk? Et Almagtens
Tegn! Hvad er Phænomenet af den Dødes Opvækkelse?
Et Almagtens Tegn! Man nævne hvilken Undergjerning,
hvilken Aabenbarelse i hele den hellige Historie, man
vil; saasnart Tanken forsøger at trænge ind i Phæno\-%
menets Væsen, støder den overalt paa den samme
Skranke, standser ved den samme Gaade, faaer paa alle
Spørgsmaal det samme eensformige Svar: det er Al\-%
magtens Tegn!

Skulde Videnskritiken nu virkelig kunne bevise, at
et Almagtens Tegn er umuligt, da maatte den kunne
bevise, at den almægtige Villie selv er umulig; men Be\-%
viser af denne Natur have kun Gyldighed for dem, der
ved at glemme den personlige Selvmodsigelse bevise, at
de i dybeste Forstand have glemt sig selv.

„Men“, indvender man, „ved slige Almagtens Tegn
gjøres der jo de allervilkaarligste Brud paa Naturens
Orden.“ Denne Indvending vilde have Noget at betyde,
dersom Tegnet i Egenskab af Aabenbaringstegn virkelig
havde Noget med Tingenes naturlige Orden at bestille;
men en saadan Forudsætning beroer paa en videnskabe\-%
lig Misforstaaelse. Hvert Tegn er et enkelt, et ene\-%
staaende Phænomen, der umulig kan indføres som Led
i Naturphænomenernes Række; men naar det ikke kan
indføres i Naturphænomenernes Række, kan det jo heller
ikke foranledige nogensomhelst for Naturlæren forstyr\-%
rende Induction: Tegnet er, hver Gang det gives, et
personligt, et udelukkende personligt Mellemværende
imellem den, der aabenbarer sig ved Tegnet, og den
eller dem, der modtage Aabenbaringen. Grændsen for
Meddelelsen af slige Tegn kan ikke bestemmes ved nogen
Naturlov; men derfor er Grændsen ikke vilkaarlig: den
bestemmes ved guddommelig Viisdom, idet den bestem\-%
mes ved guddommelig Villie.

% --- p. 168 ---
\opage{168}\setcounter{footnote}{0}%
Det følger af Troens Princip, at Naturbegrebet,
nemlig det religiøse, maa underordnes Skabelsesbegre\-%
bet; men det følger ingenlunde af Principet, at den
religiøse Bevidsthed skulde være blind for den virkelige
Natur. Den poetiske Naturbetragtning anerkjender man
dog endnu for berettiget ligeoverfor den videnskabelige,
og hvorfor? Fordi man anerkjender Skjønhedsideen for
berettiget ligeoverfor Vidensideen. Hvorfor vil man da
ikke anerkjende den religiøse Naturbetragtning for be\-%
rettiget ligeoverfor den videnskabelige? Fordi Cultur\-%
bevidstheden er kommen saaledes i Gang med at hen\-%
føre Troen under Viden, at den har overseet Troes\-%
principets Selvgyldighed! Den davidiske Naturbetragtning
f. Ex. er ikke videnskabelig, men skulde den derfor ikke
have religiøs Gyldighed? Naturbetragtningens religiøse
Charakteer er bestemt derved, at den overalt, selv hvor
der ikke er Tale om særegne Mirakler, henfører Natur\-%
tingene til den skabende og opholdende Almagt. Istedenfor
at dvæle ved Phænomenernes naturlige Sammenhæng vil
det troende Øie allevegne igjennem Phænomenerne opdage
Gud selv. „Han udfolder Himlene som et Telt; han
tømrer i Vandene sine Høisale; Skyer danner han til
sin Vogn; saa ager han paa Stormens Vinger“ (Psalm.
IV). Mon en saadan Naturbetragtning vel skulde lade
% "IV" as printed; the verses quoted are from Ps. 104 (CIV)
sig fortrænge enten af Astronomien eller af Meteorolo\-%
gien? Fordi den religiøse Naturanskuelse her udtales i
Poesiens Form, mon den derfor skulde ophøre at være
religiøs? „Jehovas Træer staae frodige, Cedrene, som
han plantede paa Libanon, hvor Smaafuglene bygge
Rede“; skulde maaskee en saadan Udtalelse ikke kunne
bestaae med en videnskabelig Botanik? „Du lader det
mørknes“, siger den hellige Sanger om Jehova, „saa
kommer Natten, i den snige alle Skovens Dyr sig om,
Løveungerne brøle efter Bytte og begjære deres Føde
% --- p. 169 ---
\opage{169}\setcounter{footnote}{0}%
af Gud“; det kan dog vist ikke være Zoologiens Opgave
at underkaste denne Betragtning af Skovens Dyr en
videnskabelig Kritik.

Vilde Nogen være saa dannet at mene, at det ikke
er Psalmistens Alvor med de Exempler paa Guds Virk\-%
somhed i Naturen, her anføres, at det kun er billedlige
Udtalelser af en flygtig Stemning: da viser en Saadan
kun, at han for lutter Dannelse har glemt, hvad Religion
er, glemt, hvilket Eftertryk den Troende lægger paa de
Ord: „Herre, hvormange ere dine Gjerninger, Du gjorde
dem alle viselig!“ „Men“, spørger man, „har Psalmisten
da virkelig selv kunnet troe, at Jehova engang imellem
ager paa Stormskyerne, at han virkelig og med egen
Haand har plantet Cedrene paa Libanon, for at Fuglene
kunne bygge Rede i dem, og at Løveungerne, naar de
brøle efter Bytte, virkelig bede til Gud? Nei, paa saa
lavt et Trin kunne Naturvidenskaberne dog umulig have
staaet i hiin vistnok iøvrigt ingenlunde oplyste Tid; det
maa være Billeder, Symboler og Billeder, intet Andet
end Symboler og Billeder!“ Ja vist er det Billeder,
men dog mere end Billeder; vist er det Symboler, men
disse Symboler have desuagtet Virkelighed, fuld Virke\-%
lighed for Aanden; i den troende Naturbetragtning er
% "fuld Virkelighed for Aanden" is underlined in pencil by a reader, with a
% pencil line in the margin; not letterspaced in the print; not transcribed
Virkeligheden Symbol, Symbolet Virkelighed. Hvad
Symbolet er i den stemningsrige Naturbetragtning, det
er igjen paa særskilt Maade Tegnet, Almagtens Tegn i
Underet. Underets Virkelighed er et Tegn, et Symbol.
„Kun et Symbol?“ Ja, men dette Symbol er tillige
Virkelighed. Grundspørgsmaalet er: Kan Troen gjælde
for et Princip eller ikke? Forkaster man Principet, saa
kan den religiøse Naturbetragtning naturligviis kun have
subjectiv, æsthetisk Gyldighed, saa giver den kun Billeder,
men ingen Virkelighed; erkjender man Principet, saa
bliver man ogsaa nødsaget til at indrømme den til Prin\-%
% --- p. 170 ---
\opage{170}\setcounter{footnote}{0}%
cipet svarende Naturbetragtning Virkelighed, fuldgyldig
Virkelighed for Aanden.

For nu nærmere at undersøge, hvorledes det egent\-%
lig forholder sig med den religiøse Naturbetragtning,
og da særlig med Opfattelsen af Underets Virkelighed,
maae vi først bestemme Forholdet mellem Troeserfaring
og Sandsningserfaring noget nøiere. Naar de Troende
i den hellige Historie modtage en Aabenbaring, da frem\-%
træder Phænomenet ofte med en saa slaaende Umiddel\-%
barhed, med et saa bestemt Præg af udvortes Virkelighed,
i en saa umiskjendelig Vexelvirkning med naturlige Om\-%
givelser, at det seer ud, som kunde et saadant Phænomen
opfattes alene ved sandselig Seen, sandselig Høren, sand\-%
selig Følen, altsaa uden Tro. Dette er ikke destomindre
en grov Vildfarelse. Ligesom Troen altid forudsætter
Aabenbaring, saaledes forudsætter Aabenbaringen igjen
Tro: Aabenbaring og Tro ere i Principet Eet, og kunne
umuligt adskilles. Den Troende seer, hører, føler; men
i al denne Sandsning er Troen med; han seer med
Troens Øie, hører med Troens Øre: al hans Erfaring
er og bliver en Troeserfaring. Hvorvidt den Enkelte i
det enkelte Tilfælde modtager en sand eller en falsk, en
virkelig eller en indbildt Aabenbaring, derfor maa den
Paagjældende selv være ansvarlig. Hvorledes en Bileam
hører og forstaaer Jehovas Bud, derfor staaer han per\-%
sonlig Jehova til Regnskab; om Sligt kan Intet afgjøres
i reen Almindelighed. Hvorledes en David hører og
forstaaer Jehovas Bud, er en personlig Sag imellem Je\-%
hova og David. I denne Henseende ere de to Beret\-%
ninger om Folketællingen høist lærerige. Ifølge 1 Sam.
(XXIV) mener David at høre Jehovas Stemme: „Gaa,
% "1 Sam." as printed (clear at 300 ppi); the census narrative is 2 Sam. XXIV
tæl Israel og Juda!“ David bilder sig altsaa ind, at
han i dette Tilfælde har faaet en sand Aabenbaring;
men det er en Indbildning. Det er ikke Jehovas Stemme,
% --- p. 171 ---
\opage{171}\setcounter{footnote}{0}%
David hører; nei, det er (1 Chr. XXI.) en anden
Stemme. Aabenbaringen er falsk; hvorfor? Fordi Da\-%
vid har vakt Jehovas Vrede, og Grunden hertil er kjen\-%
delig nok. Det var efter Davids og --- hvad man tyde\-%
ligt mærker af Joabs Ord: „hvorfor har min Herre
Konge faaet Lyst til den Ting?“ --- hans Samtidiges
religiøse Synsmaade en saare betænkelig Sag at foretage
en Folketælling; det blev af gode psychologiske Grunde
udlagt som et Tegn paa et vist Overmod; Kongen vilde
ved et saa opsigtvækkende Skridt just gjøre det vitter\-%
ligt, at Israels Folk nu var saa talrigt, at det, uden at
behøve Jehovas særdeles Bistand, nok kunde stole paa
sin egen Styrke. Saa gjærer det i Davids Indre: Ego\-%
ismen er i Strid med Fromheden, Ærgjerrigheden med
Ydmygheden, Forfængeligheden med den Jehova skyl\-%
dige Taksigelse. I en saadan Tilstand er den Troende
usikker paa sig selv og priisgiven falske Indskydelser.
Kunde Aabenbaringsordet imidlertid opfattes af et sand\-%
seligt Øre uden Tro, saa maatte jo Skylden falde paa
Jehova; men det Falske ligger i Hørelsen, og Skylden
falder paa David.

For at sikkre sig Fortolkningen af et virkeligt
Aabenbaringsphænomen, opkaster man undertiden det
Spørgsmaal: maatte en Vantro, hvis en Saadan havde
været tilstede, ikke ogsaa have seet det Overnaturlige?
Man vil ved Besvarelsen af et saadant Spørgsmaal erfare,
hvorvidt Fortolkeren anseer Aabenbarelsen for en virke\-%
lig Kjendsgjerning eller blot for et indvortes Syn, et
Foster af de Troendes egen Phantasie; men man lægger
ikke Mærke til, at Spørgsmaalet indeholder en Selvmod\-%
sigelse. Hvis det Overnaturlige virkelig kunde sandses
og sees med de samme Øine, der see det Naturlige, saa
vilde det jo ikke være overnaturligt; saa vilde en Aaben\-%
baring kunne meddeles et Menneske uden Tro, og Troen
% --- p. 172 ---
\opage{172}\setcounter{footnote}{0}%
ophøre at være Princip. Istedenfor at forbinde den
virkelige Seen og den virkelige Høren med den virkelige
Tro, vil man kaste Troen bort og gribe det Aaben\-%
barede i sandselig raa Umiddelbarhed; det er at krænke
det Hellige. Man slutter saaledes: hvis en tilstede\-%
værende Vantro ikke ogsaa kunde have seet Synet, saa
er Synet ikke virkeligt; det er en falsk Slutning. Iste\-%
denfor at lægge den fulde Forvisning ind i Troen,
tager man Troen bort og lægger Forvisningen ind i
den blotte Sandsning. Talen være om Christi Opstan\-%
delse og den Opstandnes forskjellige Aabenbarelser for
Disciplene. Evangelisterne lægge den allerstørste Vægt
paa Opstandelsens Virkelighed: „Herren er sandelig op\-%
standen og seet af Simon“, sige de Elleve i Jerusalem
til de To, der kom fra Emaus; men ikke skal man hos
Nogen af Vidnerne nogetsteds opdage fjerneste Antyd\-%
ning af, at den Opstandne er seet ogsaa af Ikketroende.
Tvertimod findes der mangfoldige Vidnesbyrd, som vise,
at endog de Troendes Øine særlig maatte oplades, før
de kunde kjende den Forvandlede. Maria Magdalene
møder ham i Urtegaarden; men hun kjender ham ikke;
hun meente jo, at det var Urtegaardsmanden (Joh. XIX.,
14---15). Om de To, der gik til Emaus, da Jesus „kom
% "Joh. XIX." as printed; the passage is John XX, 14--15
selv nær og vandrede med dem“, hedder det (Luk.
XXIV, 16): „Men deres Øine holdtes til, saa de ikke
kjendte ham“. Selv de i Jerusalem Forsamlede (Luk.
XXIV, 36---42), som dog alt fra flere Sider vare for\-%
beredte ved det glade Budskab, selv disse kunde jo ikke
strax og ligefrem gjenkjende ham, der pludselig stod
midt iblandt dem; „men de forfærdedes og betoges af
Frygt, og meente, at de saae en Aand“.

Under Forudsætning af Tro, paa Betingelse af Tro
faaer den sande Aabenbaring først fuldgyldig Virkelighed,
Virkelighed for Aanden, Aandsvirkelighed; alle Aaben\-%
% "Virkelighed for Aanden, Aandsvirkelighed" underlined in pencil by a
% reader, with a marginal stroke; not letterspaced; not transcribed
% --- p. 173 ---
\opage{173}\setcounter{footnote}{0}%
baringsphænomener ere i Virkeligheden Aandsphænome\-%
ner. At Virkeligheden her er Virkelighed for Aan\-%
den, maa ikke forstaaes, som var Aanden selv for af\-%
mægtig til at ophæve de umiddelbart naturlige Skranker.
Det er ingen Ufuldkommenhed ved Aandsphænomenet,
at dets høieste Virkelighed udelukkende er for Aanden;
det er tvertimod en Fuldkommenhed: Aabenbaringen vil
ikke stadfæste Aandens Afmagt, men dens Almagt, dens
absolute Selvmagt. Vi see det paa den Opstandne.
Efterhaanden som Disciplenes Øine oplades, efterhaanden
som det klarer sig i deres Indre, saa at de see ham i
Troen, høre ham i Troen, bliver ogsaa den udvortes
Skikkelse klar og bestemt, antager Virkelighed, ja haand\-%
gribelig Virkelighed. Hvorledes er det muligt, at en
Skikkelse, der begynder med at vise sig som Aand,
kan ende med at fremstille sig i fuld, haandgribelig
Virkelighed? Ja, hvorledes er det muligt? Ved dette
Spørgsmaal bilder Videnskritiken sig ind, at den kan
fremmane uhyre Vanskeligheder og indvikle den troende
Tænker i uhyre Modsigelser. Det er en Indbildning; fra
Troens Synspunkt er her i et saadant Spørgsmaal kun
een Modsigelse, een eneste, og den ligger ikke i Sagen,
hvorom der spørges, men i Spørgsmaalet.

Som et talende Exempel paa, hvad et Aabenbarings\-%
phænomen i fuld Virkelighed har at betyde, tjener især
% "Aabenbaringsphænomen i fuld Virkelighed" underlined in pencil by a reader; not transcribed
Beretningen om den Opstandnes Aabenbarelse for Thomas
(Joh. XX, 26---29): „Og efter otte Dage vare atter hans
Disciple inde, og Thomas med dem. Jesus kom, der
Dørene vare lukte, og stod midt iblandt, og sagde: Fred
være med Eder! Derefter siger han til Thomas: ræk
din Finger hid, og see mine Hænder, og ræk din Haand
hid, og stik den i min Side; og vær ikke vantro, men
troende. Og Thomas svarede og sagde til ham: min
Herre og min Gud! Jesus siger til ham: efterdi du
% --- p. 174 ---
\opage{174}\setcounter{footnote}{0}%
haver seet mig Thomas, haver du troet; salige ere de,
som ikke have seet, og dog troe.“ Disciplen Thomas
viser sig her i et ganske eiendommeligt Lys. Han er
vantro, siger man. Ganske vist forraader han en betæn\-%
kelig Svaghed i Troen; vi høre ham jo, før Aabenba\-%
relsen skeer, udtrykkelig at erklære: „Uden jeg faaer seet
Naglegabet i hans Hænder, og stikker min Finger i
Naglegabet, og stikker min Haand i hans Side, vil jeg
ingenlunde troe“; en Saducæer kunde sige det Samme.
Og dog er Thomas ingen Saducæer. Han negter ikke
% Rettelser (PDF 338): læs "Sadducæer" (both occurrences, ll. 9--10); a reader
% has pencilled a "d" above each; not transcribed
Opstandelsens Mulighed; men Vished vil han have, fuld
Vished; han vil have Virkeligheden, den fulde, haand\-%
gribelige Virkelighed for sig. Saa faaer han da ogsaa
Virkeligheden; han faaer den saa følelig, saa umiddel\-%
bar, som det vel kan tænkes. Men hvor, og hvorledes
faaer han den? I Disciplenes Kreds, og efter at have
hørt den Opstandnes Hilsen: „Fred være med Eder!“
For denne Stemme viger Tvivlen, ved denne Hilsen
vækkes Troen, og nu er Thomas i Aanden forberedt til
at fatte den fulde Virkelighed. Men med Aabenbarelsen
følger en Bebreidelse. Hvad er det Christus bebreider
Thomas? At han vil have fuld Forvisning om Opstan\-%
delsens Virkelighed, bliver ikke bebreidet ham; en Let\-%
troenhed, der stoler paa løse Rygter, har Christus aldrig
anbefalet. Nei, Bebreidelsen lyder paa, at den Lidet\-%
troende vil gjøre sin Forvisning afhængig af Sandsernes
Vidnesbyrd. Thomas vil, som man siger, „see det, før han
vil troe det“; han vil, med andre Ord, „forholde sig til
de hellige Kjendsgjerninger paa tilsvarende Maade som
Naturforskeren forholder sig til Naturens Kjendsgjer\-%
ninger“.

Det Feilgreb, hvori Disciplen Thomas her gjør sig
skyldig paa første Haand, er just det, de orthodoxe Dog\-%
matikere, i Protestantismen ligesaa fuldt som i Katholi\-%
% --- p. 175 ---
\opage{175}\setcounter{footnote}{0}%
cismen, paa forskjellig Viis have efterlignet, og hvori de
have gjort sig skyldige paa anden Haand. Ligesom
Thomas helst vil betragte Aabenbaringsphænomenet sand\-%
selig udvortes, empirisk, saaledes ville de objective Dog\-%
matikere endnu den Dag idag allerførst sikkre sig Kjends\-%
gjerningernes udvortes historiske Virkelighed; paa den
udvortes faste Virkelighed mene de saa, med Aandens
Bistand, at kunne bygge deres Tro som paa et muret
Grundlag. Men naar den udvortes Erfaringsvirkelighed
sættes som det Første, er Aanden krænket og negter de
objectivt Troende sit Vidnesbyrd. Kun for aandige
Kjendsgjerninger vil Aanden vidne; dens Vidnesbyrd er,
ifølge Principet, betinget af, at Aabenbaringsphænome\-%
nerne i deres høieste Virkelighed opfattes og bestemmes
som Aandsphænomener. Just fordi de orthodoxe Theo\-%
loger have fastholdt Aabenbaringsphænomenernes Virke\-%
lighed uden Aand, har den negative Kritik kunnet gaae
til den modsatte Yderlighed og i sindrige Myther byde
os hule Aandsphænomener uden Virkelighed.

At faae Aabenbaringsphænomenernes Virkelighed
grundig opfattet og ved Fortolkningen af den hellige
% "hellige": a speck of dirt above the first e (not an accent)
Historie alsidig belyst, er en væsentlig Betingelse for det
aandelige Liv i Nutiden.

\begin{center}\rule{3cm}{0.4pt}\end{center}

\chapter*{XI.}
\addcontentsline{toc}{chapter}{Forelæsning XI}
\markboth{Forelæsning XI}{Forelæsning XI}
\label{ch:xi}
% pp. 176--193
% --- p. 176 ---
\opage{176}\setcounter{footnote}{0}%
% [large initial D]
Den hellige Historie er Aabenbaringens Historie, Histo\-%
rien om den i Ordet sig aabenbarende Gud. Men den
i Ordet sig aabenbarende Gud er en talende Gud, der
adskiller sig fra Naturreligionernes stumme Guder og fra
Videnskabens tause Guddom. De mythiske Guder ere
stumme Guder; Guder af „Sølv og Guld“, et Værk af
Menneskehænder. De have vel en Mund, men kunne
ikke tale; de have vel Øine, men kunne ikke see; de
have Øren, men kunne ikke høre“ (Psalm. CXV).
% printer's defect: closing “ after "høre" with no matching opening „ (the „Sølv og Guld“ pair above is closed); transcribed as printed, quote balance -1
Hvad der fra Aabenbaringens Standpunkt siges om „Af\-%
guderne“, er bogstavelig sandt, endogsaa naar vi tage
Ideen med. Digteren paakalder Musen; han beder hende
om Stemning, om Indskydelser, om Aandens Bistand, om
Sprogets vidunderlige Gave; men hvorledes det forhol\-%
der sig med denne Paakaldelse, er let at gjennemskue.
For at Musen skal kunne give Stemning, maa Digteren
allerede være i Stemning; for at de geniale Indskydelser
kunne begynde i hans Indre, maae de allerede være be\-%
gyndte; det er Ideebevægelsen, hvorpaa Alt kommer an:
Musen som Muse er stum; den stumme Gudinde gjør
Ingen talende. At Digteren desuagtet maa beile til Mu\-%
sens Gunst, betyder, at han under Ideens Tryk maa for\-%
doble sit aandige Væsen, maa, for at bevæges, lade sig
% --- p. 177 ---
\opage{177}\setcounter{footnote}{0}%
bevæge; saa seer han det første Bevægende i Gudindens
Blik og modtager som guddommelig Gave, hvad der op\-%
rindelig boer i ham selv. Saaledes forstaaet, er Musens
Paakaldelse en psychologisk Nødvendighed, en skjøn
Tvetydighed, der klæder Poesien. Som et Væsensbillede,
der bærer Ideen, er Musen ikke Gudinde; som Gudinde,
det er: som et Væsen for sig i selvstændig Væren, er
Gudinden en Afgud og ingen Muse.

Kan Musen som Gudinde ikke tale til Digteren, da
kan Guddommen i sit rene Begreb endnu mindre tale til
Tænkeren. Vel maa Tænkeren, idet han tænker Guddom\-%
mens Idee, ogsaa fordoble sit Væsen, forsaavidt han
nødvendig maa tænke sit endelige Selv i Forskjel fra
den absolute Idee; men han maa igjen med fuld Bevidst\-%
hed kunne hæve Fordoblingen, hvis han vil holde sig
absolut i den rene almindelige Tænkning. At tænke
Gud, siger Hegel, er ikke som at tænke et Dyr
eller en Steen; thi Dyret og Stenen forestille vi os at
have en endelig Tilværelse udenfor vort Begreb; vort
Begreb om det absolute Væsen derimod, vort sande Be\-%
greb, er det absolute Væsen selv. Ihvorvel Tænken og
Talen ere saa uadskillelige, at al bestemt Tænken gaaer
i Eet med en indre Talen, saa kan det dog aldrig falde
Tænkeren ind enten at tale til det absolute Væsen eller
nogensinde vente, at det absolute Væsen skulde tale til
ham: Guddommens Talen til Mennesket er, fra Videns
Standpunkt, Menneskets Talen i sig selv.

Anderledes paa Troens Standpunkt. Troens Gud
er en Gud, der kan tale til Mennesket, en Gud, til hvem
Mennesket kan tale. Aabenbaringens Gud er en talende
Gud: Abraham hørte den talende Gud og vandrede
% "Gud:" read as a colon at 300 ppi (both witnesses give ";"); AUDIT: confirmed colon at 600 dpi
bort fra de stumme Guder; Moses hørte den talende
Gud, og Israel blev Guds Eiendoms Folk i Kraft af alle
„disse Ord, som Gud talede“.

% --- p. 178 ---
\opage{178}\setcounter{footnote}{0}%
Men for Viden og Videnskab kan der dog vel neppe
tænkes noget Begreb, der indeholder en større Modsigelse,
end Begrebet om en talende Gud. Vi maae i denne
Anledning lade den negative Kritik igjen komme til
Orde.

Blandt alle Mirakler, store og smaa, hvorom det
gamle og nye Testamente fortæller, er det største uden
Sammenligning dette, at Gud i Himlen selv tager Ordet
og giver sig til at tale. Under Forudsætning af, at han
virkelig har en Hemmelighed, er dette jo vistnok ganske
i sin Orden; Paulus siger: „Hvilket Menneske veed det,
der er i Mennesket, uden Menneskets Aand, som er i
ham? Saa veed og Ingen det, som er i Gud, uden Guds
Aand“. Ligesom det altsaa beroer paa et Menneskes
gode Villie, om han vil rykke ud med sin Hemmelighed,
eller fortie den, saaledes beroer det, i endnu strengere
Forstand, paa Guds gode og urandsagelige Villie, hvad
og hvormeget af forborgne Raadslutninger han vil betroe
sine Udvalgte. Men idet Miraklet med den talende Gud
opstilles som Aabenbaringens Grundmirakel, sees det da
ogsaa, hvorledes alle i Troens Væsen skjulte Modsigelser
uvilkaarlig springe i Øinene. Er Aabenbaringen en vil\-%
kaarlig Handling, saa maa det aabenbarede Ord nød\-%
vendig blive et Udtryk af Vilkaarligheden, altsaa et vil\-%
kaarligt Ord. At Aabenbaringsordet er et vilkaarligt
Ord, beviser da ogsaa det gamle Testamente fra Begyn\-%
delsen til Enden.

„Og Gud talede alle disse Ord og sagde“; saaledes
begynder (2. Moseb. XX) Lovgivningen fra Sinai. Det
Charakteristiske ved denne Lovgivning er altsaa ikke
nærmest dens Indhold, men dens Udspring; ikke hvad
der befales, men at det befales, at det udtrykkelig befa\-%
les af Gud. Hvad dette har at betyde, kan oplyses ved
den bekjendte Modsætning imellem to Synsmaader, der
% --- p. 179 ---
\opage{179}\setcounter{footnote}{0}%
allerede kom til Udtalelse i Middelalderens Scholastik:
\textit{det er godt, fordi Gud vil det}, og \textit{Gud vil det, fordi det
er godt}. Den sidste Sætning gjør Guds Villie afhængig
af det Godes Idee og stemmer med Videns Fordring.
Naar Gud ikkun befaler, hvad der er godt, og forbyder,
hvad der er ondt, saa behøver jeg jo kun at bruge
min Fornuft til at kjende Forskjellen mellem Ondt og
Godt, og trænger ikke til vilkaarlige Aabenbaringer.
Saadanne Bud som: Du skal ikke stjæle; Du skal ikke
sige falsk Vidnesbyrd imod din Næste; du skal ikke
slaae ihjel o. dsl., ere almeengyldige Fornuftbud, der
kunne erkjendes, uden at en talende Gud har nødig at
dictere dem. Men at opfatte de sædelige Bud saaledes,
strider imod Aanden i den mosaiske Lovgivning. Hvor\-%
for maa en Israelit ikke stjæle? Svar: fordi Jehova har
forbudt ham det. Hvis Jehova ikke havde forbudt ham
det? Saa maatte han nok stjæle. Og hvis Jehova engang
befalede ham det? Ja, saa var det hans Pligt at stjæle.
Saaledes komme vi da til at fastholde den første Sæt\-%
ning: \textit{det er godt, fordi Gud vil det}, som en ubetinget
Vilkaarlighedssætning, en uforfalsket Aabenbaringssætning.
Det kan fra dette Synspunkt ikke undre os at finde
sædelige Fornuftbud stillede i Række med aldeles speci\-%
elle, vilkaarlige Bud. Det hedder: „Du skal ikke gjøre
dig noget udskaaret Billede eller nogen Afbildning af
det, som er i Himmelen oventil, eller af det, som er paa
Jorden nedentil, eller af det, som er i Vandet under
Jorden“; et saadant Bud forudsætter ganske rigtig en
særegen Aabenbaring, thi intet fornuftigt Menneske vil
af sig selv falde paa, at Gud skulde være en saa afsagt
% a raised speck between "saa" and "afsagt" (risen space or dirt), not an apostrophe; not transcribed
Fjende af Billedhuggerkunst og Maleri.

Naar et Folk i den Grad gjøres sædelig umyndigt,
at det ikke ved egen Eftertanke er istand til at udfinde
Forskjellen mellem Godt og Ondt, men skal have alle
% --- p. 180 ---
\opage{180}\setcounter{footnote}{0}%
religiøse, sædelige og borgerlige Forskrifter udtrykkelig
dicterede, da er det en Selvfølge, at „den talende Gud“
maa tale mange Ord og foreskrive Alt indtil de mindste
Enkeltheder. De særskilte Love ere derfor ligesaa umid\-%
delbart udgaaede af Jehovas Mund som Grundloven.
Det er utroligt, hvor en lovgivende Gud kan gaae i
Detailler. „Og Jehova talede til Moses (2. Moseb. XXV),
idet han sagde: Tal til Israels Børn, at de skulle tage
mig en Offergave af hver Mand ... Og dette er Offer\-%
gaven, som I skulle tage af dem, Guld og Sølv og Kob\-%
ber. Mørkeblaat og Purpurrødt og Carmosinrødt og
fiint Linned og Gedehaar, og rødfarvet Væderskind, og
Sælhundeskind og Acacietræ, Olie til Lysning, Urter til
Salveolie“ o. s. v. o. s. v. Fremdeles (2. Moseb. XXVI)
„Og Tabernaklet skal du gjøre af ti Gardiner af tvundet
Linned og Mørkeblaat og Purpurrødt og Carmosinrødt“.
Hvis der altsaa til Tabernaklet blev brugt f. Ex. elleve
Gardiner, saa vilde det ellevte Gardin maaskee være Gud
til større Krænkelse end et falsk Vidnesbyrd imod Næsten.
„Hvert Gardin skal være otte og tyve Alen langt og otte
Alen bredt“; hvis det nu traf sig, at et Gardin blev en
halv Tomme for kort, saa vilde den Skyldige paadrage
sig Jehovas Vrede og maaskee blive udslettet af sit Folk.

Er ikke dette, spørger Culturmennesket, at drage
det Evige, det Ophøiede ned i det Endelige og Smaalige?
Og er det underligt, at de gamle Hedninger spottende
sagde: en saadan talende Gud er ingen Aandens Gud;
det er en snaksom Gud, der vil blande sig i Alt, be\-%
fatte sig med Alt, vide Besked om Alt (garrulus, inquie\-%
tus, impudenter etiam curiosus). Hvor guddommelig Vil\-%
kaarlighed hersker, er det umuligt at gjøre Forskjel
imellem Væsentligt og Uvæsentligt, Sædeligt og Usæde\-%
ligt, Tilladeligt og Utilladeligt. Der gives en Løgn, som
man pleier at kalde Nødløgn. Spørger man os, om en
% --- p. 181 ---
\opage{181}\setcounter{footnote}{0}%
saadan er tilladt, navnlig i det Store, naar det gjælder
et af Livets alvorlige Forhold, og hvor dens Anvendelse
ligefrem forraader Feighed, saa svare vi ubetinget: en
saadan Nødløgn er ikke tilladt. Gud kan umuligt bifalde
Sligt; thi af Sætningen „Gud vil det, fordi det er godt“,
følger jo omvendt: „Gud vil det ikke, fordi det er ondt.“
Seer man nu efter i første Mosebog tolvte og tyvende
Kapitel, hvad skal man saa dømme om Fader Abraham?
Det tolvte Kapitel begynder med Abrahams Udvælgelse;
Jehova taler til ham som til sin Udvalgte, og byder ham
at drage ud af Landet; han vil vise ham Veien og give
ham sin Velsignelse. Altsaa tænke man sig en Mand
der hæves saa høit, at Majestæten, den hellige Gud, vær\-%
diges at tale til ham. Saa kommer Abraham til Pharao i
Ægypten med Sara, og hvad gjør han nu? Han bliver
bange og siger til sin Hustru: „See nu, jeg veed,
at du er en deilig Qvinde at see til. Og det vil skee,
naar Ægypterne see dig og sige: Denne er hans Hustru,
at de slaae mig ihjel og lade dig leve. Kjære siig, at
du er min Søster o. s. v.“ Hvad er det Abraham her
gjør? Han gjør det, hvorover en Riddersmand vilde
rødme; han skjuler, at Sara er hans Hustru, ja hvad
mere er, han gjør sig endog en Profit deraf; thi til Ve\-%
derlag for Sara gav Pharao ham jo „Faar og Qvæg og
Æsler og Tjenere og Tjenestepiger og Aseninder og
Kameler“, saamange han vilde have; Pharao aner i sin
Troskyldighed ikke, at han er i Færd med at begaae en
stor Forbrydelse. Saa maa han faae det at vide paa en
anden Maade. For at standse de ubehagelige Følger af
Abrahams Feighed, maa Jehova til at hjælpe efter ved
Mirakler.

„Herren slog Pharao og hans Huus med store
Plager i Anledning af Sara, Abrahams Hustru.“ Hvor
maa dog Abraham ikke staae beskæmmet over for Pharao
% --- p. 182 ---
\opage{182}\setcounter{footnote}{0}%
og høre paa hans bebreidende Ord: „Hvi har du gjort
dette imod mig? Hvorfor har du ikke forkyndt mig, at
det er din Hustru? Hvi sagde du: hun er min Søster;
og jeg tog hende mig til Hustru, og nu see, der er din
Hustru, tag og gaae.“ At Jehova gaaer med i Intri\-%
guen, kunne vi til Nød forstaae, da Jehova nu engang
er Abrahams Ven; og hvad gjør man ikke for en Ven?
„Der ligger“, sige vore Alttroende, „en stor Trøst
deri, at vi see, hvorledes Gud kan bruge skrøbelige
Folk til sit Riges Støtter. Guds Førelse med Abraham
og at han anvender ham i sit Riges Tjeneste, er Forsvar
nok for Abraham.“ Men havde Abraham dog ikke for\-%
tjent en alvorlig Irettesættelse? Da han ingen Irettesættelse
faaer, er det naturligviis saa langt fra, at han fortryder,
hvad han har gjort, at han snarere i al Stilhed synes
at ønske sig selv til Lykke med at være første Op\-%
finder af Nødløgnen. I tyvende Kapitel spille saa han
og Jehova Komedien om igjen hos Abimelech i Gerar.
Og saa finder man det endda forunderligt, at engelske og
tydske Fritænkere have kunnet angribe det gamle Te\-%
stamentes positive Sædelære; og saa korser man sig over
at en Voltaire, ifølge det bekjendte \textit{taisez-vous, taisez-%
vous, Habakuk est capable de tout}, meente at kunne
lægge en Prophet de største Urimeligheder i Munden!

Hermed er Videnskritiken fyldestgjort. „Men“,
kunde Nogen spørge, „hvorfor skulle vi her sidde og
høre paa en saa forargelig Kritik, da vi nu eengang
vide, at den udspringer af den pure Vantro, og let kan
give de enfoldigt Troende Forargelse?“ Jeg svarer:
Med de enfoldigt Troendes Forargelse har det ingen
Fare; deres Tro sidder ikke saa løst, at en Smule Rai\-%
sonnementskritik strax skulde kunne rokke den. Deri\-%
mod er der al mulig Fare for, at man, under det Paaskud,
ikke at ville give de Enfoldige Forargelse, fortier Ind\-%
% --- p. 183 ---
\opage{183}\setcounter{footnote}{0}%
vendingerne, tilslører Modsigelserne, og holder det Hele
hen i en saa taaget Ubestemthed, at de Selvtænkende,
der gjerne ville oplyses om Sagens Sammenhæng, aldrig
kunne faae nogen Oplysning. Vi høre Frelserens Ord
(Mtth. XVIII, 7): „Vee Verden for Forargelser! Vel
er det fornødent, at Forargelser skulle komme; dog vee
det Menneske, ved hvilket Forargelsen kommer.“

Hvorfra kommer Forargelsen, naar Talen er om
\emph{den talende Gud}? Ikke fra Luther i Worms med
de „klare Grunde“ og den levende Troeserkjendelse,
% "Troeserkjendelse": a speck above the second e, taken for dirt, not an accent
nei ikke fra Luther, men fra Lutherolatrien, d. e. fra
den i Skoletheorien stivnede, krystalliserede Lutherdom,
fra Bibliolatrien, d. e. fra Bibeltheologiens videnskabe\-%
lige Forqvakling af Læren om Inspirationen. Man havde,
det fik dog Reformatorerne sat igjennem, omsider faaet
Bugt med den pavelige katholske Udvorteshed; men til
Erstatning for Tabet af den synlige Kirkes Ufeilbarhed,
klamrede den protestantiske Theologie sig saa igjen saa
krampeagtig til den bibelske Udvorteshed, at den er\-%
klærede hele den hellige Skrift for at være ordret dic\-%
teret af Gud selv (Solus Deus, si accurate loqui velimus,
Sacræ Scripturæ autor dicendus est, Prophetæ vero et
Apostoli autores dici non possunt, nisi per quandam cata\-%
chresin, utpote qui potius Dei autoris calami fuerint.
Quenstedt). Ved at gjøre Forfatterne af de hellige
Bøger til Guds og Christi „Haandskrivere“ (amanuenses,
manus, notarii), Guds og Christi „Penne“ (calami), ved
at lade den Hellig-Aand umiddelbart drive disse „Penne“
til at skrive (impulsus ad scribendum), indgive dem,
hvad de skulle skrive (suggestio rerum), og hvorledes de
skulle skrive (suggestio verborum), gjorde man i den
Grad det menneskelige Ord guddommeligt og det gud\-%
dommelige menneskeligt, at her end ikke blev Glimt af
% --- p. 184 ---
\opage{184}\setcounter{footnote}{0}%
Forskjel tilbage. Her er Vildfarelsernes første Grund, den
Jordbund, hvori Forargelsens Sæd saa rigelig har spiret.

Den orthodoxe Theologie opfatter Inspirationen me\-%
chanisk; Grundmodsigelsen er, at den gjør Aanden i
sin rigeste, reneste Aabenbarelse aandløs. Den gamle
Orthodoxie og den moderne Bevidsthed mødes i den
falske Forudsætning, at det Guddommelige og det Men\-%
neskelige maae falde sammen i Aanden. Forskjellen er
kun den, at Orthodoxien lader det Menneskelige gaae
under i det Guddommelige, saaledes at dette bemæg\-%
tiger sig Menneskeordets bestemte Form, medens den
moderne Bevidsthed lader det Guddommelige gaae saa\-%
ledes op i det Menneskelige, at dette i den rene Tænkning
faaer en abstract guddommelig Form. Men Eenheden
af Guds og Menneskets Aand er en Eenhed i væsentlig
Modsætning, en Modsætningseenhed. Det er Aands\-%
modsætningen mellem det Guddommelige og det Menne\-%
skelige, som fra begge Sider bliver overseet.

Vi ville, med Troesprincipet for Øie, betragte den
mosaiske Lovgivning. Culturbevidstheden kan ypper\-%
ligt forstaae, at Moses er sin Tids og sit Folks Mand,
at han bevæges af en stor Idee, at han har en Mission.
Den kan endvidere forstaae, at en af Ideen saa mægtig
bevæget Mand maa, idet han skal give sit Folk Love,
afpasse Lovgivningen efter Folkets aandelige Vilkaar, at
Lovene, udsprungne af en eneste stor Grundtanke, maae
forgrene sig vidt og udstrække sig til alle Livets
Forhold; den kan forstaae, at Grundtanken, naar den,
som i nærværende Tilfælde, er afgjørende religiøs,
maa fordre, at Lovgiveren taler, ikke i sit eget, men i
Guddommens Navn, og taler med guddommelig Myn\-%
dighed. Vi behøve blot at omskrive Begyndelses\-%
sætningen: „Og Gud talede alle disse Ord“, saaledes:
„Og Moses talede alle disse Ord i Jehovas Navn“, og
% --- p. 185 ---
\opage{185}\setcounter{footnote}{0}%
Culturbevidstheden med sit frie historiske Overblik vil
beundre Moses og see i hans Lovgivning et Aandens
Værk, et af Oldtidens største. Men videre kan Cultur\-%
bevidstheden ikke gaae. Og hvorfor ikke? Fordi den
ikke kan sætte sig ind i Gudsforholdet, ikke veed at tænke
Troens Tanke. Nogle Træk af Troespsychologien ville,
idetmindste foreløbigt, kunne oplyse Sagen.

Henføre vi Folkebefrielsens og Lovgivningens store
Grundidee, ikke til den menneskelige Bevidsthed i Moses
alene, men til Jehovas Villie, saa at det Hele oprindelig
udgaaer fra ham: saa er Forholdet imellem Jehova og
Moses ikke et almindeligt Ideeforhold, men et personligt
Aandsforhold, hvori det guddommelige Selv træder i
indre Vexelvirkning med det menneskelige. Paa Op\-%
fattelsen af denne Vexelvirkning beroer altsaa For\-%
staaelsen. Det stærkeste, mest umiddelbare Udtryk for
en saadan Vexelvirkning er en gjensidig Talen. Talen
% "Talen." : the mark sits on the baseline with no descending tail and a wide space
% follows; read as a full stop (both witnesses give a comma). Doubtful.
er Tankens umiddelbare Form; de Tvende meddele
hinanden deres Tanker derved, at de paa en eller anden
Maade tale med hinanden. Jehova taler til Moses,
Moses igjen til Jehova: hvad er det nu for en Talen?
Det er i Principet en indvortes Talen, en oprindelig
Meddelelse af Aand til Aand. Der er væsentlig For\-%
skjel mellem indvortes, aandelig og udvortes, sandselig
Talen. Naar to Individer tale til hinanden, er Talen ud\-%
vortes; de maae paa udvortes endelig Maade forstaae
hinandens Sprog. Fødte Landsmænd tale derfor lettest
og naturligst sammen, en Hebræer med en Hebræer, en
Græker med en Græker. Enhver af de Talende anvender
Sprogmidlerne paa sin Viis. Man kan opløse Talen i
saamange Slutninger, saamange Domme, saamange Sæt\-%
ninger, saamange Ord; man kan paa udvortes endelig
Maade forvisse sig om, hvad og hvormeget Enhver af de Ta\-%
lende har sagt. Anderledes med den indvortes Talen, den
% --- p. 186 ---
\opage{186}\setcounter{footnote}{0}%
Talen, hvori den guddommelige Aand meddeler sig til
den menneskelige. Her begynder det ikke umiddelbart
med Ord og Sætninger; men det begynder i Reflexion
med Tanker og Domme. De guddommelige Tanker og
Domme give sig først en indre, ideel, bevægelig Form
i Vækkelser, Indskydelser, indtrængende Grundforestil\-%
linger; Aandsmeddelelsen kan i denne Henseende siges
at gaae igjennem Ideen.
% "gaae": a small raised mark stands between the second a and the e (an
% accent-like speck, not a letter); transcribed without it.

Men nu gjør det en principiel Forskjel, at Ideen
som Idee er upersonlig, og forsaavidt selvløs, medens
Aanden som Aand er personlig et Selv. Idet Jehovas
Aand er Eet med det almægtige Selv, saa er den gud\-%
dommelige Selvbestemmelse udtrykkelig virksom i den
indre Meddelelse. Hvorledes de vækkende Indskydelser
og de stærke Grundtanker nærmere skulle bestemmes,
opfattes og tilegnes af den menneskelige Aand, kommer
ikke, saaledes som ved den blotte Ideebeaandelse, til at
beroe paa den opfattende Bevidsthed alene, men oprindelig
og afgjørende paa den guddommelige Villie. Det Næste, vi
maae fremhæve, er dette, at den, i hvis Indre Aanden
har begyndt en saa overnaturlig Virksomhed, oversætter
Indskydelserne, de primitive Forkyndelser i sit eget
menneskelige Sprog; heri ligner Aandsmeddelelsen igjen
en naturlig, f. Ex. digterisk Ideebevægelse. Og dog
er her, ifølge Principet, en afgjørende Forskjel. Den af
Ideen Bevægede har et Bevidsthedsindhold, der fordrer
og faaer sin Form af hans egen Bevidsthed; det Sprog,
hvori Ideen udtales, maa Digteren udtrykkelig erkjende
for sit Sprog: Ideen taler ikke til ham, men kun
igjennem ham, idet han taler ud af Ideen.

Med den af Aanden, Guds Aand, Bevægede for\-%
holder Sagen sig paa en anden Maade. Under Indvirk\-%
ningen af det absolute Selv antager det i Forstaaelsen
opfattede, i det menneskelige Sprog oversatte Indhold et
% --- p. 187 ---
\opage{187}\setcounter{footnote}{0}%
saadant Præg af guddommelig Gyldighed, at det ud\-%
trykkelig tager Sproget med dets Ord og Sætninger fra
Mennesket, gjør Menneskesproget til sit Sprog, Ord
og Tanker til sine Domme. I indre aandelig Mod\-%
sætning til det sig saaledes med Præg af guddomme\-%
lig Overlegenhed Fremstillende staaer den modtagende
Bevidsthed med sit reent menneskelige Indhold, sine
endelige Forestillinger, Tanker og Følelser som et
Væsen for sig, en afmægtig Aand, der tiltales af den
Almægtige.

Opgive vi Troesprincipet, saa bliver den til en saa\-%
dan Aandsfordobling svarende Umiddelbarhed øieblikke\-%
lig phantastisk; for ikke at faae to selvstændige Aander
i een Sjæl, forvandle vi da det guddommelige Selv til en
Indbildning, og fortolke Fordoblingen mythisk. I Kraft
af Principet er Fordoblingen derimod saa langt fra at
være mythisk, at den meget mere afgiver den absolute
Forudsætning for klar personlig Selverkjendelse. Staaer
% "personlig Selverkjendelse" is underlined in pencil in this copy; not transcribed.
det fast, at en alsidig, fuldstændig og tilstrækkelig be\-%
grundet Erkjendelse af det relative Selv kun er at op\-%
naae ved en indtrængende Erkjendelse af dets aandelige
Forhold til det absolute, d. v. s. af Gudsforholdet: da
er jo ogsaa den fremstillede Fordobling retfærdiggjort
med det Samme.

At det her beskrevne Inderlighedsforhold maa fra
Guds Side være et Aabenbaringsforhold, fra Menneskets
Side et Troesforhold, er klart af det Foregaaende. In\-%
derlighedsforholdet naaer sin høieste Spændkraft i In\-%
spirationen; kun paa Betingelse af Inspiration taler Gud
personlig til Mennesket, og Mennesket til Gud. Men
nu maa det Indre igjen have sit klare, umiddelbare Ud\-%
tryk i et tilsvarende Ydre; jo større Inderlighedskraften
er, desto stærkere er ogsaa den Umiddelbarhed, hvor\-%
med Udslaget skeer i det Ydre. Jehova taler ikke
% --- p. 188 ---
\opage{188}\setcounter{footnote}{0}%
umiddelbart til hele Israels Folk, thi Folket i sin Natur\-%
lighed er ikke inspireret; men han taler til Moses, thi
Moses er hans Prophet, og Propheten inspireret. Af
det Indre forstaaes det Ydre; den indre Vexelvirkning
i Aanden faaer som Aabenbaring fuld Virkelighed, idet
dens Væsen kommer tilsyne i tilsvarende Phænomener.
Den indre Talen bliver nu til en udvortes Talen; Ud\-%
vortesheden er saa stærk, at det for en overfladisk Be\-%
tragtning seer ud, som gjaldt det blot om sandselig
Seen og Høren. Det hedder (2 Moseb. XIX), at
„Moses steg op til Gud“, og at „Jehova raabte til ham
fra Bjerget.~.~.~. Og Moses kom og kaldte ad de Ældste
i Folket og han lagde alle disse Ord for dem, hvilke
Jehova havde befalet ham. Da svarede alt Folket og
sagde: alt det, som Jehova siger, det ville vi gjøre; og
Moses bragte Folkets Ord til Jehova. Og Jehova sagde
til Moses: see, jeg kommer til dig i en Skyes Mørke,
paa det Folket maa høre, naar jeg taler med dig, og de
ligeledes maae troe paa dig evindelig“. Paa en mere ud\-%
vortes og sandselig Maade kan et Aabenbaringsphænomen
vel ikke beskrives; og dog vil det, naar vi see nærmere
til, vise sig, at al denne Udvorteshed er Udslag af et
Indvortes, at Aabenbaringsphænomenet i al sin Virke\-%
lighed er et Aandsphænomen, et Phænomen, der for\-%
udsætter Tro og kun har Betydning for Troende.

Begynde vi med det mest Udvortes, det storartede
Naturoptrin, Uveiret med Lyn og Torden, da er den
umiddelbare Virkning heraf ligefrem den, at det indjager
det sandselige Menneske Skræk. I sin naturlige Raahed
er Hoben sandselig; det om Bjerget forsamlede Folk
bliver uvilkaarlig slaget med Skræk og Forfærdelse.
Men paavirket af det Ord, Moses taler, vækkes Aanden
i Folket; det stemmes i Tro til religiøs Opfattelse af
Naturphænomenet; det fornemmer, at Jehova er i Skyens
% --- p. 189 ---
\opage{189}\setcounter{footnote}{0}%
Mørke, og gribes af hellig Ærefrygt. Videre kommer
Folket ikke, og kan ikke komme videre. Den ydre Foran\-%
staltning, hvorved det (XIX, 12---13) fjernes fra Bjergets
Top, svarer ganske til dets aandelige Forfatning; det
overvældes af det Sandselige, aner det Hellige og betages
af Rædsel. „Men da Folket fornam Tordenen og Flam\-%
merne og Basunernes Lyd og Bjerget, som røg (XX, 18),
da Folket fornam dette, stod det langt borte.“ I sin For\-%
færdelse kan nu Folket forstaae, at det behøver en Tals\-%
mand, en med Aand begavet Mellemmand imellem sig
og Gud. „Og de sagde til Moses (XX, 19): tal du med
os, da ville vi høre, men lad Gud ikke tale med os, at
vi ei skulle døe.“

Følge vi i Tanken med Moses, der stiger op „til
Gud“ paa Bjerget, og spørge: hvad mon en Vidende,
f. Ex. en dristig Naturforsker, hvis han til samme Tid
havde foretaget en Bjergstigning paa Sinai, vel vilde have
seet og hørt? Et Uveir, et ganske naturligt, men forskræk\-%
keligt Uveir; intet Andet! At Uveiret her er et Almag\-%
tens Tegn, at den Almægtige selv er nærværende i
Uveiret, kan ikke sees med et naturvidenskabeligt Øie.
„Men Jehova taler jo høit, han raaber i Uveiret til
Moses?“ Ja, han raaber til Moses; men dette Raab kan
ikke høres af et naturvidenskabeligt Øre! Idet Moses
bestiger Bjerget og „gaaer nær ind til Mørket, der hvor
Gud er“, er han ikke objectiv rolig som en Naturforsker,
heller ikke betagen af Skræk som den raa Hob; men han
er bevæget i Aanden og frimodig i Troen: hvis hans Tro
vaklede og hans Aand blev svag, vilde han synke sam\-%
men i Forfærdelse. Det er i Uveirets Mørke under
umiddelbart Indtryk af den Almægtiges personlige Nær\-%
værelse, altsaa i Sjælekrisens allerhøieste Spænding, at
den indre Fordobling slaaer ud i mægtig Umiddelbarhed,
at den indvortes Talen faaer Præg af en udvortes, at Jehova
% --- p. 190 ---
\opage{190}\setcounter{footnote}{0}%
raaber til ham fra Skyen. „Og Gud talede alle disse
Ord og sagde“~.~.~.

I den mosaiske Lovgivning med de mange Bud
gaaer det Guddommelige saaledes i Eet med det Men\-%
neskelige, at Forskjellen og Modsætningen er skjult.
Denne Eensidighed charakteriserer Religionen paa det
gammeltestamentlige Trin og viser os den som en
forberedende, ikke fuldendende Aabenbaring. Hvad vi
herom kunne sige i al Korthed, er Følgende. Skulde
den til den indre Aabenbaring svarende Aandsfordobling
have været fuldkommen gjennemført, da maatte Moses
enten have været Christus selv eller en af hans Apostle.

Dette kan maaskee blive forstaaeligt, naar vi betragte
den til Fordoblingen svarende Modsætning mellem det i
Gudsforholdet absolute og relative Selv noget nærmere.
I Christus er det relative Selv evangelisk bestemt der\-%
ved, at „Menneskens Søn“ er „Guds Søn“ og Gud hans
„himmelske Fader“. Det relative Selv er altsaa i dette
sit Gudsforhold selv guddommeligt, er den paa Jorden
aabenbarede Gud selv. Sønnens Eenhed med Faderen
er udtalt i de Ord (Joh. X, 30): \textit{Jeg og Faderen vi
ere Eet}; hvorimod Sønnens Forskjel fra og Afhængighed
af Faderen --- altsaa Selvets Relativitet i Menneskesøn\-%
nen --- udtales i de Ord: \textit{Faderen er større end jeg.}
% The two sayings (Joh. X, 30; XIV, 28) are set in true italic, Danish text.

Heraf følger nu en afgjørende Forandring af For\-%
holdet mellem det Udvortes og det Indvortes i Aands\-%
fordoblingen. Christi Aabenbarelse er i en ganske anden
Forstand udvortes, end Aabenbarelsen ved Moses; thi
Christus fremtræder som den i Kjød og Blod aabenba\-%
rede Gud, som Ordet selv, der var fra Begyndelsen,
medens Moses derimod fremtræder som Menneske, som
Guds Ordfører for Israels Folk. Men just af denne
Grund er Aabenbarelsen i Christus i en ganske anden
Forstand indvortes end Aabenbarelsen ved Moses. Hvad
% --- p. 191 ---
\opage{191}\setcounter{footnote}{0}%
Christus tænker og taler som Gud, det tænker og taler
han tillige som Menneske; i Christi Indre er den ende\-%
lige Adskillelse imellem Gudstankerne og de menneske\-%
lige Tanker ophævet i personlig Eenhed. Naar Christus
taler, taler han i sit eget Navn. Det gammeltestament\-%
lige Udtryk: „Saa taler Jehova, Israels Gud“, er i det
nye Testamente afløst af det evangeliske: „Sandelig,
sandelig siger jeg Eder“. For en overfladisk Betragt\-%
% "jeg" is underlined in pencil in this copy; not transcribed.
ning tager dette sig nu ud, som var det gammeltesta\-%
mentlige Udtryk langt mere ophøiet og guddommeligt
end det nytestamentlige; men Sagen forholder sig om\-%
vendt. Det gammeltestamentlige Udtryk er en guddom\-%
melig Sanction af den prophetiske Tale; men Talen selv
i sin Form og sit endelige Indhold er en menneskelig
Tale. Moses og Propheterne tale med guddommelig
Fuldmagt i Jehovas Navn; men de tale folkeligt, men\-%
neskeligt om endelige, timelige, menneskelige Anliggen\-%
der. Den mosaiske Lov med dens mange Bud er given
af Moses, og Moses har taget alle mulige endelige Hen\-%
syn til sit Folks Vilkaar, dets Bestemmelse, dets Op\-%
dragelse; men dette Menneskelige i den mosaiske Lov\-%
givning er, vel at mærke, ikke saaledes adskilt fra de
guddommelige Grundtanker og Raadslutninger, at Moses
selv vilkaarlig har udtænkt dem og saa paa udvortes
Maade lagt Jehovas Navn til. Tvertimod! det er i Aand
og Sandhed guddommelige Tanker, som her ere menne\-%
skelig tænkte, guddommelige Bud, som her ere udtrykte
i menneskelige Ord. Ligeover for disse, med guddom\-%
melig Magt og Myndighed beseglede, Jehovatanker har
Moses i sin Bevidsthed en Kreds af andre Tanker, saa\-%
danne, som han i endelig Forstand maa kalde sine egne
Tanker. Endnu maae vi, i Conseqvens af Principet,
fremhæve, at den Adskillelse, en Moses gjør imellem sine
egne Tanker og de ham aabenbarede Jehovatanker, ud\-%
% --- p. 192 ---
\opage{192}\setcounter{footnote}{0}%
trykkelig bestemmes ved hans Tro og de til Troen sva\-%
% "til Troen": a raised, comma-like mark prints between the two words with no
% space (a risen space or stray sort, not punctuation); not transcribed.
rende Indskydelser.

Hvad endelig Lovens Indhold angaaer, vil det let
forstaaes, at Budenes saakaldte Vilkaarlighed ikke er
noget Beviis imod deres sædelige Værd. Ethvert Bud,
selv det, hvis Nødvendighed allerklarest indlyser for For\-%
nuften, maa som Bud, d. v. s. Befaling, antage Præg af
Vilkaarlighed. Sædelighedsideen er i afgjørende Forstand
Villiens Idee; Ideen har kun absolut Gyldighed derved,
% A pencil bracket in the margin marks these two lines; not transcribed.
at den sættes i Eet med den absolute Villie. Men heraf
følger igjen, at alle Sædelighedens oprindelige Fordringer
maae udtrykkes i Form af guddommelige Bud. Den
sædelige Sætning: „Gud vil det, fordi det er godt“, er
saa langt fra at udelukke den religiøse Sætning: „det er
godt, fordi Gud vil det“, at den sidste tvertimod er nød\-%
vendig til at sikkre den første. Da det ved alle Sæde\-%
lighedsbestemmelser væsentlig kommer an paa Villie, saa
kommer det i absolut Forstand først og sidst an paa den
guddommelige Villie.

Hvad her er sagt, maa være tilstrækkeligt til at
charakterisere det Troesproblem, vi kunne fastholde
under Navn af det religiøse Vilkaarlighedsproblem, et
Problem, der til sin fuldstændige Løsning vil kræve en
gjennemført Troespsychologie, Religionsphilosophie og
Ethik.

Af de to ovenfor opstillede Sætninger lyder den
sidste: „det er godt, fordi Gud vil det“, paa Optugtelse,
den første: „Gud vil det, fordi det er godt“, paa den til
Optugtelsen svarende Selvudvikling; hiin er den religiøse
Tvangssætning, denne den religiøse Frihedssætning; hiin
svarer til Lovens, denne til Evangeliets Aand. Det er
da i sin Orden, at Sætningen: „det er godt, fordi Gud
vil det“, eensidig sættes i Kraft paa det gammeltesta\-%
mentlige Trin. Med denne Grundtanke vil man forstaae,
hvi den mosaiske Lovgivning med sine mange guddom\-%
% --- p. 193 ---
\opage{193}\setcounter{footnote}{0}%
melige Bud maa gaae saa udførlig ind i Enkeltheder.
„Du skal“ og „du skal ikke“ gjentager sig i saa mang\-%
foldige Forgreninger, at alle Livets Anliggender umid\-%
delbart og middelbart inddrages derunder. Alt, ikke
blot hvad der hører til Gudsdyrkelsen med dens Offringer
og Ceremonier, men ogsaa hvad der angaaer de huus\-%
lige og borgerlige Forhold, Alt er punktlig bestemt ved
Jehovas Villie. Ved hvert Skridt, Israeliten gjør paa
Virkelighedens Vei, møder han et guddommeligt Bud;
formedelst det guddommelige Bud bliver det Timelige Led
for Led et Mellemværende imellem ham og den Evige. I
Modsætning til Præsterne, der skulde holde over Loven og
Budenes Opfyldelse, optraadte Propheterne som Talsmænd
for Aandens Ret; men Præsterne overlevede Propheterne,
Værkhelligheden med den pharisæiske Retfærdighed fik
Overhaand: hermed var Grændsen for den optugtende
Vilkaarlighed sat; Lovens Tid var omme. Ved Evan\-%
geliet blev „Budenes Lov“ forvandlet til Kjærlighedens
Lov. Med Kjærlighedsloven til Grundlag kunde Selv\-%
udviklingen forenes med den frie Lydighed, og Sætnin\-%
gen: „Gud vil det, fordi det er godt“, forklare Vilkaar\-%
ligheden af guddommelig Viisdom.

Man vil, som det synes, fra alle Sider lægge an
paa at deducere Villien af den rene, almindelige Tænk\-%
ning, ja man fordrer endogsaa, at Logiken skal begynde
med „den gode Villie“; saavidt er det kommet. Det er
maaskee endnu det mest talende Beviis for, at det sidste
Glimt af Troestanken er i Begreb med at forsvinde. At
Nutiden ikke kan fatte Aabenbaringens Ord om „den
talende Gud“ og forlige den menneskelige Tanke med
den guddommelige Vilkaarlighed, er en væsentlig Hindring
for det aandelige Liv, en Hindring, som kun Villiens
Aand kan forvandle til Betingelse.

\begin{center}\rule{3cm}{0.4pt}\end{center}

\chapter*{XII.}
\addcontentsline{toc}{chapter}{Forelæsning XII}
\markboth{Forelæsning XII}{Forelæsning XII}
\label{ch:xii}
% pp. 194--207
% --- p. 194 ---
\opage{194}\setcounter{footnote}{0}%
% [large initial I]
I andet Brev til Timotheus (III, 16---17) angives et Syns\-%
punkt for Opfattelsen og Fortolkningen af den hellige
Skrift og dermed af hele den hellige Historie. „Den
ganske Skrift“, hedder det, „er indblæst af Gud og nyttig
til Lærdom, til Overbeviisning, til Rettelse, til Optugtelse
i Retfærdighed, at det Guds Menneske maa vorde fuld\-%
komment, dygtiggjort til al god Gjerning.“ Ved dette
evangelisk-apostoliske Synspunkt er en afgjørende Skils\-%
misse mellem Troens og Videns Princip tilstrækkelig
angiven.

Hidtil har man imidlertid ikke villet agte paa denne
Skilsmisse; det vil sige, man har omskrevet Modsætnin\-%
gen saaledes, at dens Led kun afgive to Sider istedenfor at
udvise to, hinanden modsatte, Sphærer. Den Lærdom og
% "Sider", "Sphærer": pencil underlining in this copy; not transcribed.
Overbeviisning, hvorom Brevet taler, er uendelig for\-%
skjellig fra den Lærdom og Overbeviisning, hvorom der
handles i Videnskaben. Vistnok har man gjort opmærk\-%
som paa, at Udtrykkene „Lærdom og Overbeviisning“
tyde paa det Theoretiske, medens Udtrykkene „til Ret\-%
telse, til Optugtelse i Retfærdighed“ lyde paa det Prak\-%
tiske. Men ifølge denne blot formelle Forskjel har man
saa raisonneret saaledes: Den hellige Skrift fordrer Lær\-%
dom, grundig Lærdom, Overbeviisning, sand og levende
% "Lærdom, Overbeviisning": a risen space (quad) shows between the words; not typed.
% --- p. 195 ---
\opage{195}\setcounter{footnote}{0}%
Overbeviisning; men hvor skulle vi søge den grundige
Lærdom, den sande og levende Overbeviisning, uden alene
i Videnskaben, den Videnskab, der forudsætter Troen,
% "Videnskaben, den": a risen space (quad) shows after the comma; not typed.
forklarer og begrunder Troen, nemlig Theologien. Dog,
vi skulle ikke blive ved Videnskaben alene, vi skulle jo
ogsaa arbeide for Livet, til Theorien maa svare en Praxis;
i Praxis have vi det Pastorale, det Opbyggelige; det er
dette, der er betegnet ved de Ord „til Rettelse, til Op\-%
tugtelse i Retfærdighed“. Man har tilføiet, at ligesom
den gamle Orthodoxie eensidig holdt sig til Lærdommen,
saaledes har Pietismen eensidig holdt sig til det Praktiske,
„det Opbyggelige“. „Men“, hedder det, „vi oplyste
Theologer holde Ligevægten; vi have „Lærdom“ og
„Overbeviisning“ i det Objective og „Rettelse og Optug\-%
telse“ i det Subjective, hvad kan saa kræves mere?“

Imod denne Misforstaaelse opstille vi Adskillelsen
mellem to absolut ueensartede Standpunkter, det viden\-%
skabelige og det religiøse, af hvilke det første betinger
Erkjendelsen af den upersonlige, det sidste af den per\-%
sonlige Sandhed.

Naar vi kalde den videnskabelige Sandhed uperson\-%
lig, saa er Meningen naturligviis ingenlunde den, at
videnskabelig Tænkning skulde kunne foregaae uden i en
tænkende Bevidsthed, at der skulde kunne gives en Viden
uden en Vidende, en Overbeviisning uden en Subjecti\-%
vitet, der følte sig overbeviist; men Meningen er, at den
Vidende maa under sin Videnskabs Udvikling opgive
ethvert personligt Selvhensyn, maa, fordybet i Sagens, i
Tingenes, Lovenes, Kjendsgjerningernes Erkjendelse, ikke
lade sig forstyrre af nogetsomhelst Hensyn til nogensom\-%
helst personlig Interesse. En Physiolog, f. Ex., kan som
Menneske, som existerende Personlighed, være uendelig
interesseret i Spørgsmaalet om Sjælens Udødelighed; men
hvis han uden videre lader denne Interesse spille en
% --- p. 196 ---
\opage{196}\setcounter{footnote}{0}%
Rolle med i de reent videnskabelige Undersøgelser, bliver
han en maadelig Physiolog. En Astronom kan som
Menneske, som existerende Personlighed, føle Trang til
at forvisse sig om, hvorvidt der i Virkeligheden gives en
Verden af høiere Orden end denne timelige, hvori
vi leve; men hvis han, ledet af denne Trang, retter
Kikkerten imod Melkeveien og begynder at opløse Stjer\-%
netaagerne til Fordeel for Læren om et „Himmeriges
Rige“, vil han neppe ret længe bevare sin Anseelse blandt
Astronomerne. En Geolog kan som Menneske, som
existerende Personlighed, føle sig saa stærkt tiltalt, saa
dybt greben af den bibelske Fortælling om Paradiset, at
han ret kunde ønske, hvis det var muligt, at finde et
Punkt i Jordudviklingen, hvorom han kunde sige: her
er Sporet af den paradisiske Tilstand; her, paa dette
Sted, har Edens Have ligget! Men lader han slige
Længsler og Bestræbelser faae Indflydelse paa Bedøm\-%
melsen af de geologiske Kjendsgjerninger: da viser
han sig som en høist utilforladelig Forsker. Vi haabe
at blive forstaaede, ikke blot af Videnskabens Dyrkere,
men af alle Dannede, naar vi i den her angivne Mening
gjentagende erklære, at den videnskabelige Sandhed er
og maa være upersonlig.
% pencil underlining and a marginal stroke in this copy on p. 196; not transcribed.

At vi, i Modsætning til den videnskabelige, kalde
den religiøse Sandhed personlig, maa ikke forstaaes, som
var Meningen den, at den religiøse Sandhedserkjendelse
kun skulde beroe paa individuelle Forestillinger, Følelser,
Stemninger og slige ubestemmelige Forudsætninger; nei,
Meningen er, at den i Personlighedens ideale Fordringer
absolut concentrerede Selvhed betegner Standpunktet.
Hvad der fyldestgjør den absolute Personlighedsfordring,
er gyldigt; hvad der strider imod den, er ugyldigt; hvad
der hverken befordrer den eller strider imod den, er
ligegyldigt.

% --- p. 197 ---
\opage{197}\setcounter{footnote}{0}%
For at opklare, hvad heri ligger, behøve vi ikke at
see os om efter nye og forviklede Exempler; vi ere i
det Foregaaende paa vor Vei gjennem den hellige Historie
allerede stødte paa saamange Knudepunkter, saamange
endog haarde Anstødsstene, at det vel kunde være hen\-%
sigtsmæssigt at optage de allerede anførte Exempler fra
et nyt Synspunkt; om der da maaskee kunde falde et
nyt Lys paa, hvad der tidligere var gaadefuldt. Saale\-%
des saae vi i Læren om Skabelsen, at den astronomiske
Fremstilling af Verdens Bygning er i afgjørende Strid
med den Fremstilling, der gives os i første Mosebog;
vi saae, hvorledes Knuden blev knyttet navnlig ved
Christi Ord i Matthæus (XXIV, 29), hvor han forudsiger
Verdens Undergang og udtaler sig saaledes: „Strax efter
de Dages Trængsel skal Solen formørkes, og Maanen ikke
give sit Skin, og Stjernerne falde fra Himmelen.“ Ved
et Udtryk som dette, at Stjernerne engang skulle falde
fra Himmelen, bliver Striden mellem Videnskaben og
Religionen saa haard, at det synes, som om den aldrig
kunde hæves, som om et Forlig var umuligt. Sige vi:
Stjernerne kunne falde ned fra Himlen, saa modsige vi
Astronomien i et afgjørende Punkt og bryde derved
Staven over Cultur og Videnskab. Sige vi derimod, at
% "Cultur": a single faint speck above the u (not a diaeresis; 300 ppi); the layer's "Cültur" rejected.
hiint Ord af Christus ikke kan gaae i Opfyldelse, saa
krænke vi Christi Majestæt, krænke Religionen og kaste
Foragt paa det aabenbarede Ord. Paa disse Vilkaar
synes enhver Udvei til Forsoning af Tro og Viden
afspærret.

Imidlertid er Afspærringen mere tilsyneladende end
virkelig. Denne Himmel og disse Stjerner sees her fra
aldeles forskjellige Standpunkter, med ganske forskjellige
Øine; derfor tales der om dem i forskjellige Sprog.
Det er jo vitterligt for Alle, at det Synlige i Verden
skal sees af et Øie. Som Øiet er, der seer, saaledes er
% --- p. 198 ---
\opage{198}\setcounter{footnote}{0}%
Synsbilledet, saaledes er den synlige Verden; som Øret
er, der hører, saaledes er Hørelsen; som Bevidstheden
er, der betragter, saaledes er Billedet af det Betragtede.
Hvor forskjelligt Billedet af Himlen og Jorden maa blive,
eftersom Betragterens Standpunkt forandres, kunne vi
særlig indprente os i denne Sammenhæng ved at adskille
tre Standpunkter, nemlig det videnskabelige, det æsthe\-%
tiske og det absolut religiøse.

Stille vi os paa det videnskabelige Standpunkt, hvad
betragte vi da som det Første? Tingene selv, deres
Love og Sammenhæng! Det er den udvortes Himmel,
disse bestemte Stjerner, disse Bevægelses-Phænomener
med de tilsvarende Love, som ere Betragtningens Gjen\-%
stand. Betragter da den Betragtende sig selv med det
Samme? Nei, for den, der videnskabelig seer paa Stjerne\-%
himlen, er Stjernehimlen Eet og Alt. Under Iagttagel\-%
serne tænker Astronomen jo slet ikke paa sin Personlighed,
ikke paa egne Ønsker og Længsler. Slige Længsler
vilde han, hvis Nogen var saa sentimental at minde ham
om Hjertets Trang, afvise med isnende Kulde. Derfor
er Minervas Bryst saa koldt, og hendes Blik saa strengt;
derfor kaldes Videnskaben objectiv, at Betragteren, for\-%
dybet i sin Gjenstand, skal glemme sin Subjectivitet.
Astronomien er, ligesom enhver anden Videnskab, uper\-%
sonlig, og i sit upersonlige Væsen objectiv, ligegyldig mod
vore sjælelige Rørelser, vore dybeste Længsler, vore Smerter
og Glæder, vor Frygt og vort Haab, saa ligegyldig, som
var det den rene Skjæbne.

Hvad er det vel for personlige Resultater, der kunne
udledes af Verdensbygningens Betragtning? Resultaternes
Resultat er, at Mennesket svinder ind til et Intet. Vi
blive i Forhold til Jorden langt mindre end Myrerne i
deres Tue; Jorden selv forsvinder i Universet. Altsaa
% "deres Tue": a risen space (quad) shows before "Tue"; not typed.
maa Forfriskning og Befrielse søges deri, at vi, isteden\-%
% --- p. 199 ---
\opage{199}\setcounter{footnote}{0}%
for at fastholde vort eget Selv, fortabe os i det Selvløse
% "Selvløse" at line end: no comma printed (image, 300 ppi).
at vi, istedenfor at være selvstændige indenfor vor person\-%
lige Begrændsning, frigjøre os i det Grændseløse. Det er
mærkeligt, hvorledes denne Fordybelse i Himmelrummets
ydre Uendelighed, denne sublime Selvfortabelse, gjengiver
en orientalsk Anskuelse, der allerede tilhører Oldtiden;
det er Brahmaismen, hvortil vi sigte.

Som bekjendt ere Østerlænderne tilbøielige til i
dybsindige Phantasier at slaae en Streg over det Ende\-%
lige og fordybe sig i Skuet af det Evige, det Grændse\-%
løse, det Uendelige. Brahmas Præster betragte denne
Verden i dens Endelighed som et skjønt Bedrag. Dens ydre
Mangfoldighed, dens mange forskjellige Farver og Toner
betyde, at der er bredt et Tidens Slør over den evige
Sandhed; det er Maias Slør, Blendværkets, Bedragets
glimrende Slør. Det er en Opgave for den Vise, at
løfte dette Slør. Først hæver han Sandsebedraget ved
at frakjende de sandselige Ting Realitet; men det er
% "men det": a raised dot-like mark (risen space) before "det"; not typed.
endnu ikke nok: han maa ogsaa frigjøre sig fra Tanke\-%
bedraget. Tankebedraget varer saalænge, som han ad\-%
skiller bestemte Tanker, den ene begrændset ved den
anden. Saa lader han da alle Tanker flyde sammen i een
Tanke, Uendelighedens Tanke; saa nyder han i sin Selv\-%
opløsning Beskuelsens Salighed, taber sig i det Grændse\-%
løse og forsvinder i Brahm. Her er en Frigjørelse, der
i indvortes Forstand punktlig svarer til den Frigjørelse,
Astronomerne og Physikerne anbefale os i udvortes
Forstand, naar de anbefale os at fornegte hvad de kalde
vor Egoisme og forsvinde i det uendelige Verdensalt.
Dette var nu det første Standpunkt.

Forskjelligt fra det strengt videnskabelige, objective
Standpunkt er det æsthetiske. Paa det æsthetiske Trin
opgiver Personligheden vel ikke sig selv, men den op\-%
hæver i Stemningens Øieblikke sin virkelige Begrænds\-%
% --- p. 200 ---
\opage{200}\setcounter{footnote}{0}%
ning; den udvider sig phantastisk og favner i Selvud\-%
videlsen Himmel og Jord. Det Charakteristiske er, at
Menneske-Aanden seer paa Himlen ideale Skikkelser,
der engang have levet og virket paa Jorden; herved faae
de os bekjendte Stjernebilleder en ganske særegen Be\-%
% "bekjendte Stjernebilleder": a risen space (quad), not a full stop; not typed.
tydning. Vi kjende jo saadanne Stjernebilleder paa den
synlige Himmel: Cepheus, Cassiopeia, Andromeda, der
er lænket til Klippen, Perseus, Hesten Pegasus, den
knælende Herakles, den lernæiske Slange, o. s. v. Disse
Navne minde os om de mythiske Heroer og deres Be\-%
drifter. Hvad betyder det? Det betyder, at idet Men\-%
nesket seer op paa Stjernehimlen, seer han et Speil, der
gjengiver Personlighedens Billeder, en Afglands af det
personlige Liv i Tiderne. „Stjernebillederne siger J. L.
% printer's defect: no closing “ after "Stjernebillederne" (image, 300 ppi);
% transcribed as printed. Quote balance on this page +1.
Heiberg (Intelligensblade, 2det Bind, S. 41), „ere sande
Billeder, som Oldtiden har ophængt i det store Gallerie,
hvor alle Slægter skulle betragte dem. Her viser det sig
ret, at Stjernehimlen, langt fra at være fremmed for Jorden,
er Reflexen af de jordiske Skjæbner, de jordiske Tanker,
de jordiske Drømme. Man maa betragte Stjernehimlen
som en illustreret Pragtudgave af den græske Mythologie
i uhyre Format; man maa betragte den med Oldtidens
Digtere for Øiet“. Det er ikke den Enkelte, der speiler
sig selv, det er Oldtiden og dens store Skikkelser, han
seer i Speilet.

Her er altsaa Betragtningen fordoblet, og Stand\-%
punktet forvandlet; thi naar der skal tales om de Skik\-%
kelser, der sees i Stjernerne, skal der tales, ikke blot
fra Poesiens Standpunkt, men i Poesiens Sprog. Vi
kjende vist alle Heibergs deilige Digt, i hvilket han
blandt Andet lader Stjernehimlen sige til Betragteren:
„Her Fabelens Konger og Dronninger staae, med dens
Udyr, Drager og Slanger, hvis Navne lød fra den Old\-%
tids graa, to tusindaarige Sanger“. Her er det kolde
% --- p. 201 ---
\opage{201}\setcounter{footnote}{0}%
System af Phænomener og Love forsvundet; naar Dig\-%
teren taler om Fabelens Konger, dens Udyr og Slanger,
saa taler han netop i sit eget Sprog, thi Sproget maa
svare til Standpunktet. I dette Stemningens Sprog paa\-%
kaldes Erindringen. Mennesket veed sig selv begrændset,
ikke blot i Forhold til udvortes Storhed i Rummet, men og\-%
saa i Forhold til Tiden. Det lever en stakket Tid, dets Liv
er „som Græsset, som Markens Blomster; Solen gaaer
op med Hede, Blomsten visner og falder af“. Vi leve i
Døgnet det korte Liv; skal Personligheden udvide sig og
aande frit i Aandens Rige, maa den tye til Erindringen.
I Erindringen vender den sig mod det Svundne, vandrer
ad Forbigangenhedens Veie, mødes med Oldtidens Store
og forstaaer den i Digtet talende Himmel: „Hos mig Du
lever et Liv som svandt, Du lever med Guder og Helte,
hvis forklarede Skikkelser Sagnet bandt til mit brede,
funklende Belte.“

Ingen vil sige, at dette Poesiens Standpunkt kommer
i Strid med Astronomien. Her er ingen Strid; thi de
to forskjellige Standpunkter ere ganske ueensartede.
Astronomen vilde kun gjøre sig latterlig, hvis han i sit
prosaiske Sprog begyndte at føre Beviis mod, hvad
Digteren her siger om Fabelens Udyr, Drager og Slanger,
medens omvendt Digteren vilde gjøre sig latterlig, hvis
han ved Fabelens Konger og Dronninger vilde modbe\-%
vise Astronomen. Der er to Standpunkter, to Betragt\-%
ninger, to Sprog og to forskjellige Billeder af den be\-%
tragtede Verden. Hver af dem har sin Gyldighed. Vi
ville ikke undvære Astronomien, men heller ikke Poesien.

Dette kan nu Enhver forstaae; men skulde det saa
ikke ogsaa kunne blive forstaaeligt, naar vi paa tilsva\-%
rende Maade mærke os det tredie Standpunkt, det reli\-%
giøse, ikke det relativt religiøse, der er svækket i men\-%
% --- p. 202 ---
\opage{202}\setcounter{footnote}{0}%
neskelig Viden; nei det absolut religiøse, Christi eget
Standpunkt.

Hvis Christus ikke stod paa et eget Standpunkt,
kunde han jo ikke være Christus; men staaer han paa
et Standpunkt, som ingen Dødelig med ham kan ind\-%
tage, maa det vel medføre, at der til det eiendommelige
Betragtningens Standpunkt svarer et eget Billede af
Himmel og Jord. Hvad dette er for et Standpunkt, er
allerede tidligere sagt: det er Personlighedens! Religio\-%
nens Princip er Personlighedens Princip i allerdybeste
Forstand.
% marginal pencil stroke in this copy beside the last three lines; not transcribed.

Det er jo charakteristisk for den nyere Tid at have
Respect for Love, „de evige Love“; naar Nogen vil
modsige Physikeren, peger denne blot paa de evige
Love, og Modsigelserne forstumme. Naar vi nu
høre Christus og see, hvad der ligger i hans Ord,
peger ogsaa han hen til evige og uforanderlige Love;
men det er vel at mærke ikke Nødvendighedens men
Frihedens Love. Her er to Systemer af Love, fordi
her, efter den menneskelige Erkjendelses Vilkaar, maa
være to forkjellige Standpunkter. Her er Nødvendig\-%
% "forkjellige": so printed (s dropped; image, 300 ppi); transcribed as printed.
hedens Love, svarende til Naturriget; men her er ogsaa
Frihedens, Aandens og Personlighedens Love, der gjælde
i Christusriget. Det er Frihedens og Personlighedens
Love, hvorom Christus taler, naar han siger: „Himmel
og Jord skulle forgaae, men mine Ord skulle ikke for\-%
gaae“. Charakteristisk for dette Standpunkt er, at det
er den personlige Sandhed selv, der taler om Himmel
og Jord, og taler Personlighedens Sprog, myndigt af\-%
gjørende. Den Talende imponeres ikke af Verdens\-%
bygningens udvortes Storhed; thi den udvortes Verden,
dens Sol og Maane og Stjerner, underligge Nødvendig\-%
hedéns Love, men Sandheden selv veed med evig, gud\-%
% "Nødvendighedéns": so printed, a clear acute on the second e (crop at 600 dpi); printer's defect, transcribed as printed.
dommelig Vished, at Nødvendighedens Love ere grun\-%
% --- p. 203 ---
\opage{203}\setcounter{footnote}{0}%
dede paa Frihedens Love. Medens det naturlige Men\-%
neske maa føle sig at svinde ind til et Intet ligeover\-%
for Naturens Magt og Majestæt, staaer Christus, som
% "Naturens Magt": a risen space (quad) shows after "Naturens"; not typed.
Personlighedens absolute Repræsentant, urokkelig fast
paa Frihedens Grund. Lod han sig overvælde af
stærke Naturindtryk, af Verdensbygningens ydre, uende\-%
lige Storhed, saa var han ikke Christus.

Man skjønner imidlertid let, at Modsætningen imellem
den religiøse og den videnskabelige Anskuelse af Ver\-%
densbygningen er ganske anderledes haard end Modsæt\-%
ningen mellem den videnskabelige og den æsthetiske. I
den æsthetiske Anskuelse er det jo ikke fuldt Alvor med
Virkeligheden; Poesien paastaaer ikke, at „Fabelens
Konger og Dronninger“ virkelig staae paa Himmelhvæl\-%
vingen: det æsthetiske og det videnskabelige Standpunkt
ere med andre Ord, blot relativt ueensartede. I den re\-%
ligiøse Naturbetragtning er det derimod ligesaa afgjø\-%
rende Alvor med Virkeligheden som i den videnskabe\-%
lige. Spørge vi, om det kan antages, at Himmel og
Jord virkelig engang vil forgaae, saaledes som Chri\-%
stus har forudsagt os det: da svarer den paa sit Eget
reflecterende Videnskab med et udtrykkeligt Nei, den
paa sit Eget reflecterende Religion derimod med et ud\-%
trykkeligt Ja. Fra Videns Standpunkt kan Christi
Forudsigelse ikke være sand; og hvorfor ikke? Fordi
dens Opfyldelse er en Umulighed! Fra Troens Stand\-%
punkt maa Christi Forudsigelse absolut være sand; hvor\-%
for? Fordi Christus er Sandheden! Saa staaer Viden\-%
% "er Sandheden": pencil underlining of "er" in this copy; not transcribed.
skab mod Religion som Ja imod Nei, som Sandhed
imod Sandhed: vil Nogen kunne negte, at her er en
Dualisme?

De, der lægge an paa at undergrave den positive aaben\-%
barede Religion, kunne naturligviis tage sig Dualismen let;
de kunne efter Behag snart ophæve den ved træge Ab\-%
% --- p. 204 ---
\opage{204}\setcounter{footnote}{0}%
stractioner, snart igjen more sig med at gjøre den til
Gjenstand for karikeret Udpegen og fnisende For\-%
vrængning. Men de, som ville holde paa Religionen uden
at opgive Videnskaben, Theologerne f. Ex., hvad mon
de tænke paa? Kunne de virkelig indbilde sig selv og
Andre, at deres „gode Villie“, naar den blot faaer Lov
at komme med i Videnskaben, vil være istand til at
magte Dualismen?

Dualismen er ikke et vilkaarligt Paafund, ikke en
philosophisk Chicane, den er en Kjendsgjerning, en
vitterlig uimodsigelig Kjendsgjerning. Til at forsone sig
med Dualismen, eller, som det i Almindelighed udtrykkes,
forsone Troen med Viden, gives der kun een Udvei, een
eneste. Vi skulle med Hensyn paa vort opstillede Ex\-%
empel belyse Sagen lidt nærmere. Naturvidenskabens
Opfattelse af Verdensbygningen er ikke religiøs; Christi
Anskuelse af Verdensbygningen er absolut religiøs: her\-%
fra kommer Modsigelsen. Denne Modsigelse kan kun
løses derved, at Christus i den personlige Sandheds
Navn faaer ubetinget Ret, og Videnskaben med sin
upersonlige Sandhed videnskabelig Ret.

For at give Christus ubetinget Ret, maae vi søge
at forstaae hans Grundtanke, det Sprog, hvori Grund\-%
tanken udtales, og den til Udtalelsen svarende Virke\-%
lighed.

Hvad Christus med absolut personlig Vished ud\-%
taler, er dette, at Naturriget med alle udvortes sandselige
Ting er Aandens Almagt underlagt, at Aandens Rige,
det Rige, hvis Stifter han er paa Jorden, skal, naar
Timen kommer, virkeliggjøres, ikke blot i det Indvortes,
men ogsaa i det Udvortes. Denne Vished har han ikke
fra nogen umiddelbar Naturbetragtning, saa lidt som fra
astronomiske Iagttagelser af Sol, Maane og Stjerner; nei,
han har Forvisningen, den fulde, afgjørende Forvisning
% --- p. 205 ---
\opage{205}\setcounter{footnote}{0}%
i sig selv, ved sig selv, han har den personlig. Troen
paa Christi Ord er ogsaa i dette Punkt Eet med Troen
paa Christus selv; ogsaa heri viser han sig som Forbil\-%
ledet, som Personlighedens virkeliggjorte Ideal, at han
med en saa uendelig Overlegenhed kan tale om Frihedens
Magt over den i Nødvendighed bundne Natur.

De skrøbelige Individualiteter, de relative, i Skyld
og Selvmodsigelse hildede Personligheder ere, trods al
deres Higen mod Aandens Maal, idelig udsatte for at
overvældes af naturlige Hindringer og kues af overmæg\-%
tige Sandseindtryk. Det er især i mørke, tunge Øieblikke,
at Trykket føles. Forsagtheden kommer med Tvivlen;
og hvorfra kommer Tvivlen? Fra en Splid i Menneskets
Indre! Og hvorfra kommer Spliden? Fra en Strid
imellem de to absolut ueensartede Principer, Nødvendig\-%
hedsprincipet, med de tusinde Beviser fra den sandselige,
materielle Verdens Realitet, og Frihedsprincipet med sine
Fordringer og ideale Tilskyndelser fra det sjælelige Selv.
Her behøves et myndigt Ord, ikke blot til „Lærdom“,
men ogsaa til „Opbyggelse.“

I Christi Ord er der Lærdom, og Lærdommen er
denne: Vil du have fuld Forvisning om Aandens ube\-%
tingede Magt over den forkrænkelige Natur, da grib
Frihedsprincipet og vogt Dig for at stirre paa Nødven\-%
digheden! med andre Ord: vil Du finde Veien til Fri\-%
hedens, Sandhedens og Personlighedens virkelige Rige,
da forvild Dig ikke ind i de upersonlige Erkjendelsers
Egne, søg ikke Veiledning hos Astronomer og Physikere
og Naturphilosopher, men vend Dig til ham, der er Veien
og Sandheden og Livet! I Christi Ord er der „Opbyg\-%
gelse“. Det Opbyggelige er, at han udtaler Forvisningen,
ikke som en almindelig Livsbetragtning i en ved lykke\-%
lige Varsler og lyse Udsigter vakt og bevæget Stemning,
men i det tunge, forfærdelige Øieblik, i et Øieblik, da
% --- p. 206 ---
\opage{206}\setcounter{footnote}{0}%
han maatte sige til sine Fjender: denne er Eders Time
og Mørkets Magt; det er just et saadant Øieblik, der
giver Forvisningen den indre Personlighedens Vægt, som
veier op mod det udvortes Virkelige.

Som Grundtanken er, maa ogsaa Sproget være.
Hvad Christus seer i Aanden, de mægtige Syner, de store
Billeder, hvori det Tilkommende faaer aandig umiddel\-%
bar Nærværelse, det udtaler han i Ord, der endnu den
Dag idag kunne gjøre Anskuelsen levende. Hans Taler
til Disciplene om Jerusalems Forstyrrelse og Verdens
Undergang og hans Vidnesbyrd for Ypperstepræsterne i
Raadet føres i samme Sprog. „Og den Ypperste-Præst
% four dots, as printed:
.... sagde til ham: jeg besværger dig ved den levende
Gud, at du siger os, om du er den Christus, den levende
Guds Søn. Jesus sagde til ham: du haver sagt det;
dog siger jeg Eder: nu herefter skulle I see Menneskens
Søn sidde hos Kraftens høire Haand, og komme i Him\-%
melens Skyer“ (Matth. XXVI, 63---64). Der behøves
ikke Tro, men blot sund Menneskeforstand for at indsee,
hvor urimeligt det vilde være at underkaste Christi Ord
en paa meteorologiske Undersøgelser grundet Kritik.
Enhver forstaaer, at det Sprog Christus taler, er absolut
ueensartet med det Sprog, Naturlæren taler.

„Men“, indvender man --- og herved komme vi til
Spørgsmaalet om Virkeliggjørelsen --- „Christi Tale lyder
jo dog paa en virkelig forestaaende Naturbegivenhed;
det er jo netop Knuden, at dersom hans Forudsigelse
virkelig skal gaae i Opfyldelse, saa maa den komme i
afgjørende Strid med de Naturlove, hvis Uforanderlighed
Videnskaben forudsætter.“ Vi svare, at dersom en saa\-%
dan Strid virkelig skulde finde Sted, da maatte den
Virkelighed, hvorom Christus taler, være en simpel ud\-%
vortes physisk Virkelighed; men det Aabenbarelsesphæ\-%
nomen, han beskriver, er et Aandsphænomen; dets Virke\-%
% --- p. 207 ---
\opage{207}\setcounter{footnote}{0}%
lighed en Virkelighed for Aanden. „Denne Virkelighed
for Aanden har altsaa Intet med det Physiske at gjøre?“
Jo, ganske vist har den med den materielle Naturvirke\-%
lighed at gjøre; men Maaden, hvorpaa Virkeliggjørelsen
i physisk Forstand skal foregaae, beskrives ikke physisk.
Spaadommens Opfyldelse forudsætter et Mirakel. Hvor\-%
meget der af denne naturlige Verden skal omdannes og
forvandles ved et saadant Mirakel, til hvilket Tidspunkt
Forvandlingen, timelig seet, vil indtræde, og hvilke Om\-%
dannelses-Mellemtrin, det Synlige først skal gjennemløbe,
derom indeholder Talen Intet.

Stille vi os derimod paa Naturvidenskabens Stand\-%
punkt og fastholde Nødvendighedsprincipet, saa faaer det
virkelige Universum jo rigtignok et andet Udseende,
men saa bliver ogsaa Sandhedserkjendelsen af en anden
Art. Det, vi ad denne Vei skulle søge, er ikke en Fyl\-%
destgjørelse af Personlighedens ideale Fordringer, men
Tingenes reale Sammenhæng: her er Sandheden uperson\-%
lig, her stilles ingen Mirakler i Udsigt. Ja saavist som
den Sandhed, der udtaler sig i den pythagoræiske Lære\-%
sætning er upersonlig, ligesaa vist er det hele System af
Sandheder, der udtaler sig i en Laplaces \textit{Mécanique céléste}
% "céléste": so printed, with an acute on the second e as well (image, 300 ppi);
% transcribed as printed.
et System af upersonlige Sandheder. De personlige og
de upersonlige Sandheder ere ikke blot noget forskjellige,
men absolut ueensartede. Dualismen er her; den staaer
ved Magt. De, der absolut ville forkaste det aabenbarede
Ord og offre deres Personlighed til det graa Absolute,
kunne vel hæve den, men tilsløre den kunne de ikke.

At tilsløre den ved Tro og Viden bestemte Dualisme
og bedaare de Halvdannede ved dogmatiske Talemaader
og metaphysiske Tvetydigheder, er at hemme og forqvakle
det aandelige Liv i Nutiden.

\begin{center}\rule{3cm}{0.4pt}\end{center}

\chapter*{XIII.}
\addcontentsline{toc}{chapter}{Forelæsning XIII}
\markboth{Forelæsning XIII}{Forelæsning XIII}
\label{ch:xiii}
% pp. 208--224
% --- p. 208 ---
\opage{208}\setcounter{footnote}{0}%
% [large initial T]
Tildragelserne i vort personlige Liv, hvad der under
vexlende Omgivelser møder os paa vor Vei, pleie vi i
Forhold til de forskjellige Indtryk, vi modtage deraf, at
betegne med forskjellige Navne. Vi tale saaledes om
Hændelse, Tilfælde, Lykke, Ulykke, Skjæbnens Slag og
Forsynets Prøvelser; men er dette nu ogsaa i Sagen selv
forskjellige Ting? Er der i Virkeligheden nogle Til\-%
dragelser, der ligefrem maae tilskrives Skjæbnen, og ved
Siden af disse igjen andre, der umiddelbart beskikkes os
af Forsynet? Have Skjæbne og Forsyn deelt Herre\-%
dømmet imellem sig? Dette Spørgsmaal er saa person\-%
ligt som muligt.

Det Personlige vil forstaaes personligt; det vil for\-%
tolkes og anvendes i Overeensstemmelse med den Paa\-%
gjældendes eiendommelige Synsmaade: kun ved For\-%
staaelsen og Anvendelsen faaer en Begivenhed sin Be\-%
tydning for Livet. Altsaa Eet af To: enten er Alt under
Skjæbnen, og der er intet Forsyn, eller alle Ting uden
Undtagelse ere under Forsynet, og der er ingen Skjæbne.
Hvilken af de to Synsmaader der i sig selv er den rette
og sande, kan ikke afgjøres i objectiv Viden; det kan
alene bestemmes ved et personligt Valg. Den, der hol\-%
der sig til Nødvendigheden, vælger Skjæbnen; den, der
% --- p. 209 ---
\opage{209}\setcounter{footnote}{0}%
med besluttet Sind griber Friheden, vælger Forsynet.
Her staaer igjen Standpunkt mod Standpunkt, Fortolk\-%
ning mod Fortolkning; vi skulle nu see, hvorledes det
forholder sig med denne Modsætning.

At stille sig paa Naturnødvendighedens Standpunkt og
vælge Skjæbnen er med andre Ord at bygge sin Livs\-%
anskuelse paa objectiv Viden og anvende den upersonlige
Sandhed, som var den personlig. Hvad dette betyder, see
vi paa Mange af Oldtidens Vise, særlig paa Stoikerne.
For den ædle Stoiker er Naturnødvendigheden Eet med
Fornuftnødvendigheden; stolende paa Fornuftnødvendighe\-%
den har han med fast Villie og fuld Bevidsthed stillet sig
under Skjæbnen. Hans Valgsprog: „Før mig, o Skjæbne“,
vil saaledes sige: Led mig, du evige Fornuftnødvendighed,
i dig er Viisdom, i dig er Sandhed! altsaa: Før mig, o
evige Sandhed! Hvorledes føres han nu af Skjæbnen,
og hvorledes forstaaer han dens Førelse? Først gjør han
Adskillelse mellem det Væsentlige og det Ligegyldige.
Til denne Adskillelse mellem Tilværelsens Anliggender
svarer i hans egen Bevidsthed en skarp og klar personlig
Adskillelse mellem „det, der beroer paa os“, og „det,
der ikke beroer paa os“. Hvad er det nu, der ikke be\-%
roer paa os? Om vi skulle være syge eller sunde, be\-%
roer ikke paa os; det staaer i Tilfældets Magt og hører
til det Ligegyldige. Om vi i Samfundet skulle indtage
et lavt eller et høit Trin, om vi skulle høre til de Svage
eller de Stærke, de Ringe eller de Mægtige, beroer ikke
paa os; det afhænger af Tilfældet og hører til det Lige\-%
gyldige. Saaledes vedbliver han at gjennemgaae Livets
endelige Vilkaar og udjævner Forskjellighederne i det
Ene: „Det er ligegyldigt“.

At alt dette Ligegyldige dog faaer en Slags Gyl\-%
dighed, ligger ene og alene i den Betydning, det har for
Charakteren. Derfor siger Epiktet: „Verden er som en
% --- p. 210 ---
\opage{210}\setcounter{footnote}{0}%
stor Skueplads; Menneskene ere de forskjellige Skue\-%
spillere, og Guddommen tildeler Enhver hans bestemte
Rolle“. Herved vil han ingenlunde have Livet opfattet,
hverken som et blindt Spil af Kræfter, eller som et vil\-%
kaarligt Intriguespil. Naar Epiktet taler om at spille sin
Rolle godt, tager han Ordet Rolle i en anden og dybere
Betydning end en Keiser August, der paa sit Dødsleie
siger med et ironiserende Smiil til sine Fortrolige: „Ikke
sandt, jeg har spillet min Rolle godt?“ Epiktet vil ind\-%
skærpe den Sætning, at ethvert Menneske naaer sin Be\-%
stemmelse, naar han med sine Kræfter og Anlæg ud\-%
fylder den Plads, hvorpaa han efter Tingenes Orden blev
sat. Derfor bøier den Vise sig under Nødvendigheden
som under en hellig Magt. Naar du sidder i Theatret
og seer et Sørgespil, maa dit Hjerte ikke banke uroligt,
og du maa ikke græde; thi du skal ikke ønske, at det
gik Helten anderledes, end det gaaer. Det gaaer, som
det maa gaae, og dermed skal du være tilfreds. Om
du elsker din Hustru og Søn nok saa høit, saa indprent
dig bestandig, at et kostbart Leerkar kan sønderbrydes;
hvis det da sønderbrydes, vil dit Uheld ei overraske
og smerte dig. Har du itide alvorligt foreholdt dig selv,
at din Hustru og Søn ere dødelige, saa vil du kunne
see dem døe med den Ro, der er Skjæbnen og Fornuften
værdig.

Der er i den stoiske Resignation noget Storartet og
Ophøiet. En Stoiker resignerer efter en ganske anden
Maalestok, end der ellers i Hverdagslivet resigneres.
Enhver Fornuftig, hvem et Uheld rammer, maa jo vist\-%
nok med en vis Sindsro see at skikke sig i det Uund\-%
gaaelige; men man giver i endelig Resignation kun Af\-%
kald paa nogle Goder i stille Forventning om at faae
dem erstattede ved andre. Stoikerens Resignation der\-%
imod er uendelig; han resignerer een Gang for alle paa
% --- p. 211 ---
\opage{211}\setcounter{footnote}{0}%
alle endelige Goder og opgiver alle timelige Ønsker, at
han ikke skal blive en Bold for Lykkens Luner. Man
skulde troe, at dette var den mest ophøiede Ethik, som
gaves i Verden; og visselig er det ogsaa ophøiede Cha\-%
rakterer, vi møde blandt Oldtidens Stoikere. Dog er det
ved nærmere Betragtning let at indsee, at denne sæde\-%
lige Storhed har noget Opskruet, noget Umenneskeligt
ved sig, at der i dens inderste Væsen har skjult sig en
Modsigelse.

Stoikeren er kold og stærk; thi han er Skjæbnens
Mand, og Skjæbnen er kold og stærk. Han trænger til
Ingen, elsker Ingen, hader Ingen, beundrer Ingen, for\-%
agter Ingen; han haaber paa Ingen og frygter Ingen. Han
kan ikke bøies, han kan kun knuses; paa ham passe de
stærke Ord: \textit{Si fractus illabatur orbis, impavidum ferient
ruinæ}. Stoikeren lægger udelukkende au paa Charak\-%
% printer's defect: "au" printed (turned n, u-form positively visible at 600 dpi) for "an" (lægger an paa); transcribed as printed
teren. Men hvad er Charakteren Andet end udpræget
Personlighed? Og hvad er udpræget Personlighed uden
et personligt Indre? Men Principet for den stoiske Per\-%
sonlighed er upersonligt: her er Modsigelsen. Det In\-%
derste i Stoikerens Væsen er ikke et Selv, men det
kolde, selvløse Almene.

Stoicismen afgiver et mærkeligt Exempel paa, hvad
der kommer ud af at ville hævde Selvet paa et selvløst
Grundlag. I menneskelig Forstand gjør Stoikeren de
største Anstrengelser for at holde fast paa sig selv,
stemme overeens med sig selv, være sig selv nok; men
see vi nærmere til, saa beroer denne Selvbevarelse uaf\-%
brudt paa en væsentlig Opgaaen i det Selvløse. Det
stoiske Selv er uden Selvhedskjærne, er kun et Gjennem\-%
gangspunkt for Verdensfornuften, er imod Skjæbnens
Magt kun Skin og Skygge. Dets Høieste er at ville,
ikke hvad et personligt Selv fornuftigviis maa ville:
hævde sin oprindelige Gyldighed --- nei, hvad Skjæbnen
% --- p. 212 ---
\opage{212}\setcounter{footnote}{0}%
vil, gjøre sig selv til et Intet. Dette bliver især kjende\-%
ligt, hvor Conflicter indtræde, der bringe Stoikeren til
at søge Befrielse i Døden. Selvmord pleier man ellers
at ansee for en usædelig Handling, ja for en stor For\-%
brydelse. Men fra det stoiske Synspunkt kan et ved
Omstændighedernes Sammenstød motiveret, med roligt
Overlæg udført Selvmord betragtes som Dydens Fuld\-%
endelse. Cato i Utica gav sig selv Døden --- „efter at
have læst Platos Phædon og sovet rolig indtil Midnat“.
Plutarch skildrer ham som et Dydsmønster. Det er
Skjæbnelærens kolde Conseqvens, at Selvet, efter at have
øvet alle sine Kræfter og spillet sin Rolle paa Livets
brogede Skueplads, ender Rollen med at blæse Lampen
ud og forsvinde i det Selvløse.

Den „moderne Videnskab“, anstrenge sig nu, saa\-%
meget den vil, med at opsminke og forskjønne den gamle
Skjæbnelære ved at anprise „Ideen“ og „det Absolute“
og Menneskets „absolute Forhold til det Absolute“: den
kan umulig komme bort fra den stoiske Selvmodsigelse.
At et personligt Selv skal have sin Grund, sit Væsen, i
det Selvløse, er et fortvivlet Punkt, der ikke taaler
Ethikens Lys, en metaphysisk Tvetydighed, der gjør
Bevidsthedens Inderste mythisk.

Den Selvmodsigelsens Dialektik, som hertil svarer,
har S. Kierkegaard træffende paaviist i sit Skrift:
\emph{Sygdommen til Døden}, hvori han blandt Andet
(S. 65 flgd.) omhandler „den Fortvivlelse, fortvivlet at
ville være sig selv,“ og (S. 67) charakteriserer „det for\-%
tvivlede Selv“ paa følgende Maade: „Det erkjender
ingen Magt over sig, derfor mangler det i sidste Grund
Alvor, og kan kun fremkogle et Skin af Alvor, naar
det selv skjænker sine Experimenter sin allerhøieste Op\-%
mærksomhed. Dette er nu en tilløiet Alvor; som den
Ild, Prometheus stjal fra Guderne --- saaledes er dette
% --- p. 213 ---
\opage{213}\setcounter{footnote}{0}%
at stjæle fra Gud den Tanke, som er Alvoren, at Gud
seer paa En, istedetfor hvilket det fortvivlede Selv
nøies med at see paa sig selv, hvilket nu skal skjænke
hans Foretagender den uendelige Interesse og Betydning,
medens dog just dette gjør dem til Experimenter. . . .
Der er i den hele Dialektik, indenfor hvilken det handler,
intet Fast; hvad Selvet er, staaer i intet Øieblik fast,
det er, evig fast. . . . Selvet er sin egen Herre, absolut,
som det hedder, sin egen Herre, og just dette er For\-%
tvivlelsen, men ogsaa hvad det anseer for sin Lyst, sin
Nydelse. Dog forvisser man sig ved et nærmere Eftersyn
let om, at denne absolute Hersker er en Konge uden
Land, han regjerer egentligen over Intet; hans Tilstand,
hans Herredom underligger den Dialektik, at i ethvert
Øieblik Oprøret er Legitimitet. Dette beroer nemlig til
syvende og sidst vilkaarligt paa Selvet selv“.

Fra Religionens Standpunkt ere alle Tildragelser,
alle virkelige Begivenheder, alle mødende Tilfælde, for\-%
saavidt de kunne faae væsentlig, afgjørende Betydning
for Menneskets personlige Liv, uden Undtagelse at op\-%
fatte og bestemme som Tilskikkelser af det styrende
Forsyn. Denne Erkjendelse beroer ikke paa Indsigt i
Begivenhedernes naturnødvendige Sammenhæng; tvert\-%
imod! den beroer paa Indsigt i Religionens Væsen; den
fremgaaer med indre Conseqvens af Troens Princip.
Kun ved at gjennemtænke dette Princip vil det være
muligt at naae til indtrængende Forstaaelse af Aabenba\-%
ringens Tegn i det gamle og det nye Testamente. Thi
hvad er vel et Aabenbaringstegn efter sit religiøse
Begreb? Det er en Kjendsgjerning, en Begivenhed, som
ved at være udtraadt af den umiddelbare Natursammen\-%
hæng udtrykkelig betegner et personligt Mellemværende
mellem den sig aabenbarende Gud og den i Troen Ud\-%
valgte, der faaer Aabenbaringen, altsaa en Forsynets Til\-%
% --- p. 214 ---
\opage{214}\setcounter{footnote}{0}%
skikkelse. Saaledes er Phænomenet med den brændende
Tornebusk et personligt Mellemværende mellem Jehova
og Moses, en Tilskikkelse for Moses. At Moses i al
Almindelighed har troet paa Fædrenes Gud og hans
styrende Forsyn, førend Aabenbaringen skete, maa ind\-%
rømmes, ligesom det ogsaa maa forudsættes, at han i
Ørkenens Eensomhed har grundet paa sit Folks Be\-%
frielse. Men der er, ogsaa i Troeslivet, et stort Skridt
fra det Almene til det Enkelte, fra Ønske til Opfyldelse,
fra Tanke til Handling. Paa dette Overgangstrin ind\-%
træder Underet som et udtrykkeligt Tegn. Nu er Timen
kommen, da Befrielsens store Værk skal begynde; nu
skal Moses erkjende, at han og ingen Anden er den
Mand, Forsynet har kaaret til sit Redskab. „Og nu
gaa“, saa lyder Jehovas Røst, (2 Moseb. III, 10), „og
jeg vil sende dig til Pharao, og udfør mit Folk, Israels
Børn, af Ægypten“.

Men saa bestemt Kaldelsen end lyder, saa utvetydig
Tilskikkelsen synes at være, saa skal den dog opfattes
i Tro og tilegnes i Tro. Man skulde ifølge en udvortes
Betragtning antage, at Moses øieblikkelig maatte have
indseet Muligheden af, at han, trods sin menneskelige
Skrøbelighed, vilde være istand til at fuldføre, hvad der
blev ham paalagt, da Jehova jo udtrykkelig lover at
staae ham bi. Men det gaaer i Virkeligheden ikke saa
let med at fatte og forstaae de guddommelige Tilskik\-%
kelser. Det er dog ret betegnende, at Moses, skjøndt han
staaer for det almægtige Forsyn og hører Forjættelsen
om guddommelig Bistand i gjentagne Bekræftelser, dog
med Farerne for Øie viser sig forsagt og tvivlraadig.
Kan der være noget mere talende Beviis for, at Underet,
hvor udvortes det end optræder, dog væsentlig er et
Aandsphænomen, en Virkelighed for Aanden? At For\-%
synets Tilskikkelser, selv hvor de paa den mest uven\-%
% --- p. 215 ---
\opage{215}\setcounter{footnote}{0}%
tede Maade indtræde som en Hjælp i Nødens Stund, end\-%
ogsaa da, naar de vise sig som Almagtens Tegn, maae
opfattes og fortolkes i Tro, derpaa afgiver den hellige
Historie mange ret slaaende Exempler.

I anden Mosebog (XVII.) fortælles Følgende: „Og
der Folket tørstede efter Vand, da kivedes Folket med
Moses og sagde: hvorfor har du ført os op fra Ægypten,
at lade mig og mine Børn og mit Qvæg døe af Tørst?
Da raabte Moses til Jehova og sagde: hvad skal jeg
gjøre ved dette Folk, der mangler kun lidt, da stene de
mig. Men Jehova sagde til Moses: gak frem foran Folket
og tag med dig af de Ældste i Israel, og tag din Stav,
med hvilken du har slaget Floden, i din Haand og gaa.
See, jeg skal staae for dit Aasyn der paa Klippen, paa
Horeb, og du skal slaae paa Klippen, da skal der ud\-%
flyde Vand deraf, og Folket skal drikke. Og Moses
gjorde saa for Israels Ældstes Øine.“

Vil Nogen heri see et Beviis for, at et Mirakels
Virkelighed kan blive kundbart ogsaa for Ikketroende:
da er han i en stor Vildfarelse. Fortællingen beviser jo
lige det Modsatte. Den store Hob, der fortvivlede i Nø\-%
den, den tog imod Hjælpen, den drak af Kilden, (Jevnf.
1 Cor. X.), den slukkede Tørsten, den studsede maaskee
et Øieblik over en saa uventet Lykke, men her staaer
Intet om, at den fik noget Indtryk af Miraklet som Mi\-%
rakel.

Et andet, endnu mere talende Vidnesbyrd have vi
i det nye Testamente (Joh. VI.). Efter at Jesus har
bespiist 5000 Mand i Omegnen af Tiberias, drager han
til Capernaum; „(men der kom andre Skibe fra Tiberias,
nær til det Sted, hvor de aade Brødet, efter at Herren
havde gjort Taksigelse;) og da Folket nu saae, at Jesus
var der ikke, ei heller hans Disciple, traadte ogsaa de i
Skibene, og kom til Capernaum, for at søge Jesum. Og
% --- p. 216 ---
\opage{216}\setcounter{footnote}{0}%
der de fandt ham paa hiin Side Søen, sagde de til ham:
Rabbi, naar er du kommen hid? Jesus svarede dem og
sagde: sandelig, sandelig siger jeg Eder: I søge mig,
ikke fordi I saae Tegn, men fordi I aade af Brødene og
bleve mætte.“ Det er altsaa disse Mennesker ganske
ligegyldigt, hvorfra Brødet kommer; om det skeer ved
et Lykketræf, en eventyrlig Hændelse, ved Trolddom
eller ved en særegen Forsynets Foranstaltning, blot de
kunne faae Brød.

Vi lære af disse Exempler ikke alene, at Tegnet som
Tegn kun kan erkjendes af Troende, men vi lære tillige,
at de Troende, som i Tegnet see en Forsynets Tilskikkelse,
maae ville opfatte Tilskikkelsen saaledes, at den kan faae
sin Betydning for Livet: Tilegnelsen og Forstaaelsen med\-%
fører et personligt Ansvar. Dette Forhold mellem Miraklet
og Tilskikkelsen kan igjen kaste Lys over Abrahams Op\-%
træden hos Pharao og hos Abimelech.

Historien om Abraham maa, som de fleste Beret\-%
ninger i det gamle og det nye Testamente, \emph{fortolkes};
det vil sige: der maa tilføies Mellembestemmelser, Sæt\-%
ninger og Domme, som enten umiddelbart eller middel\-%
bart høre med til at bestemme det Synspunkt, hvorfra
Beretningen vil opfattes. Historien om Abraham er en
folkelig Historie; den meddeles i store Træk; Mellem\-%
bestemmelserne træde i Baggrunden: Folkets Blik er fæ\-%
stet paa Stammemoderens Skjønhed, Patriarchens Tro
og Jehovas Naade; i den folkelige Fremstilling er det
Religiøse indeholdt. Som Jehovas Udvalgte er Abraham
i ganske særegen Forstand en Forsynets Mand, hans Vei
gjennem Verden en Forsynets Vei. Men just som For\-%
synets Mand er Abraham ansvarlig for, hvorledes
han tyder og bruger Tilskikkelserne. Saa kommer
han som en værgeløs Fremmed til Ægypten, til Pharaos
Hof, og seer Faren for sig og sin Hustru. Var han nu
% --- p. 217 ---
\opage{217}\setcounter{footnote}{0}%
under de truende Omgivelser uden videre berettiget til
at stole paa overnaturlig Bistand; turde han være ganske
vis paa at forstaae Forsynets Tanke, naar han sagde til
sig selv: „Jeg har faaet Udvælgelsen; det er umuligt at
Ægypterne kunne slaae mig ihjel; før skal den Almæg\-%
tige ved et Mirakel slaae dem alle ihjel“, saa havde han
jo rigtignok let ved at spille „Riddersmand“. Men at
forstaae Tilskikkelserne er en personlig Opgave. Sagen
stillede sig for ham i et andet Lys; den Tanke, hvori
han bliver enig med Sara, kan udtrykkes saaledes: „Vi
maae give Tid, afværge den øieblikkelige Fare og lade
Begivenhederne udvikle sig; hvo veed, hvad Jehova vil
gjøre? Er dette nu med Hensyn paa Forsynslæren cor\-%
% printer's defect: the quotation opened at „Vi maae give Tid is never closed (no “ after "gjøre?", checked at 300 dpi); transcribed as printed, quote balance +1
rect, saa behøve vi kun hvad den kritiserede „Nødløgn“
angaaer, at mærke os Abrahams Ord til Abimelech
(1 Moseb. XX, 12): „Hun er dog i Sandhed min Sø\-%
ster, min Faders Datter, men ikke min Moders, og hun
er bleven mig til Hustru;“ thi det vil jo sige med andre
Ord: „jeg har fortiet dig Noget, men hvad jeg har sagt,
er sandt.“

Hvad der gjælder om de vidunderlige Tegn og Til\-%
skikkelser i den hellige Historie, det gjælder paa tilsva\-%
rende Maade om Begivenhederne i ethvert Menneskes
personlige Liv. Troen paa Forsynet fordrer en Bevidst\-%
hedsfordobling. Her forudsættes en virkelig Verden, en
Tingenes Orden, hvori Tildragelserne forholde sig som
Led i en Kjæde; dette er den ene Side. Men her for\-%
udsættes tillige, at hele den virkelige Tingenes Orden
bæres, bevæges og gjennemvirkes af Guds almægtige
Villie; hver enkelt Tildragelse viser os Modsætningen
imellem Tro og Viden paa det Allerklareste. Tildra\-%
gelsen være f. Ex. et sørgeligt Dødsfald; en Helt døer i
det Øieblik, han er i Begreb med at befrie sit Folk;
en Familiefader døer og efterlader Hustru og Børn
% --- p. 218 ---
\opage{218}\setcounter{footnote}{0}%
hjælpeløse, o. s. v. Betragtes disse Dødstilfælde fra
Videns Synspunkt, da er der vel Ingen, der betvivler,
at ethvert af dem jo maa have havt sine naturlige
Aarsager. Men blive vi staaende ved denne Opfat\-%
telse, saa kan her ikke være Tale om noget Forsyn;
hver Begivenhed er, for at minde om Stoikerne, paa
sin Viis indvævet i den store Nødvendighed: Alt falder
ind under Skjæbnen. Skal her vise sig et Forsyn, maa
Tildragelsen for sig opfattes som en Begivenhed, hvori den
guddommelige Villie udtrykkelig har indlagt en for de
Paagjældende personlig Opgave.

Det er nu tillige let at see, at den til Troen svarende
Opfattelse ingenlunde udelukker den naturlige. De na\-%
turlige Aarsager til Heltens Død kunne være Hvermand
i Landet bekjendte, og endda kan Folket i Alvor op\-%
fordres til ved den sørgelige Begivenhed at mærke sig For\-%
synets Villie. Det maa udtrykkelig fremhæves, at naar
Troen forudsætter en Eenhed af Naturaarsagerne og Guds
Villie, saa forudsætter den tillige, at den altstyrende Villie
kan i hvert Øieblik ændre de samvirkende Aarsager og
lade Begivenhederne indtræffe netop saaledes, som
det er Mennesket tjenligst. Men paa en saadan Syns\-%
maade kan Videnskaben umulig indlade sig. Naar en
Læge opfatter et Dødsfald som Læge, og en Præst det
% a risen space/quad prints as a faint mark over the "e" of "en Præst"; not typed
samme Dødsfald som Sjælesørger, saa ere de to Opfat\-%
telser absolut ueensartede. Dersom nu de to absolut
ueensartede Synsmaader vare uforenelige i een Bevidst\-%
hed, da maatte det jo være Lægen ligesaa umuligt at
forstaae Præsten, som det vilde være Præsten umuligt
at forstaae Lægen. Er dette nu Tilfældet? Ingenlunde!
Den dualistiske Betragtningsmaade er i den Grad sammen\-%
voxet med vore religiøs-sædelige Forestillinger, at det i
et Sørgehuus vilde vække den største Forargelse, om
Lægen erklærede, at det paagjældende Dødsfald ikke
% --- p. 219 ---
\opage{219}\setcounter{footnote}{0}%
vedkom Forsynet, da Patienten beviislig var død en
naturlig Død, medens Præsten i from Iver for at vise
Forsynets Finger, udtrykkelig paastod, at Gud selv bog\-%
stavelig havde slaaet Mennesket ihjel.

Den religiøst Sørgende forener Bevidstheden om na\-%
turligt virkende Aarsager med Bevidstheden om Guds
urandsagelige Villie, forener altsaa to absolut ueensartede
Synsmaader i een Bevidsthed; og saa raaber man endda
Dualisme, Dualisme! og gjør en Larm og en Blæst,
som var det noget uhørt, forfærdeligt og selvmodsigende,
hvad man dog, naar det sees i det virkelige Liv, vil hen\-%
regne til det Allersædvanligste.

Miraklet og Hverdagsbegivenheden ligne hinanden
deri, at de begge kunne opfattes som Tilskikkelser; men
saa umuligt det er uden Tro at erkjende Hverdagsbegiven\-%
heden for en Tilskikkelse, ligesaa umuligt er det uden
Tro at erkjende det Allervidunderligste for et Mirakel.
En i Raahed nedsunken Mængde kunde, ifølge den hel\-%
lige Histories Beretning, vel for et Øieblik bringes til at
studse og forfærdes, naar det Vidunderlige skete; men
den glemte snart Overraskelsen og følte ingen Til\-%
skyndelse til at takke Forsynet. Paa tilsvarende Maade
gaaer det med Tilskikkelserne endnu den Dag idag.
Der kan møde et Menneske Noget, saa glædeligt, saa
sørgeligt, at det er i Stemningens Øieblik, som saae han
% a risen space prints as a faint raised mark after "som"; not typed
Forsynets Øie; men dersom der ingen Inderlighedens
Kraft, ingen Troens Alvor er i Personligheden, hvad
skeer saa? Saa skeer det ikke Usædvanlige, at Stem\-%
ningen svinder, at Sindet taber sin Spændkraft, at det,
der begyndte med at vise sig som en høiere Tilskikkelse,
ender med at forvandle sig til en ganske naturlig Be\-%
givenhed. Forskjellen mellem Nutidens Vidende, der
ikke kunne see Tilskikkelsen i en naturlig Begivenhed,
og hine Uvidende, der ikke kunde see den i Almagtens
% --- p. 220 ---
\opage{220}\setcounter{footnote}{0}%
Tegn, er kun den, at de Vidende forblindes af Tanker,
de Uvidende af Tankeløshed.

Til hvilken Yderlighed, den nyere Theologie, bedaaret
og forhexet af „videnskabelig“ Tænkning, er gaaet i at
undergrave Forsynslæren, kan sees af Schleiermachers Dog\-%
matik. Schleiermacher siger om Opholdelsen (Der christ\-%
liche Glaube. 1ster Band. Berlin 1835. S. 222): „Den
fromme Selvbevidsthed, ifølge hvilken vi sætte Alt, hvad
der vækker os og indvirker paa os, i den ligefremme
Afhængighed af Gud, falder ganske sammen med den
Indsigt, at dette netop Altsammen er betinget og bestemt
ved Natursammenhæng.“ Saa aldeles bliver den guddomme\-%
lige Villie ifølge denne Forklaring identificeret med de vir\-%
kende Naturaarsager, at Skjæbnen selv bliver Forsyn, og
Forsynet Skjæbne. Det er den gamle stoiske Skjæbne\-%
lære, der friskes op ved lidt evangelisk Sminke. Schleier\-%
machers „fromme Selvbevidsthed“ og Epiktets fromme
Resignation bringes her til fuldstændig Congruens. Denne
Congruens har længe været forberedt: Holberg saae den
allerede.

„Intet,“ siger Holberg, „kan bevæge mere, end den
Resignation til Guds Villie, som gjøres af Epicteto.
Hans Ord ere disse: Jeg ønsker af mit ganske Hjerte,
at Døden paa mit Yderste finder mig bered, paa
det jeg uden Forvirrelse, uden Hinder og Tvang, maa
ende mit Liv, og at jeg maa sige til Gud: Herre! haver
jeg overtraadt dine Bud? haver jeg misbrugt de Gaver,
som du haver forlehnet mig med? haver jeg ikke under\-%
kastet min Villie din Villie? haver jeg nogentid klaget
over min onde Skjæbne? haver jeg nogentid beskyldet
dit guddommelige Forsyn? jeg haver været syg, efterdi
det haver været din Villie, ligesom det og haver været
min Villie; jeg haver været fattig, efterdi det saaledes
har behaget dig, og jeg haver været fornøiet i min
% --- p. 221 ---
\opage{221}\setcounter{footnote}{0}%
Fattigdom; jeg haver levet i Foragt, efterdi du haver
villet saa have det, og jeg haver ikke stræbet at reise
mig i Veiret. Du haver aldrig seet mig bedrøvet udi
min slette Tilstand. Du haver aldrig fundet mig utaal\-%
modig og knurrende. Jeg er endnu færdig til at under\-%
kaste mig alt, hvad som du finder for godt at paalægge
mig. Det ringeste Tegn, som du giver, er for mig en
hellig og ubrødelig Ordre. Du vil, at jeg skal forlade
denne glimrende Verdens Skueplads: see! jeg forlader
den, og yder dig tusinde Tak, at du haver værdiget at
sætte mig derpaa, for at vise mig dine Gjerninger, og at
stille for mine Øien den fortræffelige Orden, hvormed
du regjerer alle skabte Ting.“ (Hundrede og tyve af
Holbergs Epistler ved Fabricius. Kjøbenhavn 1858.
S. 114).

Ved et første flygtigt Øiekast maa det vel synes,
som om der i Epiktets Udtalelse indeholdtes en dybt til\-%
fredsstillende Forsynslære; men det Modsatte er netop
Tilfældet. Det er det selvløse Væsen, den almene For\-%
nuftnødvendighed, som i den fromme Betragtning her
bliver personificeret; der tillægges Skjæbnen en Viisdom,
et Selv, en Villie, som den ifølge Sagens Natur hverken
har eller kan have. Derved, at det Upersonlige i Ver\-%
densstyrelsen paa denne Maade antager Skin af Person\-%
lighed, skeer det just, at den fromme Resignation, langt
fra at nære og styrke det personlige Selv i Mennesket,
fortærer Personlighedens Kjærne. Den Fromme vil, hvad
den personificerede Skjæbne vil; ved at gjøre sin Villie
til Eet med Skjæbnevillien offrer han Personlighedens
ædleste Kræfter paa det Upersonlige.

Den fromme Skjæbnelære og den evangeliske For\-%
synslære udgaae fra absolut ueensartede Principer. Lige\-%
som Aabenbaringen i det Hele charakteriserer sig ved,
at den underordner Naturbegrebet under Skabelsesbe\-%
% --- p. 222 ---
\opage{222}\setcounter{footnote}{0}%
grebet, saaledes charakteriserer den aabenbarede Forsyns\-%
lære sig særligt derved, at den underordner alle virkende
Naturaarsager under Guds styrende Villie. Istedenfor
at identificere Guds Villie med Naturaarsagerne, forud\-%
sætter Troen tvertimod Forskjellen. (Observandum quod
Deus non solum vim agendi dat causis secundis . . ., sed
quod immediate influit in actionem et effectum creaturæ.
Qvst.): Frihedsaarsagen bestemmer selv de Begivenheder,
der, efter Phænomenet at dømme, umiddelbart fremgaae
af naturlige Aarsager.

Den, der ikke kan opfatte Miraklet som et Almag\-%
tens Tegn, kan heller ikke opfatte nogensomhelst Begiven\-%
hed som en Forsynets Tilskikkelse. Den Eenhed af
Forsynsvillie og Naturaarsager, der er betegnende for
Hverdagsbegivenheden, faaer i Troen sit Lys fra den
Modsætning mellem Forsynsvillien og Naturaarsagerne, der
anskueliggjøres i Miraklet; Eenheden maa belyses ved
Modsætningen, og Modsætningen ved Eenheden: Miraklet
og Tilskikkelsen forudsætte og forklare hinanden gjen\-%
sidig.

Den sande Forsynslære er en ubetinget Personlig\-%
hedslære; Forbilledet, den sande Personlighed, er selv
det aabenbarede Forsyn. Hvad har det aabenbarede For\-%
syn da lært os om Tilskikkelser og Prøvelser, om Mod
og Taalmod, om Frygt og Fare? Vi høre Hans Ord, som
siger (Mtth. X, 28---30): \emph{Sælges ikke to Spurve
for een Penning? og ikke een af dem falder paa
Jorden uden Eders Faders Villie. Ja endog
alle Eders Hovedhaar ere talte. Frygter derfor
ikke; I ere bedre end mange Spurve}. --- „Dersom
disse Ord kun vare Menneskeord, da maatte vi vel be\-%
tragte dem som en Overdrivelse i et Begeistringens Ud\-%
brud; thi i den Verden af Farer, der ligger imellem
Mennesket og Gud, har et Menneske jo meget at frygte.
% --- p. 223 ---
\opage{223}\setcounter{footnote}{0}%
Men naar han, der taler Ordet, er Ordet selv, da maae
vi lære af ham, hvor Faren er, og hvad vi have at frygte.
Der er i Farernes Verden kun een ubetinget Fare, kun
Eet, der er frygteligt, kun Een, der er at frygte. Der\-%
for siger Ordet: „Jeg vil vise Eder, for hvem I skulle
frygte.“ Og Forsynet er menneskeligt. Han, der taler
om Faren og viser os, for hvem vi skulle frygte; han,
der overskuer alle Endelighedens Farer og siger: „frygter
ikke!“, han veed, hvad Fare er; hans Liv var fortroligt
med Faren. Faren søgte ham i hans Vugge; den fulgte
ham til hans Grav. Overalt, i Ørkenen som i Stæderne,
paa Galilæas Sletter, mellem Judæas Bjerge, overalt,
hvor han gik hen, gik Faren efter ham og lurede paa
ham. Aarvaagen som Forsynet, frimodig som den, der
er under Forsynet, gaaer Menneskens Søn ad Farernes
Vei og frygter for Intet i Verden. Han arbeider,
og Ingen standser ham; han gaaer til sit Maal, og Ingen
lægger Haand paa ham, før hans Time kommer. Men
da hans Time var kommen, gik han frivillig Faren imøde;
den overraskede ham ikke som et Tilfælde, den kom
ikke over ham som en Skjæbne; den kom, fordi han lod
den komme. Den kom da heller ikke i udvortes Larm,
ikke i en Storm af menneskelige Lidenskaber, ikke i et
Bulder af fjendtlige Vaaben; den kom i al Stilhed een\-%
som til den Eensomme, kun saaledes lod han den komme.
Ene hiin Nat i den stille Have lod han Faren komme
til sig og stred med den, ikke som den Stærke strider,
men som den Skrøbelige, at der i Kjød og Blod ikke
skulde være en Skjæbnens Forskjel mellem Stærke og
Skrøbelige, men en Hengivenhedens Seier i Skrøbelighed.
Ene hiin Nat i den stille Have, lod han Faren gaae som
en Dødsangst igjennem sin Sjæl; han trodsede den ikke,
han foragtede den ikke; han hævede ikke stolt sit Hoved
som den, der offrer sig til Skjæbnen; han faldt paa sit
% --- p. 224 ---
\opage{224}\setcounter{footnote}{0}%
Ansigt og bøiede sig ydmyg under Faderens Villie. Ene
hiin Nat i den stille Have, lod han alle Endelighedens
Farer forvandle sig i Inderlighed til een Fare. Saa stred
han i Inderlighed og seirede i Inderlighed; og da han
saa reiste sig fra Striden, see, da var det forbi med
Faren, da var der slet ingen Fare.“ (Om Tilværelsens
ideale Magter: Skjæbne og Forsyn. Andet Oplag. Kbhvn.
1867. S. 30. flg.).

Forsynslæren er dualistisk. Ophæves Dualismen,
saa forsvinder Forsynet. Men naar Forsynet forsvinder
for Folkets Ledere og for Folket selv: hvordan vil det
da gaae med det aandelige Liv i Nutid og Fremtid?

\begin{center}\rule{3cm}{0.4pt}\end{center}

\chapter*{XIV.}
\addcontentsline{toc}{chapter}{Forelæsning XIV}
\markboth{Forelæsning XIV}{Forelæsning XIV}
\label{ch:xiv}
% pp. 225--255
% --- p. 225 ---
\opage{225}\setcounter{footnote}{0}%
% [large initial D]
Det var om Hindringer og Betingelser for Nutidens
aandelige Liv, at disse Forelæsninger væsentlig skulde
handle. Men til en Bedømmelse af det aandelige Liv,
dets Phænomener og Kræfter var det først og fornem\-%
melig nødvendigt at vælge sig et Standpunkt, hvorfra det
Hele nogenlunde kunde overskues, og sikkre sig en
Maalestok, hvorunder de indbyrdes forskjellige, hin\-%
anden krydsende Retninger lode sig henføre. Det er Skils\-%
missen mellem Tro og Viden, mellem Religion og Vi\-%
denskab, der betegner vort valgte Standpunkt. Fra for\-%
skjellige Sider er dette Standpunkt i det Foregaaende
efterhaanden blevet belyst; men vi maae tilstaae det, det
er kun en deelviis Belysning! Her kræves en alsidig
Opgjørelse, et sidste afsluttende Regnskab.

Et saadant Regnskab synes vel ved første Øiekast
at maatte blive meget indviklet, efterdi det skal aflægges
ved to Domstole, hver med sin Rettergang, Religionens og
Videnskabens; Værdierne maae ved denne Opgjørelse be\-%
regnes i to Myntsorter: det bliver, hvad man spottende har
sagt, et dobbelt Bogholderi. Ikke destomindre kan Sagen
fra Begyndelsen anlægges meget enfoldigt. Da her nemlig er
to Erkjendelseskredse, hver med sit Princip, saa ligger det
nær, at enhver af Parterne punktlig maa have Sit; men naar
% --- p. 226 ---
\opage{226}\setcounter{footnote}{0}%
Enhver faaer Sit, ophører jo al Strid og Klage. Faaer
Religionen med sine Aabenbaringsidealer, sine hellige
Kjendsgjerninger og personlige Villiestanker Lov til at
udgjøre et Aandens Rige for sig, en Region af Fore\-%
stillinger og Begreber, der med indre Conseqvens ud\-%
vikles i egen Sammenhæng af egne Forudsætninger: hvad
kan Troen da fordre mere? Og nu til den anden Side!
hvilken større Indrømmelse kan man vel gjøre Viden\-%
skaben, end naar man tillægger den et eget Standpunkt,
en egen Methode, egne Love og Conseqvenser og derhos
Frihed til at udvikle sig i alle Retninger, inddrage Alt
for sin Domstol, undersøge og prøve Alt, hvad der paa
nogen Maade kan falde ind under en menneskelig Synskreds.
Man skulde mene, at Videnskaben for sit Vedkommende
maatte finde sig tilfreds, ligesaavel som Religionen. Og
dog, vi tør ikke dølge det, har Regnskabet sine Van\-%
skeligheder. Sagen er, at de to Erkjendelseskredse,
ihvorvel de ere indbyrdes absolut ueensartede, skulle for\-%
enes i een Bevidsthed, og det saaledes, at den til Selv\-%
forstaaelse udviklede Personlighed bliver istand til at be\-%
væge sig med Sikkerhed i Viden og dog finde Hvile i
Troen. Dette kan kun skee derved, at den lærer at
gjennemskue den Grundvildfarelse, den første falske
Forudsætning, hvoraf alle Videns Indvendinger mod Troen
fremgaae, saa klart, at den overalt formaaer at hæve Mis\-%
forstaaelsen uden at krænke Videnskaben. Men under
denne Stræben er Personligheden udsat for Fristelse.
„Fristelsen er, at Tro og Viden saa let gjøres lidt eens\-%
% "lidt": underlined in pencil by a reader; not transcribed
artede, uagtet de ere absolut ueensartede. Betones det
i Ueensartetheden Absolute tilstrækkelig skarpt, saa
høres en skrigende Modsigelse: \emph{hvad der er Sand\-%
hed i Troen, er Usandhed i Videnskaben;
hvad der er Sandhed i Videnskaben er Usand\-%
hed i Troen}. Men i Modsigelsen ligger Muligheden
% --- p. 227 ---
\opage{227}\setcounter{footnote}{0}%
til dens Løsning. Det er vor Videns absolute Grændse,
det er Tilværelsens Gaade, hvorpaa saavel den Vi\-%
dendes som den Troendes Opmærksomhed her maa
fæste sig. Har al menneskelig Viden en absolut
Grændse, er her virkelig et Aandens Mysterium, en
Tilværelsens Gaade, da er her ogsaa Noget, der maa
\emph{troes}, fordi det ikke kan \emph{vides}. Naar nu det i Gaaden
Skjulte, det, der ene og alene er bestemt til at være
Gjenstand for Tro, ikke desmindre drages ind under
Viden og Videnskab, hvad saa? Saa overskrider Viden
sin absolute Grændse. Og naar Viden har overskredet sin
absolute Grændse, hvad saa? Saa giver Videnskaben
sig frisk væk ifærd med at udgrunde, begrunde og be\-%
gribe det, som ifølge sin Natur hverken kan eller skal
videnskabeligt begribes. Men naar Videnskaben giver sig
af med at begribe det Ubegribelige, hvad saa? Saa
kommer den nødvendigviis til modsigende Begreber.
Hvad der kun kan udtrykkes i modsigende Begreber
maa være usandt; Troens Mysterier lade sig kun ud\-%
trykke i modsigende Begreber: altsaa ere de usande.
At Noget kan være sandt i Troen og usandt i Viden\-%
skaben følger nu ligefrem deraf, at \emph{Videnskaben ved
at ville begribe, hvad der ligger udenfor dens
Grændse, selv bliver usand}.“

„Men hvorledes kan Noget, der virkelig er sandt i
Videnskaben, være usandt i Troen? Svar: Fordi det
ikke blot er Viden, men ogsaa Troen, der har en absolut
Grændse: at Viden har Troen til Grændse, forudsætter
nødvendigt, at Troen har Viden til Grændse. Men at
Troen har Viden til Grændse, vil sige, at et Troesbe\-%
greb umuligt kan udvikles til et Vidensbegreb. Idet
Troen støtter sig til Miraklet, forudsætter den en virkelig
Natur med virkelige Love. Men hvorledes forholder
Naturen sig til Miraklet? Ja det er jo netop dette, der
% --- p. 228 ---
\opage{228}\setcounter{footnote}{0}%
er Tilværelsens Gaade. Hvis nu Troen, forglemmende
sin absolute Forudsætning, nemlig Gaaden, overskrider
sin Grændse og begynder at ville omkalfatre det viden\-%
skabelige Naturbegreb til Bedste for Miraklet: da bliver
Naturlærens Naturbegreb ikke blot en Usandhed i Troen,
\emph{men Troen, der har indladt sig paa Noget den
ikke magter, bliver selv usand}.“ (Om den gode
Villie S. 31---33.)

Det Næste, der frembyder Vanskeligheder, er Regn\-%
skabets Form. Denne Vanskelighed angaaer ikke nær\-%
mest det Religiøse, men det Videnskabelige. Det aan\-%
delige Livs Grundtanker kunne virkelig fremstilles i en
human Meddelelsesform. Det vilde her være paa urette
Sted at gjøre Undskyldning for det Mangelfulde, saaledes
at man skjød Skylden paa Erkjendelsens Vanskelighed.
Det vilde være en Selvmodsigelse efter det valgte Stand\-%
punkt; thi det ligger, som allerede viist, i Person\-%
lighedens Natur, at dens afgjørende Livsbetingelser maae
kunne udvikles klart, fatteligt, folkeligt, og forstaaeligt
for Alle.

Anderledes med den videnskabelige Side af Regn\-%
skabet; her er den humane Meddelelsesform utilstrække\-%
lig. Et videnskabeligt Regnskab kan kun aflægges i
streng videnskabelig Form, og ethvert Forsøg paa at
gjøre det tilfulde kjendeligt og fatteligt for Udenforstaa\-%
ende, for den humane Bevidsthed i dannede Kredse, er
et Forsøg paa at omsætte det strengt Videnskabelige i
noget Populært; men hvormeget den populære Form af\-%
viger fra de videnskalige Fordringer, er allerede i det
% "videnskalige": so printed (for "videnskabelige"); printer's error, not corrected
Foregaaende tilstrækkelig udhævet. Til det Bedøvende
i Nutidens aandelige Liv kan man uden Tvivl henregne
den uforholdsmæssige Tillid til videnskabelige Medde\-%
lelser paa anden Haand, til Opfattelse i populære Former.

Jeg maa altsaa gjøre Afkald paa at fremlægge det
% --- p. 229 ---
\opage{229}\setcounter{footnote}{0}%
videnskabelige Regnskab; men for at det ikke skal see
ud, som om jeg under dette Foregivende unddrog mig et
paafølgende Regnskab, vil jeg gaae saa vidt, som det nogen\-%
lunde er muligt, idet jeg i Korthed vil paavise de Syns\-%
punkter og orienterende Udgangspunkter, hvorpaa det ved
det videnskabelige Regnskabs Aflæggelse væsentlig maa
komme an. Jeg angiver disse Synspunkter ved at henvise til
skriftlige Udtalelser, forsaavidt disse foreligge. Sagen
er jo i sin første Vorden; jeg har endnu ikke havt Lei\-%
lighed til, hvad dog er nødvendigt til det Heles fuldstæn\-%
dige Opklaring, at holde et udførligt Foredrag over
Religionsphilosophien.

Naar Talen er om Skilsmisse mellem Tro og Viden,
ligger det nær at adskille Religionen fra Videnskaben
ved at gjøre Religionserkjendelsen til et Paradox. Dette
er, hvad Kierkegaard med vidunderlig Kraft og Geni\-%
alitet har sat igjennem. Idet han resolut brød Staven
over ethvert Forsøg paa ad videnskabelig Vei at ville
retfærdiggjøre Religionens Sandheder, gjorde han Bruddet
saa afgjørende, at enhver Erkjendelse fra Religionens Side
vendte sig mod Videnskaben med en stikkende Braad, idet
Troen blev sat som Paradox. Jeg skal hertil bemærke, at
mit valgte Standpunkt ikke staaer i Strid med, men
heller ikke falder sammen med den Kierkegaardske
Anskuelse. Denne betegner et Stadium, et Standpunkt
for sig; dette Stadium maatte naaes, og Paradoxets
Knude strammes tilgavns, før vi kunde komme videre
i religiøs Erkjendelse. Hvor forskjelligt det valgte
Standpunkt er fra det Kierkegaardske, kan anskue\-%
liggjøres i faa Træk. Modsætningen mellem Tro
og Viden er hos Kierkegaard forvandlet til en Mod\-%
sætning mellem Lidenskab og Videnskab. Videnskaben
og Lidenskaben ere Grundmodsætningens Led; han klager
just over, at Theologerne i vort theocentriske 19de Aar\-%
% --- p. 230 ---
\opage{230}\setcounter{footnote}{0}%
hundrede ere ifærd med at glemme Lidenskaben over
Videnskaben. Han vil da igjen gjøre Lidenskaben
gjældende, særlig den hellige Lidenskab, Troen. Naar
Troen sættes som hellig Lidenskab, seer man strax,
hvor umuligt det er at fremstille den i en Doctrin,
i en objectiv dogmatisk Læreform. Derfor indskærper
Kierkegaard atter og atter, at Religionens Sandheder ikke
kunne doceres; Christendommen, siger han, er ingen
Lære, den er en Existensmeddelelse.

Jeg har for Overblikkets Skyld udgivet et Skrift:
„Paa Kierkegaardske Stadier, et Livsbillede“. Jeg har
ikke udgivet dette Skrift for at fritage Læseren for
at gaae til Kilderne; men den Kierkegaardske Litte\-%
ratur er saa vidtløftig, mangesidig og forgrenet, at man
let gaaer vild i de forskjellige Produktivitetsformer lige\-%
som i en Urskov. Derfor er der Mange, som læse
store Partier af Kierkegaard uden at have et sik\-%
kert Overblik over den hele Produktivitet. Nogle for\-%
tabe sig i „Enten Eller“ og tilegne sig ham som en vittig,
brillant Æsthetiker; Andre springe det Æsthetiske over,
holde sig til „Frygt og Bæven“, til „Kjærlighedens Gjer\-%
ninger“, til „Indøvelse i Christendom“ o. s. v., og tage
ham uden videre for den strenge Bodsprædikant. Det
lille Skrift fremstiller derfor Stadierne, det æsthetiske,
det ethiske og det religiøse, i deres indbyrdes Forhold
--- til Overblik.

„Kjærligheden“, siger Kierkegaard, „har dog i Dig\-%
terne sine Præster, og stundom hører man en Røst, der
veed at holde den i Hævd, men om Troen høres der
intet Ord; hvo taler til denne Lidenskabs Ære? Philo\-%
sophien gaaer videre. Theologien sidder sminket ved Vin\-%
duet og beiler til dens Gunst, falbyder sin Deilighed til
Philosophien.“ Fremstillingen af Troesidealerne bør ingen\-%
lunde være doctrinær; Kierkegaard gjør Alt for her
% --- p. 231 ---
\opage{231}\setcounter{footnote}{0}%
at bryde den videnskabelige Form. Saaledes holder han
i sin „Smule Philosophie“ det positive Christendomsind\-%
hold i Form af et Experiment og fremstiller i den viden\-%
skabelige Efterskrift det Hele saa humoristisk og let
spillende, at man skulde troe, det var kun Spøg, skjøndt
det med denne Spøg er ramme Alvor. At fremstille
Troesidealerne er for Paradoxets Forsvarer snarere en
Digteropgave end et videnskabeligt Anliggende. „Hvis
det er saa, at Du kan fremstille Troen, det beviser blot,
at Du er en Digter, og gjør Du det godt, at Du er
en god Digter, men Intet mindre, end at Du er en Troende.
Maaskee kan Du ogsaa græde, naar Du fremstiller Troen,
det vilde saa bevise, at Du var en god Skuespiller.“
At udvikle Paradoxet theologisk vilde i Kierkegaards
Øine være noget af det mest Bagvendte, man kunde fore\-%
tage sig, thi Paradoxet lader sig ikke bestemme fra et
traditionelt eller objectivt Lærestandpunkt. Hele Frem\-%
stillingsmaaden maa ifølge hans Synsmaade væsentlig
være en Fremstilling af Stemninger. Den Kierkegaardske
Dialektik er en Stemningsdialektik, hans Begreber ere
Stemningsbegreber. Saavidt jeg veed, har Ingen i Nu\-%
tiden kunnet maale sig med Kierkegaard i at udvikle de
religiøse Begreber til den høieste Pathos. Jeg vælger
et Exempel af „Frygt og Bæven“. Den pseudonyme
Forfatter taler om Abrahams Tro, og taler i Stemning.
For at charakterisere Abrahams aandelige Storhed sam\-%
menligner han ham med andre Helte og kommer derved
til at skildre os Charakterer i forskjellig Aand, forskjel\-%
lige Arter af Storhed. „Ingen“, siger han, „skal glemmes,
der var stor i Verden; men Enhver var stor paa sin
Viis, og Enhver i Forhold til dets Storhed, \emph{han elskede}.
Thi den, der elskede sig selv, blev stor ved sig selv,
og den, der elskede andre Mennesker, blev stor ved sin
Hengivenhed, men den, der elskede Gud, blev større
% --- p. 232 ---
\opage{232}\setcounter{footnote}{0}%
end Alle. Enhver skal mindes, men Enhver blev stor i
Forhold til sin \emph{Forventning}. En blev stor ved at
forvente det Mulige, en Anden ved at forvente det Evige;
men den, der forventede det Umulige, blev større end
Alle. Enhver skal mindes, men Enhver var stor alt
i Forhold til dets Størrelse, han stred med. Thi den,
% "han stred med": set in plain roman here, unlike the spaced "han elskede" on p. 231
der stred med Verden, blev stor ved at overvinde
Verden, og den, der \emph{stred} med sig selv, blev stor
% "Verden, og": a faint raised mark after the comma is a risen space/quad; not transcribed
ved at overvinde sig selv; men den, der stred med Gud,
blev større end Alle. Saaledes blev der stridt i Verden,
Mand mod Mand, Een mod Tusinde, men den, der stred
med Gud, var større end Alle. Saaledes blev der stridt
paa Jorden: der var den, der overvandt Alt ved sin
Kraft, og der var den, der overvandt Gud ved sin Af\-%
magt. Der var den, der stolede paa sig selv og vandt
Alt, der var den, der tryg ved sin Styrke offrede Alt,
men den, der troede Gud, var større end Alle. Der var
den, der var stor ved sin Kraft, og den, der var stor ved
sin Viisdom, og den, der var stor ved sit Haab, og den, der
var stor ved sin Kjærlighed; men Abraham var større end
Alle, stor ved den Kraft, hvis Styrke er Afmagt, stor ved
den Viisdom, hvis Hemmelighed er Daarskab, stor ved det
Haab, hvis Form er Vanvid, stor ved den Kjærlighed,
der er Had til sig selv.“

Her haves da Paradoxet, i den høieste Pathos, den
rigeste Stemning. Man seer let, at denne stemningsrige
Tale aldeles ikke tillader noget videnskabeligt Raisonne\-%
ment. Hvor vilde ikke enhver formel, metaphysisk Ind\-%
sigelse falde død og magtesløs til Jorden ligeoverfor
en Udtalelse som denne, at \emph{den, der forventede det
Umulige, blev større end Alle}! At forvente det
Umulige, hvad vil det sige? Naar man skal reflectere,
kan man vel indrømme, at i Troen forventer Mennesket
det Usandsynlige; men hvem kan forvente det i sig
% --- p. 233 ---
\opage{233}\setcounter{footnote}{0}%
Umulige? Og dog er det Umulige sat saa skarpt i det
Usandsynliges Sted. Meget kan synes umuligt for Men\-%
nesker, men netop det, at det er umuligt for Mennesker,
lægger et stiltiende Eftertryk paa, at det er muligt for
Gud. Enhver Akkord mellem det for Mennesket Umulige
og det for Gud Mulige er imidlertid afskaaret ved Para\-%
doxet. Den ligefremme Forstaaelse er afviist, al endelig
Raisonneren maa forstumme; Kraften og Farten ligger i
den ideale Pathos: Existensmeddelelsen giver sig ikke
Tid til theoretiske Forhandlinger.
% "Forhandlinger": a stray mark stands over the o; not an accent, not transcribed

Kierkegaard har træffende seet, at Spørgsmaalet
om Christendommens Sandhed umulig kan afgjøres i
Videnskaben, efterdi Subjectiviteten i Videnskaben altid
maa være betragtende, og den betragtende Subjectivitet
objectiv. „Den beskedne, objective Subjectivitet holder
sig med applauderet Heroisme udenfor; den er til Tje\-%
neste med at ville antage Sandheden, saasnart den er
tilveiebragt. Dog er det et fjernt Maal, hvortil der stræbes
(unegteligt, thi en Approximation kan blive ved, saalænge
det skal være) --- og medens Græsset groer, saa døer
Betragteren, rolig, thi han var objectiv. O, Du ikke for
Intet prisede Objectivitet, Du formaaer Alt; ikke den
mest Troende har været sin evige Salighed saa vis, og
fremfor Alt saa sikker paa, ikke at tabe den, som den
Objective! Det skulde da være, at denne Objectivitet og
Beskedenhed var paa urette Sted, at den var uchristelig.
Saa var det jo rigtignok betænkeligt paa den Maade at
komme ind i Christendommens Sandhed. Christendommen
er Aand, Aand er Inderlighed, Inderlighed er Subjecti\-%
vitet, Subjectivitet i sit Væsentlige Lidenskab, i sit Maxi\-%
mum uendelig personligt interesseret Lidenskab for sin
evige Salighed.“ (Afsluttende uvidenskabelig Efterskrift
til de philosophiske Smuler af Joh. Climacus. Kbhvn.
1846. S. 19---20).

% --- p. 234 ---
\opage{234}\setcounter{footnote}{0}%
Naar Kierkegaard udvikler sine Begreber, maae vi
ikke tage dem i den blot almindelige Form; thi naar
man skiller disse Begreber fra den tilsvarende Stemning
og stærke Pathos og drager dem ind under Betragtningen,
har man forvansket det Hele, og det er just det, Kier\-%
kegaard ivrer saa stærkt imod. Metaphysiske Indsi\-%
gelser, kritiske Bemærkningspillerier have her Intet at
betyde.

At dette Standpunkt ikke falder sammen med det
af os valgte, kan let sees; thi vi have jo netop, i stik
Modsætning til den Kierkegaardske Afviisning, stræbt atter
og atter at drage hele Systemet af religiøse Erkjendelser
ind under Betragtningens Niveau. Vi indrømme, at
Kierkegaard har ubetinget Ret i den Paastand, at Troen
for den Ikke-Troende maa være et Paradox. Vi ind\-%
rømme, at den, der endskjøndt han antager Troens Forud\-%
sætninger dog begjærer at klare disse i Viden og Viden\-%
skab, og mener at have Trang til en Troesvidenskab, afvises
ved Paradoxet. Vi indrømme, at den Existerende, den i
Existensens Kriser Bestedte, maa stimuleres af Modsi\-%
gelsens Braad, maa atter og atter saares af det Paradoxe.
Men nu er et Spørgsmaal tilbage, og det er det, der har
bevæget sig igjennem vore Undersøgelser, dette nemlig:
Er Troen da ogsaa for sig selv et Paradox? Er den
Troende saa blind for sine egne Tankemomenter, at de
umulig kunne udvikle sig for ham til klare, i sig sam\-%
menhængende Begreber? Dette maa benegtes. Troen
er i sig et Erkjendelsesprincip; Troens Ord er Lysets
og Livets Ord, og fra dette udgaaer al Klarhed over
Religionens Indhold. Vel er Troen „en hellig Liden\-%
% "hellig Lidenskab": a raised mark after "hellig" is a risen space/quad; not transcribed
skab“, men naar den hellige Lidenskab gjennemreflec\-%
teres, viser den sig som et pathetisk Udtryk for den efter
Hellighed stræbende Villie: Troesprincipet er i dybeste
Forstand et Villiesprincip, Troen selv en Villiesact, en
% --- p. 235 ---
\opage{235}\setcounter{footnote}{0}%
inderste Handling (nemo credit, nisi volens): Troestankerne
ere Villiestanker, alle tilsammen afgjørende Personlig\-%
hedstanker. Idet det hele, i sig sammenhængende Ind\-%
begreb af religiøse Villiestanker udvikles og klart frem\-%
stilles i et System af Troesbegreber, faae vi, ikke en
paradox Existensmeddelelse, men en, fra Troens Stand\-%
punkt, gyldig Lære om Religionen, en systematisk sam\-%
menhængende Religionslære. Dette med Hensyn til det
Kierkegaardske Stadium. Vi komme nu til det næste
Spørgsmaal.

Indrømmes det nemlig, at der gives et System af
Religionserkjendelser, at vi paany maae faae en Reli\-%
gionslære, synes det valgte Standpunkt, skjøndt vi have
erklæret os afgjørende mod al Theologie, dog at falde
tilbage under det Theologiske. Jeg vil begynde med
en dum Bemærkning, en Bemærkning der ellers fra saa\-%
mange Sider anmelder sig som et Foster af sjelden
Kløgt; Bemærkningen er denne: „at give en systematisk
Fremstilling af de religiøse Erkjendelser er at gjøre Re\-%
ligionslæren objectiv; at gjøre Religionslæren objectiv er
at stille sig paa Betragterens Standpunkt og gjøre Læren
videnskabelig, altsaa, naar man dog skal holde paa Sub\-%
jectivitetens ubetingede Gyldighed og afvise Videnskaben,
en aabenbar Modsigelse“. Hvorfor jeg kalder dette en
dum Bemærkning, kan jeg let forklare. Dersom det ikke
var muligt at udvikle og begrunde Religionslæren, fordi
det kun kan skee paa den Betingelse, at man stiller sig
paa Betragterens Standpunkt, saa vilde det jo heller ikke
være muligt, hvad dog vitterlig er skeet i „Uvidenskabe\-%
lig Efterskrift“, at give en humoristisk experimenterende
Fremstilling af de christelige Problemer. Johannes Cli\-%
macus staaer jo selv fra Først til Sidst paa Betragterens
Standpunkt; og dog er der aldeles ingen Modsigelse i, at
han, staaende paa dette Standpunkt, smiler ad den „be\-%
% --- p. 236 ---
\opage{236}\setcounter{footnote}{0}%
skedne objective“ Subjectivitets „applauderede Heroisme“,
% "applauderede": a stray acute-like mark stands over the first a; read as a blemish, not an accent
og skærper det Subjective. Her er nemlig en uhyre For\-%
skjel imellem den betragtende Bevidsthed, der argumenterer,
og den Bevidsthed, hvoraf, hvorfra og hvortil der argumen\-%
teres. Saavel i Religionslæren som ved den humoristiske
Experimenteren maa der under hele Betragtningen argu\-%
menteres fra den i sig concentrerede Selvhed, den ubetinget
gyldige Subjectivitet. Nu følger en anden, mere fornuftig
Indvending. Den systematiske Udvikling af et til Troes\-%
principet svarende Lærebegreb er, og maa ifølge sin
Natur nødvendigviis være, videnskabelig; men hvad er
en videnskabelig Udvikling af den positive Religions Ind\-%
hold andet end Theologie? „Du kan“, siger man,
„vælge dit Standpunkt, hvor du vil: dersom du virkelig
løser den Opgave, at fremstille den aabenbarede Reli\-%
gions Indhold i en sammenhængende, af eet Princip fly\-%
dende Udvikling, saa er og bliver du Theolog, og det
nytter ikke, at du protesterer“.

For at indsee, hvorledes det forholder sig med
denne Paastand, behøve vi blot at kaste et Blik paa de
theologiske Standpunkter.

Det første, som melder sig, er da det nyere specu\-%
lativ-dogmatiske. Dette Standpunkt har jo villet optage
alle Troens Forudsætninger, har villet bygge Troes\-%
erkjendelsen paa Troeslivet, har ikke villet udvikle Be\-%
greberne af den rene Viden, men af „den gode Villie“,
og har paa denne Maade villet gjøre Villiesprincipet til
Videnskabsprincip. Men at gjøre Villiesprincipet til
Videnskabsprincip er umuligt. Skulde Theologien i
Kraft af sin gode Villie virkelig være Videnskab, da
maatte den være den høieste af alle Videnskaber; den
maatte i Villiens Navn kunne regulere de øvrige, i For\-%
hold til det Hellige lavere staaende og derfor underord\-%
% --- p. 237 ---
\opage{237}\setcounter{footnote}{0}%
nede Videnskaber; den maatte kort og godt kunne re\-%
tablere sit middelalderlige Universalmonarchie.

„Men et nyt theologisk Universalmonarchie, hvilken
uhyre Tanke! Mon en Tanke af denne Høide kan, uden
at bukke sig saa dybt, at den bliver ganske kroget, gaae
ind ad Virkelighedens Dør? Ogsaa jeg har mødt den
høie Tanke; den er mange Gange traadt mig imøde,
naar jeg betragtede Dogmatiken, ikke fra Skyggesiden,
men fra Lyssiden; thi, ihvorvel jeg, ifølge min Opgave,
som oftest har maattet gjennemvandre dens Skygge\-%
partier, har jeg derfor ikke glemt, hvor jeg endnu kan
finde de solbeskinnede Steder. Men naar jeg saa, efter
at have beskuet Idealet i al dets Herlighed, alvorlig
spurgte mig selv: kan et saadant Ideal realiseres? har
jeg bestandig faaet et benegtende Svar. Det er jo vist\-%
nok saa, at det Godes Idee, opfattet som Kjærlighedens,
den hellige Villies, Idee er i guddommelig Viden Eet
med den evige Sandhedsidee; men heraf følger jo ingen\-%
lunde, at Sandhedsideen i en menneskelig Viden er Eet
% "guddommelig", "menneskelig" (first): underlined in pencil by a reader; not transcribed
med det Godes Idee i menneskelig existentiel Villen.
Det er jo vistnok saa, at der i Viisdommen selv ikke
kan tænkes nogen Kløft imellem Troen og Videnskaben:
men hvorfor er Gud ingen Videnskabsmand? Fordi
Gud ikke er nogen Troende. At tænke sig den gud\-%
dommelige Viden som en videnskabelig, være sig philo\-%
sophisk eller theologisk Viden, er ligesaa gudsbespotte\-%
ligt som at tale om en troende Gud. Hvorfor skal et
Menneske opfatte det Absolute i Troen? Fordi den re\-%
neste og høieste Udvikling af den menneskelige Viden
er, ikke guddommelig, men videnskabelig. Hvorfor skal
et Menneske opfatte det Absolute i Viden? Fordi Troen
er en menneskelig, ikke en guddommelig, Erkjendelse.
Hvorfor viser det Absolute, der opfattes i menneskelig
Viden, sig saa ueensartet med det Absolute, der opfattes
% --- p. 238 ---
\opage{238}\setcounter{footnote}{0}%
i Troen? Fordi de relative Opfattelsesformer gribe det
Absolute, ikke i dets centrale Dybde, saaledes som det
af os \textit{in abstracto} maa tænkes at være med uendelig
Concretion i det guddommelige Selv, men i to polære
Extremer, af hvilke det ene er betegnet ved den almene,
i rationel Nødvendighed grundende Tænkning, det andet
ved det punktuelle, i Samvittigheden spirende Villiesliv.
Her er i det Absolute Noget, jeg absolut skal troe, fordi
der i det Absolute er Noget, jeg absolut skal vide“. (Den
gode Villie som Magt i Videnskaben S. 6---7).

At der er et Forhold imellem Villie og Videnskab,
kan ikke benegtes; den, der skal tænke, maa ville tænke;
den, der skal dyrke en Videnskab, maa ville dyrke den.
Ja, Forholdet er endogsaa personligt, ethisk, forsaavidt
Enhver, der offrer sin Tid paa at dyrke en Videnskab,
maa være ansvarlig for, hvorledes han anvender sine
Kræfter, om der ikke indtræder Tilfælde, da han af høiere
Hensyn maa bryde af; o. s. v. Men disse og lignende
Betragtninger have aldeles Intet med vort foreliggende
Spørgsmaal at gjøre. Personligheden, der staaer til An\-%
svar for Alt, hvad den tænker og gjør, staaer ogsaa til
Ansvar for den Maade, hvorpaa den opfatter Forholdet
mellem Villien og Videnskaben. Den gode Villie er
ikke en Magt i Videnskaben, men den er en Magt i
Religionen. De religiøse Tankeproblemer ere i Grunden
% "De religiøse": a raised mark after "De" is a risen space/quad; not transcribed
alle Villiesproblemer. Ved ethvert Villiesproblem argu\-%
menteres der fra Selvet, ved Selvet, til Selvet; medens
der ved ethvert Vidensproblem argumenteres fra Sagen,
ved Sagen, til Sagen: at argumentere fra Naturen er at
argumentere fra Sagen, at argumentere fra Skabelsen er
at argumentere fra Selvet; at argumentere fra den phy\-%
sisk-organiske Individualitet er at argumentere fra Sagen,
at argumentere fra Idealmenneskets rene Hjerte og synd\-%
frie Villie er at argumentere fra Selvet; at argumentere
% pp. 238--239: "argumentere" repeatedly underlined in pencil by a reader; not transcribed
% --- p. 239 ---
\opage{239}\setcounter{footnote}{0}%
fra en indre Ideenødvendighed er at argumentere fra
Sagen, at argumentere fra Aandens Frihed er at argu\-%
mentere fra Selvet, o. s. frdl.

Confusionen i den christelig-speculative Dogmatik
forraader sig ved ethvert af Hoveddogmerne, Skabelsen,
Incarnationen, Inspirationen, idet to halve Tanker, en
% "to halve": a raised mark after "to" is a risen space/quad; not transcribed
halv fra Selvet og en halv fra Sagen, en halv fra Villien
og en halv fra Videnskaben, en halv fra det Subjective
og en halv fra det Objective, reise sig mod hinanden,
brydes med hinanden og ende med at neutralisere hin\-%
anden. Denne de ueensartede Tankers \textit{concordia discors}
er allerede beskreven paa adskillige Maader og kan end\-%
nu beskrives paa flere.

„Man veed, hvorledes det er gaaet „den chri\-%
stelige Dogmatik“, hvilke Uheld den har havt med
det Antirationelle i de store enfoldige Hoveddogmer.
Som nu Skabelsesdogmet; ja, det vilde jo være et meget
dybsindigt og særdeles videnskabeligt Dogme, dersom
der blot ikke var den Hage, det Antirationelle, derved,
at saasandt Dogmatiken rigtig oprigtig skal have det
ud, at Verden er skabt af Intet, saa er der i Verden
ingen Natur, og skal den erkjende, at her er en virkelig
Natur, saa kan den umulig faae det ud, at Verden er
skabt af Intet. (Jvfr. „Den gode Villie“ S. 96---104).
At slaae det Hele i speculativt Hartkorn, lade Natur
være Skabelse og Skabelse Natur, for at Tro kan være
Viden og Viden Tro, er en saa ydmyg Inconseqvens,
at den sikkerlig ikke vil undlade at ydmyge Theologien.
Men Incarnationen, Gudmenneskedogmet? Ja Gudmen\-%
neskedogmet, det store Dogme, som Aarhundreder igjen\-%
nem med Rationelt, Antirationelt og Irrationelt har sat saa
mangfoldige baade Tanker og Lidenskaber i Bevægelse, vilde
utvivlsomt være Facultetets Yndlingsdogme, Theologiens
Palladium, det centrale Hjertedogme, dersom „den christe\-%
% --- p. 240 ---
\opage{240}\setcounter{footnote}{0}%
lige Dogmatik“ ikke havde havt den Sorg paa det Sidste,
den store Sorg med Sønnen. Hvilken Sorg? Ak, hans
Bevidsthed, hans dogmatiske Logosbevidsthed er jo,
\textit{horribile dictu}, revnet, dualistisk revnet; den ene Søn
er --- ved Revnen --- bleven til to Sønner: en, hvis Be\-%
vidsthed er fra Evighed af, men som ikke er bleven
incarneret, og en anden, en yngre, som er bleven incar\-%
neret, men hvis Bevidsthed ikke er fra Evighed af.
(Jvfr. „Den gode Villie“ S. 91---92). Vistnok holder
Theologien, alt hvad den kan, paa „Bevidsthedens Een\-%
hed“, ogsaa i Incarnationen; men den Revne, den
Revne! Incarnationsdogmet er et meget vanskeligt Dogme.
Men Inspirationsdogmet, det er dog vel ikke saa van\-%
skeligt? Er der noget Dogme, der har Aandens Vidnes\-%
byrd, saa maa det vel være Aandsudgydelsesdogmet.
Da det nu tillige skal være ganske vist, at Dogmatiken
ogsaa, i Kraft af Aandsudgydelsesdogmet, har „Aandens
Vidnesbyrd“, saa tør vi jo vente, at Aandsudgydelses\-%
dogmet maa være ligesom skabt for Dogmatiken. Ja,
det vilde virkelig have været som skabt for dogmatisk
Speculation, dersom „den christelige Dogmatik“ ved at
snuble over det Antirationelle, ikke i Faldet havde havt det
Uheld, at dens Pintsemirakel var kommet til at staae paa
Hovedet. Men at vende op og ned paa Pintsemiraklet,
det store „kirkestiftende Mirakel“, i den Grad, at Tilhø\-%
rerne høre mange \emph{forskjellige} Tungemaal, uagtet Apost\-%
lene tilsammen kun tale et \emph{eneste} Tungemaal; at det
Øverste altsaa bliver nederst, det Nederste øverst, det er
mere end reformatorisk, det er revolutionairt. At Dog\-%
matiken med sin Mirakelrevolution har meent det godt,
inderlig godt baade med Apostlene og med Kirken og
Pintsefesten, skulle vi gjerne indrømme; men revolutio\-%
nairt er det, det er vist (Jvnfr. „Den gode Villie“ S.
84). Inspirationsdogmet er altsaa et meget vanskeligt
% --- p. 241 ---
\opage{241}\setcounter{footnote}{0}%
Dogme. Ja, og Daabsdogmet med det bevidste „mystisk
Magiske“ er et meget vanskeligt Dogme; og Dogmet om
Christi Legemes speculative Nærværelse i Nadveren er
et meget vanskeligt Dogme; og Dogmet om den evige
Fordømmelse er et vanskeligt, meget vanskeligt, Dogme.
Der er i den evige Fordømmelse noget saa antirationelt,
at „den christelige Dogmatik“ udtrykkelig har maattet
forlange en Ikke-Fordømmelse (Apokatastasis), paa det at
Fordømmelse og Ikke-Fordømmelse under en forfærdelig
Brydning kunde ende det Hele med en høitidelig An\-%
tinomie, og Antinomien (Fornuftstriden) staae som
Vidnesbyrd, et Systemet afsluttende Vidnesbyrd, om, at
der er Eenhed i den dogmatiske Bevidsthed, at Theolo\-%
gien i Sandhed magter det Antirationelle. --- Saalænge
vi letsindig smutte bort fra hver enkelt Forhandling om
et hvilketsomhelst enkelt Dogme ved at sige, naar det
kniber: ja her er naturligviis vanskelige Punkter, men
forøvrigt o. s. v., er det en let Sag at lege Forsonings\-%
leg med Tro og Viden. Men skal der holdes Stand paa
bestemte Punkter, skal der i Alvor gjøres Rede for de
enkelte bestemte Dogmer, da vil man snart komme til
at sande, at alle Dogmer uden Undtagelse ere lige van\-%
skelige, lige antirationelle, lige stridende mod videnskabe\-%
lig Begriben. Troens Tanker bære alle Troens Mærke,
ere alle forsynede med den antirationelle Braad, saa at
den, der vil kramme dem videnskabeligt, nok omsider
skal fornemme Stikket“. (Om Holbergs Kirkehistorie og
Theologie. Kbhvn. 1866. S. 35---37).

Det Samme er i mit Skrift, Den gode Villie som
Magt i Videnskaben S. 92---93, sagt paa en anden Maade:
„Det Eiendommelige ved den dogmatiske Tænkning er,
at der tænkes videnskabeligt i Troen: \textit{credo, ut intelligam.}
At der tænkes videnskabeligt i Troen, vil sige, at der
ikke tænkes, men forestilles, anskues, phantaseres paa
% --- p. 242 ---
\opage{242}\setcounter{footnote}{0}%
metaphysisk Maade til Opklaring af Mysterierne, til
Bedste for Aabenbaringen, til Glæde for Troen. Dog\-%
matiken, der objectivt bygger paa Skriftens Ord, seer i
Skriftstederne noget ganske Andet og Dybere, end enten
Christus selv eller hans Apostle nogensinde have seet,
den seer det Metaphysiske. „Han vilde“, hedder det
(Dogm. § 133) om Christi Selvfornedrelse, „ikke blot
leve i den rene Guddomsherlighed, i metaphysisk Majestæt;
men han udtømte sin Fylde i den ringe Tjenerskikkelse
for at kunne blive Hovedet i Kjærlighedens Rige, Forsoneren
og Forløseren. Men denne Selvfornedrelse maa tillige“
--- her kommer den rette positive Mirakelforklaring ---
„sees som hans Selvfuldendelse“, Selvfuldendelsen nemlig af
den i sig evig fuldendte Logos. Idet han, den ene og
samme Logos, „lever paa menneskelig Viis, og som
Menneskens Søn kun besidder sin Guddom i den men\-%
neskelige Individualitets Indskrænkning, i den menneske\-%
lige Bevidstheds begrændsede Former: maa han vistnok
siges at leve i Fornedrelse og Fattigdom, efterdi han
har givet Afkald paa den majestætiske Herlighed, i
hvilken han som den allestedsnærværende Logos gjen\-%
nemlyser Alskabningen.“ (Dogm. § 133). Lægge vi
nu Mærke til de to Udtryk: „han har givet Afkald“, og
„han gjennemlyser“, og erindre, at det er den samme,
den selvsamme Logos, det ene og selvsamme Logos-Jeg,
der her har givet Afkald paa den majestætiske Herlig\-%
hed, hvorpaa han dog ikke et eneste Øieblik har givet
Afkald, hvorledes skulle vi da kunne tvivle om, at det
netop er det høiere Phlyariske, det dybere Phlyaretiske,
som ved denne Forklaring bliver sat i sit positive
Moment.“

Undersøgelsernes Resultat kan fastholdes i følgende
lidet opbyggelige Betragtning: „Mediationen af Tro og
Viden lader sig naturligviis kun gjennemføre paa den
% --- p. 243 ---
\opage{243}\setcounter{footnote}{0}%
Betingelse, at begge Parterne nærme sig hinanden med
al mulig diplomatisk Forekommenhed, d. v. s. i al Op\-%
rigtighed, Uegennyttighed, kun besjælede af de varmeste
Ønsker om Fred og god Forstaaelse. Under den gjen\-%
sidige Tilnærmelse foregaaer der en lykkelig Forandring
med dem begge: Troen bliver gemytlig, Viden god\-%
modig. At Troen bliver gemytlig, vil sige, at den op\-%
giver sin gammeldags orthodoxe Prikkenhed, undseer
sig ved hvad der er Jøder en Forargelse og Grækere en
Daarskab, smiler til Philosophien og forelsker sig i Spe\-%
culationen. Det kan naturligviis ikke være nogen al\-%
mindelig Tro, ikke den folkelige, enfoldige Asyltro, der
teer sig saa charmant; nei, det er en dannet, velop\-%
dragen Tro, en Tro, der kjender sine Mystikere og har
hørt Forelæsninger i Institutet; kun en saadan Tro, en
subjectiv Tro med et objectivt Hjerte, en Tro, der ikke
vilde elske praktisk, dersom den ikke elskede theoretisk,
kan aabne sin Favn, sin dybe, dogmatiske Favn, for
den høiere Viden. At den Viden, der uden Qvalme
skal kunne give sig hen i en saa dannet, saa institut\-%
dannet Kjærlighed, absolut maa være \emph{godmodig}, er vel
af sig selv indlysende. Her er nemlig en Forskjel, en
meget betydelig Forskjel, imellem den godmodige og
den ikke godmodige Viden. Den ikke godmodige Viden
kunne vi med Grund kalde den drillende, nemlig den
dogmatisk drillende, fordi den altid har sin Glæde af at
drille Dogmatiken. Drilleriet bestaaer ganske simpelt
deri, at den med ubarmhjertig Conseqvens forfølger en\-%
hver Modsigelse, hvormed Dogmatiken i sin Troskyl\-%
dighed har indladt sig. Af et Par uskyldige: „Maa-%
ikke-tænkes“, fremmaner den ved Hjælp af det Logisk-%
Physiske en dæmonisk Vrimmel af „Maa-ikke-tænkes;“
den faaer ved kritisk dialektiske Hexekunster de dog\-%
matiske Modsigelser til, om vi saa maae sige, at yngle.
% --- p. 244 ---
\opage{244}\setcounter{footnote}{0}%
Ja, hvad der er det mest Oprørende, just det dybeste
Positive, det, hvorom Dogmatiken udtrykkelig siger, at
\emph{det maa tænkes}, forvandler den til et uroligt Negativt,
et i Bund og Grund Utænkeligt, et „Kan-ikke-tænkes“:
det er ondskabsfuldt. Den ondskabsfulde Viden, den
Viden, der ikke vil slippe det Logisk-Physiske, er en
syndig Viden, en Viden, der ikke kan sige \textit{credo, ut
intelligam}, en \emph{vantro} Viden. Anderledes med den god\-%
modige Viden: salige ere de, der ikke ville kritisere den
godmodige Viden. Den godmodige Viden seer vel un\-%
dertiden noget tankeløs ud; men, ved Dogmatiken, den
er ikke tankeløs, thi den gaaer i Tanker. Hvorfor gaaer
den i Tanker? Den er forelsket, subjectiv forelsket i
Troens objective Hjerte: det er jo en troende Viden.
Den troende Viden er en Viden, der længes --- ud over
det Logiske ind i det Overlogiske: Mirakelmetaphysik
er dens Længsels Maal. Naaer den ikke sin Længsels
Maal, saa døer den; men naaer den sin Længsels Maal,
saa lever den, lever af Troens Positive, blæser ad Kri\-%
tikens Negative og neddysser alle Modsigelser i absolut
Glæde over de relative Mysterier. Skulde en eller
anden Modsigelse være saa ustyrlig at skrige ud fra
Mysteriernes Rede, hvad gjør saa det? „De Franske
tage det ikke saa nøie“. Den troende Viden er god\-%
modig: godmodig synker den til Dogmatikens Bryst,
godmodig hviler den ved Troens Hjerte. Praktisk er
Godmodigheden Kjærlighed, theoretisk er Kjærligheden
Godmodighed. Maa det i det Praktiske siges, at Kjær\-%
lighed skjuler Syndernes Mangfoldighed, da maa det
her, hvor det dybeste Praktiske tillige er det høieste
Theoretiske, siges, at Godmodighed, den troende Videns
utrolige Godmodighed, skjuler Modsigelsernes Mangfol\-%
dighed. Paa disse Vilkaar er Foreningen af Tro og
% --- p. 245 ---
\opage{245}\setcounter{footnote}{0}%
Viden fuldbragt; det Phlyariske er sat i speculativ
Fuldendelse.“

Her kan altsaa ingen Tvivl være om, at der er en
principiel Strid mellem mit Standpunkt og det specula\-%
tiv-dogmatiske.

Saa opstaaer et tredie Spørgsmaal: Hvorledes for\-%
holder det valgte Standpunkt sig til den gamle ortho\-%
% a faint raised mark after "holder" and another after "et" (next line but one)
% are risen spaces, not punctuation; not transcribed.
doxe Lære? Thi naar der afviges fra Paradoxet derved,
at der danner sig et Lærebegreb, naar der afviges fra
den moderne Begrundelse derved, at Mediationen prote\-%
steres, da synes vi jo at maatte komme tilbage til den
gamle orthodoxe Fremstilling. Er dette nu Tilfældet?
Jeg svarer: det valgte Standpunkt kan ingenlunde flyde
sammen med det gamle orthodoxe; thi det orthodoxe
Standpunkt charakteriserer sig derved, at det pointerer
Troesbegreberne ganske rigtig, men motiverer dem
saare slet. Hvad der er charakteristisk for Troesbe\-%
greberne, veed den gamle Dogmatik med sjelden Skarp\-%
sindighed at udfinde, men den bedaares af Videnskaben
og optager reent menneskelig Viden i Udviklingen. Dens
Grundmodsigelse er, at ogsaa den gjør Troeslæren til
Videnskab, og Conseqvensen er, at denne Lære saa
igjen hjemfalder til videnskabelig Kritik og opløses af
Videnskaben.

For at blive Videnskab maa Theologien nødvendig
bruge Fornuften; men for at være Troeslære, rettroende
Troeslære, maa den tillige tage Fornuften fangen under
Troens Lydighed. At bruge Fornuften er at anerkjende
de Følgesætninger, som med logisk Nødvendighed frem\-%
gaae af de i Troens Navn opstillede Forudsætninger;
men at tage Fornuften fangen, med andre Ord, ikke
bruge Fornuften, er at negte en Følgesætning, uagtet
den vitterlig og med logisk uafviselig Nødvendighed
fremgaaer af sin Forudsætning. Som Troesvidenskab
% --- p. 246 ---
\opage{246}\setcounter{footnote}{0}%
skal Theologien altsaa baade anerkjende og ikke aner\-%
kjende Følgesætningens Fremgang af Forudsætningen,
skal baade bruge og ikke bruge Fornuften; dette er
Selvmodsigelsen, en Selvmodsigelse, der er ligesaa gammel
som Theologien. Der er i den hele Kirkehistorie ikke
en eneste dogmatisk Strid, som jo bærer Vidne herom.
Et mere talende Beviis for den lutherske Orthodoxies
videnskabelige Selvmodsigelse end dens Indsigelse mod
Calvins Lære om Naadevalget kan neppe tænkes. Man
indrømmer aabenlyst, at Calvin er videnskabelig conse\-%
qvent, og ligesaa aabenlyst fordømmer man hans Conse\-%
qvens; og endda er Theologien en Videnskab!

„Concordieformlens velgjørende inconseqvente Fore\-%
stillingsmaade i Forhold til den gudsbespottelig conse\-%
qvente Calvinske er“, siger en luthersk orthodox Kirke\-%
historiker, „i Korthed denne: Helliggjørelse og Salig\-%
hed ere ifølge begge kun et Værk af den guddomme\-%
lige Naade, Fordømmelsen en Følge af egen Skyld.
Men medens nu Concordieformlen ydmyg bliver staaende
ved denne sidste Sats, og derfor udtrykkelig tillægger
Mennesket alene Skylden, og ikke den manglende
Naade \ldots{}, argumenterer Calvin videre: Men hiin Skyld
finder kun Sted, fordi Gud ikke ogsaa her kraftig med\-%
deler sin Naade; men Gud meddeler ikke ogsaa her sin
Naade, fordi han fra Evighed af har besluttet Fordøm\-%
melsen, og selv Menneskets Fald har maattet tjene til
at realisere denne Raadslutning, ligesom overhovedet
Alt er besluttet fra Evighed --- Gud har altsaa selv
villet Menneskenes Synd. Denne sidste nødvendige, om
end nok saa meget besmykkede Følgesætning, afskyer
især Concordieformlen som modstridende Guds Hellig\-%
hed og det guddommelige Ord i Skriften, skjøndt den
heller ikke miskjender Vanskelighederne i denne Læres
% --- p. 247 ---
\opage{247}\setcounter{footnote}{0}%
hele Connex, hvilke ingen menneskelig Forstand for\-%
maaer at løse.“

Skulde der, efter hvad her er oplyst, være nogen
Mening i at anvende Fornuftens Love i nogle Tilfælde
og afskaffe dem i andre, da maatte der i den theolo\-%
giske Viden jo være en ufeilbar guddommelig Regel
hvorefter Anvendelsen og Afskaffelsen kunde foretages.
Men hvor skulle vi vel tye hen for at finde en saadan
Regel? Til Fornuften nytter det ikke at henvende sig;
thi dersom Fornuften selv kunde lære os, i hvilke Til\-%
fælde vi for Troens Skyld burde foretrække „den yd\-%
myge Inconseqvens“ og suspendere Fornuftloven, saa
blev Fornuftens Anvendelse jo netop stadfæstet, just
som den skulde afskaffes: Fornuften blev paa denne
Maade Dommer i sin egen Sag og fik et absolut Herre\-%
dømme over Troen. Er det derimod Troen, adskilt fra
Fornuften, altsaa en ved sin Adskillelse fra Fornuften
ufornuftig Tro, der baade skal ophæve og stadfæste For\-%
nuftloven, ja saa beholder Troen jo vistnok et ubetinget
Herredømme over Fornuften; men hvad bliver der saa
med denne i sig ufornuftige Tro af Theologiens Viden\-%
skabelighed? At tale om et vist fornuftigt Overfornuftigt,
et mystisk Noget, som ved at staae baade over For\-%
nuften og over Troen skulde regulere Forholdet imellem
dem, er at snakke hen i Taaget; thi Concordieformlens
„ydmyge Inconseqvens“ beviser jo netop, at her i Grunden
hverken er Regulativ eller Conseqvens. I Striden om
Naadevalget erklærer den lutherske Theolog høitidelig
Calvins „Følgeslutning“ for at være „modstridende Guds
Hellighed og det guddommelige Ord i Skriften“. Paa
hvilken Side er nu Conseqvensen? Naturligviis paa begge
Sider. Calvin gaaer med Lutheraneren ud fra den For\-%
udsætning, at Helliggjørelse og Salighed kun er et Værk
af den guddommelige Naade, og slutter saa paa sin Viis
% --- p. 248 ---
\opage{248}\setcounter{footnote}{0}%
conseqvent, at Salighedens Udeblivelse absolut maa være
en Følge af Naadens Udeblivelse, og denne Udeblivelse
igjen en Følge af den guddommelige Raadslutning. Con\-%
cordieformlens Forsvarer gaaer med Calvin ud fra den
Forudsætning, at Fordømmelsen er en Følge af Menne\-%
skets egen Skyld, og slutter saa deraf ganske conse\-%
qvent, at den umulig kan være en Følge af Guds evige
Raadslutning. Paa hvilken Side er Inconseqvensen?
Naturligviis paa begge Sider. Calvin er inconseqvent,
naar han sætter Fordømmelsen som en Følge af Synden
og vil gjøre Gud til Fordømmelsens evige Ophav uden
dog at gjøre ham til Syndens Ophav. Concordieformlens
Forsvarer er inconseqvent, idet han, for at gjøre Gud
alene til Salighedens Ophav, erklærer, at Mennesket vel
har Evne til at modstaae Naaden, naar det gjælder For\-%
tabelse, men derimod ingen Evne til at samvirke med
% "samvirke med / altsaa ingen Evne": printed so, no object after "med" (sense
% wants e.g. "med Naaden,"); a faint mark after "altsaa" is a risen space, not a comma.
altsaa ingen Evne til at modstaae Naaden, naar det
gjælder om Frelse og Salighed. Lutheranerens Incon\-%
seqvens bestaaer deri, at han snart retter Principet efter
Udfaldet, snart igjen Udfaldet efter Principet: er Ud\-%
faldet Fordømmelse, saa har Mennesket kunnet modstaae
Naaden, thi Mennesket maa selv være Skyld i sin For\-%
dømmelse; er Udfaldet derimod Frelse og Salighed, saa
har Mennesket ikke kunnet modstaae Naaden. Havde
Menneskets Villie kunnet modstaae den saliggjørende
Naade, saa havde dette Menneskes Villie netop, ved
ikke at modvirke Naaden, samvirket med Naaden. Men
Tanken om, at Menneskets Villie skulde paa nogen
Maade kunne samvirke med den guddommelige Naade,
er en Synergisme, hvorover Concordieformlen kun kan
udtale sig fordømmende.

Ere nu begge, den strenge Calvinist og den strenge
Lutheraner, lige conseqvente og lige inconseqvente, saa
er det jo Videnskaben umuligt at tilkjende den ene For\-%
% --- p. 249 ---
\opage{249}\setcounter{footnote}{0}%
trinet fremfor den anden. Hvorpaa begrunder vel den
lutherske Theologie sit Fortrin for den reformeerte? Ikke
paa sin conseqvente Charakteer (thi den indrømmer selv
sin Inconseqvens); heller ikke ligefrem paa sin Inconse\-%
qvens, men paa det i Inconseqvensen Ydmygende, altsaa
paa Ydmygheden. Naar Calvinisten er inconseqvent, saa
har Lutheraneren Lov til at bruge sin Fornuft, appellere
til Logiken og holde paa Conseqvensen, thi hvad er den
calvinistiske Inconseqvens vel Andet end en uvidenskabelig
Selvmodsigelse, en dogmatisk Prostitution? Men naar Lu\-%
theraneren ikke blot bliver greben, men aabenbart griber
sig selv i vitterlig Inconseqvens: hvad er det saa? Ja,
saa er det ikke Selvmodsigelse, ikke videnskabelig Pro\-%
stitution, men Ydmyghed, reen, skjær, theologisk Ydmyg\-%
hed. Og hvad er saa Calvinistens Conseqvens? Natur\-%
ligviis Hovmod, det rene Hovmod, den uomskaarne For\-%
nufts, den uigjenfødte Forstands ugudelige Indbildskhed.

Dersom vi endda til Slutning kunde øine et Tegn,
et Mærke, et ubedrageligt Kjendemærke, hvorefter man
paa videnskabelig Maade var istand til at skjelne imellem
de uvidenskabelige, altsaa prostituerende, og de ydmyge,
altsaa ikke prostituerende, Inconseqvenser: saa kunde
den skibbrudne Theologie maaskee endnu haabe paa en
Redningsplanke; men hvor finde vi et saadant Kjende\-%
mærke? „I den hellige Skrift,“ svarer Lutheraneren; „thi
den hellige Skrift er Guds Ord, og der staaer udtrykkelig
i den hellige Skrift: Gud vil, at Alle skulle blive salige
og komme til Sandheds Erkjendelse.“ Calvinistens Hov\-%
mod sees altsaa deraf, at hans Conseqvens staaer i Strid
med Guds Ord. Men der staaer jo ogsaa udtrykkelig i
den hellige Skrift, at „Gud forbarmer sig over den, som
han vil, og forhærder den, som han vil.“ Gaae vi nu
ud fra det sidste Skriftsted, saa bliver Calvinistens Con\-%
seqvens jo ydmyg, medens Lutheranerens Inconseqvens,
% --- p. 250 ---
\opage{250}\setcounter{footnote}{0}%
ved at komme i aabenbar Strid med Skriftens Ord: „Men,
o Menneske, hvo er Du, at Du vil gaae i Rette med Gud?
Mon Noget, som er dannet, kan sige til den, som dannede
det: hvi gjorde Du mig saaledes? Eller haver Pottemageren
ikke Magt over Leret, af det samme Stykke at gjøre et
Kar til Ære, men et andet til Vanære?“ --- forvandler sig
og bliver et Tegn, ikke paa Ydmyghed, men paa Hov\-%
% "dét": a short acute-like stroke stands over the e at 300 ppi; set as printed.
% Doubtful: it may be a speck, in which case read "det".
mod, dét mest oprørende Hovmod, det Hovmod at ville
gaae i Rette med Gud og negte ham Ret til fra Evighed
af at udvælge --- „et Kar til Ære, et andet til Vanære“.
Dette maa være nok til at belyse den orthodoxe Theo\-%
logies oprindelige Grundmodsigelse.

Denne Grundmodsigelse maa haves skarpt for Øie,
naar man vil forstaae Forholdet mellem den orthodoxe
Theologie og den paafølgende Rationalisme.

Rationalismen har, saa stridende den end er mod
den orthodoxe Lære, udviklet sig som en videnskabelig
Conseqvens af Orthodoxiens egen Methode. Hvis denne
ei havde taget, hvad den kaldte Fornuften, den fangne
Fornuft, tilhjælp, saa at den kom i den Modsigelse at
betjene sig af Fornuften og dog at holde den fangen, var Ra\-%
tionalismen aldrig trængt ind i Kirken. Dette har jeg, som
allerede antydet, søgt at klare i et lidet Skrift „Om Holbergs
Kirkehistorie og Theologie, et Bidrag fra Fortiden til at
belyse Nutiden“, af hvilket Skrift jeg endnu vil fremhæve
Følgende: Al Viden og Videnskab er i sit Væsen rationel;
her taales ingen confessionelle Haandfæstninger, intet yp\-%
perstepræsteligt Formynderskab \ldots{} uden Fornuft kunde
Kirken umulig opføre sin Lærebygning; Fornuften var
altsaa den egentlige Bygmester, og Bygmesteren den
Eneste, der kjendte Techniken og kunde vurdere Byg\-%
ningens Varighed. Vistnok havde Fornuften fra Begyn\-%
delsen ikke opført den store Bygning enten af egen Til\-%
skyndelse eller for egen Regning; det var tvertimod Tvangs\-%
% --- p. 251 ---
\opage{251}\setcounter{footnote}{0}%
arbeide og paa Troens Bekostning; Kirketroen var, om
vi kunne bruge dette Billede, den egentlige Bygherre.
Skulde Bygherren nu for bestandig have sikkret sig mod
Bygningens Nedrivelse, da maatte Troen, ligerviis som
hiin barbariske Despot, der ved at lade Mesteren myrde
sikkrede sig imod, at der ikke blev opført et andet endnu
herligere Mesterværk, have sikkret sig imod, at Fornuften
ikke nedrev det gamle Bindingsværk, ryddede Pladsen
og opførte en ny og grundmuret Bygning i en ny Stiil,
ved at faae Fornuften qvalt. Men til en saadan Ugjer\-%
ning havde Kirketroen, hvor despotisk den end iøvrigt
teede sig, hverken Mod eller Kraft. Saa skete det, at
Fornuften, saasnart den var bleven færdig med Lærebyg\-%
ningen for Troens Regning, øieblikkelig begyndte at rive
den ned igjen for egen Regning.

Da der nu som bekjendt ligesaavel arbeides indeni
en Bygning, naar der rives ned, som naar der bygges
op: saa fik det for udenfor Staaende let Udseende af, at
det Nedrivningsarbeide, som Fornuften i Slutningen af
18de og Begyndelsen af 19de Aarhundrede foretog med
den gamle theologiske Bygning, var, ikke blot Theolo\-%
gernes, men Theologiens Værk, at den Opløsningsproces,
der ødelagde Theologien, var en Reformation, som Theo\-%
logien med indre Conseqvens og af egen videnskabelig
Drift fuldbyrdede paa sig selv: et stort litterært Sandse\-%
bedrag! Theologiens Tid ligger bagved Holberg; de travle
Tanker, vi møde hos ham, ere et Nedrivningscorps, han
just har samlet og sat i Arbeide. Da saa den Fornuft,
der paa egen Haand havde begyndt sit Arbeide udenfor
det theologiske Facultet, efterhaanden, skjøndt ikke uden
adskillig Modstand og allehaande Bryderier, fik Ansæt\-%
telse i Facultetet, hvad skete saa? Endte det maaskee
med, at Theologerne, efter at have ladet Fornuften i noget
over en Menneskealder fortsætte Opløsningsarbeidet i Fa\-%
% --- p. 252 ---
\opage{252}\setcounter{footnote}{0}%
cultetet selv, kom til mere positive Resultater end Hol\-%
berg? nei, det endte med en Udrensning af Orthodoxiens
gamle Suurdeig indtil sidste Levning; det endte med en
reen, klar, aldeles vandklar Rationalisme.

I Rationalismens aandløse Forstandighed gik Troes\-%
livet tilgrunde. Derfor gjaldt det fra Troens Side om
et Brud, ikke blot med en Verdslighed, hvori en i Tri\-%
vialiteter udtværet Nyttighedslære havde udslettet ethvert
Spor af det Ideale, men først og sidst med en Kirkelig\-%
hed, hvis opbyggelige Foredrag vel kunde udmærke sig
endogsaa ved at anbefale Kirketaarnes Afbenyttelse til
Veirmøller og ret con amore variere Mester Mannegrims
Yndlingsthemaer: Overtro og Vexeldrift og Staldfodring
med Kartofler. Paa disse Vilkaar maatte Indsigelserne
blive hvasse, og selv Dommen over Holberg falde streng
og skarp; thi, det kan ikke negtes: med Holberg træder
Rationalismen op i Danmark.

Det gjaldt ved Begyndelsen af vort Aarhundrede om
noget Mere end at blive lidt aandrig, i en aandløs Tid;
det gjaldt at gribe Troens Banner, bryde med Vantroens
Væsen og gjøre Bruddet saa dybt og ulægeligt, at det
aldrig mere kunde heles. Saa blev der brudt med alle
Tidens Magter, brudt for Troens Skyld i Menighedens
Navn; det kom til Kamp, begeistret Kamp for Livets
Ret, for Aandens Lys, „for Gaaden i det Guddomsord,
som skaber, hvad det nævner, som fylder Dale trindt paa
Jord, og Klipperne udjævner“. Det gjaldt ikke om at
gaae paa Akkord med den formelle Oplysning, smykke
sig med Fornuftens Yndigheder og faae sig en Glorie
ved at anbefale en ny Middelvei; nei, Middelveistiden
var forbi; det gjaldt om et Gjennembrud. Havde det
været Opgaven, paa en sindrig, smidig, for de Dannede
behagelig Maade at indsmugle et Qvantum satis af den
% --- p. 253 ---
\opage{253}\setcounter{footnote}{0}%
gamle Tro uden at støde den nye Fornuft: da skulde en
smagfuld oratorisk Smeltning af meget Subjectivt med
noget Objectivt, lidt Mystik og lidt Philosophie, lidt
Johannes og lidt Jakobi, have været fuldkommen til\-%
strækkelig. Men her handledes om noget Mere; det gjaldt
netop om ikke at glide hen over knudrede Steder, men
trykke paa Ploven og pløie saa dybt, at det spirende
Sædekorn kunde slaae Rod, og Menigheden see sin frugt\-%
bare Grund med de hellige Kilder. Den oplyste Tids\-%
alder meente nu rigtignok, at en Gjenopvækkelse af Fæ\-%
drenes enfoldige Tro var en ligesaa folkelig Umulighed
som videnskabelig Urimelighed; men den oplyste Tids\-%
alder tog feil. Den gamle Tro stod op „som Ax i Vang,
som Mai i Bøgeskove“; den stod virkelig op igjen til
Alles Forundring. En Paaskemorgen, ja hvor ubegribe\-%
ligt, en Paaskemorgen gav den Døde Liv, „en Paaske\-%
morgen kom og tog --- det sorte Kors fra Graven“.

Ved de sidste Ord ville Kjendere høre en Efterklang
af de grundtvigske Toner og see den mægtige Skikkelse
i Baggrunden. Saaledes er Bruddet skeet, væsentlig og
fra Grunden af. Ved Bruddet er Dualismen sat; men
derfor er den personlige Bevidsthedseenhed ingenlunde
brudt. Langtfra! Det er just Personlighedens høieste
intellectuelle Opgave at erkjende det Værd, der maa til\-%
lægges Videnskaben, og dog fastholde de ideale For\-%
dringer, der udgaae fra Troen. De herhen hørende Un\-%
dersøgelser blive at gjennemføre, ikke i nogen Theologie
--- thi Theologien maa bort --- men vel i en Religions\-%
philosophie.

„Men hvad er da Religionsphilosophie? til hvilken
Sphære maa den henregnes? Forsaavidt Religionsphiloso\-%
phien er en Videnskab, maa den naturligviis høre under
Videns Kreds; men forsaavidt den har Religionen til Ind\-%
% --- p. 254 ---
\opage{254}\setcounter{footnote}{0}%
hold, maa den orientere os i Troens Sphære. Ved denne
dobbeltsidige Opgave bestemmes dens dobbeltsidige Cha\-%
rakteer som \emph{Grændsevidenskab}. Den philosopherende
Bevidsthed er, hvad der da forstaaer sig af sig selv, det
subjective Medium, hvori de ueensartede Begreber mødes;
heraf følger, at Bevidsthedseenheden maa forudsætte en
formel, til Grændsebestemmelsen svarende Begrebseenhed,
at, med andre Ord, de absolut ueensartede Begreber maae
i Grændsen selv blive formelt eensartede. Hvad en saa\-%
dan formel Eensartethed har at betyde, kan let paavises.
Vi tale om absolut ueensartede Principer, men Princip
er jo \textit{in abstracto} altid Princip. Vi tale om ueensartede
Erkjendelser, ueensartede Begreber o. s. v.; men idet vi
abstrahere fra det Eiendommelige, bliver Erkjendelse $=$
Erkjendelse, Begreb $=$ Begreb.

Som Grændsevidenskab faaer Religionsphilosophien
sine Troesbegreber fra Troeslivet og sine Vidensbegreber
fra Videnskaben. I sidste Henseende viser det sig, hvad
en logisk Theisme har at betyde; i første Henseende
kommer det an paa religiøs Psychologie. Det nytter
ikke, at Religionsphilosophen vil hjælpe paa Videnskaben
ved Forsikkringer om sin egen Tro, sine egne individu\-%
elle Erfaringer. Her spørges ikke, hvorfra Philosophen
har sine Erfaringer; her spørges om Erfaringernes psy\-%
chologiske Værdi i deres Forhold til Troesprincipet. Den
tilsyneladende Modsigelse, at Religionsphilosophien paa
de opstillede Vilkaar maa blive en Videnskab, der handler
baade om det, der falder \emph{udenfor}, og det, der falder
\emph{indenfor} Videnskaben, er let at hæve. Vil Nogen paa\-%
staae, at der umulig kan gives en Viden om en Troes\-%
erkjendelse, naar Troeserkjendelsen efter sin Natur skal
falde udenfor Viden, lad ham da begynde med Begyn\-%
delsen og paastaae Umuligheden af, at der kan gives en
% --- p. 255 ---
\opage{255}\setcounter{footnote}{0}%
Bevidsthed om Jorden, naar den virkelige Jord efter sin
Natur skal falde udenfor Bevidstheden. En Indvending
som den, at Troesbegreberne ved at reflectere sig i Viden
absolut maae blive Vidensbegreber, har Intet at betyde;
Videnskabens Opgave er netop at vide det Ueensartede
som ueensartet.“ (Den gode Villie. S. 136---37).

\begin{center}\rule{3cm}{0.4pt}\end{center}

\chapter*{XV.}
\addcontentsline{toc}{chapter}{Forelæsning XV}
\markboth{Forelæsning XV}{Forelæsning XV}
\label{ch:xv}
% pp. 256--278
% --- p. 256 ---
\opage{256}\setcounter{footnote}{0}%
% [large initial S]
Saa ofte der tales om det aandelige Liv, dets Vilkaar,
Betingelser og Hindringer, særlig om Striden mellem Tro
og Viden, saa ofte maae vi tænke paa Luthers Ord:
Troen er en urolig Ting. Og hvilken af Kirkens store
Lærere kunde vel ogsaa sige dette med stærkere Efter\-%
tryk end Luther selv? Det var Troen, der gav ham Fred!
Ja, men Troen beredte ham ogsaa Kamp, en Kamp, der
varede Livet ud, en Kamp, først i det Indvortes og der\-%
paa tillige i det Udvortes. I Nattens Mørke, angst for
Verden og for sig selv, tyer den unge Mand til Augu\-%
stinerklosteret i Erfurt, banker paa og indlades. Man
skulde nu mene, at han, der, med et Tordenveir efter sig,
gik ind ad Klosterets Dør, slap fri for Livets Storme,
for Verdens Uro og indgik til Fred; men saaledes var
det ikke. Han gik ikke ind til Fred, men til Strid, den
forfærdeligste Strid, Strid med sig selv og med Gud,
Strid i Samvittigheden, Strid med Tvivlen for at vinde
Troen. Og da han havde vundet Troen, da han erkjendte,
at det ikke var ved pavelige gode Gjerninger, men ved
Troen alene, at et Menneske bliver retfærdiggjort, hvor\-%
ledes charakteriserer han da Troen? Siger han maaskee,
at i Troen er der lutter Fred, og at det kun koster Strid
og Uro, inden man naaer til Troen? nei, om Troen, der
% --- p. 257 ---
\opage{257}\setcounter{footnote}{0}%
giver Fred, just om den Tro, siger han, at den er en
urolig Ting. Og den vedblev at være det for ham. Thi
det var ham ikke nok at have vundet Fred i sit eget
Indre. Hvad der havde givet ham indvortes Fred, det
vilde frem paa hans Læber, vilde forkyndes med hans mæg\-%
tige Røst; saa begyndte Striden i det Udvortes. Da saa
Uveiret kom i det Udvortes, da han for Troens Skyld
fandt sig midt paa Begivenhedernes Strøm, midt i den vilde
Kamp, ombruset af Storme, da maatte han vel sige med
forstærket Eftertryk: Troen er en urolig Ting.

Luthers Ord om Troen som en urolig Ting stad\-%
fæstes i Kirkens Historie fra Først til Sidst. Det er Sag\-%
kyndige bekjendt, hvorledes den gamle Kirke under blo\-%
dige Forfølgelser stred sig frem til Seier. I Farernes Tid,
naar Forfølgelserne ret rasede mod de Christne, da voxte
Troeskraften, da meldte sig Flere og Flere med stedse
stigende Begeistring, ikke blot kraftige Mænd, men Yng\-%
linge, Qvinder og Oldinge, alle med frimodig Bekjendelse,
alle beredte til at offre Livet i den Korsfæstedes Navn.
Havde derimod Kirken i nogen Tid havt Fred, saa mær\-%
kedes Slapheden. Jo længere den Mellemtid havde været,
hvori Troen ligesom var i Begreb med at blive en rolig
Ting, saa at man ogsaa med den kunde indrette sig paa
Hygge og Velvære, desmindre forberedt var man, naar
% a raised mark before "desmindre" is a risen space/quad, not a hyphen; not transcribed
Stormen pludselig brød løs. Under de Decianske og se\-%
nere under de Diocletianske Forfølgelser faldt just Saa\-%
mange fra, fordi Kirken havde havt Fred forlænge.

At Troen er og altid maa være en urolig Ting,
følger ligefrem deraf, at den er det høieste personlige An\-%
liggende og som det høieste personlige Anliggende grun\-%
det i en dualistisk Bevidsthed. Eller skulde maaskee
Modsætningen imellem Verdsligt og Helligt, mellem Kjødet
og Aanden, imellem det gamle og det nye Menneske
nogentid i Verden kunne ophøre at være dualistisk? En
% --- p. 258 ---
\opage{258}\setcounter{footnote}{0}%
anden Sag er det, at Dualismen kan til forskjellige Tider
vise sig forskjellig, og den personlige Strid som Følge
heraf faae et forskjelligt Udseende.

Det gjør, hvad Uroligheden angaaer, en meget be\-%
tydelig Forskjel, om Striden er udadgaaende eller blot
indadgaaende, om den føres paa barbariske, halvbarbariske
eller humane Vilkaar.

Forfølgelsernes Tid var for de første Bekjendere en
barbarisk Tid; under den første umiddelbare Modsætning
mellem Christendom og Hedenskab maatte den afgjørende
Troesprøve nødvendigviis være: at kunne offre Livet for
Christi Navns Bekjendelse; Idealet var at blive Blod\-%
vidne. Saa glødende var Aanden, saa flammende den
hellige Lidenskab, at Higen efter Martyrkronen endog
blev en Slags forfængelig Æressag. Forløserens Ord:
„Frygter ikke for dem, som slaae Legemet ihjel, men
kunne ikke ihjelslaae Sjælen“, fik rigelig Opfyldelse:
Troens unge „Soldater“ gik modigt i Ilden. Men saa
storartede Vidnesbyrd om Aandens Magt og Frihedens
Seier end herved fremkom, saa var Opfyldelsen af den
første elementære Grundfordring dog ikke tilstrækkelig til
alsidig Udvikling af det aandelige Liv i det Hele. En
uophørlig Fortsættelse med Forfølgelser og Blodvidner
skulde sikkerlig have ført, ikke til høiere personlig Selv\-%
forstaaelse, men til vildt Sværmeri og mørk Fanatisme.
Montanisternes og Donatisternes Udskeielser vidne derom.

Det var da for Aandsudviklingens Skyld ligesaa\-%
meget Religionen som Culturen, ligesaavel det nye som
det gamle Menneske til Baade, at Kirken omsider fik
udvortes Fred.

Men den blot udvortes Fred kunde ikke bringe Ro,
thi Troen er en urolig Ting. Den udvortes Modstand
i det ophidsede Hedenskab forvandlede sig nu til en ind\-%
% a raised dot between "i" and "det" is a risen space/quad; not transcribed
vortes Modstand i Kjætterier og Sektvæsen. Havde det
% --- p. 259 ---
\opage{259}\setcounter{footnote}{0}%
under Forfølgelserne været den væsentligste Betingelse
for den evige Salighed at kunne offre Livet for Troen, saa
blev det nu, efter at Forfølgelserne vare ophørte, en uafvise\-%
lig Betingelse at sikkre sig Adskillelsen mellem den sande
og den falske Tro. Herved kom ikke blot Lidenskaberne,
men ogsaa Tankerne i Bevægelse. De heftige dogma\-%
tiske Stridigheder i den græske Kirke fra det fjerde
til det ottende Aarhundrede førtes alle paa halvbarba\-%
riske Vilkaar. Bispernes Skinsyge, Munkenes Fanatisme,
Folkets Lidenskaber og Keiserens Magtsprog giver Ud\-%
viklingens Gang en saa verdslig Retning, at det kirkelige
Lærebegreb synes at maatte blive et simpelt Produkt af
menneskelig Vilkaarlighed; og dog foregaaer der en Idee-%
Udvikling. Idet de mange Stemmer høres, og alle Midler
anvendes, falder Vilkaarligheden bort: gjennem Stridens
Larm bevæger Kirke-Aanden sig med indre Sikkerhed
fremad til conseqvent Udformning af Troens Indhold.

Det er et første Skridt paa Selvtænkningens Vei, at
Kirken faaer sine Troessætninger opstillede og bragte i
System. Men just ved dette første Skridt begik den et
Feiltrin, der bragte den til Affald fra den apostoliske
Tankegang: den kom under Lærebegrebets videre Ud\-%
% a reader's pencilled marginal note beside this line; not transcribed
dannelse ind paa den objective Videns Vei, den fik Theo\-%
logien grundlagt.

Ved Theologien blev Læren udvortes objectiv i samme
Forstand som en positiv Lovgivning kan siges at være
udvortes objectiv; ved den Begunstigelse, Christendommen
fra Constantin den Stores Tid i det Hele nød af Kei\-%
serne, blev Kirkens Stilling i Staten ligeledes udvortes
objectiv: Eftertrykket falder i alle Retninger paa det Ob\-%
jective, det udvortes Objective.

Hovedsagen er i det udvortes Objective, at Kirken
faaer et Lærebegreb af een Gang for alle fastslaaede
Dogmer; thi ved et saadant Lærebegreb er dens ideale
% --- p. 260 ---
\opage{260}\setcounter{footnote}{0}%
Subsistens „igjennem skiftende Tidsfølger“ langt bedre be\-%
trygget end ved Menighedernes „usikkre og fluctuerende,
slet garanterede“ Tilegnelse af Evangeliets Aand efter den
blot apostoliske Troesbekjendelse. Hovedsagen er --- i
det udvortes Objective, „at de kirkelige Indkomster ere
i Kornværdi efter Capitelstaxt, og at Kirken besidder
Jordegodser“; thi derved, at Kirken besidder Jordegodser
„er dens materielle Subsistens gjennem skiftende Tids\-%
følger langt bedre betrygget end ved usikkre og fluctu\-%
erende, slet garanterede Pengesager“. Hovedsagen er,
at det hele Kirkevæsen, aandeligt saavel som timeligt,
kommer til Ro i det Objective. Men Troen er en urolig
Ting.

Hvorfor var det nu netop nødvendigt, at der til den
apostoliske Troesbekjendelse blev føiet den nicæniske om
Sønnens Væsenseenhed med og evige Fødsel af Faderen,
til denne igjen den constantinopolitanske om de tre Hy\-%
postaser i eet Væsen, til denne igjen den chalcedonen\-%
siske om Christi to Naturer, som hverken maatte gaae
sammen eller holde sig adskilte, til denne igjen Bekjen\-%
delsen af Christi to Villier o. s. v., hvorfor, sige vi, var
en saadan Opdyngen af subtile Dogmebestemmelser, hvis
Udtydning og Begrundelse maatte udkræve overmenne\-%
skelig Skarpsindighed, vel nødvendig, uden for at give
Troen et ret indviklet, forviklet og lærd Grundlag.
Hvorfor skulde Troen vel have et saa indviklet, forviklet
og lærd Grundlag? Naturligviis for Objectivitetens Skyld,
at Lærerne, staaende paa dette objectivt faste Grundlag,
gjennem skiftende Tidsfølger kunde med Fynd og
Kraft fortolke de enfoldigt Troende Christi Ord: Jeg
priser dig Fader, at du haver skjult det for de Vise og
Forstandige og aabenbaret det for de Umyndige! Og
hvorfor maatte Kirken nu absolut have sine Indkomster
i Kornværdi efter Capitelstaxten? hvorfor maatte den
% --- p. 261 ---
\opage{261}\setcounter{footnote}{0}%
til sin materielle Subsistens nødvendigviis besidde Jorde\-%
godser (Patrimonium Petri)? Naturligviis for Objectivi\-%
tetens Skyld, at Christi Statholdere gjennem skiftende
Tidsfølger kunde vide sig betryggede! Hvorledes skulde
vel Ordets Forkyndere være istand til at yde de Troende
Trøst mod Bekymringer for de timelige Ting, dersom de ikke
selv følte sig betryggede mod timelige Bekymringer; men
hvorledes skulde de vel føle sig betryggede, hvorledes
skulde de med fuld Fortrøstning tolke Evangeliets Ord:
„Bekymrer Eder ikke for den anden Morgen, thi den
Dag imorgen skal bekymre sig for sit Eget“, dersom de
ikke vidste med sig selv, at de gjennem skiftende Tids\-%
følger fik deres kirkelige Indkomster i Kornværdi efter
Capitelstaxten? „Betragter Lilierne paa Marken!“ Ja,
men naar Kirken ingen Jordegodser har, saa har den jo
ingen Marker, og naar den ingen Marker har, saa har
den jo heller ingen Lilier. Skulde en Evangeliets For\-%
kynder gjennem skiftende Tidsfølger, uden Marker og
uden Lilier, betrygges alene ved usikkre og fluctuerende,
slet garanterede Pengesager, hvor skulde han da med
roligt Sind kunne gaae ud og betragte Lilierne paa
Marken!

Der gives en forskjellig Betragtning af Lilierne paa
Marken. Et er det at betragte dem fra Evangeliets
Standpunkt, et Andet at betragte dem --- efter Capitels\-%
taxten; en er det nye, en anden det gamle Menneskes
Betragtning; de to Betragtninger ere absolut ueensartede;
de absolut ueensartede Betragtninger gjøre Bevidstheden
dualistisk! den dualistiske Bevidsthed gjør Troen til en
urolig Ting.

Men er det da ikke en Civilisationens Lov, at det
menneskelige Samfund skal arbeide sig frem til Orden
og Ro? Vi svare: det er ikke alene en Civilisationens,
men ogsaa en Religionens Lov; kun maa man fra
% --- p. 262 ---
\opage{262}\setcounter{footnote}{0}%
Grunden af gjøre Forskjel imellem Orden og Ro i det
% "af gjøre": printed thus, the second letter of "af" a clear f (full ascender), where
% the sense wants "at gjøre"; wrong sort, transcribed as printed; not in the Rettelser
borgerlige og Orden og Ro i det religiøse Samfund.
Just fordi der i det religiøse Samfund maa herske en
anden Aand end den, der paa menneskelige Vilkaar her\-%
sker i det borgerlige, maa Orden og Ro ogsaa tilveie\-%
bringes ved andre Midler og paa andre Betingelser i det
religiøse end i de borgerlige Anliggender. I det borger\-%
lige Samfund gjælde Tvangslove, i det religiøse Friheds\-%
love. Det er en raa Forvexling af Frihedslove med
Tvangslove, der gjør den middelalderlige Kirke-Aand
halvbarbarisk.

Paa Hierarchiet og Pavemagten skulde den store,
saavel borgerlige som religiøse, Tingenes Orden be\-%
grundes i Occidenten. Alle, Store og Smaa, Læge og
Lærde, vare religiøst og borgerligt interesserede i, at der
op over Forvirringen kunde hæve sig en ordnende sty\-%
rende Magt, et Bestaaende, som var istand til at holde
Uroen nede og paa menneskelig Viis i det Helliges Navn
qvæle Tankernes Oprør; men Troen er en urolig
Ting. I Troens Navn og for Troens Skyld fik Pave\-%
kirken de store Landstrækninger, det meget timelige
Gods med de mange Herligheder, at Christi Rige ret
skulde herliggjøres i det Synlige; men Nydelsen af
disse Herligheder havde for Geistlighedens Tro saa be\-%
tænkelige Følger, at en Arnold af Brescia tog sig det
inderlig nær, og Fratricellerne fortvivlede derover. Uroen
kom desuden fra en anden Side. I Troens Navn og
for Troens Skyld fik Kirken omsider en saadan Over\-%
høihed over Staten, at Innocens III kunde sammenligne
Pavedømmet med Solen og Keiserdømmet med Maanen,
der faaer sit Lys fra Solen. Men den samme Aands\-%
retning, der gav Kirken politisk Magt, fristede ogsaa
Hierarchiet til at misbruge Magten. Da Paven i liden\-%
skabelig Kamp mod Tydsklands middelalderlige Keiser
% --- p. 263 ---
\opage{263}\setcounter{footnote}{0}%
havde slaaet Hohenstauferne til Jorden, kom han som
bekjendt i Afhængighed af det mod en ny Tingenes
Orden politisk fremadstræbende Frankrig, og maatte
som Følge heraf beqvemme sig til at forlade sit evige Rom
og drage til Avignon. Saaledes indtraadte, hvad Kirke\-%
historien kalder Pavens babyloniske Fangenskab eller
Exilet. Følgen af Exilet var det store pavelige Schisma;
Følgen af Schismaet var en uhyre Forargelse. Pave stod
op imod Pave, Pave fordømte Pave; til hvem skulde
Christendommen holde sig? Ja Troen er en urolig Ting.
Uroen gik ikke blot udad i Stridigheder om det Verds\-%
lige, den gik ogsaa indad i Stridigheder om det Aande\-%
lige, i Tankernes indbyrdes Strid om de evige Sandheder.
Videnstankerne fra en classisk, hedensk Cultur skulde for\-%
liges med Troestankerne i de kirkelige Dogmer, den
græske Philosophie skulde i Theologiens Navn brede
Fornuftens Lys over Aabenbaringens Mysterier, Aristo\-%
teles, en Christen før Christus (Christianus ante Christum),
skulde allerede have begrebet Alt, hvad der af Kirke\-%
læren kunde begribes. Saa blev der raisonneret og dis\-%
puteret; i Scholastikens og Dialektikens Høiskoler stred
og fægtede man paa ridderlig Viis som ved en Turne\-%
ring. Under denne fortsatte Feide vare Mange urolige,
Nogle paa Videns, Andre paa Troens Vegne.

Uroen tiltog, da man omsider efter Constantino\-%
pels Erobring under den stadig voxende Interesse for
Videnskaben, som charakteriserer det Afsnit af Italiens
Historie, man kalder Humanisternes Tidsalder, mere og
mere fordybede sig i Classikerne, altsaa ogsaa i Plato og
Aristoteles. Da blev det klart, at her maatte være to
Erkjendelsessystemer, et i Philosophien, det græsk clas\-%
siske, grundet paa Fornuften, et andet i Theologien, det
orthodox-kirkelige, hvilende paa Pavens Autoritet. Disse
To kunde ikke forliges; derfor søgte visse Tænkere, der
% --- p. 264 ---
\opage{264}\setcounter{footnote}{0}%
ikke turde bryde med Kirken, medens de paa den anden
Side heller ikke vilde opgive Philosophien, at undvige den
overhængende Fare ved at forkynde en Lære om to
Sandheder, en philosophisk og en theologisk. Naar de
blot maatte have deres philosophiske Sandheder i Fred,
vare de villige til i alle Maader, hvad angaaer Autoritet
og Tro, at bøie sig under Kirkelæren. Saaledes Pom\-%
ponatius i det 16de Aarhundrede og Vanini i det 17de.
% a mark after "Pomponatius" stands at mid-height, not on the baseline: a risen
% space/quad, not a full stop; not transcribed
Men det gik ikke. Pomponatius opdagede, at den clas\-%
siske Aristoteles havde en ganske anden Lære om Sjælens
Væsen og Udødelighed end den scholastiske. Den scho\-%
lastiske Aristoteles skulde stemme med Kirkelæren og
give Beviis for Sjælens Udødelighed; den classiske der\-%
imod gav netop Beviis for den naturlige Sjæls Døde\-%
lighed. Inqvisitionen kom i Bevægelse, Pomponatius blev
frelst ved mægtige Venners (Cardinalers) Mellemkomst,
men Vanini blev brændt. Man brændte Bekjenderen af
de to Sandheder, og det var i hiint Tidspunkt ganske
naturligt.

Det er begribeligt, at den katholske Kirke ikke kunde
tage tiltakke med hiin Sandhedsdeling; vi for vort Ved\-%
kommende kunne det ikke heller. Skulde Delingen be\-%
tyde Noget, maatte den forstaaes saaledes, at kun Philo\-%
sophien var Videnskab, og Theologien et selvmodsigende
Monstrum; men til et saadant Brud med den dogmatiske
Mirakelvidenskab var Tiden endnu ikke kommen. Hiin
Strid blev da ogsaa snart afløst af en langt mere dybt\-%
gaaende. Det kom saa vidt, at Troessandheden blev
splidagtig med sig selv ved Reformationen. Hvad skeer?
Ud af den gamle Kirkes Skjød fremgaaer en ny Kirke,
og dog døer den gamle Kirke ikke. Saa staaer Kirke
mod Kirke, hver med sin Troeslære. Hvad der er sandt
i Protestantismen, er usandt i Katholicismen og omvendt.
% --- p. 265 ---
\opage{265}\setcounter{footnote}{0}%
Var Katholikens Tro sand, saa var Protestantens usand.
Saa kunde man da see, at Troen er en urolig Ting.

Uroen gik endnu videre. Idet den evangelisk, pro\-%
testantiske Kirke brød med den pavelig katholske, spal\-%
tedes den igjen dualistisk i to store Modsætninger: den
lutherske og den reformeerte Typus med derunder hø\-%
rende Traditioner. Ja, Spaltningen gik saavidt, at den
katholske Kirke endog blev splidagtig med sig selv, idet
den jesuitisk pelagianiserende Retning lagde an paa
at qvæle den augustinske. Kræfterne kom i Bevægelse:
Paastand stod mod Paastand, Lidenskab mod Lidenskab,
Magtsprog mod Magtsprog; men Troesprincipet selv blev
ikke klaret, Troestanken ikke udviklet.

Den evangelisk protestantiske Grundtanke, saaledes som
den blev opfattet og fastholdt af Luther selv, skulde vist\-%
nok være bleven gjennemført langt conseqventere, og Stri\-%
den saavel med de Pavelige som med de Reformerte have ført
til et andet Resultat, dersom ikke den theologiske Fixidee, at
Dogmerne dog, idetmindste til en vis Grad, burde forklares og
begrundes videnskabeligt, var traadt imellem. Vi see det af
Nadverstridighederne. Imod den katholske Transsubstan\-%
tiationslære satte Luther Indstiftelsens Ord. Men saa træder
Carlstadt op og fordreier Indstiftelsesordene, fordi han ikke
kan begribe, hvorledes det er muligt, at Christi Legeme og
Blod kan være i Brød og Viin, medens han derimod vel kan
begribe, at Christus, pegende paa sit Legeme, kan have
sagt: dette er mit Legeme, og derhos, pegende paa Brødet,
tilføiet: tager dette --- Brødet, som ikke er mit Legeme
--- hen o. s. v. Hvad Luther hertil vilde svare, var
allerede udtalt i hans Skrift om Kirkens babyloniske
Fangenskab, hvori han siger: „Sandelig, om jeg ikke
kan begribe, hvorledes Brødet er Christi Legeme, saa
vil jeg dog tage min Forstand fangen og holde mig en\-%
foldig og alene til hans Ord.“ Hvad han tilføier til Op\-%
% --- p. 266 ---
\opage{266}\setcounter{footnote}{0}%
klaring: „Naar Jernet og Ilden, to Substanser, ere saa\-%
ledes forbundne i det ene, glødende Jern, at enhver Deel
er Jern og Ild, hvi formaaer da ikke meget mere det
almægtige Gudmenneskes, den forklarede Herres, Legeme
at være i det indviede Brøds Substans?“ udgives visselig
ikke for nogen dogmatisk Begrundelse. „Ingen Fornuft“,
siger Luther i sit „Manifest mod de himmelske Propheter“,
„er saa ringe, at den ikke heller skulde troe, at slet og
ret Brød og Viin her ere tilstede, end at Christi Kjød
og Blod her er skjult. Dertil behøver man ingen Aand,
det er let at troe for Enhver, og her behøves ikke Andet,
end at En, som har en Smule Anseelse, er saa dristig
at prædike det, saa har han allerede nok af Disciple.
Men naar man vil omgaaes saaledes med vor Tro, at vi
først lægge vor egen forudfattede Mening ind i Skriften,
og derpaa dreie den efter vort Sind og alene see hen til,
hvad der er ligefrem for den almindelige formørkede For\-%
stand, saa vil ingen Troesartikel blive staaende; thi der
findes ikke een i Skriften, som ikke af Gud er stillet
over Fornuften.“ --- Saa er her altsaa, ifølge Luthers
udtrykkelige Vidnesbyrd, Dualisme af Tro og Fornuft:
Troen er en urolig Ting.

Luther havde Ret, han følte det dybt; det var denne
Følelse, der gjorde ham saa stærk i Bekjendelsen, saa
bitter i Striden. Ja visselig havde Luther Ret; den
moderne Christenhed er sig neppe bevidst, hvad Kirken
skylder den Mand, der i dette Punkt stod saa haardt
paa Textens Ord, og trodsede en Verden med Troens
Evangelium. Ja Luther havde Ret, dog ikke som Theo\-%
log, men som Præst. Hans Opfattelse af Nadveren for\-%
klarer sig just deraf, at han oprindelig saae Sakramentet,
ikke med Theologens, men med Præstens Øie. Hvad
Luther som Theolog har fremsat om Nadverlæren, er
kun betydningsfuldt for hans Tid; hans Theologumena
% --- p. 267 ---
\opage{267}\setcounter{footnote}{0}%
ere forsvindende, men hans præstelige Vidnesbyrd ville
forblive, saalænge Troen er --- Tro.

Luther selv var ikke klar over det principielle For\-%
hold imellem Religionen og Theologien. Den nye Lære
trængte til Videnskabens Bistand; Melanchton var Theo\-%
log, og Theologen meente saa, at Calvin havde Ret, thi
Calvin var jo ogsaa Theolog. Men den lutherske Kirke,
der ubetinget havde Ret i Troen, vilde da nu ogsaa med
det Samme have ubetinget Ret i Theologien, saa gik
det jo længere jo værre. Den lutherske Theologie ar\-%
beidede sig træt i uspeculative Speculationer, og har alle\-%
rede længe maattet see sine Yndlingstheorier om \textit{com\-%
munication idiomatum og ubiquitas corporis Christi} givne
% "communication": printed thus (for "communicatio"); transcribed as printed
til Priis for opløsende Kritik. (Evangelietr. og Theolog.
S. 117---118.)

Allerede i Slutningen af det 17de Aarhundrede havde
der dannet sig \textit{to aldeles modsatte Anskuelser; ifølge den
ene af disse skulde man reentud opgive Forbindelsen af Tro
og Fornuft, Religion og Videnskab, ifølge den anden skulde
man ved alle mulige videnskabelige Midler forsvare og
hævde den}. Den førstnævnte Synsmaade gik igjen i to
Retninger, idet man paa den ene Side indskrænkede sig
til at operere videnskabeligt ved at paavise Modsigelser
og vække Tvivl, paa den anden Side opfattede Religionen
praktisk som Liv og stolede paa Følelsen: Repræsen\-%
tanten for den første Retning er Bayle, for den sidste
Spener.

I Bayles historiske og kritiske Ordbog (Dictionaire
historique et critique) er Striden mellem Tro og Viden,
% "Dictionaire" (one n) as printed; roman
Religion og Videnskab med stor Skarpsindighed paaviist.
Dogmet vil troes, men ikke begribes, \emph{det er og maa
være den menneskelige Fornuft et Mysterium}.
„Aabenbaringens Hemmeligheder tilhøre et overnatur\-%
ligt Rige, de støtte sig paa den høieste Autoritet,
% --- p. 268 ---
\opage{268}\setcounter{footnote}{0}%
Guds Autoritet, og han har ikke forelagt os dem, for at
vi skulle begribe dem, men for at vi skulle troe dem
med al den ydmyge Underkastelse, som vi ere det ufeil\-%
bare høieste Væsen skyldige. Aabenbaringens og Philo\-%
sophiens Uforenelighed gjør det nødvendigt at vælge
mellem begge og at decidere sig for en af dem. Den,
som ikke vil troe Andet end det, som er klart og over\-%
eensstemmende med de almindelige Begreber, han gribe
Philosophien og lade Christendommen fare; thi det er
umuligt paa een Gang at besidde Klarheden og Ubegri\-%
beligheden; Foreningen af disse to Dele er ikke mindre
ugjørlig, end det er at finde Cirklens Qvadratur. Den,
for hvem et rundt Bord ikke er beqvemt, han lade sig
gjøre et fiirkantet deraf; men han maa ikke forlange, at
eet og samme Bord skal yde ham baade et rundt og et
fiirkantet Bords Beqvemmeligheder.“ Bayle taler med
lige Færdighed snart Troens, snart Fornuftens Sag. Til
Fordeel for Troen anfører han: „For ikke at gribe blindt
til ved Valget maa man gjøre sig nøie bekjendt med
begge Siders, Troens og Videnskabens Natur. Ved første
Øiekast skulde det, naar man hører, at Troen er en
Vished, hvis Gjenstand forbliver dunkel, men at Viden\-%
skaben derimod, foruden det at den giver Vished, ogsaa
gjør det, som vides, klart og tydeligt, synes, som om det
afgjorte Fortrin var paa Videnskabens Side. Men Viden\-%
skaben kommer istand ved Fornuften, og hvad om For\-%
nuften nu ikke skulde være en saadan Opgave voxen?
Men dette er i Virkeligheden Tilfældet; Fornuften er
kun et Ødelæggelsens, ikke et Opbyggelsens Princip, den
veed bedre, hvad Tingene ikke ere, end hvad de ere, den
er kun skikket til at vække Tvivl og til at snoe sig til
Høire og Venstre for at udstrække Striden i en Uende\-%
lighed. Den er en Landløberske, som ikke kan komme
til Ro, en Penelope, der bestandig igjen opløser sin egen
% --- p. 269 ---
\opage{269}\setcounter{footnote}{0}%
Væv. Den duer kun til at bringe Menneskene det Mørke,
hvori de famle, deres Uduelighed til at arbeide sig ud
deraf, og Trangen til en Aabenbaring, til Bevidsthed.“
For Fornuften anfører han: „Der gives Grundsætninger, mod
hvilke nok saa udtrykkeligt og bestemt udtalte Lærdomme
i den hellige Skrift ikke vilde kunne udrette Noget, som
f. Ex. de Grundsætninger, at det Hele er større end en
Deel af det, at to hinanden modsatte Sætninger umulig
begge kunne være sande o. dsl. Dette tilstaae ogsaa
de rettroende Theologer af alle Confessioner, idet de ikke
indrømme, at Treenigheden og Menneskevordelsen ere
Dogmer, der indeholde en Selvmodsigelse; hvorved de
i Gjerningen anerkjende Fornuften som høieste Instans,
idet de ikke troe, at et Dogme er tilstrækkelig sikkert,
naar det ikke saa at sige er indregistreret og lovformelig
stadfæstet ved Fornuftens Høiesteret. Gud vilde skjænke
% "til at": a faint raised mark between the words is a risen space; not transcribed.
Sjælen et ufeilbart Middel til at skjelne det Sande fra
det Falske, og dette Middel er det naturlige Lys, de
metaphysiske Principer, i hvilke vi besidde en oprindelig
Regel, hvorefter vi kunne skjelne, om de forskjellige Lær\-%
domme, som vi finde i Bøger eller høre foredrages af
vore Lærere, ere rigtige eller urigtige. Det eneste Kri\-%
terium for Sandheden have vi derfor i Gjenstandenes
Overeensstemmelse med dette almindelige og oprindelige
Lys, som af Gud er udgydt i alle Menneskers Sjæle.“

Man skulde nu formode, at slige dybt indgribende
Forhandlinger især maatte have tiltrukket sig Theolo\-%
gernes Opmærksomhed, da det jo aabenbart er Theologiens
videnskabelige Anseelse, der her staaer paa Spil. Dette
var nu vel ogsaa for en Deel Tilfældet i andre Lande,
navnlig i Holland, hvor en Clericus optraadte mod Bayle;
men det var ingenlunde Tilfældet i Danmark. I Dan\-%
mark havde Theologerne jo nok i det Daglige; hvad
skulde de vel med almeenvidenskabelige Principunder\-%
% --- p. 270 ---
\opage{270}\setcounter{footnote}{0}%
søgelser? Bayles Skriverier var jo et ugudeligt Væsen,
noget philosophisk Tøieri, der ikke kom gudfrygtige
Mænd ved. Vistnok havde de en uhyggelig Følelse af,
at Tiden ikke var dem gunstig, at her var Noget i Veien;
men hvori det stak, begreb de ikke. Det maatte være
Verdens Ondskab, Slægtens og Tidens Vantro, Fornuf\-%
tens Hovmod; saa bleve de let enige om at betragte de
mange uhyggelige Phænomener, hvor forskjelligartede
disse end iøvrigt vare, som Aabenbarelser af eet og
samme ugudelige Væsen. Hvad der ikke stemte med
deres theologiske Anskuelser, blev under de allerforskjel\-%
ligste Benævnelser slaaet sammen i een eneste stor Kjæt\-%
terbunke. Naturalister, Indifferentister, Synkretister, Athe\-%
ister, Pietister --- Pontoppidan tør af lutter Gudfrygtighed
end ikke skrive de gudsbespottelige Navne heelt ud, men
% "A . . .ister og P . .ister": the abbreviating points as printed (three after A, two after P,
% the last point close against "ister").
forkorter dem til A\,.\,.\,.ister og P\,.\,.ister --- Fanatikere, En\-%
thusiaster, Deister og Rationalister ryge En om Ørene,
naar man giver sig til at røre i Bunken. Altsaa Theo\-%
logerne gik i den gamle Trædemølle, og den gamle Træde\-%
mølle gik, som den kunde. (Om Holbergs Kirkehistorie
og Theol. S. 41---43.)

Dersom Vanskelighederne kunde bortveires derved,
at Theologerne lukkede Øinene til, for ikke at see dem,
da vilde de forlængst være bortveirede; saaledes mener
vel Agerhønen, at den ved at stikke Hovedet i et Stage\-%
hul bliver usynlig for Jægeren; men Kjendsgjerning er
altid Kjendsgjerning. Har Theologien ikke sin Modstand
udenfor sig, saa har den den inden i sig. Den gamle
Orthodoxie holdt paa Troen ved at stride mod Fornuften;
det var Dualisme. Den theologiske Rationalisme holdt
paa Fornuften ved at stride imod Troen; det var ogsaa Dua\-%
lisme. Dualismen gaaer igjen under de allerforskjelligste
Skikkelser; selv Adskillelsen mellem Tro og Tro, mellem
% --- p. 271 ---
\opage{271}\setcounter{footnote}{0}%
Aabenbaringstro og Fornufttro er jo dualistisk. Vi see
det hos Strausz.

% The quotation opened here („Uagtet …) closes on p. 273 after "uforenelige."; the
% closing mark after "videre." on this page closes the inner quotation of Strauss.
„Uagtet Strausz ved at antage en Fornufttro endnu
ikke er ganske paa det Rene med Forholdet mellem Ra\-%
tionalisme og Philosophie, mellem Fornuftreligion og
Videnskab, saa er det dog klart, at hans Fornufttro kun
passer for hvad han kalder de Vidende, ikke for de
Troende. „Her er“, siger han, „en Kløft befæstet mellem
to Klasser af det menneskelige Selskab, de Vidende og
Folket; d. e. de Ikke-Philosopherende saavel i de høiere
som i de lavere Stænder, og den vil maaskee aldrig
fyldes\,... Saa lade den Troende da den Vidende, ligesom
denne hiin, roligen vandre sin Vei, vi lade dem beholde
deres Tro, saa maae de ogsaa lade os beholde vor Philo\-%
sophie; og hvis det skulde lykkes de overvættes Fromme
at udelukke os fra deres Kirke, saa ville vi agte det for
en Vinding; af falske Foreningsforsøg er der nu skeet
nok; kun Skærpelse af Modsætningerne kan føre videre.“
Hermed er Dualismen, den af Bayle paapegede Dua\-%
lisme mellem Religion og Videnskab, og med den Theo\-%
logiens Ugyldighed, udtalt som Resultat af vort Aar\-%
hundredes videnskabelige Prøvelse.

Og dermed, sige maaskee vore lærde Theologer, skal
det saa være beviist, at Theologiens Dødsdom uigjenkal\-%
delig er fældet. Hvo seer nu ikke, at Sammenstillingen
mellem Bayle og Strausz netop beviser det Modsatte.
Bayle og med ham mangfoldige Fritænkere vare i deres
Tid ligesaa sikkre paa den saakaldte Dualisme, d. e. paa
den uoverstigelige Kløft mellem Tro og Viden, som Strausz
med mangfoldige andre, baade skjulte og aabenbare, Fri\-%
tænkere kunne være det i vor Tid. Desuagtet skete det
dog, at Tro og Viden under en bestandig fremadskridende
theologisk Videnskab vedbleve at arbeide sammen et godt
Aarhundrede efter Bayle: hvad skulde nu, spørge vi, være
% --- p. 272 ---
\opage{272}\setcounter{footnote}{0}%
til Hinder for, at Tro og Viden, naar Døgnets paradoxe
Ilinger med „Strøtanker, Aphorismer, Indfald og Glimt“
ere dragne forbi, og det igjen klarer op paa Videnska\-%
bernes Himmel, kunne fremdeles arbeide sammen paa nye
Vilkaar i et nyt Aarhundrede og give Theologien en ny
Skikkelse? Erfaringen lærer jo dog, at man ikke skal
forhaste sig med at erklære et videnskabeligt Standpunkt,
endnu mindre en Videnskab, der har Aarhundreders Hævd,
% The two Latin tags on this page are in TRUE ITALIC (Horace, Ars poetica 70--71; and the
% maxim "qui nimium probat"); "quæ" printed with the æ ligature; "etc." is roman.
for et absolut Overvundet: \textit{multa renascuntur, quæ jam
cecidere, caduntque, quæ nunc sunt} etc. --- Betragtninger
af denne Natur minde stærkt om det bekjendte: \textit{qui ni\-%
mium probat, nihil probat}. Ved at anføre begrændsede
Sandheder i grændseløs Almindelighed ender man enten
i det Meningsløse eller i det Trivielle. Forsaavidt det
gaaer op og ned med alle Ting i Verden, maa det na\-%
turligviis gaae op og ned med Theologien. --- Theologiske
Trøstegrunde af denne Art minde stærkt om Folkets kloge
Nar, der altid loe, naar det gik opad Bakke, fordi der
efter en Modgang opad maatte følge en Medgang nedad,
eller, maaskee endnu bedre, om den drukne Poet, som
da Alting løb rundt, blev staaende med Nøglen og ven\-%
tede paa, at Døren nok skulde komme til ham, uden at
han, i sin Ufeilbarhed, behøvede at skifte Standpunkt.
Det er længe gaaet med Theologien op og ned, ned og
op som i en Gynge; men just fordi Theologien har siddet
saalænge i Gyngen, er det trods de mange Afvexlinger
i Bevægelser og Svingninger dog nu omsider blevet mu\-%
ligt at tegne dens Bane og bestemme Curvens Ligning.

Sammenligne vi Dualismen af Religion og Videnskab,
saaledes som den optræder hos Bayle, med Dualismen,
saaledes som den er kommen til Gjennembrud i den nyere
Tid, da er det klart, at den sidste er mere end en simpel
Gjentagelse af den første. Naar Bayle stillede Religion
og Philosophie imod hinanden, tog han uden videre Re\-%
% --- p. 273 ---
\opage{273}\setcounter{footnote}{0}%
ligionen, som den dengang existerede i de kirkelige Be\-%
kjendelsesskrifter, i theologisk positiv bestemt Form. Det
var paa hans Tid endnu ikke blevet Tænkeren klart, hvad
der vilde skee, naar man brød de krystalliserede Former,
% "udsmeltede · det", "antog · Begrebets", "støbte · Indholdet": faint raised marks between
% the words are risen spaces; not transcribed.
udsmeltede det Religiøse og støbte Indholdet om i en
anden Form; det var, med andre Ord, endnu ikke gaaet
op for den videnskabelige Bevidsthed, hvorledes Sagen
vilde see ud, naar Troens Indhold antog Begrebets Form.
Til at udføre Arbeiderne ved dette store Smeltnings- og
Omstøbningsforsøg krævedes ikke mindre end et heelt
Aarhundrede, Tidsrummet mellem Bayle og Hegel. He\-%
gels Religionsphilosophie er Rationalismens dybeste Con\-%
seqvens, dens Forvandling til egentlig Videnskab. Men
neppe var Forvandlingen gjennemført, før det stod klart
for den Vidende som for den Troende, at Tro ikke er
Viden, Religion ikke Videnskab. Det sidste Skjær af
Fornufttro, et Efterskin af den gamle Rationalisme, som
Strausz i sin Troeslære endnu lod blive tilbage, er svundet
i Videns Navn hos L. Feuerbach, i Troens Navn hos
S. Kierkegaard. Kierkegaard og Feuerbach, saa diame\-%
tralt modsatte de end ere, ere begge enige i at erklære
Religionens og Videnskabens Charakterer for uforenelige.“
--- Om end Viden kan komme til Ro i Ideen, saa kan
Troen dog ikke komme til Ro i Viden: Troen er en
urolig Ting.

Men --- spørger man --- vil ikke hele den kirke\-%
lige Lærebygning falde sammen; ville ikke de religiøse
Forestillinger komme i Gjæring og megen Forvirring;
vil ikke det for Præsternes grundigere Dannelse, høiere
Aandsudvikling og friere Overblik over Folkets religiøse
og statsborgerlige Vilkaar nødvendige Grundlag under\-%
graves, naar den videnskabelige Forskole lukkes, naar
der, med andre Ord, kastes Vrag paa Theologien? Nei,
lad os tænke paa en Reform af det theologiske Studium,
% --- p. 274 ---
\opage{274}\setcounter{footnote}{0}%
hvis en saadan for Tiden virkelig behøves, det er der
Mening i; men at ville afskaffe Theologien som Theo\-%
logie, det er et forrykt Indfald, et Project, der maaskee
kan finde Indgang hos dem, der ville Forstyrrelse og Uro,
men ikke hos dem, der ville Fred og Orden.

For at imødegaae slige Indvendinger, maae vi først
og sidst være enige med os selv om, hvad Theologie er,
eller rettere, giver sig ud for at være: at Theologie ud\-%
trykkelig giver sig ud for at være en Videnskab, en lige\-%
frem Troesvidenskab, en Videnskab om det Overnatur\-%
lige, det Overvidenskabelige, Miraklet, at den altsaa er en
sig selv modsigende Mirakelvidenskab; at det i den theolo\-%
giske Videnskab Selvmodsigende kun tilhylles og skjules
derved, at den har et Apparat af historiske, philologiske,
dogmatiske, metaphysiske, halvmetaphysiske Bestand\-%
dele til Afbenyttelse og kritisk Bearbeidelse, forsaavidt
der kan være Tale om objectiv Afbenyttelse efter subjec\-%
% "subjectivt · Skjøn": a raised mark between the words is a risen space; not transcribed.
tivt Skjøn og kritisk Bearbeidelse uden Princip. Dette
er Modsigelsen, en ældgammel, indgroet Grundmodsigelse,
der fuldstændig maa blottes og fremdrages, naar vi skulle
forstaae tilgavns, hvad Theologie er.

„Uagtet Philologie, Kritik og Historie, som i og for
sig høre under Videnskaben, udgjøre Bestanddele af Theo\-%
logien, udgjøre de dog ingenlunde det i Theologien Theo\-%
logiske. Forsaavidt de nævnte Discipliner behandles vi\-%
denskabeligt, behandles de utheologisk; forsaavidt de
derimod behandles theologisk, behandles de uvidenskabe\-%
ligt. Det i Theologien Theologiske kjendes nemlig der\-%
paa, at den Aarsag, hvoraf alle Virkninger i Naturens
som i Aandens Rige deels umiddelbart, deels middelbart
afledes og forklares, ikke søges i Tingenes Natur eller
Fornuftens Væsen, men i den absolute Villie, i den Al\-%
mægtige, der skaber af Intet. Men selv den reneste,
skarpsindigste og conseqventeste Anvendelse af et saadant
% --- p. 275 ---
\opage{275}\setcounter{footnote}{0}%
Princip (Occasionalismen; Malebranche) gjør Videnskaben
uvidenskabelig. --- Naar der i Videnskabslæren handles
om et Princip, kan Sagen ikke afgjøres ved et Mere eller
Mindre, som om En ved at følge „Middelveien“ vilde
mene, at man endda nok kunde komme videnskabeligt
igjennem ved at antage lidt af det Overnaturlige, naar
man kun ikke gjorde for meget af det. Spørgsmaalet
om et Princip er et Spørgsmaal om en ubetinget, Alt
betingende Forudsætning; et saadant Spørgsmaal maa
derfor altid besvares med Ja eller Nei. Enten maa da
Theologien een Gang for alle vedkjende sig den over\-%
naturlige Causalitet, og har saa med det Samme bortstødt
det objectivt rationelle Grundlag, eller ogsaa een Gang
for alle vedkjende sig dette Grundlag, og har saa med
det Samme fraskrevet sig ethvert Tilhold i det Over\-%
naturlige. I første Tilfælde vil Theologien altsaa ophøre
at være Videnskab, i sidste Tilfælde at være Theologie.“
(Jvfr. Philosoph. Propædeutik. Kbhvn. 1857. S. 70---80).

Et Blik paa Nutidens aandelige Vilkaar lærer, at Theo\-%
logernes mangefarvede Dogmatiker ikke længere ere istand
til at magte de i stedse voxende Antal og med fornyet
Styrke indbrydende Modsigelser eller give den religiøse
Folkeunderviisning noget fast almeengjældende Grundlag.
En Revision af Troeslæren i dens Heelhed, en Udvikling
og Fremstilling af Christendommens Væsen, der kunde
vise, hvorledes de personlige Sandheder i indbyrdes Sam\-%
menhæng fremgaae af eet eneste, i sig selvgyldigt Prin\-%
cip, maa vel, hvad den religiøse Tænkning angaaer, siges
at være fornøden. Men skal en saadan Revision gjennem\-%
føres, da maa Theologien bort; thi Theologien forqvakler
Principet. Skal en yngre Slægt af folkelige Religions\-%
lærere finde den rette Tone, ramme det rette Maal og
% "aandige" (not "aandelige") as printed; a good word, not a defect.
med den aandige Overlegenhed, som kun Sikkerhed i
Tænkning kan give Ordets Forkyndere, træde op mod
% --- p. 276 ---
\opage{276}\setcounter{footnote}{0}%
% This page carries many faint raised marks between words (risen spaces/quads,
% e.g. "Løsning · udtrykkelig", "er · en", "at · gjennemtænke"); none transcribed.
en reflecterende Culturbevidsthed: da er en strengere, paa
Livsopgavernes Løsning udtrykkelig beregnet Forskole
aldeles nødvendig. Paa theologiske Vilkaar er en saa\-%
dan Forskole umulig: skulle Præsterne komme frem, maae
Theologerne bort; thi Theologien forqvakler Principet.

Med Dualismen af Tro og Viden har det ingen Fare,
naar de, hvis Kald det er at forkynde Troens Ord, først
have lært at tænke Troens Tanker. Men at tænke Troens
Tanker er ikke muligt uden væsentlig Indsigt i de tilsva\-%
rende Videnstanker. Have de Vidende end hidtil forsømt
at gjennemtænke Troesprincipet, saa maae Troens Tals\-%
mænd for al Ting see, at de ikke til Gjengjæld forsømme
at gjennemtænke Vidensprincipet. Først naar det kommer
saavidt, at den troende Tænker med Sandhed kan sige
til den Vidende: Jeg forstaaer nok Dig, men Du for\-%
staaer endnu ikke mig, da først vil den Overlegenhed
være naaet, som ad Reflexionens Vei kan naaes til Bedste
for Troen. Vil saa Troen med det Samme være sikkret?
Nei, thi Troen er mere end en Tænkning, dens Uro er
mere end en Tankernes Uro, den er ogsaa en Hjerternes
Uro: Troen er en urolig Ting.

Troens Uro hidrører imidlertid ikke fra dens Princip,
men fra de troende Individualiteters naturlige Skrøbelighed.
Naar selv en Paulus (Philip. III, 13---14) maa sige til Philip\-%
penserne: „Brødre, jeg agter mig ikke at have grebet
det. Men Eet gjør jeg: forglemmende, hvad der er bag\-%
ved, og rækkende efter det, som er foran, iler jeg mod
Maalet,“ saa taler han i „Frygt og Bæven“ om Troens
Uro og Personlighedens Kamp for det hellige Ideal; men
naar han saa, saaledes som i Brevet til Romerne (VIII,
37---39), tilføier: „Men i alt dette mere end seire vi ved
ham, som os elskede. Thi jeg er vis paa, at hverken
Død, ei heller Liv, ei heller Engle, ei heller Fyrsten\-%
dømmer, ei heller Magter, ei heller det Nærværende, ei
% --- p. 277 ---
\opage{277}\setcounter{footnote}{0}%
% Faint raised marks between words on this page (e.g. "fra · Guds", "Sind · finder")
% are risen spaces; not transcribed.
heller det Tilkommende, ei heller det Høie, ei heller det
Dybe, ei heller nogen anden Skabning skal kunne skille
os fra Guds Kjærlighed i Christo Jesu, vor Herre,“ ---
da taler han om den Forvisning, hvori Uroen hæves, den
Inderlighedens Grund, hvori det troende Sind finder
Hvile. Til ingen Tid i det Timelige kan Fordoblingen
% "forsvinder men": no punctuation is printed after "forsvinder" at the line end (sense
% wants a comma or semicolon); transcribed as printed.
af Ro og Uro forsvinde, uden at Troen selv forsvinder
men den sig uophørlig fortsættende Strid kan føres paa
stedse humanere Vilkaar.

Fra mange Sider og paa mange Maader ere Troes\-%
kræfterne prøvede; lige fra den første Christenforfølgelse
indtil de sidste jesuitisk reactionære Forsøg paa at ud\-%
rydde de sidste Spirer af evangelisk Frihed er Prøven
bestaaet, ofte heroisk og under de mest barbariske Mis\-%
handlinger; men først i den nærværende Tidsalder, Tids\-%
alderen med de humane Betingelser, er Aandsudviklingen
naaet saavidt, at ikke blot Troeskræfterne, men, hvad der
i Forhold til vort Culturliv er afgjørende, ogsaa Troes\-%
tankerne kunne fra Grunden af underkastes en Prøvelse.

Vil Nogen maaskee under Paaskud af, at Christen\-%
% "et Liv" and "en Theorie" letterspaced (including the short words "et", "en"; "ikke" is not).
dommen er \emph{et Liv}, ikke \emph{en Theorie}, afvise den in\-%
tellectuelle Prøvelse som en upraktisk, spidsfindig Skole\-%
strid, slaae om sig med Talemaader og forsikkre, at der
ikke kommer Noget ud af at strides om Troen, at man
strax skal erklære sig for Troende og stride for Troen
o. s. v.: da veed jeg ikke at slutte paa en mere passende
Maade end ved at minde om en Udtalelse, hvormed jeg for
% Title letterspaced: "Evangelietroen" and "Theologien"; "og" between them is not spaced.
sytten Aar siden indledede mit Skrift: \emph{Evangelietroen}
og \emph{Theologien}: „Hvis hine Dage, da Spørgsmaalet om
det salige Liv var Slægtens høieste Spørgsmaal, umiddel\-%
bart vendte tilbage; hvis Kirkens store Jubelaar nu var
saa nær, at det sang for vore Øren, som hørte vi allerede
Klangen af de Klokker, der skulle ringe det ind: ja da
kunde vore Opvakte jo sagtens afvise „en blot theoretisk
% --- p. 278 ---
\opage{278}\setcounter{footnote}{0}%
Skolestrid“, og give dem Anviisning paa Livets „Ild\-%
prøve“, som --- „istedenfor at gjøre derefter!“ --- kun
skrive derom, forsaavidt de skrive om Troen. „Men“
--- tilføier en christelig Psycholog --- „mon Troen findes
paa Jorden?“ .\,.\,. Skulde det omsider komme saavidt,
% "strides om" and "stride for" letterspaced (including "om", "for"); "Troen" is not.
at man ikke engang fandt det Umagen værd at \emph{strides
om} Troen, mon det da ogsaa var et godt Tegn? mon det
ikke meget mere var et Tegn paa, at man foragtede den
menneskelige Strid, fordi man i Grunden heller ikke ag\-%
tede den guddommelige; et Tegn paa at man endnu mindre
fandt det Umagen værd at \emph{stride for} Troen: et Tegn
paa, at Troen ikke mere fandtes paa Jorden?“

Saa faaer det da, naar Talen er om Nutidens aande\-%
lige Vilkaar, at blive ved Luthers Ord: Troen er en
urolig Ting!

\begin{center}\rule{3cm}{0.4pt}\end{center}

\chapter*{XVI.}
\addcontentsline{toc}{chapter}{Forelæsning XVI}
\markboth{Forelæsning XVI}{Forelæsning XVI}
\label{ch:xvi}
% pp. 279--295
% --- p. 279 ---
\opage{279}\setcounter{footnote}{0}%
% [large initial D]
Det er „en nærmere videnskabelig Samvirken mellem de
nordiske Universiteter“, jeg ved de nu holdte Foredrag
% a comma-like speck before "nordiske" at the line start is a risen space/quad; not transcribed
har havt den Ære at indlede. Lægger man imidlertid
særlig Vægt paa, at det er en Samvirken mellem Uni\-%
versiteterne, hvorom Talen er, og at denne skal være vi\-%
denskabelig, reiser der sig en naturlig Indvending. Hvor\-%
for, vil man spørge, just hvor det gjælder et Bidrag til
en videnskabelig Samvirken, vælge en Fremstillingsform,
der ikke er videnskabelig? Det er, som det synes, en
Modsigelse at ville indlede noget Videnskabeligt paa en
ikke videnskabelig Maade. Jeg skylder til Slutning den
høitærede Forsamling Regnskab for mit Valg.

Skal det betyde Noget med en Samvirken mellem
de forskjellige Universiteter, skulle Lærerne fra disse Uni\-%
versiteter kunne gjøre sig Haab om at give Bidrag til
en saadan Samvirken, da maae, saavidt jeg skjønner, to
Ting haves for Øie. Det er for det Første naturligt og
nødvendigt, at enhver Universitetslærer vælger sit Emne
af sit Fag og ifølge sin videnskabelige Stilling. Det er
naturligt, at Juristen vælger et juridisk, Physikeren et
physisk, Historikeren og Philologen et historisk og phi\-%
lologisk Emne o. s. v. Dette er mit Første. Det Næste
er, at ingen hospiterende Lærer vil vælge sit Emne saa\-%
% --- p. 280 ---
\opage{280}\setcounter{footnote}{0}%
ledes, at han meddeler Noget, der maa forudsættes at
være Alle ved det andet Universitet bekjendt. Thi det
kan dog i Sandhed ikke kaldes at indlede nogen høiere
Samvirken, at en Docent fra et andet Universitet kom\-%
mer og meddeler de Studerende, hvad de allerede
vide. Jeg maatte saaledes ansee det for mindre passende
at foredrage de norske Studenter et Overblik over et eller
andet almindeligt Parti af Philosophien eller et Brudstykke
af Propædeutiken, kort sagt Noget, som de Studerende
ifølge den Grundighed, Reenhed og Strenghed, hvor\-%
med den philosophiske Videnskab dyrkes her, kunde
høre ligesaa godt hos de beskikkede Lærere. Nei, at
byde de opvakte unge Mænd, hvis levende Interesse
for en friere Aandsudvikling særlig har sat den her om\-%
handlede Sag i Bevægelse, noget Saadant, vilde jeg ikke
blot betragte som et Tegn paa Illusion for mit eget Ved\-%
kommende, men som en Fornærmelse.

Skulde jeg have valgt et philosophisk Emne, maatte
jeg have valgt det fra et andet Stade og paa en saadan
Maade, at det for mig Eiendommelige deri kom til Udtalelse,
og et saadant Emne kunde jeg let have valgt; thi saavidt
jeg veed, afvige mine Bestræbelser, jeg troer paa en tem\-%
melig afgjørende Maade, fra Philosophernes i de forskjellige
Litteraturer, naar Talen er om Methode og særlig om Vexel\-%
virken mellem Philosophien og de forskende Videnskaber.

Jeg har i en lang Aarrække gjort mig det til sær\-%
lig Opgave at studere Videnskabernes Encyclopædie, idet
jeg har anseet det for aldeles nødvendigt, at den philo\-%
sophiske Videnskab, hvis den skal føres et Skridt videre
frem, først maa gjøre sig Rede for sit Forhold til de øv\-%
rige Videnskaber, de saakaldte Fagvidenskaber. Skulde
jeg altsaa have valgt et philosophisk, mig eiendomme\-%
ligt Emne, maatte jeg have begrændset det i en saadan
Grad, at Emnet selv upaatvivlelig kun vilde have inter\-%
% --- p. 281 ---
\opage{281}\setcounter{footnote}{0}%
esseret en meget lille Kreds af de Studerende. Man
tænke sig, hvad det vil sige at skulle i et forholdsviis
ringe Antal Timer gjennemføre en strengt videnskabelig
Analyse, hvoraf det kunde sees, hvorledes Philosophien
f. Ex. forholder sig til Mathematiken. Hvor mange Til\-%
hørere vilde vel have fulgt disse Analyser med udeelt
Interesse, hos hvor mange turde jeg vel forudsætte, at
her var ligelig Kundskab paa den ene Side til Philoso\-%
phien, paa den anden til Mathematiken? Havde jeg valgt
et andet, hermed beslægtet Emne, f. Ex. at fremstille Phi\-%
losophiens Forhold til Chemien, særlig med Hensyn paa
den saa vigtige Atomlære, og da ogsaa med Hensyn til
Forandrings-Problemet og Forholdet mellem de to Be\-%
greber Forandring og Forvandling, da var jeg atter bleven
henviist til en meget begrændset Tilhørerkreds.

Maatte jeg nu af anførte Grunde opgive at vælge et
strengt videnskabeligt Emne, var jeg paa den anden
Side henviist til den store Kreds af humane Ideer, hvor\-%
over Aandsvidenskaben saa rigelig har at disponere, og
idet jeg saae mig om i Kredsen af disse Ideer, manglede
jeg sandelig ikke Emner, der kunde komme paa Valg.
Men blandt alle Emner var der eet, der atter og atter
traadte i Forgrunden, Spørgsmaalet om Tro og Viden,
et Spørgsmaal, der er paa Dagsordenen, ja, jeg kunde
sige, for Tiden et brændende Spørgsmaal i mit Fædre\-%
land. Dette Spørgsmaal har just en eiendommelig Hoved\-%
side, hvorved dets Behandling falder ind under det hu\-%
mane Synspunkt, og jeg var i denne Henseende særlig
fristet til netop at lade Valget falde paa dette Emne.

Men, her reiser sig igjen en naturlig Indvending.
Hvorfor, vil man spørge, saa ikke uden videre anmelde
Forelæsninger over Forholdet mellem Tro og Viden og
holde disse, da de jo skulle sigte til at befordre en viden\-%
skabelig Samvirken mellem Universiteterne, i den akade\-%
% --- p. 282 ---
\opage{282}\setcounter{footnote}{0}%
miske, strengt videnskabelige Form? Jeg svarer: det var
% "var": the r stands displaced to the right ("va r") at the line end -- a loose sort, not letterspacing; set as one word
mig umuligt! Hvad Tro angaaer, kunde jeg vel, som jeg
har forsøgt at vise, udvikle de store Grundbegreber ogsaa
for dem, som ikke ganske vare paa det Rene med sig
selv om Vidensprinciperne; men naar det paa den anden
Side gjaldt om at anvise Viden og Videnskab sin strenge
Grændse, var det mig aldeles umuligt her at give noget
Overblik, der kunde gjøre Fordring paa videnskabelig
Stringens; jeg maatte da --- det ligger i Sagens Natur ---
udførlig have gjort Rede for mit philosophiske Vidensbe\-%
greb. Den eiendommelige Begrændsning, jeg, ifølge min
Synsmaade, maa give Videnskaben, er netop motiveret
ved de før omtalte encyclopædiske Studier og ved Under\-%
søgelsen af Philosophiens Forhold til de exacte og posi\-%
tive Videnskaber. Da jeg nu, som sagt, ikke kunde vente
at forebringe Noget om dette sidste Punkt, der vilde
blive indtrængende forstaaeligt, kunde jeg ikke faae Vi\-%
densbegrebet med tilstrækkelig Grundighed behandlet, og
maatte da, hvad dette Emne angaaer, opgive den i snev\-%
rere Forstand videnskabelige Fremstillingsmaade.

Hertil kom endnu en anden Overveielse: det, som
er mig det Vigtigste for Øieblikket --- thi med Viden og
Videnskab troer jeg nok, jeg skal blive færdig --- er ikke
Videns, men Troens Anliggende. Det er ikke den reent
videnskabelige, men den folkelige Udvikling af de mod\-%
satte Begreber, hvorpaa jeg har at lægge særligt Eftertryk.
For nu ikke at forskylde den Inconseqvens, naar jeg ud\-%
trykkelig sammenstillede de to Begreber, Viden og Tro,
og derved forpligtede mig til at oplyse Viden ligesaa
klart som Troen, at lade Videnskaben staae i Skygge,
omdannede jeg Titelen paa mine Forelæsninger og an\-%
meldte dem, ikke under Overskriften Tro og Viden, men
under det almindeligere og mere humane Synspunkt:
Hindringer og Betingelser for det aandelige Liv i Nu\-%
% --- p. 283 ---
\opage{283}\setcounter{footnote}{0}%
tiden. Under dette Synspunkt kunde jeg paa en friere
Maade fremdrage Meget, der deels hører under Civili\-%
sationen, deels under Folkelivet og særlig under Reli\-%
gionen. Jeg kunde her bevæge mig frit paa de frie Er\-%
kjendelsers Omraade. Saavidt om Valget.

Var Valget mig nu --- som jeg troer at have viist
--- paa en Maade foreskrevet, saa kom endnu hertil det
ingenlunde uvæsentlige Hensyn, at jeg paa disse Vilkaar
kunde gjøre mig Haab om at samle en Kreds af dan\-%
nede Tilhørere; ja, ikke blot af Tilhørere, men,
% "dannede . Tilhørere": a raised dot between the words is a risen space/quad, not a full stop; not transcribed
jeg maa særlig lægge Vægt paa denne Omstændighed,
ogsaa af Tilhørerinder. Det er ikke smigrende Ord; til
smigrende Ord er her ikke Sted, ikke Stemning; det er min
fulde Overbeviisning, at Grundtanker, der bevæge sig i
Troeslivet, Grundtanker med et personligt og folkeligt
Præg, maae paa Dannelsens øverste som paa dens nederste
Trin høres ligeligt af Mand og Qvinde. Det er ved Be\-%
handlingen af saadanne Opgaver af Vigtighed at have den
qvindelige Individualitet med, navnlig for Fremstillingens
Skyld, hvor det just gjælder om at skjelne skarpt imellem
det Videnskabelige, det Populære og det Folkelige. Kun
for en qvindelig Domstol ere vi ganske sikkre paa ikke
at blande de ueensartede Former mellem hverandre.
Hvad for det Første den videnskabelige Behandling an\-%
gaaer, da vil vist Ingen, der har foresat sig at tale om
Dannelsens almene Anliggender til Dannede, være udsat
for at forvilde sig ind paa de lærde Enemærker; men skulde
det hændes ham, vilde han med en vis Undseelse blive
Feiltagelsen vaer, naar han saae sig omgivet af opmærk\-%
somme Tilhørerinder. Det er vel kun sjeldnere, at en
begavet Qvinde føler sig hendragen til lærde Undersøgelser.
Det forstaaer sig: ingen Regel uden Undtagelse, og hvor Ta\-%
lentet er, være sig hos Qvinde eller Mand, maae vi holde
det i Ære. Vi kjende enkeltviis den eiendommelige Be\-%
% --- p. 284 ---
\opage{284}\setcounter{footnote}{0}%
gavelse hos ædle og skjønne Qvinder. Madame Chatelet
f. Ex. havde et saa eminent Talent for Mathematik, at hun
ved Siden af at spille den elskværdige fornemme Dames Rolle
% this line is very tightly set ("Sidenaf", "spilleden" nearly touch); word spaces supplied
i Salonerne, gjaldt for en videnskabelig Autoritet blandt sin
Tids store Mathematikere. Men, hvor sjeldent, ja, til
Held for Samfundet, saare sjeldent er ikke et saadant
Phænomen. Holde vi os til Reglen, da maa et Foredrag,
der skal kunne gjøre Regning paa qvindelig Opmærksom\-%
hed, helst være befriet fra lærd Oppakning og kunstige
Formler; men Foredrag af denne Natur egne sig jo aller\-%
bedst for Meddelelse af humant, æsthetisk, ethisk og re\-%
ligiøst Indhold.

Det Næste, hvorpaa jeg maa lægge Vægt, er Adskillel\-%
sen mellem det Populære og det Folkelige. For det Popu\-%
lære har Qvinden i en vis Forstand ligesaa megen Sands
som Manden; hun har og maa have sin Andeel i Cultur
og Civilisation: uden hendes Deeltagelse vilde det see
tarveligt ud med den saakaldte høiere Dannelse. Men
det blot Populære vil ikke lettelig bedaare hende; hvis
hun ikke ved altfor encyclopædiske Overhalinger i In\-%
stitutet er bleven seminaristisk-gouvernantisk forskruet,
vil hun altid med fiin Takt mærke sig den Grændse, hvor
det Besjælede hører op, og det Uforstaaelige begynder;
den flygtige, lette, tilsyneladende overfladiske Maade, hvor\-%
paa hun antyder en Opfattelse, er ofte et talende Beviis
% "et talende": a faint mark over the e (layer reads "ét"); taken for a speck, not an accent -- doubtful; rejected alternative "ét"
for, at hun med Instinktets Sikkerhed i Grunden har
truffet det Rette. Kan hun end ikke gaae slet saa vidt
ud i den populære Astronomie som den dannede Mand,
saa kan hun til Gjengjæld som oftest gaae nok saa dybt
% "Det Folkelige" and "det Personlige": pencil underlining by a reader; not transcribed. Not letterspaced.
ind i Poesie og Psychologie. Det Folkelige, det Reli\-%
giøse, kort: alt det Personlige opfatter en Qvinde dybt
og væsentligt, ligesaa dybt og væsentligt som en Mand.
Men naar vi sige, at Qvinden er istand til at forstaae
Personlighedens høieste Sandheder ligesaa oprindelig som
% --- p. 285 ---
\opage{285}\setcounter{footnote}{0}%
Manden, saa sige vi derfor ikke, at hun forstaaer dem
paa samme Maade som han. Nei! Forskjellen mellem
Mand og Qvinde er væsentlig en sjælelig Forskjel. Man
har ofte, naturligviis ikke uden lærd Selvvigtighed, spurgt,
om Qvinden, dette svagere Kjøn, nu ogsaa i Virkelig\-%
heden kunde tænke grundigt, reflectere grundigt, argu\-%
mentere grundigt o. s. v. Vi svare Ja! men hun reflec\-%
terer ikke, hun tænker ikke, hun argumenterer ikke som
en Mand; hun er jo en Qvinde.

Dette fører os et Skridt videre; thi netop ved at
berøre Forskjellen mellem Mandens og Qvindens Til\-%
egnelse af de høieste Sandheder, komme vi med det
Samme til at besinde os paa den dobbelte Opfattelse af
den folkelige Underviisning, af hvad der i en tidligere
Sammenhæng er charakteriseret som Folkeoplysning.
Skulde det virkelig forholde sig saaledes, som „den mo\-%
derne Bevidsthed“ forsikkrer, at Viden tilsidst maa op\-%
løse Troen, da vil Qvinden igjen synke ned paa et
meget underordnet Trin i Samfundet; thi det er vist:
paa Videns Vilkaar kan Qvinden aldrig blive Mandens
Lige. Skulde det ved nærmere Eftertanke vise sig
umuligt for en alsidig dannet Mand at forene Viden og
Tro i een Bevidsthed, da vil der aabne sig en sjælelig
Kløft mellem Mand og Qvinde; thi det er vist: en
Qvinde, der for Videns Skyld vil opgive at have sit
Liv i Troens Mysterium, maa sætte Qvindeligheden til
og visne som en Plante, der afskæres fra sin Rod.
En Mand, om han ikke er egentlig Videnskabsmand,
kan dog til Nød hutle sig igjennem ved Hjælp af det
Populære; men en Qvinde, hvis hele Dannelse er bygget
paa Raisonnementer og Grundsætninger, hentede fra
populære Skrifter, er en prosaisk Størrelse, en sjælelig
forfjasket Individualitet. En Mand kan til Nød blive
Kosmopolit og maaskee komme nogenlunde vel der\-%
% --- p. 286 ---
\opage{286}\setcounter{footnote}{0}%
fra; men naar en Qvinde bliver kosmopolitisk, hvad
bliver hun saa? Et goldt Abstractum, et i mange Ru\-%
brikker bemalet Schema, et Omrids af det Almeenqvinde\-%
lige uden Lys og Liv. Selv paa Dannelsens øverste
Trin staaer den qvindelige Individualitet paa Vagt saa\-%
vel for det Religiøse som for det Folkelige. Dette har
Grundtvig træffende seet, og hvad Grundtvig i djærv
Umiddelbarhed gjør gjældende, har S. Kierkegaards
subtile Reflexion gjennem æsthetisk, ethisk og religiøst
nuancerede Overgange tilfulde stadfæstet.

Som betegnende Fællesmærke for det Qvindelige
og det Folkelige nævne vi det dybere Enfoldige. Fra
det Enfoldige maa al sand personlig Dannelse, al
Aandsudvikling begynde; til det Enfoldige maa den
hele Udvikling, hvor rig og mangfoldig den end iøvrigt
kan være, igjen vende tilbage. Den folkelige Underviis\-%
ning begynder enfoldigt; thi den begynder med Bør\-%
nene, og den begynder med Almuen. Heri ligger, at
den just maa begynde saaledes, at den i sine Begyndel\-%
sesgrunde har Muligheder nok, har hele Anlægets Rig\-%
dom i de første Spirer. Thi ligesaa lidt som et Barn
er bestemt til at henleve sin Tid i en bestandig Barn\-%
dom, ligesaa lidt er Almuen bestemt til for bestandig at
forblive som Almue. Almuen er Folket i Tilstand af
aandelig Barndom; i Almuen har det folkelige Liv sin
Rod, sine Anlæg, sine Instinkter; i Almuen røre de
aandelige Kræfter sig paa dunkel, men oprindelig og
natureiendommelig Maade. Almuen er Folket selv, ikke
i fuldstændig Væren, men i Vorden og Opvæxt. Glemme
vi Opvæxten og sætte en stadig Væren istedenfor en be\-%
standig Vorden, da maae vi nødvendig strande paa een
af to Yderligheder: enten ligefrem gjøre Folket til Al\-%
mue eller betragte Almuen blot som Masse. De, som
paa en raa Maade identificere Folket med Almuen, ud\-%
% --- p. 287 ---
\opage{287}\setcounter{footnote}{0}%
mærke sig især ved et eget Had til Alt, hvad der tyder
paa høiere Dannelse, rigere Aandsudvikling, intellectuel
Overlegenhed: Videnskab, Kunst, Poesi betragtes med
skjæve Blik af Massernes Førere. De, der i aandsari\-%
stokratisk Forblindelse see ned paa Almuen som paa
den Deel af Befolkningen, Skjæbnen een Gang for
alle har bestemt til at trælle i Umyndighed, og gotte
sig ved de bekjendte Ord: „Zuschlagen kann die Masse,
da ist sie respectabel, urtheilen aber gelingt ihr miserabel“ ---
% German quotation set in the body roman, as the book does; not marked
forglemme Historiens dyrt kjøbte Erfaringer, hvor for\-%
færdelig en saadan Ringeagt dog har hævnet sig, og
fremdeles vil hævne sig. Folkeunderviisningen tør ikke
blive staaende ved det Almueagtige; den tør ligesaa lidt
indskrænke sig til det Religiøse alene som til det Borger\-%
lige. Der er i Folket, vi gjentage det, en dobbelt Trang:
Trang til Religion og Trang til Oplysning, Trang
til Sjælefrelse og Trang til Cultur. Trin for Trin maa
Almuen hæves til menneskelig Indsigt i det Menneske\-%
lige, til stedse høiere Selverkjendelse, til klarere For\-%
stand paa Virkeligheden, til dybere personlig Overbe\-%
viisning.

Uden personlig Overbeviisning gives der ingen indre
Selvstændighed, uden indre Selvstændighed ingen egent\-%
lig Frihed. Der gives og kan gives et frit Folk, men der
gives ikke og kan ikke gives en fri Almue; thi Almuens
Trin er, aandeligt seet, Umyndighedens, altsaa Ufrihedens
Trin. Det gjælder ved folkelig Underviisning om at
bringe den menige Mand saavidt, at han vil kunne:
tænke klart sine egne Tanker, gjøre sig tilstrækkeligt
Rede for sine egne Handlinger; indsee og forstaae, at
Ingen i Samfundet kan have Rettigheder, uden at han
tillige har Pligter; indsee og forstaae, at Enkeltmands
Interesse kun fremmes paa fornuftig Maade derved, at
Samfundsinteresserne varetages i det Store; indsee og
% --- p. 288 ---
\opage{288}\setcounter{footnote}{0}%
forstaae, at Samfundsinteresserne i det Store umulig
kunne fremmes, med mindre de forskjellige Arbeids\-%
klasser, ikke blot de, der arbeide i det Legemlige, men
ogsaa de, der arbeide i det Aandelige, nyde fuld Aner\-%
kjendelse, hver efter de i Forholdenes Natur begrundede
Vilkaar. Kun den, som er istand til at give sig selv
Grunde, for hvad han vil, og hvad han gjør, vil være
istand til at forstaae de Grunde, der gives ham af Andre.
Vanskeligheden med Menigmands Frigjørelse, er ikke
den, at han mangler enten ædel Selvfølelse eller sund
Forstand eller redelig Villie; men Vanskeligheden er, at
hans Tænkning er uklar, hans Synskreds indskrænket,
hans Sind mistænksomt. Vanskeligheden er, at han
ikke ret veed enten at give Grunde eller tage imod
Grunde, fordi han er væsentlig ufri. Menigmand er
ufri, saalænge hans Indre er dunkelt; men at hans Indre
er dunkelt, giver sig især tilkjende derved, at han i de
mest personlige Anliggender er mørk, indesluttet. Indre
Dunkelhed og Indesluttethed giver Charakteren, under
Livets Tryk, et kjendeligt Anstrøg af Tungsindighed.
Paa Overfladen kan Folkelivet vel i visse Solskins-Øie\-%
blikke være som Festligheden selv, et Billede paa natur\-%
lig Sorgløshed, umiddelbar Glæde og uudtømmelig
Livslyst; men inderst inde er der noget Mørkt, og i
Mørket boer det Tungsindige. En grundig Kjender af
Folkemelodierne fra Verdens forskjellige Egne har sagt,
at de i det Hele bære Præg af, at det ikke er Glæden,
men Sorgen, der har lært Folkene at synge. Dette
tunge, trykkede, indesluttede Væsen vil først forsvinde,
efterhaanden som Varme og Lys fra Frihedens Sol
% "efterhaanden som": a raised mark after "efterhaanden" is a risen space/quad; not transcribed
trænger mere igjennem i Masserne.

Og hvem skal saa undervise Almuen, sørge for dens
Opvæxt, opdrage og danne den til Folk? Der kan, hvor\-%
mange Forandringer Samfundsformerne end undergaae,
% --- p. 289 ---
\opage{289}\setcounter{footnote}{0}%
ikke være Tvivl om, at Præster og Skolelærere frem\-%
deles maae vedblive at være de egentlige Folkelærere.
Hvorledes disse Lærere opfatte og røgte deres Kald, er
da af den allerstørste Betydning ikke blot for Staten, at
den kan holde over Fred og Orden, men for Folket
selv, at det kan sikkre sig aandelig Frihed. Man hører
undertiden Yttringer som disse: til at lære en Videnskab
hører der et godt Hoved, til at blive Kunstner hører
der Talent; selv til at blive Berider hører der naturlige
Anlæg. Men til at blive Landsbypræst? Ja, Præst paa
% "Anlæg. Men": a raised mark after the full stop is a risen space/quad; not transcribed
Landet kan man da blive med et lille Hoved, uden
Talent. Og Skolelærer? Dertil behøves neppe naturlige
Anlæg! Slige Domme ere i høi Grad skjæve og forvir\-%
rende. Naar der disputeres om, hvor megen Begavelse
der hører til for at prædike for Almuen, er det strax
% "hører til": spacing.py reports a RUN; the gap is a risen space/quad (raised mark) after "hører", not letterspacing -- not \emph
en grundfalsk Begyndelse at søge Maalestokken i det
gode Hoved. Maalestokken er theologisk, og just derfor
falsk. Man gaaer theologisk ud fra den Forudsætning,
at en Præst i Grunden bør være en lærd Mand; eller
dog en Slags populær Videnskabsmand at sige, naar han
skal være Præst med Anseelse. Dette er nu Idealet.
Men da Idealerne saare sjelden realiseres i Verden, maa
man i mange Tilfælde see igjennem Fingre med de vi\-%
denskabelige Mangler og trøste sig med, at en Præst,
uagtet han hverken er egentlig lærd eller Videnskabs\-%
mand, dog, naar han for Resten er en troende Mand,
der kjender sin Christendom, meget vel kan være en
god Sjælesørger. Forstaaer sig, han kan i Reglen ikke
faae det store Kald, det med mange og rige Ind\-%
komster, naar han ikke har den store theologiske Cha\-%
rakteer; men derfor kan han jo nok faae et betydeligt,
byrdefuldt Kald, nemlig et Kald med mange fattige
Sjæle, naar han, som sagt, er god nok til at være Sjæle\-%
sørger. Der kan i det praktiske Livs mangfoldige Til\-%
% --- p. 290 ---
\opage{290}\setcounter{footnote}{0}%
fælde neppe paavises nogen større Modsigelse mellem
Maalestokken og det, der skal maales, end den, der
falder i Øinene, naar Talen er om at bedømme de af
Staten lønnede Sjælesørgeres Qvalificationer. Ingen kan,
derom holder jeg mig overbeviist, føle dette dybere end
Præsterne selv.

Thi, lad os tale reent ud, hvormange Præster blive
vel egentlig Videnskabsmænd? Ikke saare mange --- til
Held for Præstestanden, ja til Held for Menigheden!
Det er dog i Sandhed ikke Præstens Kald at fordybe
sig i lærde videnskabelige Studier; han har en anden
Gjerning at gjøre, en Aandens Gjerning, hvortil der,
ikke alene i guddommelig, men ogsaa i menneskelig
Forstand, udkræves en egen Begavelse, en særegen For\-%
beredelse, Opdragelse og Dannelse. Lad os spørge
Præsterne omkring i Landet, om ikke de fleste ville til\-%
staae, at de Uddrag af lærd Theologie, der paa en saa\-%
kaldet videnskabelig Maade blev dem indprentet, just
høre til det, de helst maatte see at glemme jo før jo
heller, for at kunne begynde forfra? Mon de Fleste
iblandt dem ikke skulde have mærket sig en Forskjel,
en væsentlig Forskjel, mellem den theologiske og den
religiøse Tankegang, mellem det lærde og det opbygge\-%
lige Foredrag? Istedenfor at hjælpe sig med lærde
Brudstykker og videnskabelige Overblik ere de netop
henviste til Iagttagelser af det virkelige Liv, til person\-%
lig Opfattelse og Tilegnelse af det apostoliske Ord, til
at øse af de indre Erfaringers Kilde.

Vistnok kan der behøves en objectiv Regel, hvor\-%
efter Folkelærernes Qvalificationer kunne bedømmes, saa
sandt Staten skal, ikke blot lønne sine Embedsmænd,
men ogsaa saavidt muligt afpasse Lønnen efter Tjenesten;
thi en subjectiv, pietistisk Mening om den Enkeltes „Be\-%
segling“ er ganske vist utilforladelig. Men just i denne
% --- p. 291 ---
\opage{291}\setcounter{footnote}{0}%
Henseende viser det Urimelige sig, naar Theorien gaaer
i en Retning og Praxis i en anden. Var den nu be\-%
staaende Præsteskole virkelig skikket til at give sine
Studerende den videnskabelige Dannelse, den religiøse
Erkjendelse, den aandelige Indvielse, der i særegen
% "Erkjendelse,": a raised mark before the comma is a risen space/quad; not transcribed
Grad gjør dem skikkede til at opdrage Folket og for\-%
kynde Evangeliet i Menigheden: saa vilde et Vidnesbyrd
om, hvorledes den Enkelte havde bestaaet sin Afgangs\-%
prøve, ganske vist kunne være orienterende; men er dette
nu Tilfældet? Der kan ikke være Tvivl om, at jo Præ\-%
sterne selv, naar det bliver dem klart, at den Vei, der
skal føre dem til Maalet, er en ganske anden end den
objectivt dogmatiske, ville være de Ivrigste til at modar\-%
beide det theologiske Uvæsen. Thi det tynger paa dem,
dette Uvæsen; det sløver Sandsen og forvirrer Blikket;
det lammer Aandskræfterne ved at proppe Hukommelsen
med Stykker og Stumper af lærde Traditioner, litterære Re\-%
liqvier, hvori end ingen Dæmon skal kunne indblæse enten
Liv eller Aand, med mindre det skeer ved et speculativt
Mirakel. At danne Præster i en Skole, der udvisker
Aabenbaringens Charakterer, overseer Troesprincipet og
fornegter den religiøse Tænknings oprindelige Gyldighed,
hvad fører det til? Det fører til en stille, planmæssig,
jævnt fremadskridende Underminering af de Geistliges Stil\-%
ling i det civiliserede Samfund. Hvad var det vel, der
egentlig begrundede Geistlighedens Magt og Myndighed
i Middelalderen? Det var dens klarere Blik for Aanden,
dens dybere Indsigt, dens ideelle Forspring. I Nu\-%
tiden, hvis aandelige Retning i det Hele er overveiende
intellectuel, reflecterende, kritiserende, maa den Geistlige
ikke blot see at følge med, men, hvad det Religiøse an\-%
gaaer, udtrykkelig hævde sin intellectuelle Overlegenhed.
Samtiden maa faae et Indtryk af, at han netop er Manden,
den rette Mand, der veed Besked, naar Spørgsmaalet er
% --- p. 292 ---
\opage{292}\setcounter{footnote}{0}%
om det personlige Livs afgjørende Grundtanker, den
Mand, der i religiøse Reflexioner kan magte Fritænkeren,
gjøre Regnskabet op med de Vidende og orientere de
Troende. Kan dette Indtryk blive det overveiende i Sam\-%
fundet, saa kan Geistligheden tilvisse være ubekymret for
alt det Øvrige. Blandt krigsførende Magter vil den, der
har de største Kræfter, den strengeste Disciplin og de
bedste Váaben, altid være sikkrest paa Seiren. Det skal i
% printer's defect: "Váaben" -- an acute stands plainly over the first a at 300 ppi (a wrong sort; cf. "Vaaben" twice below); transcribed as printed. Rejected alternative: a speck, "Vaaben"
Aandernes Strid vel stadfæste sig, at der paa Historiens Vei
ikke gives nogen Aand saa mægtig som Evangeliets Aand,
ingen Tanke saa gjennemtrængende som Villiestanken,
intet Ord saa kraftigt som Troens Ord, og at da heller
ingen Stridende, der kjender sine Vaaben, er saa vel ud\-%
rustet som den, der strider for Personlighedens Ret, dens
sande Liv i det evige Rige. Men det præstelige Vaa\-%
benskjold er ikke Theologie; at iføre den Geistlige et
theologisk Harnisk er i vore Dage ligesaa bagvendt, som
det er bagvendt at skrue ham op til „Sandhedsvidne“ og
give ham til Priis for Satiren.
% this page is curved at the gutter edge; raised marks beside the commas after "Fritænkeren", "Aand", "Villiestanken", "fordre", "Høiskole", "Elementer" and between "Grundlag" and "efter" are risen spaces/quads; not transcribed

Endnu et Punkt skal her berøres. Det gjælder i
Nutidens aandelige Liv om Vexelvirkning mellem Viden\-%
skab og Folkeoplysning, en høiere Vexelvirkning, der
upaatvivlelig vil føre med sig, at hvad der allerede i mit
Fædreland har været Gjenstand for meget omfattende og
alvorlige Overveielser, ogsaa snart i Norge og Sverrig
vil tildrage sig Opmærksomheden, dette nemlig, om
ikke den nærværende Tids aandelige Vilkaar skulde fordre,
at der ligeoverfor Universitetet, den lærde Høiskole,
reiser sig en anden Høiskole paa folkeligt Grundlag efter
en stor Maalestok. Det er ikke nok med de smaa Bonde\-%
høiskoler; her behøves en Læreanstalt, hvilende paa et
saa bredt Grundlag, at den kan omfatte de forskjellige
Dannelsens og Videnskabens Elementer, vel at mærke
saaledes, at Folkeoplysningens Grundtanke gjennem\-%
% --- p. 293 ---
\opage{293}\setcounter{footnote}{0}%
trænger hvert enkelt Element og behersker det Hele.
Det er mit Haab, at det vil komme dertil, idet jeg til\-%
lige forudsætter, at en saadan Modsætning af Folkehøi\-%
skole og Universitet langt fra at trykke Universitetet og
sætte det i Skygge som en forældet Indretning, tvert\-%
imod vil mægtig bidrage til at hæve dets videnskabelige
Anseelse, indgive det nye Impulser, befrie det for uved\-%
kommende Byrder, sætte Ideerne i friere Omløb og for\-%
ynge Blodet i de gamle Lemmer. Der er dem, der mene,
at Universiteterne ere forældede Institutioner, der væsentlig
have udspillet deres Rolle. Denne Mening kan jeg ikke
tiltræde; hvorimod jeg vel kan forstaae, at Videnskabernes
høieste Foreningspunkt ikke kan gjøre Fyldest som Folke\-%
oplysningens eneste Brændpunkt, naar det staaer fast,
at denne Oplysning maa have et dobbelt Udgangspunkt,
et i Troen og et i Viden.

Der kan jo dog ikke være Tvivl om, at de særlige
Fagvidenskaber have bedst af at gruppere sig til et
stort Hele, ligesom det heller ikke kan betvivles, at
Folkeoplysningen med sine særegne Opgaver helst
maa have sit eget Midtpunkt. Hvad Philosophien
angaaer, vil den som Aandsvidenskab forhandle lige\-%
lig med begge; til den ene Side vil den fordybe
sig i de strenge Videnskaber og optage videnskabe\-%
lige Elementer i Ideelæren, til den anden Side gaae ind
i den folkelige Underviisnings Grundtanke og optage
religiøse Problemer i Erkjendelseslæren, altsaa danne et
Bindeled mellem de to Skoler. Er det nu --- efter hvad
vi kunne skjønne --- saa langt fra, at de to Skoler een\-%
sidig ville indskrænke, forhutle og fortrædige hinanden,
at de tvertimod ved Modsætningen ville styrke og frem\-%
hjælpe hinanden, da viser det sig jo ogsaa fra dette Syns\-%
punkt, hvad den Fordobling, vi her have søgt at gjøre
gjældende, væsentlig har at betyde. Dette med Hensyn
% --- p. 294 ---
\opage{294}\setcounter{footnote}{0}%
til Hovedsagen. Hvad Grundtankens nærmere Gjennem\-%
førelse angaaer, maae vi, saavel i det Praktiske som i
det Theoretiske, henvise til Fremtiden selv. Det gaaer
nemlig med ideelle Opgaver saaledes, at den, der løser
een, altid faaer to igjen. Dette, at jo mere man arbeider
med en Aandsopgave, desto mangfoldigere bliver den,
er netop det, som forskrækker Saamange, baade dem, der
% "dem, der": a raised mark before "der" is a risen space/quad; not transcribed
vælge Theorien, og dem, der betræde Erfaringens prak\-%
tiske Vei. Og dog er dette netop Aandsudviklingens Lov.
Den ene omfattende Opgave, at fremme og befordre det
aandelige Liv, indeholder utallige Opgaver; de færreste
blive vel løste i Nutiden; de fleste tilhøre Fremtiden.
Om Fremtidens Opgaver bør der helst tales kort, da
Talen ellers let bliver trættende, idet den bliver svævende,
famlende, ubestemt: at bruge det Nærværende paa bedste
Maade er at forberede det Tilkommende.

% [short centred rule between sections, mid-page; the lecture continues below it]
\begin{center}\rule{3cm}{0.4pt}\end{center}

Det var en stor og, som vi haabe, frugtbar Frem\-%
tidstanke, der samlede denne Hæderskreds af Tilhørere
og Tilhørerinder. Hvorvidt der nu i de holdte Foredrag
er hørt et Ord, udviklet en Livstanke, som Fremtiden
selv vil godkjende, det er et Spørgsmaal, der vel i nogen
Tid endnu maa vente paa sin Besvarelse. Kun dette
kan jeg sige, og jeg siger det aabenlyst, at jeg stoler
paa Grundenes Vægt. Jeg har talt uforbeholdent,
af indre Tilskyndelse, bestemt; men jeg har aldrig
% "indre Tilskyndelse": pencil underlining by a reader; not transcribed. "Jeg" here and below is loosely justified, not letterspaced
tænkt eller meent, at nogen af mine Sætninger skulde
gjælde for Troesartikler. Jeg har talt for at indstille
Sagen til Sagkyndiges Overveielse; jeg har talt mod
Civilisationens Overgreb, men ikke mod Civilisationen,
mod Fornuftens Misgreb, men ikke mod Fornuften,
mod videnskabeligt Blendværk, men ikke mod Viden\-%
% --- p. 295 ---
\opage{295}\setcounter{footnote}{0}%
skaben. Jeg har talt om Folkets Trang og Folkets Ret
til en særegen af Videnskab og lærd Forskning uafhængig
Oplysning; jeg har talt frit og aabent, thi jeg var mig
bevidst, at jeg talte til en Forsamling, i hvis Midte
Mænd af høi Dannelse, grundig Lærdom, praktisk og
theoretisk Sagkundskab, saadanne blandt Rigets første
Mænd, til hvilke Folket med Høiagtelse og Tillid maatte
see op, havde for Sagens Skyld bestemt sig til at vælge
en Plads. Jeg holdt mig overbeviist om, at man til
Videnskabsmænd kunde tale, ikke blot til Priis for Viden\-%
skaben, men ogsaa om videnskabelige Overgreb, uden
at blive misforstaaet; jeg holdt mig overbeviist om, at
man til folkelige Mænd kunde tale om Folkets Trang og
Savn og Ret, uden at foranledige Misstemning. Jeg
har talt frit; men jeg tager med mig den Overbeviisning,
at hvormeget Mangelfuldt der end kan have været i det
Enkelte, saa kan min frie Tale dog neppe have for\-%
dunklet eller forqvaklet nogen god Sag.

Nu staae vi ved Maalet; jeg tilstaaer med dyb Er\-%
kjendtlighed, at disse Foredrag ville være i mere end
een Henseende betegnende for min Lærervirksomhed, at
de for den angive et væsentligt Vendepunkt. Det er et
Blad af mit Livs Bog, et Mindeblad, jeg her fik beskrevet;
blandt mange Blade er det i Sandhed et af de rigeste,
jeg tilstaaer det, et af de skjønneste. Saa vil jeg da
prente dette dybt i mit Sind, vie det til en taknem\-%
melig Erindring; saa være da Sagen indledet, saa være
hermed de indledende Forelæsninger sluttede: den ærede
Forsamling af Tilhørere og Tilhørerinder modtage min
varmeste, min ærbødigste Tak!

\begin{center}\rule{3cm}{0.4pt}\end{center}

% --- p. 296 ---
% [p. [296] is BLANK: the verso after the end of Forelæsning XVI. No folio,
%  no text, so no \opage.]

\chapter*{Efterskrift.}
\addcontentsline{toc}{chapter}{Efterskrift}
\markboth{Efterskrift}{Efterskrift}
\label{ch:efterskrift}
% pp. [297]--326  (p. 297 carries no folio)
% --- p. 297 ---
\opage{297}\setcounter{footnote}{0}%
% [large initial H]
Hvormeget jeg end i disse Forelæsninger har bestræbt mig
for at udtale min Opfattelse af Forholdet mellem Tro og
Viden ret tydelig og bestemt, har det dog ikke kunnet und\-%
gaaes, at adskillige, endog meget væsentlige Punkter ere blevne
udsatte for Misforstaaelse. Jeg skjønner det, deels af enkelte
Forespørgsler fra Tilhørere, deels af Artikler i norske Blade.
Navnlig er der eet Punkt, til hvis yderligere Opklaring jeg
ved denne Anledning allerførst finder mig opfordret, det angaaer
% a raised mark before "det angaaer" is a risen quad; not transcribed
den \emph{Virkelighed}, der maa tillægges de hellige Kjendsgjer\-%
ninger, særlig Christi Opstandelse.

Jeg kan, hvad denne Sag angaaer, henholde mig til en
Artikel, der findes i \emph{Aftenbladet} for 24de December 1867
under Overskrift: Professor C. P. Caspari om „Tro og Viden“.
Det er i „Bekjendelsens og Vidnetsbyrdets Form“, at Hr.
% "Vidnetsbyrdets" (for "Vidnesbyrdets") as printed, confirmed at 600 dpi (audit); not emended
Prof. angiver, hvad han „med et lidet passende Ord“ maa
kalde sit „Standpunkt“. „Jeg gjør det,“ siger han, „fordi
jeg paa Grund af Anskuelser og Tanker, som i den senere
Tid efterhaanden have begyndt at røre sig og gjøre sig gjæl\-%
dende iblandt os, . . . maa ansee det for at være ikke saa
ganske afveien.“ Da de „Anskuelser og Tanker“ Hr. Prof.
Caspari ved at aflægge sin Troesbekjendelse her vil imøde\-%
gaae, ikke ligefrem ere satte i Forbindelse med mine Fore\-%
drag, kan det naturligviis ikke falde mig ind at ville paabyrde
% --- p. 298 ---
\opage{298}\setcounter{footnote}{0}%
ham nogen videnskabelig Forhandling med mig; jeg afbenytter
blot hans Udtalelse, fordi den giver mig en beqvem Anled\-%
ning til at anbringe endnu nogle Bestemmelser af, hvad jeg
forstaaer ved \emph{Aandsphænomen, Aandsvirkelighed,
Virkelighed for Aanden}.

Hr. Prof. Caspari bekjender sig til den Tro, „at den
Frelsesaabenbaring, hvorom de hellige Skrifter berette, tilhører
den ydre historiske Virkelighed, at den ikke blot er sand, men
ogsaa virkelig, at den ikke blot eier en indre Virkelighed for
Troen --- hvad der vilde være det samme som, at den kun
er til i Ideen, eller at den i Grunden kun er en Indbildning
--- men ogsaa en ydre Virkelighed i sig selv, og det ikke
blot i sine Hovedtræk og Grundkjendsgjerninger, men ogsaa
i alle sine Enkeltheder. Og paa denne Frelsesaabenbarings
ydre Virkelighed beroer vor Frelse, uden den vilde vi være
og i al Evighed forblive i vore Synder, være og blive evig
fortabte. Blotte Frelsesideer kunne ikke skjænke det i Syn\-%
den urolige Menneskehjerte Fred og Salighed. . . . Men“ ---
vedbliver Hr. Professoren --- „Videnskabens Indsigelser mod
Frelsesaabenbaringens Grundkjendsgjerninger og Enkeltheder?
--- Jeg er selv, af Guds Naade, en Videnskabsmand, men
jeg maa aabent bekjende, at de aldrig have forvoldt mig nogen
søvnløs Nat. Hvad navnlig Indsigelserne mod Frelsesaaben\-%
baringens Mirakler angaaer, maa jeg derimod stille disse Mi\-%
raklers og fremfor Alt Christi Opstandelses historiske Virke\-%
lighed. Gives der nogen Kjendsgjerning, der er tilstrækkelig,
ja overflødig historisk bevidnet, saa er det Christi Opstan\-%
delse; skal den ikke kaldes historisk, da kunne vi gjerne slaae
en Streg over Alt, hvad der kaldes Historie. . . . Og saa
gives der jo et Mirakel, som skeer den Dag idag, som er en
Frugt af hine Frelsesaabenbaringens gamle Mirakler og der\-%
for ogsaa giver dem Vidnesbyrd, og som Enhver kan erfare
paa sig selv, jeg mener Gjenfødelsens Mirakel, det store,
indre Mirakel, et Mirakel som --- hvis det er tilladt paa Mi\-%
% --- p. 299 ---
\opage{299}\setcounter{footnote}{0}%
raklernes Omraade at tale om Større eller Mindre --- er lige\-%
saa stort som Christi Opstandelse. Enhver sand Christen
veed sig selv som et levende Mirakel, Miraklernes Virkelighed
er ham derfor ogsaa ligesaa vis som hans egen Tilværelse
som nyt Menneske, som Christen.“

I denne Udtalelse er der Noget, hvorom jeg strax maa
erklære mig ubetinget enig med Hr. Prof. Caspari, og det er:
at Kjendsgjerningerne i den hellige Historie ikke ere „Ind\-%
bildninger“, at „blotte Frelsesideer ikke kunne skjænke det i
Synden urolige Menneskehjerte Fred og Salighed“; med hvor\-%
megen Grund jeg her fremhæver dette, vil Enhver, der med
uhildet Blik har gjennemlæst, ikke blot nærværende Forelæs\-%
ninger, men et hvilketsomhelst af mine øvrige Skrifter, kunne
sige sig selv. Det er just de hellige Kjendsgjerningers Vir\-%
kelighed, deres religiøst udprægede Charakteer, deres fulde,
uforvanskede Realitet, deres usvækkede Gyldighed i sig, hvis
Anerkjendelse jeg fra alle Sider og paa enhver mulig Maade
søger at indskærpe. Men, naar det for Alvor gjælder om
at hævde Virkeligheden, den fulde, selvgyldige Virkelighed,
saa paatrænger sig jo det Spørgsmaal: hvad er Virkelighed,
den sande, i sig fuldgyldige Virkelighed? Det er i Besvarelsen
% "fuldgyldige": an ink blot on the i gives it a tail like a j ("fuldgyldjge");
% (both OCR witnesses read "fuldgyldjge"); at 300 ppi the i keeps its dot and
% the tail is a blot, not a j descender, so the sense-reading is kept
% AUDIT doubtful: at 600 dpi the i carries a dot AND a descending tail; cannot rule out a j sort
af dette vigtige Spørgsmaal, at jeg ikke veed, om jeg tør er\-%
klære mig enig med Hr. Caspari.

For det Første forekommer det mig, som om Hr. Pro\-%
fessoren, hvor Talen er om Aabenbaringsphænomenerne, lægger
vel megen Vægt paa det sandselige Ydre, naar han lader
vor Frelse og Salighed være afhængig deraf og indskærper,
at dette Virkelighedens Ydre skal gjælde, ikke blot om visse
% a dot on the baseline after "Ydre" (inked quad, no sentence ends); not transcribed
Kjendsgjerninger, men om alle, ikke blot om „Grundkjends\-%
gjerningerne“, men om „deres Enkeltheder“. Lad mig vælge
nogle Exempler. Esaias beskriver sin Kaldelse til Prophet
(Es. VI.): „I det Aar, Kong Usias døde, da saae jeg Herren
sidde paa en høi og ophøiet Throne, og hans Flige opfyldte
Templet. Serapher stode over ham, hver med sex Vinger“
% --- p. 300 ---
\opage{300}\setcounter{footnote}{0}%
% AUDIT doubtful: a faint, lightly inked dash-like mark between "o. s. v." and "Propheten"; not transcribed
o. s. v. Propheten Amos beskriver sine Syner (Am. VII.):
„Dette Syn lod Jehova mig see: han var ifærd med at danne
Græshopper. . . . Dette Syn lod Jehova mig see: See den
Herre Jehova kaldte paa Luen som Sagfører; den opslikkede
den mægtige Dybde og fortærede Lodden. . . . Dette Syn
lod Jehova mig see: Herren staaende paa en afloddet Muur,
med et Blylod i Haanden. Og Jehova sagde til mig: hvad
seer du Amos? Jeg sagde et Blylod. Saa sagde Herren:
% "Jeg sagde": a speck above the a of "sagde"; not an accent, not transcribed
nu vil jeg kaste Blyloddet hen i mit Folk Israels Midte“
o. s. v. Til Ezechiel lyder den guddommelige Stemme (Ez.
III.): Menneskesøn! æd hvad du her annammer, æd denne
% printer's defect: no opening „ before "Menneskesøn!", though the quotation
% is closed after "Hon-ning." below; transcribed as printed, quote balance -1
Bogrulle; gaa derpaa hen og tal til Israels Huus. Og jeg
oplod min Mund, og han gav mig denne Bogrulle at æde.
Og han sagde til mig: Menneskesøn; dit Liv lade du fortære,
og dit Indre fylde du med denne Rulle, som jeg giver dig;
og jeg aad, og den blev i min Mund til Sødme som Hon\-%
ning.“ . . .

At disse prophetiske Syner ikke ere tomme Indbildninger,
men sande Aabenbaringer, at disse Aabenbaringer have Vir\-%
kelighed, fuld Virkelighed for Aanden, maa fastholdes; men,
spørge vi, er det nu tilladt ved slige Phænomener at sætte
Skilsmisse imellem Prophetens Syn, og hvad Propheten seer?
Tør vi virkelig antage, at den Throne, hvorpaa Jehova lader
sig see for Esaias i Templet, har samme udvortes Virkelig\-%
hed som Jerusalems Tempel? at de Græshopper, Jehova
danner i Synet for Amos, have samme udvortes Virkelighed,
som naturhistoriske Græshopper? at det Blylod, Jehova kaster
hen i Israels Midte, har en udvortes Virkelighed som det
Blylod, der hænger ved et bornholmsk Stueuhr? at den Bog\-%
rulle, Ezechiel æder --- i Synet, har samme udvortes Virke\-%
lighed som den Bibel, der ligger paa Bordet?

De prophetiske Syner sees i Aanden. Den, som siger,
% a dot on the baseline after "sees" (inked quad); not transcribed
at de Syner, Propheterne see, have samme udvortes Virke\-%
lighed, som de Ting, der sees med legemlige Øine, hvad gjør
% --- p. 301 ---
\opage{301}\setcounter{footnote}{0}%
han? Han gjør Aabenbaringen aandløs. Kan en aandløs
Aabenbaring kaldes „Frelsesaabenbaring“? Kan Nogen, uden
at drive Spot med det Hellige, virkelig paastaae, at al Frelse
og Salighed beroer paa Frelsesaabenbaringens ydre Virkelig\-%
hed, at den, der, hvad de nævnte Syner angaaer, ikke lægger
Sandheden ind i det Ydre af „Thronens“, af „Græshoppernes“,
af „Blyloddets“, af „Bogrullens“, kort: af „alle Enkelthe\-%
dernes“ ydre Virkelighed, maa „i al Evighed forblive i sine
Synder, være og blive evig fortabt“?

Det er ingenlunde min Hensigt at ville paatvinge Hr.
Prof. Caspari Conseqvensen af hans egen Udtalelse; jeg vil
kun gjøre tænkende Læsere opmærksom paa den farlige Mis\-%
viisning, der ligger i et saa eensidigt „Vidnesbyrd“. Et er det
at anerkjende det Ydre, forsaavidt det i Aabenbaringsphæno\-%
menet hører med til at bestemme det Indre, et Andet at for\-%
see sig paa det Ydre, det sandselige Ydre, i den Grad, at
man gjør al Frelse og Salighed afhængig af Enkelthedernes
ydre Virkelighed. At Hr. Professoren ingenlunde vil fornegte
det Indre, giver han tydelig nok tilkjende ved det Eftertryk,
han lægger paa Gjenfødelsens Mirakel, „det store indre Mi\-%
rakel, som --- hvis det er tilladt paa Miraklernes Omraade
% a dot on the baseline after "er" (inked quad); not transcribed
at tale om Større eller Mindre --- er ligesaa stort, som Christi
Opstandelse“. Ved dette „Vidnesbyrd“ er der nu ganske vist
lagt afgjørende Vægt paa det Indre; men, hvordan forholder
det Sidste sig nu til det Første? Om Gjenfødelsens Mirakel
gjælder det jo, at det har indre Virkelighed, Virkelighed for
Troen. Tag Troen bort, hvad bliver det saa til med Gjen\-%
fødelsens „store Mirakel“? Forvisningen om dette Mirakels
% a raised mark after "Mirakels" is a risen quad, not an apostrophe; not transcribed
Gyldighed er i Troen alene; eller skulde den Troende maa\-%
skee kunne anskue sin egen Gjenfødelse i en ydre Virkelig\-%
hed, en Virkelighed, hvormed det forholder sig, for at bruge
Hr. Professorens Udtryk, „som med Solen, der vel først bliver
til for mig derved, at jeg seer og føler den, men dog ikke
derved først overhovedet bliver til, men er til før al min Seen
% --- p. 302 ---
\opage{302}\setcounter{footnote}{0}%
og Følen, har en af den uafhængig Tilværelse og vil være
til, om aldrig noget menneskeligt Øie havde seet eller noget
menneskeligt Legeme følt den“?

% AUDIT doubtful: "Maa" may be letterspaced (emphasis?); left unmarked
Maa ethvert tænkende Menneske --- følgelig ogsaa Hr.
Professor Caspari --- indrømme, at Gjenfødelsens Mirakel
har sin sande Gyldighed i sig formedelst Aanden ved Troen
ja for Troen og Aanden \emph{alene}: saa er jo hermed ind\-%
rømmet, at der gives en indre, og dog i sig gyldig Virkelig\-%
hed, en Aandsvirkelighed, en Virkelighed for Aanden. Men
hvorledes skulle vi nu forstaae Hr. Professorens „Bekjendelse
og Vidnesbyrd“, thi han bekjender og vidner, at den Frelses\-%
aabenbaring, der blot eier en indre Virkelighed i Troen,
„kun er til i Ideen, eller at den i Grunden kun er en Ind\-%
bildning“. I Sandhed, et forfærdeligt Vidnesbyrd --- imod
„Gjenfødelsens store Mirakel“! Da nu Hr. Professor Cas\-%
pari her har aflagt, ikke eet, men to Vidnesbyrd om Virke\-%
ligheden af Gjenfødelsens Mirakel, nemlig eet \emph{for} og eet
\emph{imod}, begge lige kraftige, saa spørges: hvorledes skulle vi,
der gjerne ville annamme hans „Vidnesbyrd“, finde Rede i
en med sig selv saa splidagtig „Bekjendelse“?

Vi blive vist nødsagede til at indrømme: 1) at Aaben\-%
baringens Virkelighed først og sidst er Aandsvirkelighed, Vir\-%
kelighed for Aanden; 2) at Aandsvirkeligheden altid gaaer
fra det Indre gjennem det Ydre tilbage til det Indre, saa at
dens afgjørende Kjendemærke altid maa søges i det Indre;
3) at det Ydre, hvilket Aanden, den sig aabenbarende Guds
Aand, magter, gjennemvirker og optager i de hellige Kjends\-%
gjerningers eiendommelig bestemte Virkelighed, er af meget
forskjellig Art, snart legemligt, snart ulegemligt, snart natur\-%
ligt, snart overnaturligt, snart umiddelbart sandseligt, saa at
Aabenbaringsphænomenerne just af den Grund maae, hvis
vi virkelig skulle kunne „give os hen til dem og fordybe os i
dem“, klare sig for os ved en, ikke til „vore egne vilkaarlige
% --- p. 303 ---
\opage{303}\setcounter{footnote}{0}%
Forestillinger“, men til Aabenbaringens og Troens selvgyldige
Princip, svarende Kritik.

Jeg skal i Korthed belyse enhver af de opstillede Sæt\-%
ninger. Altsaa:

% the first thesis carries no "1)" (cf. "2)" p. 304, "3)" p. 306), as printed
\emph{Aabenbaringens Virkelighed er Aandsvirke\-%
lighed, Virkelighed for Aanden}. Heri ligger, at Intet
kan have Gyldighed for Troen uden det, som den Troende
i Kraft af Aandens Vidnesbyrd maa tillægge Gyldighed i sig.
Ifølge Hr. Professor Casparis „Vidnesbyrd“ skulde man jo
antage, at Noget kunde have indre Virkelighed for Troen
og endda maaskee være en puur og bar Indbildning; dette
maa benegtes. At Noget har indre Virkelighed for Troen,
vil sige, at det for den Troende har Gyldighed i sig; saa\-%
snart den Troende bliver sig bevidst, at det, hvorpaa han
troer, ingen Gyldighed har i sig selv, saa ophører han ogsaa
at troe derpaa. Den, der troer paa Gud, tillægger Gud
Gyldighed i sig. Vilde Nogen sige: Gud er vel i og for
sig ikke til; men just derfor troer jeg paa ham, at jeg kan
give ham en indre Virkelighed i min Tro: hvad maatte man
da dømme om et saadant Menneskes mentale Tilstand? Vilde
Nogen sige: Jeg troer paa Gud, det er: jeg giver ham en
foreløbig indre Virkelighed i min Tro; men, om den Gud,
paa hvem jeg troer, nu ogsaa har Virkelighed i sig, det maa
jeg lade staae hen, indtil jeg paa en eller anden Maade kan
faae ham at see i en ydre Virkelighed, f. Ex. siddende i
Templet paa en Throne eller staaende paa en Muur med en
Græshoppe eller med et Blylod i Haanden: hvad maatte man
da dømme om en saadan Tro? Nei, at troe paa Gud er at
tillægge Gud Gyldighed, absolut Gyldighed i sig selv; den,
for hvem Guds personlige Væsen ikke har Gyldighed i sig
selv, han troer ikke paa Gud. Naar Hr. Professor Caspari
istedenfor „blotte Frelsesideer“ vil have „Personer, guddom\-%
melige Personer, den personlige treenige Gud“ at holde sig
til, saa kunne vi heri være ganske og aldeles enige med ham;
% --- p. 304 ---
\opage{304}\setcounter{footnote}{0}%
men den treenige Guds Virkelighed er jo Aandsvirkelighed,
Virkelighed for Troen, indre Virkelighed for Aanden. Det
forholder sig ikke med Forvisningen om Guds Tilværelse
som med Forvisningen om Solens Tilværelse. Et Menneske
kan, om han end bliver legemlig blind og derved sat ud af
Stand til at see Solen, dog vedblive at føle sig forvisset om
Solens Tilværelse; men et Menneske kan ikke i aandelig
Forblindelse ophøre at troe paa Gud, og dog vedblive at føle
sig forvisset om Guds Tilværelse: den oprindelige Forvisning
om „de guddommelige Personers“ Tilværelse gives ikke ved
ydre Virkelighed for Sandserne, men ved indre Virkelighed
for Troen.

2) \emph{Aandsvirkeligheden gaaer fra det Indre
gjennem det Ydre tilbage til det Indre, saa at dens
afgjørende Kjendemærke altid maa søges i det
Indre}. Naar Spørgsmaalet er om den ydre Virkelighed, i
hvilken Gud personlig har aabenbaret sig, da maae vi særlig
tænke paa Christi Aabenbarelse i Kjødet. Med Hensyn paa
denne Kjendsgjerning har Hr. Professor Caspari da ogsaa
ganske rigtig henviist til 1 Joh. 1, 1 flgd. Med denne ydre
Virkelighed er det i den Grad Alvor, at Evangelisten er\-%
klærer enhver Aand, som vidner imod den, for en falsk
Aand. „I Elskelige! troer ikke hver Aand, men prøver
Aanderne, om de ere af Gud; thi mange falske Propheter
ere udgangne i Verden. Derpaa kjende I Guds Aand: hver
Aand, som bekjender Jesum Christum at være kommen i
Kjødet, er af Gud. Og hver Aand, som ikke bekjender
Jesum Christum at være kommen i Kjødet, er ikke af Gud,
og denne er Antichristens Aand, om hvilken I hørte, at
han kommer, og han er allerede nu i Verden“. (1 Joh. 4,
1---3). Men lad os nu see, hvorledes det Ydre i denne
Aabenbarelse stiller sig til det Indre. Er Forholdet dette:
at Visheden om Guds Søns Aabenbarelse begynder med det
Ydre og fra det Ydre gaaer over i det Indre? er Forholdet
% --- p. 305 ---
\opage{305}\setcounter{footnote}{0}%
ikke meget mere dette, at Visheden begynder fra det Indre,
dvæler i nogen Tid ved det Ydre, og saa tilsidst, efter at
dette Ydre er forsvundet, vender med forstærket Eftertryk til\-%
bage til det Indre? Dersom Erkjendelsen af, at Jesus af
Nazareth er Christus, oprindelig beroede paa det Ydre, da
maatte jo Alle, der med legemlige Øine saae Mennesket
Jesus, ogsaa af hans Ydre kunne have seet, at han virkelig
var Guds Søn. Men at dette ikke var Tilfældet, lærer den
evangeliske Historie. Den sandselig nærværende Jesus af
Nazareth var i Virkeligheden kun Guds Søn for dem, der
saae ham med Troens Øie, hørte hans Ord med Troens
Øre. Hvad det vil sige, at den i Kjød og Blod aabenbarede
Christi Virkelighed er en Virkelighed for Aanden, sees blandt
Andet af Matth. 16, 13---17. „Men der Jesus var kommen
i Egnen af Cæsarea Philippi, spurgte han sine Disciple ad,
og sagde: hvem sige Menneskene mig, Menneskens Søn, at
være? Men de sagde: Nogle sige, at du er Johannes den
Døber; men Andre Elias, men Andre Jeremias, eller En af
Propheterne. Han siger til dem: men I, hvem sige I mig at
være? Da svarede Simon Petrus og sagde: du er Christus,
den levende Guds Søn. Og Jesus svarede og sagde til ham:
salig er du, Simon Jonas Søn, thi Kjød og Blod haver ikke
aabenbaret dig det, men min Fader, som er i Himlen“.
Havde Christus været ligefrem kjendelig paa det Ydre, da
vilde der jo ikke have været saa afvigende Meninger om
ham, og der havde ikke været Grund til at prise Disciplen
salig, naar han kun saae, hvad Enhver, selv en Herodes og
en Pilatus, maatte kunne see. Hvorledes det Ydre i Christi
Aabenbarelses Virkelighed forholder sig til det Indre, lære vi
af Jesu egne Ord ved Afskeden med Disciplene (Joh. 16, 7 og
13---15). „Men jeg siger Eder Sandhed: det er Eder gavn\-%
ligt, at jeg gaaer bort; thi gaaer jeg ikke bort, skal Tals\-%
manden ikke komme til Eder“ . . . Dette maae vi jo forstaae
saaledes: det er Eder gavnligt, at min ydre Virkelighed for\-%
% --- p. 306 ---
\opage{306}\setcounter{footnote}{0}%
svinder for Eders Blik, thi kun under denne Betingelse kan
Aanden i indre Virkelighed berige Eder med den fulde Sandhed.
Det Ydres Omsætning i det Indre bliver udtrykkelig betegnet
paa en saadan Maade, at vi i Jesu egne Ord have en sikker
Veiledning til at opfatte og forstaae den hellige Historie.
Her siges, at Aanden „ikke skal tale af sig selv, men hvad\-%
somhelst han hører“ (fra Faderens og Sønnens indbyrdes Talen
det Indre) „skal han tale“. Aandens Fremstilling af Christi
% "Talen det Indre" as printed (sense perhaps "Talen i det Indre"); not emended
% (audit, 600 dpi: "det" starts slightly in from the margin, no ink where an i would stand)
Liv paa Jorden skal ikke indskrænke sig til en endelig Co\-%
pieren og Indregistreren af ydre Tildragelser; Sandhedens
Aand skal forklare, belyse og fremstille Frelserens Historie
saaledes, at det Ydre faaer sin Betydning fra det Indre; i dette
Indre er hans Herlighed, en Herlighed, som kun Troens Øie
kan see i Aandens Lys. „Han skal herliggjøre mig; thi han
skal tage af mit og forkynde Eder.“

3) \emph{Det Ydre, hvilket Aabenbaringens Aand
optager, gjennemvirker og forklarer, er af meget
forskjellig Art, saa at Kjendsgjerningerne i den
hellige Historie paa Grund heraf antage et meget
forskjelligt Virkelighedspræg}. „Gives der“, siger Hr.
Prof. Caspari, „nogen Kjendsgjerning, der er tilstrækkelig, ja
overflødig historisk bevidnet, saa er det Christi Opstandelse;
skal den ikke være historisk, da kunne vi gjerne slaae en
Streg over Alt, hvad der kaldes Historie“. Vi sige det
Samme, dog paa det Vilkaar, at denne Kjendsgjerning ikke
drages ned i Kredsen af de Begivenheder, der tilhøre Pro\-%
fanhistorien og ere Videnskritiken underkastede, men aner\-%
kjendes for hvad den i Sandhed er, en overnaturlig Kjends\-%
gjerning, en Kjendsgjerning, hvis Virkelighed med alt sit Ydre,
er Aandsvirkelighed, Virkelighed for Aanden.

At den hellige Historie optager og forener absolut ueens\-%
artede Virkelighedselementer, indlyser strax, naar vi besinde
os paa, at den i Forbindelse med naturlige Begivenheder inde\-%
holder Mirakler. Korsfæstelsen f. Ex. er en naturlig Begi\-%
% --- p. 307 ---
\opage{307}\setcounter{footnote}{0}%
venhed, Opstandelsen derimod en overnaturlig; dette giver en
væsentlig Forskjel, naar Talen er om en Vurdering af de
historiske Vidnesbyrd. At Jesus af Nazareth virkelig blev
korsfæstet under Pontius Pilatus, er en Kjendsgjerning, som
ingen enten jødisk eller hedensk Historieskriver har draget i Tvivl;
tvertimod, det blev med Haan og Spot saavel af Jøder som
Hedninger forekastet de første Christne, at de troede paa en
korsfæstet Misdæder. Anderledes med Opstandelsen; denne
Kjendsgjerning er hverken af Jøder eller Hedninger bleven an\-%
tagen for at være historisk virkelig; den er, siger den profan\-%
historiske Kritiker, ingenlunde tilstrækkelig bevidnet. Hvor\-%
ledes forholder det sig hermed?

Opstandelsen er et Mirakel, et af de stærkest bevidnede
Mirakler; men af hvem bliver et Mirakel bevidnet? Kun af
Troende, aldrig af Ikke-Troende! Dette er ingen Tilfældig\-%
hed; det ligger i Sagens Natur, det følger med uafviselig
Conseqvens af Miraklets Charakteer og Betydning.

Fra Magthandlingens Side kan der, ifølge Troens Forud\-%
sætning, vistnok Intet indvendes imod, at et Mirakel blev ud\-%
ført saaledes, at selv de mest Forstokkede maatte indrømme,
at de virkelig havde seet et Almagtens Tegn. Havde Christus,
f. Ex. da han vilde opvække Lazarus, ladet det høitidelig
bekjendtgjøre nogle Dage iforveien, saa at de, der vilde anstille
Iagttagelser, Pharisæere, Sadducæere, en Annas, en Cajaphas,
en Herodes, en Pilatus, kunde være tilstrækkelig forberedte;
havde han endvidere sørget for, at Enhver af de Iagttagende
kunde gaae ind i Graven og forvisse sig om, at Liget virkelig
var gaaet i Forraadnelse; havde han dernæst stillet disse Iagt\-%
tagere i Kreds omkring Graven og i deres Paahør udtalt
Ordene: Lazarus kom herud! med en saadan Virkning, at den
Døde virkelig saaes at reise sig og komme ud, bunden med Jord\-%
klæderne om Fødder og Hænder osv.; havde han saa endelig, efter
at dette var skeet, formaaet Iagttagerne, hvoriblandt maatte
være, ikke blot Tilhængere, men ogsaa Modstandere, dannede
% --- p. 308 ---
\opage{308}\setcounter{footnote}{0}%
Mænd, hvis Vidnesbyrd havde Vægt og Betydning i alle
Kredse, til, Een for Alle og Alle for Een, høitidelig at be\-%
vidne, hvad de havde seet og hørt: da maatte den profan\-%
historiske Kritik absolut have erklæret sig tilfredsstillet, da
maatte enhver forskende Historiker, han være forøvrigt tro\-%
ende eller vantro, indrømme, at dersom nogen Kjendsgjerning
„er tilstrækkelig, ja overflødig historisk bevidnet“, saa er det
Lazari Opvækkelse. Men saaledes er Historien om Lazari
Opvækkelse ikke bevidnet, og saaledes er Historien om Christi
Opstandelse heller ikke bevidnet.

Hvorfor hviler der nu et for profane Øine mystisk Slør
over det Mirakuløse i ethvert af den hellige Histories Mi\-%
rakler? Fordi det strider mod Miraklets Øiemed, at træde
saa udvortes øiensynlig frem, at en Betragter uden Tro skulde
ved det blotte Ydre kunne forvisse sig om, at det er en
Villie med overnaturlig Magt, som i en saadan Gjerning giver
sig tilkjende. Dersom sand Underkastelse under Guds Villie
kunde bevirkes ved et uomtvisteligt Almagtens Tegn i det
Ydre, da vilde Christus ved et Par Mirakler let have bragt
baade en Annas og en Cajaphas, baade en Herodes og
en Pilatus til at kaste sig i Støvet for den guddommelige
Majestæt. Pilatus, der er saa feig, saa angst for at falde
i Unaade hos Keiseren, fordi han ved at holde sig til det
Ydre, gaae ud fra det Ydre, og sætte Sandheden til paa det
Ydre, er forvisset om, at Keiseren har Magten, den Magt, som
gjør fattig, og som gjør rig, den Magt, som ophøier, og som
nedtrykker, den samme Pilatus skulde visselig have krøbet som
en Hund for Fødderne af den Mand, hvis almægtige Ord
kunde slaae en Legion af romerske Soldater til Jorden. Ved
at træde frem i en saadan ydre Virkekreds, vilde et Mirakel
af Christus ganske vist være blevet en verdenshistorisk uom\-%
tvistelig Kjendsgjerning; men dersom et eneste saadant Mi\-%
rakel var skeet, saa vilde Almagtsaabenbaringen med eet Slag
have knust al indre Selvstændighed, al menneskelig Person\-%
% --- p. 309 ---
\opage{309}\setcounter{footnote}{0}%
lighed og ladet en Messias efter Jødernes Sind, uendelig mæg\-%
tigere end Keiseren, men i Lighed med Keiseren, oprette et
jordisk Rige og herske som guddommelig Despot over Slaver.
Ganske vist har det virkelige Mirakel sit Ydre, ja, efter dets
forskjellige Bestemmelse, et meget forskjelligt Ydre; men det
Ydre har ikke her sin Betydning i og ved sig selv, det har
og maa have sin sande Betydning ved det Indre, ene og
alene ved det Indre: det forudsætter Troen, opfattes i Troen,
styrker Troen; det forudsætter Aanden, bevidnes af Aanden,
styrker Aanden.

At fordre et Mirakel for at grunde sin Tro paa det Ydre
alene som Ydre er Aandsforbrydelse, det er \emph{at friste Gud}. At
vove det Yderste for Troens Skyld og saa midt i Farerne stole
paa Guds vidunderlige Hjælp og Bistand er ikke --- hvad blød\-%
agtige Prædikanter stundom ved Anpriisning af „den christe\-%
lige Viisdom“ lade Menigheden underhaanden forstaae --- at
friste Gud; men at fordre Mirakler med et saa fremtrædende
historisk Ydre, at al Frelse og Salighed kunde bygges paa
umiddelbar Forvisning om alle historiske Enkeltheder i dette
historiske Ydre, fordi man bilder sig ind, at vi uden al
denne ydre Virkelighed vilde være og i al Evighed forblive
% a dot on the baseline after "være" (inked quad); not transcribed
i vore Synder, være og blive evig fortabte: det er i Sandhed at
friste Gud. Eet Mirakel er der, men kun eet eneste, om hvilket
den Troende maa sige, at dets Ydre, naar det virkelig ind\-%
træder, vil, ogsaa for Ikke-Troende, være hævet over al Tvivl,
det er Christi Aabenbarelse til Dommen; men denne Aaben\-%
barelse kan da heller ikke blive Gjenstand for historisk Kritik,
thi med den vil Historien ende.

Kun paa eet Vilkaar kan det siges, men da ogsaa med
fuld Ret, at Christi Opstandelse er tilstrækkelig, om end ikke
„overflødig“, historisk bevidnet, dette nemlig, at dens Virke\-%
lighed ikke eensidig sættes i det Ydre, men med alle Bestem\-%
melser af et Ydre henføres til og forenes med sit Indre, at
den, med andre Ord, ikke nedværdiges til at staae paa lige
% --- p. 310 ---
\opage{310}\setcounter{footnote}{0}%
Trin med anden historisk Virkelighed, men hæves til sin rette
Høide ved at opfattes og bestemmes som Aandsvirkelighed,
Virkelighed for Aanden. Den sande, fuldstændige Aands\-%
virkelighed er hverken et Ydre for sig eller et Indre for sig, men
et til sit Indre svarende Ydre, et ved sit Ydre bestemt Indre;
om en saadan i Kjendsgjerning udtalt Aandsvirkelighed kan
der kun vidnes i Troen: alle Opstandelsesvidnerne ere Troes\-%
vidner.

Efter disse Oplysninger vil det vel kunne skjønnes, at
er det ikke Hr. Prof. Caspari selv, saa er det heller ikke mig,
der sætter Virkeligheden af Christi Opstandelse for lavt; jeg
er mig bevidst, at den her fremsatte Anskuelse i Eet og Alt
stemmer med den i Forelæsningerne udtalte, og kan derfor
slutte med den Forsikring, at de Forestillinger om „blotte
Frelsesideer“, som Hr. Prof. vil bekæmpe, kun ved Misfor\-%
staaelse kunne tilskrives mine Ideer. Forstaaer sig, Ideer
kunne jo vel gjøre et Menneske søvnløst. Beklageligt vilde det
imidlertid være, om Ideerne, var det end i Form af „Frelses\-%
ideer“, nu med Eet skulde komme op og volde Uro og foruro\-%
lige en Mand, der hidtil har havt det saa roligt, at han som
„Videnskabsmand af Guds Naade“ aabent maa bekjende, at
„Videnskabens Indsigelser mod Frelsesaabenbaringens Grund\-%
kjendsgjerninger og Enkeltheder“ endnu aldrig „have voldt
ham nogen søvnløs Nat“.

Endnu førend hiin Udtalelse af Hr. Prof. Caspari blev
ndrykket i \emph{Aftenbladet}, havde \emph{Morgenbladet} leveret en
% printer's defect: "ndrykket" (for "indrykket"), the i missing at the line start; transcribed as printed
Række læseværdige Artikler under Overskrift: \emph{Epilog} til R. N.s
\emph{Forelæsninger}. Disse Artikler, der alle bære Spor af en øvet
Pen, ere affattede med frit videnskabeligt Overblik, i en vel\-%
villig Stemning, en velklingende, human Tonart og et let Sprog;
men, spørge vi, i hvilket Øiemed? Er det for at \emph{veilede}
Publikum eller for at \emph{vildlede} det? Aabenbart det Sidste!
Idet jeg uden videre Omstændigheder antager dette, vil en
Kyndig strax forstaae, at det ikke er nogen Dadel, men
% --- p. 311 ---
\opage{311}\setcounter{footnote}{0}%
tvertimod en Roes, jeg hermed udtaler; thi som \emph{veiledende}
vilde Artiklerne være temmelig maadelige, som \emph{vildledende}
ere de fortræffelige. „Men vildledende? Det er jo en For\-%
nærmelse mod Forfatteren!“ Paa ingen Maade; det er tvert\-%
imod en Compliment. Den ærede Forfatter har ved sin Op\-%
træden givet mig saamange Complimenter, at jeg maatte mangle
al Dannelse og Belevenhed, hvis jeg undlod pligtskyldigst at
gjøre ham min Contracompliment: til en Paradefægtning bør
man jo hilse med Kaarden.

\emph{Vildlede} og \emph{vildlede} er nemlig i philosophiske For\-%
handlinger to høist forskjellige Ting. Der gives en Vildleden,
der skader Sandheden, den plumpe, uhyggelige, som frem\-%
kommer derved, at en Umoden vil orientere Publikum uden
selv at være orienteret; under denne Kategorie tør den ærede For\-%
fatter paa ingen Maade henregnes. Men der gives ogsaa en
Vildleden, som ved dialektiske Streiflys lokker Publikum med
i Undersøgelsen, fører Læseren ud paa Glatiis og giver ham
en Forestilling om, at den, der nu virkelig vil ind i Sagen,
maa see med egne Øine og tænke selv. Forfatterens Vild\-%
leden er, hvad jeg haaber at kunne bevise, en høiere intel\-%
lectuel Veileden; det er den gamle, velbekjendte Gjordemoder\-%
kunst, han anvender: Methoden --- Enhver, der kjender lidt til
den græske Philosophie, seer det med et halvt Øie --- Me\-%
thoden er \emph{sokratisk}.

Strax i Begyndelsen, hvor Forf. (Epil. II) udtaler sig
om Indledningsforedraget, aner man, at her maa stikke Noget
under. I Anledning af et billedligt Udtryk, jeg i Talens Løb
skal have brugt, nemlig „\emph{Aandens Pust}“, yttrer han sig
paa følgende Maade: „Naar Aanden saaledes ved sin blotte
Blæsen forvandler Hindringer til Betingelser, saa kan det vel
paa den ene Side synes næsviist at spørge atter om Betin\-%
gelsen for, at Aanden gjør dette, eller om Grunden, hvorfor
den ikke altid blæser. Thi det kunde jo siges, at Aanden
blæser, hvor den vil, og vi høre dens Røst, men Ingen veed,
% --- p. 312 ---
\opage{312}\setcounter{footnote}{0}%
hvor den gaaer hen. Og hiint Spørgsmaal kan desuden,
som man let seer, gaae i det Endeløse. Men naar paa den
anden Side Aanden er en saadan Troldmand, der ubetinget,
naar og hvor den vil, forvandler Hindringer til Betingelser
eller omvendt, hvortil saa det hele Spørgsmaal om Hindringer
og Betingelser for det aandelige Liv? Begge Dele synes kun
at have Noget at gjøre, hvor der er en fornuftig begrundet
Sammenhæng og en lovbunden Virksomhed; ligeoverfor det
absolut Magiske ere baade Hindringer og Betingelser tomme
Ord. Altsaa paa den ene Side stiller sig i Udsigt en ende\-%
løs Ophoben af Betingelse paa Betingelse, paa den anden en
Bortvisen af det hele Spørgsmaal! Dette synes visselig at
være en Hindring for den begyndte Undersøgelse. Vil Ta\-%
lerens Aand maaskee blæse paa den og forvandle den til en
Betingelse, eller vil han blæse ad den og springe den forbi?“
% the „ that opens this quotation from the Epilogue is on p. 311
--- See, det var nu Udbyttet af den første Forelæsning, alt\-%
saa den første Veiledning.

De Læsere, som uden at have gjort sig bekjendte med
mit indledende Foredrag betroede sig til en saadan Veiled\-%
ning, maatte, ved at tage den ligefrem og for ramme Alvor, rig\-%
tignok faae en underlig Forestilling om min Optræden ved
det fremmede Universitet. Man tænke sig en hospiterende
Docent komme blæsende med nogle Talemaader om „Aan\-%
dens Blæst“, og det uden engang at have saameget Omløb i
Hovedet, at det midt i Ordstrømmen kunde falde ham ind at
spørge, „hvorfor Aanden ikke altid blæser“! Man tænke sig
en philosophisk Professor, der dog ifølge sin Stilling maa være
forpligtet til at vide lidt Besked med Forholdet imellem Be\-%
tingelserne og det Ubetingede, i den Grad mangle Reflexion,
at man bagefter skal minde ham om, at „ligeoverfor det ab\-%
solut Magiske ere baade Hindringer og Betingelser tomme
Ord“! Man tænke sig en Kreds af dannede Tilhørere, der\-%
iblandt endog høit ansete Videnskabsmænd, der i hele sexten
Timer har kunnet holde ud at høre paa noget saa Ugrundet,
% --- p. 313 ---
\opage{313}\setcounter{footnote}{0}%
i sig selv saa lidet Sammenhængende, som et saadant Fore\-%
drag over „Aandens Blæst“ lader forvente! Man tænke sig
den Uvillie, der nu bagefter maa vækkes i Publikum, naar
det dog, veiledet af Epilogen, skal tilstaae sig selv, at det
har ladet sig overraske, at de Talemaader, hvoraf det i et
ubevogtet Øieblik lod sig bedaare, kun indeholdt vindige
Aandrigheder, Luftbilleder i Skyer, som dreves af Aandens
Blæst. „Beklagelige Blæst, ubehagelige --- men dog gavnlige
--- Veiledning!“

At Veiledningen er gavnlig, skal jeg villig indrømme,
naar vi uden Pust og Blæst først faae det Veiledende for\-%
vandlet til et Vildledende. Dette skeer ganske simpelt ved
at agte paa, hvad vor Veileder i Epilogen kalder „en for\-%
nuftig begrundet Sammenhæng“. Der maa jo dog være en
fornuftig begrundet Sammenhæng imellem det, der virkelig er
blevet sagt i et Foredrag, og det, der bagefter siges derom.
At mit Indledningsforedrag, saaledes som det nu foreligger, i
alt Væsentligt er det samme som det, der blev holdt, kunne
mine Tilhørere bevidne; men, naar dette staaer fast, saa vil
Enhver, endog ved en flygtig Gjennemlæsning, kunne overbe\-%
vise sig om, at, hvad der i Epilogen staaer som et Veile\-%
dende, netop er vildledende. Men dette Vildledende er saa
igjen som ved et Smiil fra salige Guder blevet os til en sand
Veiledning. Den Aand, hvorom der tales i Indledningen, er
„den absolut guddommelige Aand“, altsaa den almægtige
Aand. Nu veed Epilogens Forfatter meget vel, at naar der
skal tales om Aandens Almagt, saa kan en Tilhører eller en
Tilhørerinde let høre feil, og bagefter mene, at der blev talt
om en Aand, der virkede magisk og derfor hverken behøvede
at ændse Hindringer eller Betingelser, og at dette just er en
forstyrrende Misforstaaelse, naar Foredraget udtrykkelig skal
handle om Hindringer og Betingelser. For at forebygge slige
intellectuelle Efterveer har den ærede Forfatter i bagvendt Stiil
villet indskærpe det Fornødne. Husk paa, hvad her er sagt:
% --- p. 314 ---
\opage{314}\setcounter{footnote}{0}%
\emph{Aanden er i sin Virken ingenlunde ligegyldig mod
Betingelser, ikke betingelsesløs, heller ikke skaber
den Betingelserne ligefrem af Intet; nei, den danner
dem af naturlige Forudsætninger, den danner dem
ved en Forvandling}. (S. 8). Denne Forvandling er ikke
magisk, thi husk paa, hvad her er sagt! Det gamle forslidte
Spørgsmaal: \emph{naar Aanden selv maa beaande sine
Tjenere og selv bringe Betingelserne med, er det
da ikke ubegribeligt, at den ikke altid og overalt
vil virke med lige kraftig Nærværelse}? er ikke flygtig
afviist, men udførlig besvaret ved en Belysning af Friheden
og Frihedens uafviselige Betingelser (S. 9---16), og at denne
for Grundtanken nødvendige Belysning netop udgjør Hoved\-%
indholdet af det indledende Foredrag --- see dette har Epilo\-%
gens snilde Forfatter aabenbart villet indskærpe. Men hvor\-%
dan har han saa baaret sig ad? For at indprente Tilhøreren,
hvor farligt det er at begynde med en complet Misforstaaelse,
har han selv begyndt med en complet Misforstaaelse; for at
vise det Komiske i, at en Tilhører glemmer de Indvendinger,
der i et Foredrag blive, ikke blot opstillede, men ogsaa imøde\-%
gaaede, har han selv --- komisk nok --- begyndt med at
glemme dem; for at forebygge, at Ingen med kort Hukom\-%
melse skal gjøre sig vigtig ved at have glemt det Vigtigste,
begynder Forf. med at lade som om det var ham, der blev
vigtig ved at have glemt det Vigtigste. I Sandhed, en Vei\-%
ledning, der udkræver Resignation og Selvfornegtelse: man
advarer mod Misforstaaelser ved at gjøre sig skyldig i Mis\-%
forstaaelse; det er, som om en Fader vilde advare sine Sønner
imod Drukkenskab ved selv at tage sig en Ruus.

I Epilog IV forekomme følgende ret vildledende Vink
til Veiledning. „Ingen skal, fordi han siger A i Videnskaben,
have Ret til at sige B i Troen. \ldots{} Bare nu ikke, strax dette
indrømmes, ialfald Logiken stikker sin Næse ind i dette Sy\-%
stem af Troesbegreber, og altsaa den magiske Kreds, hvor\-%
% --- p. 315 ---
\opage{315}\setcounter{footnote}{0}%
ved al Videnskab skulde holdes borte, bliver gjennembrudt.“
Herved giver Forf. alle dem en sokratisk Snært, som, uagtet
det atter og atter paavises, at Tænkningen, følgelig ogsaa
Logiken, ligesaa vel har Gyldighed i Troeslæren som i Viden\-%
% [a raised space/quad prints as a faint mark between "i" and "Troeslæren"; not typed]
skaben, ere hildede i den gamle indgroede Fordom, at Troen
som Tro er tankeløs og maa laane alle Tanker fra Viden\-%
skaben. For at drive denne Misforstaaelse indtil Karikatur,
udmaler han nu ret sindrigt, hvad der kommer ud af at give
sig hen i en Slags videnskabelig Phantaseren over Troens og
Videns ueensartede Tankekredse. „Jeg veed ikke,“ siger han,
„til hvilken Grad jeg tør fordre indbyrdes Overeensstem\-%
melse paa Troens Omraade. Hvorledes forholder Troen sig
til den hellige Skrift? Ere disse maaskee ogsaa ueensartede,
saa at, hvad der gjælder i den ene, ikke maa overføres til
den anden? --- Og den hellige Skrift selv? Der er jo paa\-%
viist indbyrdes Modsigelser, f. Ex. mellem 2 Samuels og 1
Krønikernes Bog. Kunne disse da fjernes ved den meddeelte
Trylleformel, at A og B ikke maa tilhøre ueensartede Ge\-%
beter? Er da Bibelen selv indbyrdes saa ueensartet, at det
Ene gjerne kan være sandt i Samuels Bøger, medens det Mod\-%
satte er sandt i Krønikernes Bøger?“

Af den iøinefaldende Overlegenhed, hvormed Forf. her
giver det Vildledende frit Spillerum, mærker man, hvor let
det vilde have været ham at holde sig til hvad der er udtalt
i Forelæsningerne, og give os en ligefrem Veiledning. At
det maa have kostet ham Selvovervindelse ikke at gjøre det,
tør jeg vel antage; thi Fristelsen har sandelig været stor.
Hvorvidt en Slutning af givne Forudsætninger lader sig drage,
beroer jo ikke paa dit og mit Forgodtbefindende; det beroer
paa Tankelovene. Skulde Veiledningen ikke, ifølge Methoden,
absolut have været vildledende, saa vilde Epilogens Forfatter
naturligviis stærkt have fremhævet, at der i Forelæsningerne
slet ikke er opstillet nogen Trylleformel, men i bestemte Ex\-%
empler paaviist, hvad det fører til, naar man uddrager Videns\-%
% --- p. 316 ---
\opage{316}\setcounter{footnote}{0}%
conseqvenser af Troens og Troesconseqvenser af Videns For\-%
udsætninger. Spørgsmaalet er: lader det sig gjøre eller ikke?
Hvis det virkelig lader sig gjøre, saa vise man Docenten,
hvorledes det lader sig gjøre, og hvori hans Feiltagelse be\-%
staaer; hvis det ikke lader sig gjøre, saa har Docenten jo
Ret, og hvad hjælper det saa med de mange svævende, over\-%
svævende, udsvævende Overveielser angaaende Trylleformler,
hvor der ingen Trylleformel er?

Den Taktik, Forf. med saameget Held har anvendt til
at vildlede Læserne angaaende Troesbegrebernes indre Over\-%
eensstemmelse i deres Kreds, den samme Taktik anvender han
med ligesaa glimrende Virkning til at forvirre Forestillingerne
angaaende Vidensbegrebernes indre Sammenhæng. „Der er“,
siger han, „saavidt jeg husker, ikke talt Noget om indbyrdes
Modsigelser paa Videnskabens Omraade. Maaskee forudsættes
det, at de ikke ville forekomme. Men om de alligevel fore\-%
kom? Jeg kan ialfald ikke føle mig ganske tryg. Thi jeg
har dog hørt Tale om Stridigheder mellem forskjellige Viden\-%
skaber, der, naar de fremsættes med Skarphed, let kunne
synes ligesaa radicale og utilgængelige for Mægling, som
Striden mellem Samuel og Krønikerne, ja som flere af de
Modsigelser, der synes at finde Sted mellem Troen og Viden\-%
skaben overhoved.\ldots{} Jeg er overhoved slet ikke uden Frygt
for, at den ved den nye Lære satte absolute Skilsmisse mellem
Tro og Videnskab conseqvent --- hvis Conseqvens ikke tilsidst
er et ganske tomt Ord --- vil føre til at opløse baade Troens
og Videnskabens indre Eenhed og søndersplitte begge Dele i
Atomer, saa at der hverken bliver nogen Tro eller nogen
Videnskab mere, men kun en ubestemmelig Mængde af løs\-%
revne Troesforestillinger og videnskabelige Begreber eller Sæt\-%
ninger. Virkelig synes der allerede at vise sig en herskende
Tilbøielighed til idetmindste at opfatte Videnskaberne udeluk\-%
kende enkeltviis, medens Tanken om en almindelig Videnska\-%
bens Organismus, hvoraf de enkelte Fag ere integrerende Led
% --- p. 317 ---
\opage{317}\setcounter{footnote}{0}%
forholdsviis træder i Baggrunden.“ Her have vi igjen det
sokratiske Skjelmeri.

Den ironiserende Forfatter vil bilde troskyldige Læsere
ind, at fordi Religionen er ueensartet, absolut ueensartet med
Videnskaben, saa skulde deraf følge, at samtlige Videnskaber,
naar de mærkede, at Troen ikke længere var istand til at holde
Styr paa dem, kunde faae isinde at gjøre sig afspænstige og
erklære sig for indbyrdes absolut ueensartede. Han maler et
mørkt Billede. Det kunde da skee, at endog hver enkelt Viden\-%
skab, naar den først havde smagt det absolut Ueensartede, fik
Værk i Lemmerne og begyndte at partere sig selv; at Arithmetik
og Geometrie f. Ex., naar de fik Nys om, at Noget kunde
være sandt i 2 Samuel (XXIV), som var usandt i 1 Krø\-%
niker (XXI), og omvendt, toge den Beslutning, at Noget for
Fremtiden skulde kunne være sandt i Arithmetiken, som var
usandt i Geometrien, og omvendt. Det er truende Udsigter!
En eller anden „Videnselsker“, der endnu ikke ret er orien\-%
teret i den lærde Republik, kunde ved slige ulykkespaaende
Varsler betages af Angst, faae Skræk i Blodet, bilde sig ind,
at vi allerede havde Anarchiet, at Videnskaberne vare i fuldt
Oprør, at der var Ildløs i det philosophiske Qvarteer, og give
sig til at raabe Brand!

Jeg maa til Gemytternes Beroligelse gjøre opmærksom
paa, at Sagen er ikke saa farlig. Videnskaberne ere jo alle
af en sindig, rolig og i Grunden fredsommelig Charakteer;
vistnok have de adskillige indbyrdes Stridigheder; disse Stri\-%
digheder kunne vel ogsaa nu og da blive noget heftige og
lidenskabelige; men det Heftige og Lidenskabelige udgaaer
fra de menneskelige „Videnselskere“, ikke fra Videnskaberne.
Det er netop ved de videnskabelige Sandheders stille uperson\-%
lige Magt over de lidenskabelige Personligheder, at den tabte
Ligevægt tilveiebringes, og Harmonien bevares. Skulde Noget
bringe Videnskaberne i Oprør, da maatte det være, at de
igjen bragtes ind under Troens Formynderskab, ind under
% --- p. 318 ---
\opage{318}\setcounter{footnote}{0}%
Dogmatikens „gode Villie“, ind i det theologiske Universal\-%
monarchie. Jeg har virkelig undersøgt Sagen, jeg har været
omkring og forhørt mig hos de enkelte Videnskaber, og kan fra
hver og een (Jvnfr. „Den gode Villie som Magt i Videnskaben“
S. 9---108) skaffe Attest for, at de alle ubetinget ere enige i
ikke at ville optage mirakuløse Elementer. Derimod skulde
jeg hilse og sige, at naar de blot maatte være frie for
Aabenbaringer, Mirakler og Spaadomme, skulde de al\-%
drig sige et ondt Ord enten imod Troen eller imod de Tro\-%
ende, men i alle personlige Anliggender opføre sig sædelig og
fornuftig, og vel vide at iagttage den Værdighed, de skylde
sig selv ligeover for Religionen. Vor Sokratikers Frygt for,
at „Videnskab uden Christendom i sin yderste Conseqvens
maatte ophæve sig selv og gaae tilgrunde i Skepticisme“, er
altsaa en vildledende Veiledning.

For at det dog ikke skal see ud, som om jeg i Anled\-%
ning af Forskrækkelsen nu først kom i Tanker om, at Viden\-%
skaberne udgjøre en Organisme, skal jeg her tillade mig at
henvise til „Philosophisk Propædeutik“ i Grundtræk. Kbhvn.
1857. S. 68---117, og særlig fremhæve Følgende (S. 71):
„Af Forholdet mellem Erkjendelseslæren og Videnskabslæren
fremgaaer, at den sande omfattende Videnskabslære kun kan
komme istand derved, at de forskjelligartede Videnskabers
Grundinteresser repræsenteres i de Forhandlinger, hvoraf den
fremgaaer. Under „den omfattende Discussion og Debat af
Alt imod Alt“ (Sibbern) har altsaa hver Videnskab sit Votum,
og yder sit Bidrag. At Philosophien sammenarbeider, klarer
og organiserer disse Bidrag til et intellectuelt Hele, beviser,
at den i saa Henseende ikke afgjør Sagen paa egen Haand,
men virker som Videnskabernes Fællesorgan, som deres egen
Almeenbevidsthed, det Væsen, der er deres Slægtskab, det
Lys, hvori de kjende og erkjende hverandre indbyrdes.“

Den Fixidee, hvorover vor Sokratiker paa det Aller\-%
fineste ironiserer, og det med Rette, da den endnu spøger
% --- p. 319 ---
\opage{319}\setcounter{footnote}{0}%
i Saamanges Hjerne, er netop Forestillingen om Christendommens
Nødvendighed til at begrunde, opretholde og fremme Viden\-%
skaberne, hvoraf igjen følger Videnskabernes Nødvendighed
til at begrunde, opretholde og fremme Christendommen. „For\-%
udsat,“ siger han (Epil. V), „at der var nogen Sandhed i den
Tanke, jeg ikke endnu har kunnet befrie mig fra, nemlig, at
der dog er en indre og væsentlig Sammenhæng mellem Viden\-%
skab og Tro, saa vilde deri altsaa ligge et rimeligt Svar paa
det \ldots{} Spørgsmaal, hvorfor dog --- i det Store og Hele ---
Videnskab og Tro ikke have kunnet lade hinanden være, men
enten søgt at forenes eller bekæmpet hinanden. Der kan vist\-%
nok have været meget Sophisteri baade i Striden og i For\-%
eningsforsøgene og navnlig mangt et tilsigtet Overgreb fra
Videnskabens Side; men Sophisteriet og Overgrebet vilde være
uden Motiv, hvis Troen kunde være Videnskaben i det Hele
ligegyldig. Er det derimod saa, hvad ogsaa Historien synes
at vise, at visse Troesformer ville udelukke Videnskaben, og
at dennes Trivsel overhoved er betinget af et vist herskende
Religions-Standpunkt --- saa at f. Ex. Videnskaben og den
hele dermed forbundne Cultur aldrig har kunnet faae den Dybde
og Udstrækning paa muhamedansk, som paa christelig Bund --- saa
har jo Videnskaben god Grund til ikke at lade Troens Indhold
og Form være sig uvedkommende. Kun den Tro, som er
væsentlig reen og sand og fri, kan derfor ogsaa lade Viden\-%
skaben frit ndvikle sig; men netop den sande Tro behøver
% "ndvikle" as printed (for "udvikle"; the first letter is plainly n, probably a turned u)
heller ikke --- f. Ex. ved en tvivlsom exceptio fori eller andre
Midler --- at unddrage sig den modne Videnskabs Kritik.
Thi den fulde Sandhed taaler fra alle Sider det fulde Lys.“

Historisk og Uhistorisk, Sandt og Halvsandt, Løst og
Fast har Forf. her paa en saare genial Maade kjørt imellem
hverandre, idet han under allehaande Bemærkninger snurrer
Læseren een Gang, to Gange, tre Gange rundt og hjælper
ham paa det Behændigste til at see feil af det, hvorpaa det
væsentlig kommer an --- naturligviis for at sætte hans Selv\-%
% --- p. 320 ---
\opage{320}\setcounter{footnote}{0}%
tænkning paa Prøve. Hvad Historien lærer om de mange
Stridigheder mellem Tro og Viden og de tilsvarende Forso\-%
ningsforsøg, har jeg jo netop søgt at vise i saa bestemte
Træk, at Ingen, som vil gjøre sig den Uleilighed at gjennem\-%
gaae mine Skrifter, let vil tage feil af Resultaterne. Den, der
% [a pencil mark over the n of "Den"; not transcribed]
endnu kan mene, at her efter saamange strandede Mæglings\-%
forsøg er mindste Udsigt til enten at gjøre Videnskaben tro\-%
ende eller Toen videnskabelig, maa jo være uden al Kund\-%
% "Toen" as printed (for "Troen"; r dropped, confirmed at 300 ppi)
skab til det sidste Aarhundredes theologiske og philosophiske
Litteratur.

Det er vist ogsaa dette, hvorom Forfatteren paa sin
sokratisk-bagvendte Maade vil minde. „Der er,“ siger han,
„en indre væsentlig Sammenhæng mellem Videnskab og Tro,
Troen kan ikke være Videnskaben ligegyldig“. Den, der tager
disse Sætninger bogstavelig, bliver fixeret, thi --- at vi skulle
gjentage det endnu en Gang --- der er i vore Dage ingen
sand Videnskab, være sig empirisk eller apriorisk, der paa
nogen Maade enten retter eller vil rette sig efter Troen, nem\-%
lig den positive Aabenbaringstro, og det er om den vi tale;
derimod presses der Mening ind i Talemaaderne, naar man
forstaaer dem saaledes, at \emph{det ingenlunde kan være Per\-%
sonligheden ligegyldigt, hvorledes de religiøse og
de videnskabelige Sandheder forholde sig til hin\-%
anden}.

„Den sande Tro behøver ikke at unddrage sig den modne
Videnskabs Kritik,“ denne Sætning er vildledende, naar den
opfattes som ligefrem veiledende. Hvad bliver der af Evan\-%
geliet om Christi Fødsel, om Christi Opstandelse, om Christi
Himmelfart, om Aandens Udgydelse o. s. v., naar den hellige
Historie virkelig skal underkastes „den modne Videnskabs
Kritik“? Giv mig en Prøve paa et saadant Stykke Kritik,
en Kritik, som ligelig tilfredsstiller baade Troen og Videnska\-%
ben, og jeg skal sige med Sokrates: ἀγαιμην ἀν θαυμαστως,
ω Ζηνων. Forstaae vi derimod Sætningen som vildledende,
% Greek (Plato, Parm. 129c) set as printed: smooth breathings on ἀγαιμην and ἀν are
% visible at 400 dpi; no accents, and no breathing on ω, are visible (standard text:
% ἀγαίμην ἂν θαυμαστῶς, ὦ Ζήνων). θ is printed in the curly form ϑ.
% AUDIT: re-read letter by letter at 600 dpi: α γ α ι μ η ν / α ν / ϑ α υ μ α σ τ ω ς , / ω Ζ η ν ω ν .
% Both breathings are smooth (psili); no acute, grave or circumflex anywhere, none on ω or Ζηνων.
% --- p. 321 ---
\opage{321}\setcounter{footnote}{0}%
saa er der Mening i den; Meningen er da denne: den sande
Tro behøver ikke at frygte den modne Videnskabs Kritik,
thi dersom den Troende tillige er en Vidende, da maa han
jo vide, at Videnskritiken just opløser de hellige Kjendsgjer\-%
ninger, fordi disse Kjendsgjerninger have deres Sandhed, ikke
i Viden, men i Troen.

Til Stadfæstelse af min Synsmaade, at Epilogens For\-%
fatter veileder bagvendt ved at vildlede, kunde jeg endnu anføre
mangfoldige, endog glimrende Exempler; men Tiden tillader
det ikke. For at afvise den dumme Snik-Snak, hvormed
Theologerne saa ofte opvarte os, at Visheden om det Sandse\-%
lige beroer paa Tro, at Astronomien beroer paa Tro, at In\-%
ductionen og dermed al empirisk Viden og Videnskab beroer
paa Tro, og at en saadan Tro saa igjen er beslægtet med
religiøs Tro, har jeg (navnlig i IV Forel.) stærkt udhævet
\emph{Visheden} af den empiriske Viden. For at gjøre denne
Vished indlysende for Udenforstaaende, har jeg beraabt mig
paa et naturligt \emph{Sandhedsinstinkt}; dette faaer vor So\-%
kratiker vittigt bøiet om til den bagvendte Paastand, at jeg
skulde være ifærd med at ville grunde Videnskaben paa „det
uformidlede Instinkt“; han minder os om den Sætning af Des\-%
cartes, at al menneskelig Vished beroer paa Tilliden til Guds
Sanddruhed. Herved har han paa en indirect og meget fiin
Maade henviist de Læsere, der misforstaae min Fremstilling,
som om jeg byggede Videnskaben paa det raa Instinkt, til
min Videnskabslære, navnlig til Grundideernes Logik II (S.
176---275).

Hvilket ironisk Skjelmeri der gaaer igjennem Forfatterens
majeutiske Fordreininger, derpaa anfører jeg endnu et Ex\-%
empel, hvori Lune og Vid næsten stige indtil Overgivenhed.
Jeg har (Forel. V. S. 81) udtalt mig saaledes: Spørges der f. Ex.
om, hvad Granit og Flint er, da kunne disse Begreber med tilsva\-%
rende Kjendsgjerninger fastsættes saa bestemt, at Enhver, som har
% --- p. 322 ---
\opage{322}\setcounter{footnote}{0}%
sund Forstand, Fordannelse og Sands for videnskabelige Un\-%
dersøgelser, han være forresten Jøde, Græker, Muhamedaner,
uden Hensyn til folkelige Forskjelligheder og Tro, bliver
tvungen til at indrømme Begrebernes Gyldighed: „Dette maa
være Granit, dette maa være Flint“. \ldots{} Hvad har vor So\-%
kratiker ikke vidst at bringe ud heraf? „Der foreholdes mig“,
siger han, „endnu et Kjendetegn, hvorpaa jeg skal kjende
Troesbegrebet fra det videnskabelige, nemlig, at det sidste kan
gjøres ethvert Menneske begribeligt, medens det første forud\-%
sætter et vist religiøst Standpunkt. Suurstoffets Begreb og
Geologiens Sætninger maae kunne fattes og tilegnes af Christne,
Muhamedanere og Brahminer uden Forskjel, medens man for\-%
gjæves vil søge at overbevise Andre end Christne om Christi
Opstandelse“. Den behændige Fordreielse bestaaer nu deri,
at medens jeg hos ethvert Menneske gjør „sund Forstand, For\-%
dannelse og Sands for videnskabelige Undersøgelser“ til udtrykke\-%
lige Betingelser for at forstaae en videnskabelig Læresætning,
springer vor Sokratiker rask væk Betingelserne over, og lader mig
i al Almindelighed tale om „ethvert Menneske“. Den Over\-%
drivelse, at \emph{ethvert Menneske} uden videre skulde kunne
overbevises om en videnskabelig Sandhed, bliver saa (Epil.
VI) illustreret ved Exempler som disse, „at det maaskee skulde
være ligesaa vanskeligt at faae en Chineser til at fatte euro\-%
pæisk Ret og europæisk Historie, som til at fatte Christen\-%
dommen; at det er vanskeligt at faae Tydskerne til at begribe,
at den største Deel af Slesvig er dansk, og at en Neger, der
troer ved Hjælp af en Fetisch at kunne vække eller stille
Storm, vil have vanskeligt nok ved at fatte Tanken om ufor\-%
anderlige Naturlove“.

Hermed slutter jeg disse Forhandlinger, og takker Epi\-%
logens ærede Forfatter for den væsentlige Tjeneste, han ved
denne Anledning har ydet mig. Epilogen har en bred Strøm;
Strømmen gaaer igjennem hele 9 Artikler, indtil den i en
% --- p. 323 ---
\opage{323}\setcounter{footnote}{0}%
10de omsider udmunder i en „Epilogens Epilog“. Just dette
er det sokratisk Fortræffelige. Skulde jeg udførlig og Punkt
for Punkt oplyse enhver af de subtilt fingerede Misforstaaelser,
da maatte jeg skrive en tyk Bog. Og hvad blev saa Resul\-%
tatet? At Publikum, der havde faaet nok af Artiklerne, be\-%
takkede sig for at læse Bogen. Dette har Epilogens For\-%
fatter meget snildt beregnet og just derfor villet give mig An\-%
ledning til at faae det constateret.

Der kan, som bekjendt, altid siges Adskilligt for og imod
enhver Ting, og da især for og imod en ny philosophisk Lære.
At der reiser sig Modsigelser mod min Philosophie, er
saa langt fra at overraske mig, at det Modsatte tvertimod
vilde have beredt mig en ingenlunde behagelig Overraskelse.
„\emph{Hvad er Philosophie}? En Videnskab fuld af Modsigel\-%
ser! Denne Definition, hvormeget den end, fra een Side be\-%
tragtet, synes anlagt paa at bryde Staven over al Philosophie,
er ikke desto mindre begrundet saavel i Videnskabens Historie,
som i Sagens Natur. For enhver anden Videnskab vilde en
saadan Tilstaaelse vel være tilintetgjørende, for Philosophien
angiver den et væsentligt Udgangspunkt, et paalideligt Støtte\-%
punkt, et fast, urokkeligt Grundlag. At Philosophien maa
indeholde Modsigelser, er jo en Selvfølge, eftersom det netop
er den Videnskab, der overalt, saavel i Tilværelsen som i
den sunde Forstands Opfattelse af Tilværelsen (Propæd.
S. 40---44) opdager og paaviser Modsigelser; Philosophien
bliver just med Hensyn hertil at opfatte som en \emph{almindelig
Modsigelseslære}. Som Modsigelseslære maa Philosophien
imidlertid ikke blot opdage og paavise Modsigelserne, den maa
ogsaa prøve Modsigelserne; thi der er Forskjel paa Modsi\-%
gelser. Nogle Modsigelser ere nemlig overfladiske, unødven\-%
dige, lette at undgaae, andre derimod paatrængende og uaf\-%
viselige; nogle Modsigelser indsnige sig paa Grund af Tanke\-%
løshed, andre fremkunstles ved Spidsfindighed, andre afsløres
% --- p. 324 ---
\opage{324}\setcounter{footnote}{0}%
ved Skarpsindighed; nogle Modsigelser spille som i en Tan\-%
keleg, andre angive væsentlige Knuder, vanskelige, tungt op\-%
løselige Problemer.“ (Forelæsninger over „Philosophisk Pro\-%
pædeutik“ fra Universitetsaeret 1860---61. Kbhvn. 1862.
S. 1---2).
% "Universitetsaeret" as printed (for "Universitetsaaret"; the second letter after s is plainly an e, unlike the a before it; AUDIT 600 dpi)

Den Omstændighed, at Philosophien bevæger sig igjen\-%
nem Modsigelser af forskjellig Art og meget forskjellig Dybde,
gjør det saa saare let at føre en Opinionspolemik imod et
% [a thin acute-like stroke over the second a of "saa"; a stray mark, not typed]
hvilketsomhelst philosophisk Lærebegreb; thi Publikum blendes
let, naar det, uden dybere Indsigt i Sagens Sammenhæng,
seer de til løsrevne Citater støttede, i faa Sætninger opstillede
Modsigelser. Faaer den sunde Forstand en Modsigelse under
een Belysning, kan den øieblikkelig gjennemskue dens Intethed,
og bliver ikke bedaaret; faaer den derimod den samme Mod\-%
sigelse under en anden Belysning og i en anden Forbindelse,
saa imponeres den, og troer at det er Noget. Siger jeg:
„Gud og Verden kunne ikke staae i Modsætning til hinanden
uden at begrændse hinanden; de kunne ikke begrændse hinan\-%
den uden at indskrænke hinanden. Gud maa altsaa ligesaa
vel være indskrænket af Verden, som Verden af Gud; men
en indskrænket Gud er en Modsigelse, altsaa er Gud ikke til“:
saa mærker den sunde Forstand Uraad. Man behøver blot
at minde om Forskjellen mellem Begrændsning og Indskrænk\-%
ning, saasom, at ethvert Hoved maa være begrændset, men
at deraf ingenlunde følger, at ethvert Hoved er et indskrænket
Hoved; at en Kunst kan have sin Grændse i en anden, uden
at den ene derfor indskrænker den anden, eller \emph{gjør den
anden Afbræk} o. s. v. --- og Enhver, der har naturlig Forstand,
fatter strax Forskjellen mellem det at \emph{begrændse} d. e. \emph{be\-%
stemme}, og det at \emph{indskrænke} d. e. \emph{gjøre Afbræk}. Er
Talen derimod om Tro og Viden, altsaa om en Modsætning, som
for den almene Bevidsthed endnu er temmelig uklar, saa er
det let at imponere ved Modsigelsen: „Tro og Viden be\-%
% --- p. 325 ---
\opage{325}\setcounter{footnote}{0}%
grændse hinanden, thi Tro er ikke Viden, Viden ikke Tro, og
dog skulle de ikke indskrænke hinanden; men at begrændse er jo
at indskrænke; altsaa: Tro og Viden skulle indskrænke hinan\-%
den derved, at de ikke indskrænke hinanden. Hvilken Mod\-%
sigelse, ja hvilken uhyre Modsigelse!“ Læserens Forstand staaer
stille, Opinionspolemiken florerer, og Publikum troer, at det
er Noget.
% [pencil: a stroke after "altsaa:", and a marginal note beside ll. 1--6,
% "Jvf. Høgaard, R. Nielsen om Tro og Viden, p. 16" (reading uncertain); not transcribed]

Til at oplyse alle de Huulheder, Fordreielser og plumpe
Misforstaaelser, der kunne rummes i en Piece paa nogle faa
Ark, maatte jeg, som sagt, hver Gang skrive en voluminøs
Bog; det er et ufrugtbart og uoverkommeligt Arbeide. Her
kommer En i den gøglende galante Vittighedsmaneer med
„Dualisme“ og „Genialitet“ og „Kulsviertro“ og præsenterer
Publikum halvanden gesvindt fabrikeret Modsigelse, og saa
skal jeg svare paa det! En Anden i den dogmatisk trum\-%
fende Høitryksmaneer med et theologisk Universalmonarchie
paa Spidsen af „den gode Villie“ træder op og foreviser Publi\-%
kum en Blind, der ikke rider paa en Halt, en billedlig,
sindbilledlig Modsigelse, og saa skal jeg svare paa det! En
Tredie i speculativ-kalkunsk Pluddermaneer melder sig med
tolv Stød i en mystisk Trompet; hvert Stød er en Modsigelse,
en kalkunsk Modsigelse, og saa skal jeg svare paa det! En
Fjerde i den selvvigtige Bornertheds velbekjendte Pænemands\-%
maneer kommer anstigende med Noget paa en kritisk Liimstang,
han kalder en Overbeviisning; Overbeviisningen er en Modsigelse,
og Liimstangen er ogsaa en Modsigelse; han slaaer til min Lære\-%
bygning med Overbeviisningen og Modsigelserne og Liim\-%
stangen og siger Ruin! saa skal jeg svare paa det! Her
kommer en Femte, en Sjette \ldots. --- det overgaaer mine
% four points, as printed
Kræfter. Maaskee de Herrer Piecemagere nok vilde leve af
Svar; men det er vist, at skulde jeg svare dem Alle og paa
Alt, saa maatte jeg døe af Svar.

Fra disse Vidtløftigheder, Bryderier og Besværligheder
% --- p. 326 ---
\opage{326}\setcounter{footnote}{0}%
har nu min sokratiske Velgjører ved sine fingerede Mistyd\-%
ninger, Mistydninger, der paa en let og behagelig Maade
karikere, hvad Publikum ellers saa gjerne tager for alvorlige
Indsigelser, indtil videre befriet mig. Epilogens Forfatter har
rakt mig en hjælpsom Haand; jeg takker ham rørt og trykker
hans Haand med Erkjendtlighed.

\begin{center}\rule{3cm}{0.4pt}\end{center}

\backmatter
% --- Rettelser (PDF 338), unnumbered leaf after p. 326: no page marker.
% --- PDF 338: Rettelser ---
\chapter*{Rettelser.}
% the head "Rettelser." is letterspaced; a short rule stands below it
\addcontentsline{toc}{chapter}{Rettelser}

\noindent
\begin{tabular}{@{}l r l l@{}}
S. & 50  & Linie 1 f.\ o. & vente at finde den \emph{læs}: vente det\\
-  & 70  & --- 1 - -      & stridt mod \emph{læs}: stridt med\\
-  & 109 & --- 3 f.\ n.   & det absolute Selv \emph{læs}: det Absolute selv\\
-  & 174 & --- 9 og 10 f.\ o. & Saducæer \emph{læs}: Sadducæer\\
\end{tabular}
% "læs" is letterspaced in all four rows; "-" and "---" are the printer's ditto marks.
% Row 3: the final v of "selv" is imperfectly printed (heavy, filled-in blot, checked at 600 dpi); read as v.

\begin{center}\rule{3cm}{0.4pt}\end{center}

\end{document}
